% Generated mechanically from frozen input inputs/p03/ja/tex/cohomology.tex.
% Do not edit this generated file; edit the bound locale source instead.

% OK, start here.
%

\setcounter{chapter}{19}
\chapter{層コホモロジー}



\phantomsection
\label{cohomology-section-phantom}


\section{序論}
\label{cohomology-section-introduction}

\noindent
本章では，位相空間上の層コホモロジーに関するいくつかの話題を扱う。
主として環の層上の加群という一般性で論じ，環付き空間の射を扱う。
サイト上の層の場合については，《サイト上のコホモロジー》第
\ref{sites-cohomology-section-introduction} 節を参照せよ。
基本文献は \cite{Godement} および \cite{Iversen} である。





\section{層のコホモロジー}
\label{cohomology-section-cohomology-sheaves}

\noindent
$X$ を位相空間とし，$\mathcal{F}$ をアーベル群の層とする。
$X$ 上のアーベル群の層の圏は十分多くの入射対象をもつ。《入射対象》補題
\ref{injectives-lemma-abelian-sheaves-space} を参照せよ。
したがって，入射分解
$\mathcal{F}[0] \to \mathcal{I}^\bullet$ を選べる。慣例に従い，
\begin{equation}
\label{cohomology-equation-cohomology}
H^i(X, \mathcal{F}) = H^i(\Gamma(X, \mathcal{I}^\bullet))
\end{equation}
と定め，これをアーベル群の層 $\mathcal{F}$ の
{\it 第 $i$ コホモロジー群}と呼ぶ。関手族 $H^i(X, -)$ は
$\textit{Ab}(X) \to \textit{Ab}$ の普遍 $\delta$-関手をなす。

\medskip\noindent
$f : X \to Y$ を位相空間の連続写像とする。上の
$\mathcal{F}[0] \to \mathcal{I}^\bullet$ を用いて，
\begin{equation}
\label{cohomology-equation-higher-direct-image}
R^if_*\mathcal{F} = H^i(f_*\mathcal{I}^\bullet)
\end{equation}
と定め，これを $\mathcal{F}$ の {\it 第 $i$ 高次順像}と呼ぶ。
関手族 $R^if_*$ は $\textit{Ab}(X) \to \textit{Ab}(Y)$ の
普遍 $\delta$-関手をなす。

\medskip\noindent
$(X, \mathcal{O}_X)$ を環付き空間とし，$\mathcal{F}$ を
$\mathcal{O}_X$-加群とする。$X$ 上の $\mathcal{O}_X$-加群の圏は
十分多くの入射対象をもつ。《入射対象》補題
\ref{injectives-lemma-sheaves-modules-space} を参照せよ。
したがって，入射分解
$\mathcal{F}[0] \to \mathcal{I}^\bullet$ を選べる。慣例に従い，
\begin{equation}
\label{cohomology-equation-cohomology-modules}
H^i(X, \mathcal{F}) = H^i(\Gamma(X, \mathcal{I}^\bullet))
\end{equation}
と定め，これを $\mathcal{F}$ の {\it 第 $i$ コホモロジー群}と呼ぶ。
関手族 $H^i(X, -)$ は
$\textit{Mod}(\mathcal{O}_X) \to \text{Mod}_{\mathcal{O}_X(X)}$ の
普遍 $\delta$-関手をなす。

\medskip\noindent
$f : (X, \mathcal{O}_X) \to (Y, \mathcal{O}_Y)$ を環付き空間の射とする。
上の $\mathcal{F}[0] \to \mathcal{I}^\bullet$ を用いて，
\begin{equation}
\label{cohomology-equation-higher-direct-image-modules}
R^if_*\mathcal{F} = H^i(f_*\mathcal{I}^\bullet)
\end{equation}
と定め，これを $\mathcal{F}$ の {\it 第 $i$ 高次順像}と呼ぶ。
関手族 $R^if_*$ は
$\textit{Mod}(\mathcal{O}_X) \to \textit{Mod}(\mathcal{O}_Y)$ の
普遍 $\delta$-関手をなす。





\section{導来関手}
\label{cohomology-section-derived-functors}

\noindent
分解関手を用いて右導来関手を得る方法を簡単に説明する。
非有界導来関手については，第 \ref{cohomology-section-unbounded} 節を参照せよ。

\medskip\noindent
$(X, \mathcal{O}_X)$ を環付き空間とする。圏
$\textit{Mod}(\mathcal{O}_X)$ はアーベル圏である。《加群》補題
\ref{modules-lemma-abelian} を参照せよ。本章では
$$
K(\mathcal{O}_X) = K(\textit{Mod}(\mathcal{O}_X))
\quad
\text{および}
\quad
D(\mathcal{O}_X) = D(\textit{Mod}(\mathcal{O}_X)).
$$
と書き，《導来圏》定義 \ref{derived-definition-complexes-notation} および
定義 \ref{derived-definition-unbounded-derived-category} で導入された
三角圏の有界版についても同様の記法を用いる。
《導来圏》注意 \ref{derived-remark-big-abelian-category} により，分解関手
$$
j = j_X :
K^{+}(\textit{Mod}(\mathcal{O}_X))
\longrightarrow
K^{+}(\mathcal{I})
$$
が存在する。ここで $\mathcal{I}$ は，入射的な層からなる
$\textit{Mod}(\mathcal{O}_X)$ の厳密充満加法部分圏である。
任意のアーベル圏 $\mathcal{B}$ への左完全関手
$F : \textit{Mod}(\mathcal{O}_X) \to \mathcal{B}$
に対し，《導来圏》第 \ref{derived-section-right-derived-functor} 節で
述べられ，今記した分解関手 $j_X$ を用いて構成される右導来関手を
$RF$ と書く。
\begin{equation}
\label{cohomology-equation-RF}
RF = F \circ j_X' : D^{+}(X) \longrightarrow D^{+}(\mathcal{B})
\end{equation}
記法については《導来圏》補題
\ref{derived-lemma-right-derived-functor} を参照せよ。状況に応じて
$RF$ は
$\textit{Mod}(\mathcal{O}_X)$,
$\text{Comp}^{+}(\textit{Mod}(\mathcal{O}_X))$,
$K^{+}(X)$, or $D^{+}(X)$
上に定義されたものとみなせる。《導来圏》定義
\ref{derived-definition-higher-derived-functors} に従い，第 $i$ 右導来関手
\begin{equation}
\label{cohomology-equation-RFi}
R^iF = H^i \circ RF : \textit{Mod}(\mathcal{O}_X) \longrightarrow \mathcal{B}
\end{equation}
を得る。このとき $R^0F = F$ であり，
$\{R^iF, \delta\}_{i \geq 0}$ は普遍 $\delta$-関手である。
《導来圏》補題 \ref{derived-lemma-higher-derived-functors} を参照せよ。

\medskip\noindent
この構成の二つの特別な場合を挙げる。
環 $R$ に対して $K(R) = K(\text{Mod}_R)$ および
$D(R) = D(\text{Mod}_R)$ と書き，有界版についても同様とする。
任意の開集合 $U \subset X$ に対し，左完全関手
$
\Gamma(U, -) :
\textit{Mod}(\mathcal{O}_X)
\longrightarrow
\text{Mod}_{\mathcal{O}_X(U)}
$
があり，上の議論から
\begin{equation}
\label{cohomology-equation-total-derived-cohomology}
R\Gamma(U, -) :
D^{+}(X)
\longrightarrow
D^{+}(\mathcal{O}_X(U))
\end{equation}
が得られる。$H^i(U, -) = R^i\Gamma(U, -)$ と置く。
$U = X$ ならば (\ref{cohomology-equation-cohomology-modules}) が得られる。
$f : X \to Y$ を環付き空間の射とすると，左完全関手
$
f_* :
\textit{Mod}(\mathcal{O}_X)
\longrightarrow
\textit{Mod}(\mathcal{O}_Y)
$
があり，そこから {\it 導来順像}
\begin{equation}
\label{cohomology-equation-total-derived-direct-image}
Rf_* :
D^{+}(X)
\longrightarrow
D^{+}(Y)
\end{equation}
が得られる。$Rf_*\mathcal{F}^\bullet$ の第 $i$ コホモロジー層を
$R^if_*\mathcal{F}^\bullet$ と書き，
(\ref{cohomology-equation-higher-direct-image-modules}) に従って
第 $i$ {\it 高次順像}と呼ぶ。上に表示した二つの関手はいずれも
導来圏の完全関手である。

\medskip\noindent
{\bf 記法の濫用：} 関手 $Rf_*$ または他の導来関手を $X$ 上の層
$\mathcal{F}$ や層の複体に適用するときは，$\mathcal{F}$ を適切な
分解で置き換えたものと了解する。この種の操作を容易にするため，対象
$\mathcal{F}^\bullet \in D(\mathcal{O}_X)$,
が与えられたとき，
$\textit{Mod}(\mathcal{O}_X)$
の入射対象からなる下に有界な複体 $\mathcal{I}^\bullet$ について，
擬同型 $\mathcal{F}^\bullet \to \mathcal{I}^\bullet$ が存在するならば，
$\mathcal{I}^\bullet$ は {\it 導来圏において
$\mathcal{F}^\bullet$ を表す}という。同様に，「
$\alpha : \mathcal{F}^\bullet \to \mathcal{G}^\bullet$ を
$D(\mathcal{O}_X)$ の射とする」という句は，$\alpha$ が複体の射で
表されていることを意味しない。実際の複体の射がある場合には，その旨を明記する。









\section{第1コホモロジーとトーサー}
\label{cohomology-section-h1-torsors}

\begin{definition}
\label{cohomology-definition-torsor}
$X$ を位相空間とし，$\mathcal{G}$ を $X$ 上の
（非可換でもよい）群の層とする。{\it 擬トーサー}，より正確には
{\it 擬 $\mathcal{G}$-トーサー}とは，作用
$\mathcal{G} \times \mathcal{F} \to \mathcal{F}$ を備えた
$X$ 上の集合の層 $\mathcal{F}$ で，次を満たすものをいう。
\begin{enumerate}
\item $\mathcal{F}(U)$ が空でないときは常に，作用
$\mathcal{G}(U) \times \mathcal{F}(U) \to \mathcal{F}(U)$
は単純推移的である。
\end{enumerate}
{\it 擬 $\mathcal{G}$-トーサーの射}
$\mathcal{F} \to \mathcal{F}'$ とは，$\mathcal{G}$-作用と両立する
集合の層の射をいう。{\it トーサー}，より正確には
{\it $\mathcal{G}$-トーサー}とは，さらに次を満たす
擬 $\mathcal{G}$-トーサーをいう。
\begin{enumerate}
\item[(2)] 任意の $x \in X$ に対して茎 $\mathcal{F}_x$ は空でない。
\end{enumerate}
{\it $\mathcal{G}$-トーサーの射}とは，擬 $\mathcal{G}$-トーサーの射をいう。
{\it 自明な $\mathcal{G}$-トーサー}とは，明らかな左
$\mathcal{G}$-作用を備えた層 $\mathcal{G}$ のことである。
\end{definition}

\noindent
トーサーの射が自動的に同型射であることは明らかである。

\begin{lemma}
\label{cohomology-lemma-trivial-torsor}
$X$ を位相空間とし，$\mathcal{G}$ を $X$ 上の
（非可換でもよい）群の層とする。$\mathcal{G}$-トーサー
$\mathcal{F}$ が自明であるための必要十分条件は
$\mathcal{F}(X) \not = \emptyset$.
\end{lemma}

\begin{proof}
省略する。
\end{proof}

\begin{lemma}
\label{cohomology-lemma-torsors-h1}
$X$ を位相空間とし，$\mathcal{H}$ を $X$ 上のアーベル群の層とする。
$\mathcal{H}$-トーサーの同型類の集合と $H^1(X, \mathcal{H})$ の間には
標準的な全単射がある。
\end{lemma}

\begin{proof}
$\mathcal{F}$ を $\mathcal{H}$-トーサーとする。
$\mathcal{F}$ 上の自由アーベル群の層
$\mathbf{Z}[\mathcal{F}]$ を考える。これは，開集合 $U \subset X$ に
$n_i \in \mathbf{Z}$，$s_i \in \mathcal{F}(U)$ なる有限形式和
$\sum n_i[s_i]$ の集まりを対応させる規則の層化である。自然な写像
$$
\sigma : \mathbf{Z}[\mathcal{F}] \longrightarrow \underline{\mathbf{Z}}
$$
があり，局所切断 $\sum n_i[s_i]$ に $\sum n_i$ を対応させる。
$\sigma$ の核は $[s] - [s']$ の形の局所切断で生成される。標準的な写像
$a : \Ker(\sigma) \to \mathcal{H}$
があり，$[s] - [s']$ を $h \cdot s' = s$ を満たす
$\mathcal{H}$ の局所切断 $h$ に移す。押し出し図式
$$
\xymatrix{
0 \ar[r] &
\Ker(\sigma) \ar[r] \ar[d]^a &
\mathbf{Z}[\mathcal{F}] \ar[r] \ar[d] &
\underline{\mathbf{Z}} \ar[r] \ar[d] &
0 \\
0 \ar[r] &
\mathcal{H} \ar[r] &
\mathcal{E} \ar[r] &
\underline{\mathbf{Z}} \ar[r] &
0
}
$$
を考える。ここで $\mathcal{E}$ は押し出しで得られる拡大である。
下の短完全列に付随するコホモロジー長完全列において
$1 \in H^0(X, \underline{\mathbf{Z}})$ に境界準同型を適用し，元
$\xi = \xi_\mathcal{F} \in H^1(X, \mathcal{H})$
を得る。

\medskip\noindent
逆に，$\xi \in H^1(X, \mathcal{H})$ が与えられたとき，次のように
$\xi$ にトーサーを対応させる。$\mathcal{H}$ から入射的な
アーベル群の層 $\mathcal{I}$ への埋め込み
$\mathcal{H} \to \mathcal{I}$ を選ぶ。
$\mathcal{Q} = \mathcal{I}/\mathcal{H}$ と置くと，短完全列
$$
\xymatrix{
0 \ar[r] &
\mathcal{H} \ar[r] &
\mathcal{I} \ar[r] &
\mathcal{Q} \ar[r] &
0
}
$$
を得る。$H^1(X, \mathcal{I}) = 0$ なので（《導来圏》補題
\ref{derived-lemma-higher-derived-functors} を参照），元 $\xi$ は
大域切断 $q \in H^0(X, \mathcal{Q})$ の像である。
$\mathcal{Q}$ において $q$ に移る切断からなる
$\mathcal{I}$ の部分層（集合の層として）を
$\mathcal{F} \subset \mathcal{I}$ とする。
$\mathcal{F}$ がトーサーであることは容易に確かめられる。

\medskip\noindent
上の二つの構成が互いに逆であることの確認は省略する。
\end{proof}







\section{第1コホモロジーと拡大}
\label{cohomology-section-h1-extensions}

\begin{lemma}
\label{cohomology-lemma-h1-extensions}
$(X, \mathcal{O}_X)$ を環付き空間とし，$\mathcal{F}$ を
$\mathcal{O}_X$-加群の層とする。標準的な全単射
$$
\Ext^1_{\textit{Mod}(\mathcal{O}_X)}(\mathcal{O}_X, \mathcal{F})
\longrightarrow
H^1(X, \mathcal{F})
$$
があり，これは拡大
$$
0 \to \mathcal{F} \to \mathcal{E} \to \mathcal{O}_X \to 0
$$
に，$1 \in \Gamma(X, \mathcal{O}_X)$ の $H^1(X, \mathcal{F})$ における像を
対応させる。
\end{lemma}

\begin{proof}
補題で与えた写像の逆を構成しよう。$\xi \in H^1(X, \mathcal{F})$ とする。
$\textit{Mod}(\mathcal{O}_X)$ において $\mathcal{I}$ が入射的となるような
単射 $\mathcal{F} \subset \mathcal{I}$ を選び，
$\mathcal{Q} = \mathcal{I}/\mathcal{F}$ と置く。
コホモロジー長完全列により，$\xi$ は切断
$\tilde \xi \in \Gamma(X, \mathcal{Q}) =
\Hom_{\mathcal{O}_X}(\mathcal{O}_X, \mathcal{Q})$
の像である。そこで引き戻し
$$
\xymatrix{
0 \ar[r] &
\mathcal{F} \ar[r] \ar@{=}[d] &
\mathcal{E} \ar[r] \ar[d] &
\mathcal{O}_X \ar[r] \ar[d]^{\tilde \xi} &
0 \\
0 \ar[r] &
\mathcal{F} \ar[r] &
\mathcal{I} \ar[r] &
\mathcal{Q} \ar[r] &
0
}
$$
をとればよい。《ホモロジー》第 \ref{homology-section-extensions} 節を参照せよ。
\end{proof}








\section{第1コホモロジーと可逆層}
\label{cohomology-section-invertible-sheaves}

\noindent
環付き空間の Picard 群は《加群》第
\ref{modules-section-invertible} 節で定義した。

\begin{lemma}
\label{cohomology-lemma-h1-invertible}
$(X, \mathcal{O}_X)$ を環付き空間とする。すべての茎
$\mathcal{O}_{X, x}$ が局所環ならば，アーベル群の標準的な同型
$$
H^1(X, \mathcal{O}_X^*) = \Pic(X).
$$
がある。
\end{lemma}

\begin{proof}
$\mathcal{L}$ を可逆 $\mathcal{O}_X$-加群の層とする。規則
$$
U \longmapsto \{s \in \mathcal{L}(U)
\text{ であって } \mathcal{O}_U \xrightarrow{s \cdot -} \mathcal{L}_U
\text{ が同型であるもの}\}
$$
で定める前層 $\mathcal{L}^*$ を考える。この前層は層条件を満たす。
さらに，$f \in \mathcal{O}_X^*(U)$ かつ $s \in \mathcal{L}^*(U)$ ならば，
明らかに $fs \in \mathcal{L}^*(U)$ である。同様に，
$s, s' \in \mathcal{L}^*(U)$ ならば，$fs = s'$ を満たす
$f \in \mathcal{O}_X^*(U)$ が一意に存在する。また，《加群》補題
\ref{modules-lemma-invertible-is-locally-free-rank-1} により，層
$\mathcal{L}^*$ は局所的に切断をもつ。言い換えれば，
$\mathcal{L}^*$ は $\mathcal{O}_X^*$-トーサーである。したがって写像
$$
\begin{matrix}
\text{環付き空間 }(X, \mathcal{O}_X)\text{ 上の可逆層} \\
\text{の同型類}
\end{matrix}
\longrightarrow
\begin{matrix}
\mathcal{O}_X^*\text{-トーサー} \\
\text{の同型類}
\end{matrix}
$$
を得る。これがアーベル群の準同型であることの確認は省略する。
補題 \ref{cohomology-lemma-torsors-h1} により，右辺は
$H^1(X, \mathcal{O}_X^*)$ と標準的に全単射である。
したがって，この写像が単射かつ全射であることを示せばよい。

\medskip\noindent
単射性。トーサー $\mathcal{L}^*$ が自明ならば，補題
\ref{cohomology-lemma-trivial-torsor} により $\mathcal{L}^*$ は大域切断をもつ。
これはちょうど $\mathcal{L} \cong \mathcal{O}_X$，すなわち
$\mathcal{L}$ が $\Pic(X)$ の単位元であることを意味する。

\medskip\noindent
全射性。$\mathcal{F}$ を $\mathcal{O}_X^*$-トーサーとする。
集合の前層
$$
\mathcal{L}_1 : U \longmapsto
(\mathcal{F}(U) \times \mathcal{O}_X(U))/\mathcal{O}_X^*(U)
$$
を考える。ここで $f \in \mathcal{O}_X^*(U)$ の $(s, g)$ への作用は
$(fs, f^{-1}g)$ である。$s'/s$ を，$fs = s'$ を満たす
$\mathcal{O}_X^*$ の局所切断 $f$ とし，
$(s, g) + (s', g') = (s, g + (s'/s)g')$ と定め，また
$\mathcal{O}_X$ の局所切断 $h$ に対して $h(s, g) = (s, hg)$ と定めると，
$\mathcal{L}_1$ は $\mathcal{O}_X$-加群の前層となる。その層化
$\mathcal{L} = \mathcal{L}_1^\#$ が可逆 $\mathcal{O}_X$-加群の層であり，
それに付随する $\mathcal{O}_X^*$-トーサー $\mathcal{L}^*$ が
$\mathcal{F}$ と同型であることの確認は省略する。
\end{proof}













\section{コホモロジーの局所性}
\label{cohomology-section-locality}

\noindent
次の補題は，開集合上で層 $\mathcal{F}$ のコホモロジーを定義する際に
曖昧さがないことを述べる。

\begin{lemma}
\label{cohomology-lemma-cohomology-of-open}
$X$ を環付き空間とし，$U \subset X$ を開部分空間とする。
\begin{enumerate}
\item $\mathcal{I}$ が入射的な $\mathcal{O}_X$-加群の層ならば，
$\mathcal{I}|_U$ は入射的な $\mathcal{O}_U$-加群の層である。
\item 任意の $\mathcal{O}_X$-加群の層 $\mathcal{F}$ に対して
$H^p(U, \mathcal{F}) = H^p(U, \mathcal{F}|_U)$.
\end{enumerate}
\end{lemma}

\begin{proof}
$j : U \to X$ を開埋め込みとする。$U$ への制限関手 $j^{-1}$ は，
零による延長関手 $j_!$ の右随伴であった。《層》補題
\ref{sheaves-lemma-j-shriek-modules} を参照せよ。さらに $j_!$ は完全である。
したがって (1) は《ホモロジー》補題
\ref{homology-lemma-adjoint-preserve-injectives} から従う。

\medskip\noindent
定義により $H^p(U, \mathcal{F}) = H^p(\Gamma(U, \mathcal{I}^\bullet))$
である。ここで $\mathcal{F} \to \mathcal{I}^\bullet$ は
$\textit{Mod}(\mathcal{O}_X)$ における入射分解である。上で見たことから，
$\mathcal{F}|_U \to \mathcal{I}^\bullet|_U$ は
$\textit{Mod}(\mathcal{O}_U)$ における入射分解である。したがって
$H^p(U, \mathcal{F}|_U)$ は
$H^p(\Gamma(U, \mathcal{I}^\bullet|_U))$ に等しい。もちろん，$X$ 上の
任意の層 $\mathcal{F}$ に対して
$\Gamma(U, \mathcal{F}) = \Gamma(U, \mathcal{F}|_U)$ である。
ゆえに (2) の等式を得る。
\end{proof}

\noindent
$X$ を環付き空間，$\mathcal{F}$ を $\mathcal{O}_X$-加群の層，
$U \subset V \subset X$ を開集合とする。このとき標準的な
{\it 制限写像}
\begin{equation}
\label{cohomology-equation-restriction-mapping}
H^n(V, \mathcal{F})
\longrightarrow
H^n(U, \mathcal{F}), \quad
\xi \longmapsto \xi|_U
\end{equation}
があり，$\mathcal{F}$ に関して関手的である。実際，任意の入射分解
$\mathcal{F} \to \mathcal{I}^\bullet$ を選ぶ。層 $\mathcal{I}^p$ の
制限写像から，複体の射
$$
\Gamma(V, \mathcal{I}^\bullet)
\longrightarrow
\Gamma(U, \mathcal{I}^\bullet)
$$
を得る。左辺は $R\Gamma(V, \mathcal{F})$ を表す複体，右辺は
$R\Gamma(U, \mathcal{F})$ を表す複体である。関手 $H^n$ を適用して
コホモロジー群上の写像を得る。上に示したとおり，この写像を
$\xi \mapsto \xi|_U$ と書く。したがって規則
$U \mapsto H^n(U, \mathcal{F})$ は $\mathcal{O}_X$-加群の前層である。
この前層は慣例的に $\underline{H}^n(\mathcal{F})$ と書かれる。
補題 \ref{cohomology-lemma-include} でこの前層の別の解釈を与える。

\begin{lemma}
\label{cohomology-lemma-kill-cohomology-class-on-covering}
$X$ を環付き空間，$\mathcal{F}$ を $\mathcal{O}_X$-加群の層，
$U \subset X$ を開部分空間とする。$n > 0$ かつ
$\xi \in H^n(U, \mathcal{F})$ とする。このとき，任意の $i \in I$ に対して
$\xi|_{U_i} = 0$ となる開被覆 $U = \bigcup_{i\in I} U_i$ が存在する。
\end{lemma}

\begin{proof}
$\mathcal{F} \to \mathcal{I}^\bullet$ を入射分解とする。このとき
$$
H^n(U, \mathcal{F}) =
\frac{\Ker(\mathcal{I}^n(U) \to \mathcal{I}^{n + 1}(U))}
{\Im(\mathcal{I}^{n - 1}(U) \to \mathcal{I}^n(U))}.
$$
上の表示においてコホモロジー類を表す元
$\tilde \xi \in \mathcal{I}^n(U)$ を選ぶ。$\mathcal{I}^\bullet$ は
$\mathcal{F}$ の入射分解で $n > 0$ だから，複体 $\mathcal{I}^\bullet$ は
次数 $n$ で完全である。したがって層として
$\Im(\mathcal{I}^{n - 1} \to \mathcal{I}^n) =
\Ker(\mathcal{I}^n \to \mathcal{I}^{n + 1})$ である。
$\tilde \xi$ は $U$ 上の核の層の切断なので，開被覆
$U = \bigcup_{i \in I} U_i$ で，$\tilde \xi|_{U_i}$ が切断
$\xi_i \in \mathcal{I}^{n - 1}(U_i)$ の $d$ による像となるものが存在する。
制限 $\xi|_{U_i}$ は $\tilde \xi|_{U_i}$ の類に対応するものとして
定義したから，結論を得る。
\end{proof}

\begin{lemma}
\label{cohomology-lemma-describe-higher-direct-images}
$f : X \to Y$ を環付き空間の射とし，$\mathcal{F}$ を
$\mathcal{O}_X$-加群の層とする。層 $R^if_*\mathcal{F}$ は，前層
$$
V \longmapsto H^i(f^{-1}(V), \mathcal{F})
$$
に付随する層である。ここで制限写像は式
(\ref{cohomology-equation-restriction-mapping}) のものである。下に有界な複体
$\mathcal{F}^\bullet$ に $R^if_*$ を適用する場合にも同様の主張が成り立つ。
\end{lemma}

\begin{proof}
$\mathcal{F} \to \mathcal{I}^\bullet$ を入射分解とする。定義により
$R^if_*\mathcal{F}$ は複体
$$
f_*\mathcal{I}^0 \to f_*\mathcal{I}^1 \to f_*\mathcal{I}^2 \to \ldots
$$
の第 $i$ コホモロジー層である。$\mathcal{O}_Y$-加群の層のアーベル圏構造の
定義により，このコホモロジー層は前層
$$
V
\longmapsto
\frac{\Ker(f_*\mathcal{I}^i(V) \to f_*\mathcal{I}^{i + 1}(V))}
{\Im(f_*\mathcal{I}^{i - 1}(V) \to f_*\mathcal{I}^i(V))}
$$
に付随する層である。これは明らかに
$$
\frac{\Ker(\mathcal{I}^i(f^{-1}(V)) \to \mathcal{I}^{i + 1}(f^{-1}(V)))}
{\Im(\mathcal{I}^{i - 1}(f^{-1}(V)) \to \mathcal{I}^i(f^{-1}(V)))}
$$
に等しく，これは $H^i(f^{-1}(V), \mathcal{F})$ に等しい。これでよい。
\end{proof}

\begin{lemma}
\label{cohomology-lemma-localize-higher-direct-images}
$f : X \to Y$ を環付き空間の射，$\mathcal{F}$ を
$\mathcal{O}_X$-加群の層，$V \subset Y$ を開部分空間とする。
$g : f^{-1}(V) \to V$ を $f$ の制限とする。このとき
$$
R^pg_*(\mathcal{F}|_{f^{-1}(V)}) = (R^pf_*\mathcal{F})|_V
$$
が成り立つ。$\mathcal{F}^\bullet$ が $\mathcal{O}_X$-加群の層の
下に有界な複体であるとき，導来順像 $Rf_*\mathcal{F}^\bullet$ についても
同様の主張が成り立つ。
\end{lemma}

\begin{proof}
第1の証明。補題 \ref{cohomology-lemma-describe-higher-direct-images} と
\ref{cohomology-lemma-cohomology-of-open} を適用すれば，表示した等式を得る。
第2の証明。入射分解 $\mathcal{F} \to \mathcal{I}^\bullet$ を選び，
$\mathcal{F}|_{f^{-1}(V)} \to \mathcal{I}^\bullet|_{f^{-1}(V)}$ も
入射分解であることを用いる。
\end{proof}

\begin{remark}
\label{cohomology-remark-daniel}
上の補題 \ref{cohomology-lemma-kill-cohomology-class-on-covering} と
\ref{cohomology-lemma-describe-higher-direct-images} の証明に対する別の方法を述べる。
$(X, \mathcal{O}_X)$ を環付き空間とする。
$i_X : \textit{Mod}(\mathcal{O}_X) \to \textit{PMod}(\mathcal{O}_X)$
を包含関手，$\#$ を層化関手とする。$i_X$ は左完全で，$\#$ は完全であることを
想起せよ。
\begin{enumerate}
\item まず，$i_X$ の右導来関手が
$R^pi_X\mathcal{F} = \underline{H}^p(\mathcal{F})$ で与えられることを述べる
下の補題 \ref{cohomology-lemma-include} を証明する。別の証明は次のとおりである。
$p = 0$ に対して等式は明らかである。
$(R^pi_X)_{p \geq 0}$ と $(\underline{H}^p)_{p \geq 0}$ はともに
入射的対象上で消える $\delta$-関手であり，したがってともに普遍的であるから，
互いに同型である。《ホモロジー》第
\ref{homology-section-cohomological-delta-functor} 節を参照せよ。
\item 補題 \ref{cohomology-lemma-kill-cohomology-class-on-covering} は，任意の
$\mathcal{O}_X$-加群の層 $\mathcal{F}$ と $p > 0$ に対して
$(\underline{H}^p(\mathcal{F}))^\# = 0$ であると言い換えられる。
これを見るには，${}^\#$ が完全であることを用いれば
$$
(\underline{H}^p(\mathcal{F}))^\# =
(R^pi_X\mathcal{F})^\# =
R^p(\# \circ i_X)(\mathcal{F}) = 0
$$
を得る。実際，$\# \circ i_X$ は恒等関手である。
\item $f : X \to Y$ を環付き空間の射，$\mathcal{F}$ を
$\mathcal{O}_X$-加群の層とする。前層
$V \mapsto H^p(f^{-1}V, \mathcal{F})$ は
$R^p (i_Y \circ f_*)\mathcal{F}$ に等しい。これは，上の (1) の議論と同様に，
両者が普遍 $\delta$-関手を与えることから証明できる。したがって補題
\ref{cohomology-lemma-describe-higher-direct-images} は
$R^p f_* \mathcal{F}= (R^p (i_Y \circ f_*)\mathcal{F})^\#$ を述べている。
再び，$\#$ が完全であり，かつ $\# \circ i_Y$ が恒等関手であることを用いると
$$
R^p f_* \mathcal{F} =
R^p(\# \circ i_Y \circ f_*)\mathcal{F} =
(R^p (i_Y \circ f_*)\mathcal{F})^\#
$$
を得る。これが所望の等式である。
\end{enumerate}
\end{remark}






\section{Mayer-Vietoris}
\label{cohomology-section-mayer-vietoris}

\noindent
以下では {\v C}ech-コホモロジー・スペクトル系列を構成する。補題
\ref{cohomology-lemma-cech-spectral-sequence} を参照せよ。そのスペクトル系列の特別な場合が
Mayer-Vietoris 長完全列である。これはスペクトル系列の非常に基本的かつ有用で，
理解しやすい変種なので，ここで別に扱う。

\begin{lemma}
\label{cohomology-lemma-injective-restriction-surjective}
\begin{slogan}
入射的加群の層はフラスクである。
\end{slogan}
$X$ を環付き空間とし，$U' \subset U \subset X$ を開部分空間とする。
任意の入射的 $\mathcal{O}_X$-加群の層 $\mathcal{I}$ に対し，制限写像
$\mathcal{I}(U) \to \mathcal{I}(U')$ は全射である。
\end{lemma}

\begin{proof}
$j : U \to X$ と $j' : U' \to X$ を開埋め込みとする。
$j_!\mathcal{O}_U$ は $\mathcal{O}_U = \mathcal{O}_X|_U$ の
零による延長であった。《層》第 \ref{sheaves-section-open-immersions} 節を
参照せよ。$j_!$ は制限の左随伴なので，任意の $\mathcal{O}_X$-加群の層
$\mathcal{F}$ に対して
$$
\Hom_{\mathcal{O}_X}(j_!\mathcal{O}_U, \mathcal{F})
=
\Hom_{\mathcal{O}_U}(\mathcal{O}_U, \mathcal{F}|_U)
=
\mathcal{F}(U)
$$
である。《層》補題 \ref{sheaves-lemma-j-shriek-modules} を参照せよ。
同様に，層 $j'_!\mathcal{O}_{U'}$ は関手
$\mathcal{F} \mapsto \mathcal{F}(U')$ を表現する。さらに，明らかな標準的
$\mathcal{O}_X$-加群の層の写像
$$
j'_!\mathcal{O}_{U'} \longrightarrow j_!\mathcal{O}_U
$$
があり，Yoneda の補題（《圏》補題 \ref{categories-lemma-yoneda}）を通じて
制限写像 $\mathcal{F}(U) \to \mathcal{F}(U')$ に対応する。
層 $j'_!\mathcal{O}_{U'}$ と $j_!\mathcal{O}_U$ の茎の記述から，
上に表示した写像は単射である（上で引用した補題を参照）。したがって
$\mathcal{I}$ が入射的 $\mathcal{O}_X$-加群の層ならば，写像
$$
\Hom_{\mathcal{O}_X}(j_!\mathcal{O}_U, \mathcal{I})
\longrightarrow
\Hom_{\mathcal{O}_X}(j'_!\mathcal{O}_{U'}, \mathcal{I})
$$
は全射である。《ホモロジー》補題
\ref{homology-lemma-characterize-injectives} を参照せよ。以上を合わせて補題を得る。
\end{proof}

\begin{lemma}[Mayer-Vietoris]
\label{cohomology-lemma-mayer-vietoris}
$X$ を環付き空間とし，$X = U \cup V$ が二つの開集合の和であると仮定する。
任意の $\mathcal{O}_X$-加群の層 $\mathcal{F}$ に対して，コホモロジー長完全列
$$
0 \to
H^0(X, \mathcal{F}) \to
H^0(U, \mathcal{F}) \oplus H^0(V, \mathcal{F}) \to
H^0(U \cap V, \mathcal{F}) \to
H^1(X, \mathcal{F}) \to \ldots
$$
が存在する。この長完全列は $\mathcal{F}$ に関して関手的である。
\end{lemma}

\begin{proof}
層条件によれば，任意のアーベル群の層 $\mathcal{F}$ に対して
$(1, -1) : \mathcal{F}(U) \oplus \mathcal{F}(V) \to \mathcal{F}(U \cap V)$
の核は，第1の写像による $\mathcal{F}(X)$ の像に等しい。上の補題
\ref{cohomology-lemma-injective-restriction-surjective} によれば，
$(1, -1) : \mathcal{I}(U) \oplus \mathcal{I}(V) \to \mathcal{I}(U \cap V)$
は $\mathcal{I}$ が入射的 $\mathcal{O}_X$-加群の層ならば全射である。
したがって $\mathcal{F} \to \mathcal{I}^\bullet$ が $\mathcal{F}$ の
入射分解ならば，複体の短完全列
$$
0 \to
\mathcal{I}^\bullet(X) \to
\mathcal{I}^\bullet(U) \oplus \mathcal{I}^\bullet(V) \to
\mathcal{I}^\bullet(U \cap V) \to
0.
$$
を得る。コホモロジーをとれば結果を得る（《ホモロジー》補題
\ref{homology-lemma-long-exact-sequence-cochain} を用いよ）。
この列の関手性の証明は省略する。
\end{proof}

\begin{lemma}[相対 Mayer-Vietoris]
\label{cohomology-lemma-relative-mayer-vietoris}
$f : X \to Y$ を環付き空間の射とする。
$X = U \cup V$ が二つの開集合の和であると仮定する。
$a = f|_U : U \to Y$，$b = f|_V : V \to Y$，および
$c = f|_{U \cap V} : U \cap V \to Y$ と書く。
任意の $\mathcal{O}_X$-加群の層 $\mathcal{F}$ に対して長完全列
$$
0 \to
f_*\mathcal{F} \to
a_*(\mathcal{F}|_U) \oplus b_*(\mathcal{F}|_V) \to
c_*(\mathcal{F}|_{U \cap V}) \to
R^1f_*\mathcal{F} \to \ldots
$$
が存在する。この長完全列は $\mathcal{F}$ に関して関手的である。
\end{lemma}

\begin{proof}
$\mathcal{F} \to \mathcal{I}^\bullet$ を $\mathcal{F}$ の入射分解とする。
複体の短完全列
$$
0 \to
f_*\mathcal{I}^\bullet \to
a_*\mathcal{I}^\bullet|_U \oplus b_*\mathcal{I}^\bullet|_V \to
c_*\mathcal{I}^\bullet|_{U \cap V} \to
0.
$$
を得ると主張する。実際，任意の開集合 $W \subset Y$ と任意の $n \geq 0$ に
対して，対応する $W$ 上の切断群の列
$$
0 \to
\mathcal{I}^n(f^{-1}(W)) \to
\mathcal{I}^n(U \cap f^{-1}(W))
\oplus \mathcal{I}^n(V \cap f^{-1}(W)) \to
\mathcal{I}^n(U \cap V \cap f^{-1}(W)) \to
0
$$
が短完全であることは，補題 \ref{cohomology-lemma-mayer-vietoris} の証明で示した。
コホモロジー層をとり，$\mathcal{I}^\bullet|_U$ が
$\mathcal{F}|_U$ の入射分解であること，および
$\mathcal{I}^\bullet|_V$，$\mathcal{I}^\bullet|_{U \cap V}$ についても同様で
あることを用いれば補題が従う。補題 \ref{cohomology-lemma-cohomology-of-open} を参照せよ。
\end{proof}




















\section{{\v C}ech 複体と {\v C}ech コホモロジー}
\label{cohomology-section-cech}

\noindent
$X$ を位相空間とし，
$\mathcal{U} : U = \bigcup_{i \in I} U_i$ を開被覆とする。
《位相空間》基本概念 (\ref{topology-item-covering}) を参照せよ。
慣例に従い，$\mathcal{U}$ の元の $(p + 1)$ 重共通部分を
$U_{i_0\ldots i_p} = U_{i_0} \cap \ldots \cap U_{i_p}$ と書く。
$\mathcal{F}$ を $X$ 上のアーベル群の前層とし，
$$
\check{\mathcal{C}}^p(\mathcal{U}, \mathcal{F})
=
\prod\nolimits_{(i_0, \ldots, i_p) \in I^{p + 1}}
\mathcal{F}(U_{i_0\ldots i_p}).
$$
と置く。これはアーベル群である。
$s \in \check{\mathcal{C}}^p(\mathcal{U}, \mathcal{F})$ に対し，
$\mathcal{F}(U_{i_0\ldots i_p})$ におけるその値を
$s_{i_0\ldots i_p}$ と書く。
$s \in \check{\mathcal{C}}^1(\mathcal{U}, \mathcal{F})$ かつ
$i, j \in I$ ならば，$s_{ij}$ と $s_{ji}$ はともに
$\mathcal{F}(U_i \cap U_j)$ の元だが，両者の間には何の関係も課していないことに
注意せよ。言い換えれば，ここでは交代余鎖を扱って{\it いない}（交代余鎖は第
\ref{cohomology-section-alternating-cech} 節で定義する）。写像
$$
d : \check{\mathcal{C}}^p(\mathcal{U}, \mathcal{F})
\longrightarrow
\check{\mathcal{C}}^{p + 1}(\mathcal{U}, \mathcal{F})
$$
を公式
\begin{equation}
\label{cohomology-equation-d-cech}
d(s)_{i_0\ldots i_{p + 1}}
=
\sum\nolimits_{j = 0}^{p + 1}
(-1)^j
s_{i_0\ldots \hat i_j \ldots i_{p + 1}}|_{U_{i_0\ldots i_{p + 1}}}
\end{equation}
で定める。$d \circ d = 0$ は直ちに確かめられる。言い換えれば，
$\check{\mathcal{C}}^\bullet(\mathcal{U}, \mathcal{F})$ は複体である。

\begin{definition}
\label{cohomology-definition-cech-complex}
$X$ を位相空間，$\mathcal{U} : U = \bigcup_{i \in I} U_i$ を開被覆，
$\mathcal{F}$ を $X$ 上のアーベル群の前層とする。複体
$\check{\mathcal{C}}^\bullet(\mathcal{U}, \mathcal{F})$ を，
$\mathcal{F}$ と開被覆 $\mathcal{U}$ に付随する
{\it {\v C}ech 複体}という。そのコホモロジー群
$H^i(\check{\mathcal{C}}^\bullet(\mathcal{U}, \mathcal{F}))$ を，
$\mathcal{F}$ と被覆 $\mathcal{U}$ に付随する
{\it {\v C}ech コホモロジー群}といい，
$\check H^i(\mathcal{U}, \mathcal{F})$ と書く。
\end{definition}

\begin{lemma}
\label{cohomology-lemma-cech-h0}
$X$ を位相空間とし，$\mathcal{F}$ を $X$ 上のアーベル群の前層とする。
次は同値である。
\begin{enumerate}
\item $\mathcal{F}$ はアーベル群の層である。
\item 任意の開被覆 $\mathcal{U} : U = \bigcup_{i \in I} U_i$ に対して
自然な写像
$$
\mathcal{F}(U) \to \check{H}^0(\mathcal{U}, \mathcal{F})
$$
は全単射である。
\end{enumerate}
\end{lemma}

\begin{proof}
層条件とは，任意の開被覆に対して
$\mathcal{F}(U) \to \check{H}^0(\mathcal{U}, \mathcal{F})$ が
全単射であることにほかならないからである。
\end{proof}

\begin{lemma}
\label{cohomology-lemma-cech-trivial}
$X$ を位相空間，$\mathcal{F}$ を $X$ 上のアーベル群の前層，
$\mathcal{U} : U = \bigcup_{i \in I} U_i$ を開被覆とする。
ある $i \in I$ に対して $U_i = U$ ならば，拡張 {\v C}ech 複体
$$
\mathcal{F}(U) \to \check{\mathcal{C}}^\bullet(\mathcal{U}, \mathcal{F})
$$
であって，$\mathcal{F}(U)$ を次数 $-1$ に置き，その微分を
$\mathcal{F}(U)$ から $\check{\mathcal{C}}^0(\mathcal{U}, \mathcal{F})$ への
標準的な写像としたものは，$0$ とホモトピー同値である。
\end{lemma}

\begin{proof}
$U = U_i$ を満たす元 $i \in I$ を固定する。$i_j = i$ ならば
$U_{i_0 \ldots i_p} = U_{i_0 \ldots \hat i_j \ldots i_p}$ であることに注意せよ。
ホモトピー
$$
h :
\prod\nolimits_{i_0 \ldots i_{p + 1}} \mathcal{F}(U_{i_0 \ldots i_{p + 1}})
\longrightarrow
\prod\nolimits_{i_0 \ldots i_p} \mathcal{F}(U_{i_0 \ldots i_p})
$$
を規則
$$
h(s)_{i_0 \ldots i_p} = s_{i i_0 \ldots i_p}
$$
で定める。言い換えれば，
$h : \prod_{i_0} \mathcal{F}(U_{i_0}) \to \mathcal{F}(U)$ は因子
$\mathcal{F}(U_i) = \mathcal{F}(U)$ への射影であり，一般に写像 $h$ は因子
$\mathcal{F}(U_{i i_1 \ldots i_{p + 1}}) =
\mathcal{F}(U_{i_1 \ldots i_{p + 1}})$ への射影である。計算すると
\begin{align*}
(dh + hd)(s)_{i_0 \ldots i_p}
& =
\sum\nolimits_{j = 0}^p
(-1)^j
h(s)_{i_0 \ldots \hat i_j \ldots i_p}
+
d(s)_{i i_0 \ldots i_p}\\
& =
\sum\nolimits_{j = 0}^p
(-1)^j
s_{i i_0 \ldots \hat i_j \ldots i_p}
+
s_{i_0 \ldots i_p}
+
\sum\nolimits_{j = 0}^p
(-1)^{j + 1}
s_{i i_0 \ldots \hat i_j \ldots i_p} \\
& =
s_{i_0 \ldots i_p}
\end{align*}
となる。これにより，恒等写像が零写像とホモトピックであることが示された。
\end{proof}






\section{前層上の関手としての {\v C}ech コホモロジー}
\label{cohomology-section-cech-functor}

\noindent
警告：この節ではほとんど専ら前層と前層の圏を扱う。以下の結果は，
層および層の圏という設定ではまったく誤りである！

\medskip\noindent
$X$ を環付き空間，$\mathcal{U} : U = \bigcup_{i \in I} U_i$ を開被覆，
$\mathcal{F}$ を $\mathcal{O}_X$-加群の前層とする。$\mathcal{F}$ を
アーベル群の前層とみなすだけで，その {\v C}ech 複体
$\check{\mathcal{C}}^\bullet(\mathcal{U}, \mathcal{F})$ を得る。
しかし，各項 $\check{\mathcal{C}}^p(\mathcal{U}, \mathcal{F})$ は自然な
$\mathcal{O}_X(U)$-加群構造をもち，微分は $\mathcal{O}_X(U)$-加群の写像で
与えられる。さらに，構成
$$
\mathcal{F} \longmapsto \check{\mathcal{C}}^\bullet(\mathcal{U}, \mathcal{F})
$$
は明らかに $\mathcal{F}$ に関して関手的である。実際，これは関手
\begin{equation}
\label{cohomology-equation-cech-functor}
\check{\mathcal{C}}^\bullet(\mathcal{U}, -) :
\textit{PMod}(\mathcal{O}_X)
\longrightarrow
\text{Comp}^{+}(\text{Mod}_{\mathcal{O}_X(U)})
\end{equation}
である。記法については《導来圏》定義
\ref{derived-definition-complexes-notation} を参照せよ。アーベル圏における
下に有界な複体の圏はアーベル圏であることを想起せよ。《ホモロジー》補題
\ref{homology-lemma-cat-cochain-abelian} を参照。

\begin{lemma}
\label{cohomology-lemma-cech-exact-presheaves}
式 (\ref{cohomology-equation-cech-functor}) で与えた関手は完全関手である
（《ホモロジー》補題 \ref{homology-lemma-exact-functor} を参照）。
\end{lemma}

\begin{proof}
任意の開集合 $W \subset U$ に対して，関手
$\mathcal{F} \mapsto \mathcal{F}(W)$ は
$\textit{PMod}(\mathcal{O}_X)$ から $\text{Mod}_{\mathcal{O}_X(U)}$ への
加法的完全関手である。複体の各項
$\check{\mathcal{C}}^p(\mathcal{U}, \mathcal{F})$ はこれらの完全関手の積なので
完全である。さらに，複体の列が完全であるための必要十分条件は，各次数の項の
列が完全であることである。したがって補題が従う。
\end{proof}

\begin{lemma}
\label{cohomology-lemma-cech-cohomology-delta-functor-presheaves}
$X$ を環付き空間，$\mathcal{U} : U = \bigcup_{i \in I} U_i$ を開被覆とする。
関手 $\mathcal{F} \mapsto \check{H}^n(\mathcal{U}, \mathcal{F})$ は，
$\mathcal{O}_X$-加群の前層のアーベル圏から $\mathcal{O}_X(U)$-加群の圏への
$\delta$-関手をなす（《ホモロジー》定義
\ref{homology-definition-cohomological-delta-functor} を参照）。
\end{lemma}

\begin{proof}
補題 \ref{cohomology-lemma-cech-exact-presheaves} により，$\mathcal{O}_X$-加群の
前層の短完全列
$0 \to \mathcal{F}_1 \to \mathcal{F}_2 \to \mathcal{F}_3 \to 0$
は $\mathcal{O}_X(U)$-加群の複体の短完全列に移る。したがって
《ホモロジー》補題 \ref{homology-lemma-long-exact-sequence-cochain} を用いて
境界写像
$\delta_{\mathcal{F}_1 \to \mathcal{F}_2 \to \mathcal{F}_3} :
\check{H}^n(\mathcal{U}, \mathcal{F}_3) \to
\check{H}^{n + 1}(\mathcal{U}, \mathcal{F}_1)$
と，対応する長完全列を得る。これらの写像が，前層の短完全列の間の写像と
両立することの確認は省略する。
\end{proof}


\noindent
次の補題の定式化では，開埋め込み $j : U \to X$ に関する加群の前層の
零による延長関手 $j_{p!}$ を用いる。《層》第
\ref{sheaves-section-open-immersions} 節を参照せよ。任意の開集合
$W \subset X$ と任意の $\mathcal{O}_X|_U$-加群の前層 $\mathcal{G}$ に対して
$$
(j_{p!}\mathcal{G})(W) =
\left\{
\begin{matrix}
\mathcal{G}(W) & \text{if } W \subset U \\
0 & \text{else.}
\end{matrix}
\right.
$$
である。さらに，関手 $j_{p!}$ は制限関手の左随伴である。《層》補題
\ref{sheaves-lemma-j-shriek-modules} を参照せよ。特に公式
$$
\Hom_{\mathcal{O}_X}(j_{p!}\mathcal{O}_U, \mathcal{F})
=
\Hom_{\mathcal{O}_U}(\mathcal{O}_U, \mathcal{F}|_U)
=
\mathcal{F}(U).
$$
を得る。関手 $\mathcal{F} \mapsto \mathcal{F}(U)$ は前層の圏上の完全関手
なので，前層 $j_{p!}\mathcal{O}_U$ は圏
$\textit{PMod}(\mathcal{O}_X)$ の射影対象である。《ホモロジー》補題
\ref{homology-lemma-characterize-projectives} を参照せよ。

\medskip\noindent
対応する開埋め込みを $j_U,j_V$ とする開集合 $U \subset V \subset X$ が
与えられたとき，標準的な写像
$(j_U)_{p!}\mathcal{O}_U \to (j_V)_{p!}\mathcal{O}_V$ があることに注意せよ。
これは任意の開集合 $W \subset U$ 上の切断では恒等写像で，それ以外では
$0$ である。同一視
$\Hom_{\mathcal{O}_X}((j_U)_{p!}\mathcal{O}_U, (j_V)_{p!}\mathcal{O}_V) =
(j_V)_{p!}\mathcal{O}_V(U) = \mathcal{O}_V(U)$
のもとで，これは元 $1 \in \mathcal{O}_V(U)$ に対応する。

\begin{lemma}
\label{cohomology-lemma-cech-map-into}
$X$ を環付き空間，$\mathcal{U} : U = \bigcup_{i \in I} U_i$ を開被覆とする。
$j_{i_0\ldots i_p} : U_{i_0 \ldots i_p} \to X$ を開埋め込みとする。
$\mathcal{O}_X$-加群の前層からなる鎖複体 $K(\mathcal{U})_\bullet$
$$
\ldots
\to
\bigoplus_{i_0i_1i_2} (j_{i_0i_1i_2})_{p!}\mathcal{O}_{U_{i_0i_1i_2}}
\to
\bigoplus_{i_0i_1} (j_{i_0i_1})_{p!}\mathcal{O}_{U_{i_0i_1}}
\to
\bigoplus_{i_0} (j_{i_0})_{p!}\mathcal{O}_{U_{i_0}}
\to 0 \to \ldots
$$
を考える。ここで最後の非零項は次数 $0$ に置き，写像
$$
(j_{i_0\ldots i_{p + 1}})_{p!}\mathcal{O}_{U_{i_0\ldots i_{p + 1}}}
\longrightarrow
(j_{i_0\ldots \hat i_j \ldots i_{p + 1}})_{p!}
\mathcal{O}_{U_{i_0\ldots \hat i_j \ldots i_{p + 1}}}
$$
は標準的な写像の $(-1)^j$ 倍で与える。このとき同型
$$
\Hom_{\mathcal{O}_X}(K(\mathcal{U})_\bullet, \mathcal{F})
=
\check{\mathcal{C}}^\bullet(\mathcal{U}, \mathcal{F})
$$
があり，$\mathcal{F} \in \Ob(\textit{PMod}(\mathcal{O}_X))$ に関して関手的で
ある。
\end{lemma}

\begin{proof}
補題の直前の議論で
$$
\Hom_{\mathcal{O}_X}(
(j_{i_0\ldots i_p})_{p!}\mathcal{O}_{U_{i_0\ldots i_p}},
\mathcal{F})
=
\mathcal{F}(U_{i_0\ldots i_p}).
$$
を見た。したがって直和
$$
\bigoplus\nolimits_{i_0 \ldots i_p}
(j_{i_0 \ldots i_p})_{p!}\mathcal{O}_{U_{i_0 \ldots i_p}}
$$
は確かに関手
$$
\mathcal{F}
\longmapsto
\prod\nolimits_{i_0\ldots i_p} \mathcal{F}(U_{i_0\ldots i_p}).
$$
を表現する。よって《圏》Yoneda の補題 \ref{categories-lemma-yoneda} により，
上で与えた項をもつ複体 $K(\mathcal{U})_\bullet$ が存在する。
写像が補題で与えたものになることは容易に確かめられる。
\end{proof}

\begin{lemma}
\label{cohomology-lemma-homology-complex}
$X$ を環付き空間，$\mathcal{U} : U = \bigcup_{i \in I} U_i$ を開被覆とする。
$\mathcal{O}_\mathcal{U} \subset \mathcal{O}_X$ を写像
$\bigoplus j_{p!}\mathcal{O}_{U_i} \to \mathcal{O}_X$ の像前層とする。
上の補題 \ref{cohomology-lemma-cech-map-into} の前層の鎖複体
$K(\mathcal{U})_\bullet$ のホモロジー前層は
$$
H_i(K(\mathcal{U})_\bullet) =
\left\{
\begin{matrix}
0 & \text{if} & i \not = 0 \\
\mathcal{O}_\mathcal{U} & \text{if} & i = 0
\end{matrix}
\right.
$$
である。
\end{lemma}

\begin{proof}
次数 $-1$ に $\mathcal{O}_\mathcal{U}$ を置き，明らかな写像
$K(\mathcal{U})_0 =
\bigoplus_{i_0} (j_{i_0})_{p!}\mathcal{O}_{U_{i_0}} \to
\mathcal{O}_\mathcal{U}$ を備えた拡張複体 $K^{ext}_\bullet$ を考える。
任意の開集合 $W \subset X$ 上でこの拡張複体の切断をとると完全複体になることを
示せば十分である。実際，任意の $W \subset X$ に対して複体
$K^{ext}_\bullet(W)$ は零複体とホモトピー同値であると主張する。
$W \subset U_i$ であることと $i \in I_1$ であることが同値となるように
$I = I_1 \amalg I_2$ と書く。

\medskip\noindent
$I_1 = \emptyset$ ならば複体 $K^{ext}_\bullet(W) = 0$ なので，示すことはない。

\medskip\noindent
$I_1 \not = \emptyset$ ならば
$\mathcal{O}_\mathcal{U}(W) = \mathcal{O}_X(W)$
であり，また
$$
K^{ext}_p(W) =
\bigoplus\nolimits_{i_0 \ldots i_p \in I_1} \mathcal{O}_X(W).
$$
これは前層 $(j_{i_0 \ldots i_p})_{p!}\mathcal{O}_{U_{i_0 \ldots i_p}}$
の簡単な記述から従う。さらに，複体 $K^{ext}_\bullet(W)$ の微分は
$$
d(s)_{i_0 \ldots i_p} =
\sum\nolimits_{j = 0, \ldots, p + 1} \sum\nolimits_{i \in I_1}
(-1)^j s_{i_0 \ldots i_{j - 1} i i_j \ldots i_p}.
$$
で与えられる。元 $s$ は有限台をもつので，この和は有限である。
元 $i_{\text{fix}} \in I_1$ を固定し，写像
$$
h : K^{ext}_p(W) \longrightarrow K^{ext}_{p + 1}(W)
$$
を規則
$$
h(s)_{i_0 \ldots i_{p + 1}} =
\left\{
\begin{matrix}
0 & \text{if} & i_0 \not = i_{\text{fix}} \\
s_{i_1 \ldots i_{p + 1}} & \text{if} & i_0 = i_{\text{fix}}
\end{matrix}
\right.
$$
で定める。これを略して
$h(s)_{i_0 \ldots i_{p + 1}} = (i_0 = i_{\text{fix}}) s_{i_1 \ldots i_p}$
と書くことにする。すると
\begin{eqnarray*}
& & (dh + hd)(s)_{i_0 \ldots i_p} \\
& = &
\sum_j \sum_{i \in I_1} (-1)^j h(s)_{i_0 \ldots i_{j - 1} i i_j \ldots i_p}
+
(i = i_0) d(s)_{i_1 \ldots i_p} \\
& = &
s_{i_0 \ldots i_p} +
\sum_{j \geq 1}\sum_{i \in I_1}
(-1)^j (i_0 = i_{\text{fix}}) s_{i_1 \ldots i_{j - 1} i i_j \ldots i_p}
+
(i_0 = i_{\text{fix}}) d(s)_{i_1 \ldots i_p}
\end{eqnarray*}
を得る。これは所望のとおり $s_{i_0 \ldots i_p}$ に等しい。
\end{proof}

\begin{lemma}
\label{cohomology-lemma-cech-cohomology-derived-presheaves}
$X$ を環付き空間とし，$\mathcal{U} : U = \bigcup_{i \in I} U_i$ を
$U \subset X$ の開被覆とする。{\v C}ech コホモロジー関手
$\check{H}^p(\mathcal{U}, -)$ は，$\delta$-関手として，関手
$$
\check{H}^0(\mathcal{U}, -) :
\textit{PMod}(\mathcal{O}_X)
\longrightarrow
\text{Mod}_{\mathcal{O}_X(U)}.
$$
の右導来関手と標準的に同型である。さらに，関手的な擬同型
$$
\check{\mathcal{C}}^\bullet(\mathcal{U}, \mathcal{F})
\longrightarrow
R\check{H}^0(\mathcal{U}, \mathcal{F})
$$
がある。ここで右辺は，左完全関手
$\check{H}^0(\mathcal{U}, -)$ の右導来関手
$$
R\check{H}^0(\mathcal{U}, -) :
D^{+}(\textit{PMod}(\mathcal{O}_X))
\longrightarrow
D^{+}(\mathcal{O}_X(U))
$$
を表す。
\end{lemma}

\begin{proof}
$\mathcal{O}_X$-加群の前層の圏は十分多くの入射対象をもつことに注意せよ。
《入射的対象》命題 \ref{injectives-proposition-presheaves-modules} を参照。
また，$\check{H}^0(\mathcal{U}, -)$ は $\mathcal{O}_X$-加群の前層の圏から
$\mathcal{O}_X(U)$-加群の圏への左完全関手である。したがって導来関手と
右導来関手が存在する。《導来圏》第
\ref{derived-section-right-derived-functor} 節を参照せよ。

\medskip\noindent
$\mathcal{I}$ を入射的 $\mathcal{O}_X$-加群の前層とする。この場合，関手
$\Hom_{\mathcal{O}_X}(-, \mathcal{I})$ は
$\textit{PMod}(\mathcal{O}_X)$ 上で完全である。補題
\ref{cohomology-lemma-cech-map-into} により
$$
\Hom_{\mathcal{O}_X}(K(\mathcal{U})_\bullet, \mathcal{I})
=
\check{\mathcal{C}}^\bullet(\mathcal{U}, \mathcal{I}).
$$
である。補題 \ref{cohomology-lemma-homology-complex} により，
$K(\mathcal{U})_\bullet$ は $\mathcal{O}_\mathcal{U}[0]$ と擬同型である。
したがって，上で述べた $\mathcal{I}$ への Hom の完全性から，すべての
$i > 0$ に対して $\check{H}^i(\mathcal{U}, \mathcal{I}) = 0$ である。
ゆえに $\delta$-関手 $(\check{H}^n, \delta)$（補題
\ref{cohomology-lemma-cech-cohomology-delta-functor-presheaves} を参照）は
《ホモロジー》補題 \ref{homology-lemma-efface-implies-universal} の仮定を満たし，
したがって普遍 $\delta$-関手である。

\medskip\noindent
《導来圏》補題 \ref{derived-lemma-higher-derived-functors} により，列
$R^i\check{H}^0(\mathcal{U}, -)$ も普遍 $\delta$-関手をなす。
普遍 $\delta$-関手の一意性（《ホモロジー》補題
\ref{homology-lemma-uniqueness-universal-delta-functor} を参照）から
$R^i\check{H}^0(\mathcal{U}, -) = \check{H}^i(\mathcal{U}, -)$ を得る。
大半の応用にはこれで十分であり，読者には証明の残りを読み飛ばすことを勧める。

\medskip\noindent
$\mathcal{F}$ を任意の $\mathcal{O}_X$-加群の前層とし，圏
$\textit{PMod}(\mathcal{O}_X)$ における入射分解
$\mathcal{F} \to \mathcal{I}^\bullet$ を選ぶ。項
$\check{\mathcal{C}}^p(\mathcal{U}, \mathcal{I}^q)$ をもつ二重複体
$\check{\mathcal{C}}^\bullet(\mathcal{U}, \mathcal{I}^\bullet)$ と，
それに付随する全複体
$\text{Tot}(\check{\mathcal{C}}^\bullet(\mathcal{U}, \mathcal{I}^\bullet))$,
を考える。《ホモロジー》定義
\ref{homology-definition-associated-simple-complex} を参照せよ。写像
$\check{\mathcal{C}}^p(\mathcal{U}, \mathcal{F})
\to \check{\mathcal{C}}^p(\mathcal{U}, \mathcal{I}^0)$ から生じる複体の写像
$$
\check{\mathcal{C}}^\bullet(\mathcal{U}, \mathcal{F})
\longrightarrow
\text{Tot}(\check{\mathcal{C}}^\bullet(\mathcal{U}, \mathcal{I}^\bullet))
$$
があり，また写像
$\check{H}^0(\mathcal{U}, \mathcal{I}^q) \to
\check{\mathcal{C}}^0(\mathcal{U}, \mathcal{I}^q)$ から生じる複体の写像
$$
\check{H}^0(\mathcal{U}, \mathcal{I}^\bullet)
\longrightarrow
\text{Tot}(\check{\mathcal{C}}^\bullet(\mathcal{U}, \mathcal{I}^\bullet))
$$
がある。《ホモロジー》補題
\ref{homology-lemma-double-complex-gives-resolution} を適用すると，これら二つの
写像はいずれも擬同型である。実際，関手としての {\v C}ech 複体が完全なので
（補題 \ref{cohomology-lemma-cech-exact-presheaves}），二重複体の列は正の次数で完全である。
また，今見たように入射的前層 $\mathcal{I}^q$ の高次
{\v C}ech コホモロジー群は零なので，二重複体の行も正の次数で完全である。
$D^{+}(\mathcal{O}_X(U))$ において擬同型は可逆になるから，これにより補題で
最後に表示した射を得る。この射が関手的であることの確認は省略する。
\end{proof}





\section{{\v C}ech コホモロジーとコホモロジー}
\label{cohomology-section-cech-cohomology-cohomology}

\begin{lemma}
\label{cohomology-lemma-injective-trivial-cech}
$X$ を環付き空間，$\mathcal{U} : U = \bigcup_{i \in I} U_i$ を開被覆，
$\mathcal{I}$ を入射的 $\mathcal{O}_X$-加群の層とする。このとき
$$
\check{H}^p(\mathcal{U}, \mathcal{I}) =
\left\{
\begin{matrix}
\mathcal{I}(U) & \text{if} & p = 0 \\
0 & \text{if} & p > 0
\end{matrix}
\right.
$$
である。
\end{lemma}

\begin{proof}
入射的 $\mathcal{O}_X$-加群の層は，圏
$\textit{PMod}(\mathcal{O}_X)$ の対象としても入射的である（例えば，層化は
包含関手の完全な左随伴だからである。《ホモロジー》補題
\ref{homology-lemma-adjoint-preserve-injectives} を用いよ）。したがって補題
\ref{cohomology-lemma-cech-cohomology-derived-presheaves}（またはその証明）を適用すれば
結果を得る。
\end{proof}

\begin{lemma}
\label{cohomology-lemma-cech-cohomology}
$X$ を環付き空間，$\mathcal{U} : U = \bigcup_{i \in I} U_i$ を開被覆とする。
関手の自然変換
$$
\check{\mathcal{C}}^\bullet(\mathcal{U}, -)
\longrightarrow
R\Gamma(U, -)
$$
がある。これらは $\textit{Mod}(\mathcal{O}_X)$ から
$D^{+}(\mathcal{O}_X(U))$ への関手である。特に，これにより
$\mathcal{F}$ が $\textit{Mod}(\mathcal{O}_X)$ を動くときの標準的な写像
$\check{H}^p(\mathcal{U}, \mathcal{F}) \to H^p(U, \mathcal{F})$ が得られる。
\end{lemma}

\begin{proof}
$\mathcal{F}$ を $\mathcal{O}_X$-加群の層とし，入射分解
$\mathcal{F} \to \mathcal{I}^\bullet$ を選ぶ。項
$\check{\mathcal{C}}^p(\mathcal{U}, \mathcal{I}^q)$ をもつ二重複体
$\check{\mathcal{C}}^\bullet(\mathcal{U}, \mathcal{I}^\bullet)$ を考える。
写像 $\mathcal{I}^q(U) \to \check{H}^0(\mathcal{U}, \mathcal{I}^q)$ から
生じる複体の写像
$$
\alpha :
\Gamma(U, \mathcal{I}^\bullet)
\longrightarrow
\text{Tot}(\check{\mathcal{C}}^\bullet(\mathcal{U}, \mathcal{I}^\bullet))
$$
と，写像 $\mathcal{F} \to \mathcal{I}^0$ から生じる複体の写像
$$
\beta :
\check{\mathcal{C}}^\bullet(\mathcal{U}, \mathcal{F})
\longrightarrow
\text{Tot}(\check{\mathcal{C}}^\bullet(\mathcal{U}, \mathcal{I}^\bullet))
$$
がある。《ホモロジー》補題
\ref{homology-lemma-double-complex-gives-resolution} を適用すると，
$\alpha$ は擬同型である。実際，補題 \ref{cohomology-lemma-injective-trivial-cech} により，
二重複体 $\check{\mathcal{C}}^\bullet(\mathcal{U}, \mathcal{I}^\bullet)$ の
第 $q$ 行は $\Gamma(U, \mathcal{I}^q)$ の分解である。したがって
$\alpha$ は $D^{+}(\mathcal{O}_X(U))$ において可逆となり，補題の自然変換は
$\beta$ に続けて $\alpha$ の逆を合成したものである。これが関手的であることの
確認は省略する。
\end{proof}

\begin{lemma}
\label{cohomology-lemma-cech-h1}
$X$ を位相空間，$\mathcal{H}$ を $X$ 上のアーベル群の層，
$\mathcal{U} : X = \bigcup_{i \in I} U_i$ を開被覆とする。写像
$$
\check{H}^1(\mathcal{U}, \mathcal{H}) \longrightarrow H^1(X, \mathcal{H})
$$
は単射であり，補題 \ref{cohomology-lemma-torsors-h1} の全単射を通じて
$\check{H}^1(\mathcal{U}, \mathcal{H})$ を，各 $U_i$ 上への制限が自明となる
$\mathcal{H}$-トーサーの同型類の集合と同一視する。
\end{lemma}

\begin{proof}
これを見るため逆写像を構成する。すなわち，$\mathcal{F}$ を，各 $U_i$ への
制限が自明である $\mathcal{H}$-トーサーとする。補題
\ref{cohomology-lemma-trivial-torsor} により，切断
$s_i \in \mathcal{F}(U_i)$ が存在する。$U_{i_0} \cap U_{i_1}$ 上には
$s_{i_0i_1} \cdot s_{i_0}|_{U_{i_0} \cap U_{i_1}} =
s_{i_1}|_{U_{i_0} \cap U_{i_1}}$ を満たす $\mathcal{H}$ の切断
$s_{i_0i_1}$ が一意に存在する。計算により，$s_{i_0i_1}$ は
{\v C}ech コサイクルであり，その類は矛盾なく定まる（すなわち，切断
$s_i$ の選び方に依存しない）。逆写像は $\mathcal{F}$ の同型類を
コサイクル $(s_{i_0i_1})$ のコホモロジー類へ移す。
この写像が実際に逆であることの確認は省略する。
\end{proof}

\begin{lemma}
\label{cohomology-lemma-include}
$X$ を環付き空間とし，関手
$i : \textit{Mod}(\mathcal{O}_X) \to \textit{PMod}(\mathcal{O}_X)$.
を考える。これは左完全関手であり，その右導来関手は
$$
R^pi(\mathcal{F}) = \underline{H}^p(\mathcal{F}) :
U \longmapsto H^p(U, \mathcal{F})
$$
で与えられる。第 \ref{cohomology-section-locality} 節の議論を参照せよ。
\end{lemma}

\begin{proof}
$i$ が左完全であることは明らかである。入射分解
$\mathcal{F} \to \mathcal{I}^\bullet$ を選ぶ。定義により，$R^pi$ は複体
$\mathcal{I}^\bullet$ の第 $p$ コホモロジー{\it 前層}である。
言い換えれば，開集合 $U$ 上の $R^pi(\mathcal{F})$ の切断は
$$
\frac{\Ker(\mathcal{I}^p(U) \to \mathcal{I}^{p + 1}(U))}
{\Im(\mathcal{I}^{p - 1}(U) \to \mathcal{I}^p(U))}.
$$
で与えられる。これは $H^p(U, \mathcal{F})$ の定義である。
\end{proof}

\begin{lemma}
\label{cohomology-lemma-cech-spectral-sequence}
$X$ を環付き空間，$\mathcal{U} : U = \bigcup_{i \in I} U_i$ を開被覆とする。
任意の $\mathcal{O}_X$-加群の層 $\mathcal{F}$ に対して，
$$
E_2^{p, q} = \check{H}^p(\mathcal{U}, \underline{H}^q(\mathcal{F}))
$$
をもち $H^{p + q}(U, \mathcal{F})$ に収束するスペクトル系列
$(E_r, d_r)_{r \geq 0}$ がある。このスペクトル系列は
$\mathcal{F}$ に関して関手的である。
\end{lemma}

\begin{proof}
これは関手
$$
i :  \textit{Mod}(\mathcal{O}_X) \to \textit{PMod}(\mathcal{O}_X)
\quad\text{and}\quad
\check{H}^0(\mathcal{U}, - ) : \textit{PMod}(\mathcal{O}_X)
\to \text{Mod}_{\mathcal{O}_X(U)}.
$$
に対する Grothendieck スペクトル系列である（《導来圏》補題
\ref{derived-lemma-grothendieck-spectral-sequence} を参照）。実際，補題
\ref{cohomology-lemma-cech-h0} により
$\check{H}^0(\mathcal{U}, i(\mathcal{F})) = \mathcal{F}(U)$ である。補題
\ref{cohomology-lemma-injective-trivial-cech} により $i(\mathcal{I})$ は
{\v C}ech 非輪状である。また，補題
\ref{cohomology-lemma-cech-cohomology-derived-presheaves} により，
$\textit{PMod}(\mathcal{O}_X)$ 上の関手として
$\check{H}^p(\mathcal{U}, -) = R^p\check{H}^0(\mathcal{U}, -)$ である。
以上を合わせれば補題を得る。
\end{proof}

\begin{lemma}
\label{cohomology-lemma-cech-spectral-sequence-application}
$X$ を環付き空間，$\mathcal{U} : U = \bigcup_{i \in I} U_i$ を開被覆，
$\mathcal{F}$ を $\mathcal{O}_X$-加群の層とする。すべての $i > 0$，
すべての $p \geq 0$，およびすべての $i_0, \ldots, i_p \in I$ に対して
$H^i(U_{i_0 \ldots i_p}, \mathcal{F}) = 0$ と仮定する。このとき
$\mathcal{O}_X(U)$-加群として
$\check{H}^p(\mathcal{U}, \mathcal{F}) = H^p(U, \mathcal{F})$ である。
\end{lemma}

\begin{proof}
補題 \ref{cohomology-lemma-cech-spectral-sequence} のスペクトル系列を用いる。
仮定は，$q \not = 0$ を満たすすべての $(p, q)$ に対して
$E_2^{p, q} = 0$ であることを意味する。したがってスペクトル系列は
$E_2$ で退化し，結果が従う。
\end{proof}

\begin{lemma}
\label{cohomology-lemma-ses-cech-h1}
$X$ を環付き空間とし，
$$
0 \to \mathcal{F} \to \mathcal{G} \to \mathcal{H} \to 0
$$
を $\mathcal{O}_X$-加群の層の短完全列とする。$U \subset X$ を開集合とする。
$\check{H}^1(\mathcal{U}, \mathcal{F}) = 0$ を満たす $U$ の開被覆
$\mathcal{U}$ の共終系が存在するならば，写像
$\mathcal{G}(U) \to \mathcal{H}(U)$ は全射である。
\end{lemma}

\begin{proof}
元 $s \in \mathcal{H}(U)$ をとる。(a)
$\check{H}^1(\mathcal{U}, \mathcal{F}) = 0$ であり，かつ (b)
$s|_{U_i}$ が切断 $s_i \in \mathcal{G}(U_i)$ の像であるような開被覆
$\mathcal{U} : U = \bigcup_{i \in I} U_i$ を選ぶ。(b) を満たす
$\mathcal{U}$ は確かに存在するので，補題の仮定から (a) と (b) をともに満たす
$\mathcal{U}$ を選べる。切断
$$
s_{i_0i_1} = s_{i_1}|_{U_{i_0i_1}} - s_{i_0}|_{U_{i_0i_1}}.
$$
を考える。$s_i$ は $s$ の持ち上げなので，
$s_{i_0i_1} \in \mathcal{F}(U_{i_0i_1})$ である。
$\check{H}^1(\mathcal{U}, \mathcal{F})$ の消滅により，
$$
s_{i_0i_1} = t_{i_1}|_{U_{i_0i_1}} - t_{i_0}|_{U_{i_0i_1}}.
$$
を満たす切断 $t_i \in \mathcal{F}(U_i)$ が存在する。このとき切断
$s_i - t_i$ は明らかに層条件を満たし，$s$ に移る $U$ 上の
$\mathcal{G}$ の切断へ貼り合わさる。これでよい。
\end{proof}

\begin{lemma}
\label{cohomology-lemma-cech-vanish}
\begin{slogan}
アーベル群の層の高次 {\v C}ech コホモロジーがすべての開被覆について
消滅するならば，高次コホモロジーも消滅する。
\end{slogan}
$X$ を環付き空間とし，$\mathcal{F}$ を次を満たす
$\mathcal{O}_X$-加群の層とする：
$$
\check{H}^p(\mathcal{U}, \mathcal{F}) = 0
$$
が，すべての $p > 0$ と，$X$ の開集合の任意の開被覆
$\mathcal{U} : U = \bigcup_{i \in I} U_i$ に対して成り立つ。このとき，
任意の $p > 0$ と任意の開集合 $U \subset X$ に対して
$H^p(U, \mathcal{F}) = 0$ である。
\end{lemma}

\begin{proof}
$\mathcal{F}$ を補題の仮定を満たす層とする。このことを「$\mathcal{F}$ の
高次 {\v C}ech コホモロジーは任意の開被覆について消滅する」と表す。
$\mathcal{F}$ から入射的 $\mathcal{O}_X$-加群の層 $\mathcal{I}$ への埋め込み
$\mathcal{F} \to \mathcal{I}$ を選ぶ。補題
\ref{cohomology-lemma-injective-trivial-cech} により，$\mathcal{I}$ の高次
{\v C}ech コホモロジーは任意の開被覆について消滅する。
$\mathcal{Q} = \mathcal{I}/\mathcal{F}$ と置くと，短完全列
$$
0 \to \mathcal{F} \to \mathcal{I} \to \mathcal{Q} \to 0.
$$
を得る。補題 \ref{cohomology-lemma-ses-cech-h1} と仮定により，この列は実は前層の列として
完全である！ 特に，任意の開被覆 $\mathcal{U}$ に対して
{\v C}ech コホモロジー群の長完全列を得る。例えば補題
\ref{cohomology-lemma-cech-cohomology-delta-functor-presheaves} を参照せよ。
したがって $\mathcal{Q}$ も，すべての開被覆について高次
{\v C}ech コホモロジーが消滅する $\mathcal{O}_X$-加群の層である。

\medskip\noindent
次に，任意の開集合 $U \subset X$ に対するコホモロジー長完全列
$$
\xymatrix{
0 \ar[r] &
H^0(U, \mathcal{F}) \ar[r] &
H^0(U, \mathcal{I}) \ar[r] &
H^0(U, \mathcal{Q}) \ar[lld] \\
&
H^1(U, \mathcal{F}) \ar[r] &
H^1(U, \mathcal{I}) \ar[r] &
H^1(U, \mathcal{Q}) \ar[lld] \\
&
\ldots & \ldots & \ldots \\
}
$$
を見る。$\mathcal{I}$ は入射的なので，$n > 0$ に対して
$H^n(U, \mathcal{I}) = 0$ である（《導来圏》補題
\ref{derived-lemma-higher-derived-functors} を参照）。上で見たことから写像
$H^0(U, \mathcal{I}) \to H^0(U, \mathcal{Q})$ は全射であり，したがって
$H^1(U, \mathcal{F}) = 0$ である。$\mathcal{F}$ は高次
{\v C}ech コホモロジーが消滅する任意の $\mathcal{O}_X$-加群の層だったから，
$\mathcal{Q}$ もそのような層であること（上を参照）から
$H^1(U, \mathcal{Q}) = 0$ も従う。長完全列により，今度は
$H^2(U, \mathcal{F}) = 0$ が従う。同様に繰り返せばよい。
\end{proof}

\begin{lemma}
\label{cohomology-lemma-cech-vanish-basis}
（補題 \ref{cohomology-lemma-cech-vanish} の変種。）
$X$ を環付き空間，$\mathcal{B}$ を $X$ の位相の基，$\mathcal{F}$ を
$\mathcal{O}_X$-加群の層とする。次の性質をもつ開被覆の集合
$\text{Cov}$ が存在すると仮定する。
\begin{enumerate}
\item $\mathcal{U} : U = \bigcup_{i \in I} U_i$ なる任意の
$\mathcal{U} \in \text{Cov}$ に対して，$U, U_i \in \mathcal{B}$ であり，
すべての $U_{i_0 \ldots i_p} \in \mathcal{B}$ である。
\item 任意の $U \in \mathcal{B}$ に対して，$\text{Cov}$ に現れる
$U$ の開被覆は $U$ の開被覆の共終系をなす。
\item 任意の $\mathcal{U} \in \text{Cov}$ とすべての $p > 0$ に対して
$\check{H}^p(\mathcal{U}, \mathcal{F}) = 0$ である。
\end{enumerate}
このとき，すべての $p > 0$ と任意の $U \in \mathcal{B}$ に対して
$H^p(U, \mathcal{F}) = 0$ である。
\end{lemma}

\begin{proof}
$\mathcal{F}$ と $\text{Cov}$ を補題のとおりとする。このことを
「任意の $\mathcal{U} \in \text{Cov}$ に対して $\mathcal{F}$ の高次
{\v C}ech コホモロジーが消滅する」と表す。$\mathcal{F}$ から入射的
$\mathcal{O}_X$-加群の層 $\mathcal{I}$ への埋め込みを選ぶ。補題
\ref{cohomology-lemma-injective-trivial-cech} により，任意の
$\mathcal{U} \in \text{Cov}$ に対して $\mathcal{I}$ の高次
{\v C}ech コホモロジーは消滅する。$\mathcal{Q} = \mathcal{I}/\mathcal{F}$
と置くと，短完全列
$$
0 \to \mathcal{F} \to \mathcal{I} \to \mathcal{Q} \to 0.
$$
を得る。補題 \ref{cohomology-lemma-ses-cech-h1} と仮定 (2) により，この列から
任意の $U \in \mathcal{B}$ に対する完全列
$$
0 \to \mathcal{F}(U) \to \mathcal{I}(U) \to \mathcal{Q}(U) \to 0.
$$
が生じる。したがって任意の $\mathcal{U} \in \text{Cov}$ に対して
{\v C}ech 複体の短完全列
$$
0 \to
\check{\mathcal{C}}^\bullet(\mathcal{U}, \mathcal{F}) \to
\check{\mathcal{C}}^\bullet(\mathcal{U}, \mathcal{I}) \to
\check{\mathcal{C}}^\bullet(\mathcal{U}, \mathcal{Q}) \to 0
$$
を得る。実際，仮定 (1) により {\v C}ech 複体の各項は
$\mathcal{B}$ の元上での値の積からなる。特に，任意の開被覆
$\mathcal{U} \in \text{Cov}$ に対して {\v C}ech コホモロジー群の長完全列を
得る。したがって $\mathcal{Q}$ も，すべての
$\mathcal{U} \in \text{Cov}$ に対して高次 {\v C}ech コホモロジーが消滅する
$\mathcal{O}_X$-加群の層である。

\medskip\noindent
次に，任意の $U \in \mathcal{B}$ に対するコホモロジー長完全列
$$
\xymatrix{
0 \ar[r] &
H^0(U, \mathcal{F}) \ar[r] &
H^0(U, \mathcal{I}) \ar[r] &
H^0(U, \mathcal{Q}) \ar[lld] \\
&
H^1(U, \mathcal{F}) \ar[r] &
H^1(U, \mathcal{I}) \ar[r] &
H^1(U, \mathcal{Q}) \ar[lld] \\
&
\ldots & \ldots & \ldots \\
}
$$
を見る。$\mathcal{I}$ は入射的なので，$n > 0$ に対して
$H^n(U, \mathcal{I}) = 0$ である（《導来圏》補題
\ref{derived-lemma-higher-derived-functors} を参照）。上で見たことから写像
$H^0(U, \mathcal{I}) \to H^0(U, \mathcal{Q})$ は全射であり，したがって
$H^1(U, \mathcal{F}) = 0$ である。$\mathcal{F}$ は，すべての
$\mathcal{U} \in \text{Cov}$ に対して高次 {\v C}ech コホモロジーが消滅する
任意の $\mathcal{O}_X$-加群の層だったから，$\mathcal{Q}$ もそのような層で
あること（上を参照）から $H^1(U, \mathcal{Q}) = 0$ も従う。長完全列により，
今度は $H^2(U, \mathcal{F}) = 0$ が従う。同様に繰り返せばよい。
\end{proof}

\begin{lemma}
\label{cohomology-lemma-pushforward-injective}
$f : X \to Y$ を環付き空間の射とし，$\mathcal{I}$ を入射的
$\mathcal{O}_X$-加群の層とする。このとき次が成り立つ。
\begin{enumerate}
\item $\check{H}^p(\mathcal{V}, f_*\mathcal{I}) = 0$
が，すべての $p > 0$ と，$Y$ の任意の開被覆
$\mathcal{V} : V = \bigcup_{j \in J} V_j$ に対して成り立つ。
\item すべての $p > 0$ と任意の開集合 $V \subset Y$ に対して
$H^p(V, f_*\mathcal{I}) = 0$ である。
\end{enumerate}
言い換えれば，任意の開集合 $V \subset Y$ に対して
$f_*\mathcal{I}$ は $\Gamma(V, -)$ に関して右非輪状である
（《導来圏》定義 \ref{derived-definition-derived-functor} を参照）。
\end{lemma}

\begin{proof}
$\mathcal{U} : f^{-1}(V) = \bigcup_{j \in J} f^{-1}(V_j)$ と置く。
これは $X$ の開被覆であり，
$$
\check{\mathcal{C}}^\bullet(\mathcal{V}, f_*\mathcal{I}) =
\check{\mathcal{C}}^\bullet(\mathcal{U}, \mathcal{I}).
$$
である。実際，
$$
f_*\mathcal{I}(V_{j_0 \ldots j_p})
= \mathcal{I}(f^{-1}(V_{j_0 \ldots j_p})) =
\mathcal{I}(f^{-1}(V_{j_0}) \cap \ldots \cap f^{-1}(V_{j_p}))
= \mathcal{I}(U_{j_0 \ldots j_p}).
$$
である。したがって補題の第1の主張は補題
\ref{cohomology-lemma-injective-trivial-cech} から従う。第2の主張は第1の主張と補題
\ref{cohomology-lemma-cech-vanish} から従う。
\end{proof}

\noindent
次の補題から，特に $f_* : \textit{Ab}(X) \to \textit{Ab}(Y)$ は入射的な
アーベル群の層を入射的なアーベル群の層に移すことが従う。

\begin{lemma}
\label{cohomology-lemma-pushforward-injective-flat}
$f : X \to Y$ を環付き空間の射とし，$f$ は平坦であると仮定する。
このとき，任意の入射的 $\mathcal{O}_X$-加群の層 $\mathcal{I}$ に対して
$f_*\mathcal{I}$ は入射的 $\mathcal{O}_Y$-加群の層である。
\end{lemma}

\begin{proof}
この場合，関手 $f^*$ は単射を単射に移す（《加群》補題
\ref{modules-lemma-pullback-flat}）。したがって結果は《ホモロジー》補題
\ref{homology-lemma-adjoint-preserve-injectives} から従う。
\end{proof}

\begin{lemma}
\label{cohomology-lemma-cohomology-products}
$(X, \mathcal{O}_X)$ を環付き空間，$I$ を集合とする。各 $i \in I$ に対し，
$\mathcal{F}_i$ を $\mathcal{O}_X$-加群の層とする。$U \subset X$ を
開集合とする。標準的な写像
$$
H^p(U, \prod\nolimits_{i \in I} \mathcal{F}_i)
\longrightarrow
\prod\nolimits_{i \in I} H^p(U, \mathcal{F}_i)
$$
は $p = 0$ のとき同型であり，$p = 1$ のとき単射である。
\end{lemma}

\begin{proof}
$p = 0$ に対する主張は，層の積がその台となる前層の積に等しいことから従う。
《層》第 \ref{sheaves-section-limits-sheaves} 節を参照せよ。
$p = 1$ に対する証明。$\mathcal{F} = \prod \mathcal{F}_i$ と置く。
$\xi \in H^1(U, \mathcal{F})$ が $\prod H^1(U, \mathcal{F}_i)$ において
零に移るとする。コホモロジーの局所性（補題
\ref{cohomology-lemma-kill-cohomology-class-on-covering} を参照）により，すべての $j$ に
対して $\xi|_{U_j} = 0$ となる開被覆
$\mathcal{U} : U = \bigcup U_j$ が存在する。補題 \ref{cohomology-lemma-cech-h1} により，
$\xi$ は元
$\check \xi \in \check H^1(\mathcal{U}, \mathcal{F})$.
から来る。すべての $i$ に対して写像
$\check H^1(\mathcal{U}, \mathcal{F}_i) \to H^1(U, \mathcal{F}_i)$ は単射で
あり（補題 \ref{cohomology-lemma-cech-h1} による），$\xi$ の像は
$\prod H^1(U, \mathcal{F}_i)$ において零なので，その像は
$\check H^1(\mathcal{U}, \mathcal{F}_i)$ において
$\check \xi_i = 0$ である。一方，$\mathcal{F} = \prod \mathcal{F}_i$ なので，
$\check{\mathcal{C}}^\bullet(\mathcal{U}, \mathcal{F})$ は複体
$\check{\mathcal{C}}^\bullet(\mathcal{U}, \mathcal{F}_i)$ の積である。
したがって《ホモロジー》補題
\ref{homology-lemma-product-abelian-groups-exact} により，所望のとおり
$\check \xi = 0$ を得る。
\end{proof}







\section{フラスク層}
\label{cohomology-section-flasque}

\noindent
定義は次のとおりである。

\begin{definition}
\label{cohomology-definition-flasque}
$X$ を位相空間とする。集合の前層 $\mathcal{F}$ が
{\it フラスク}または {\it フラビー}であるとは，$X$ の任意の開集合
$U \subset V$ に対して制限写像
$\mathcal{F}(V) \to \mathcal{F}(U)$ が全射であることをいう。
\end{definition}

\noindent
この用語を，アーベル群の層，および $X$ が環付き空間であるときの加群の層にも
用いる。任意の開集合 $U \subset X$ に対して制限写像
$\mathcal{F}(X) \to \mathcal{F}(U)$ が全射であると仮定するだけで十分なのは
明らかである。

\begin{lemma}
\label{cohomology-lemma-injective-flasque}
$(X, \mathcal{O}_X)$ を環付き空間とする。このとき任意の入射的
$\mathcal{O}_X$-加群の層はフラスクである。
\end{lemma}

\begin{proof}
これは補題 \ref{cohomology-lemma-injective-restriction-surjective} の言い換えである。
\end{proof}

\begin{lemma}
\label{cohomology-lemma-flasque-acyclic}
$(X, \mathcal{O}_X)$ を環付き空間とする。任意のフラスク
$\mathcal{O}_X$-加群の層は，$R\Gamma(X, -)$，および $X$ の任意の開集合
$U$ に対する $R\Gamma(U, -)$ に関して非輪状である。
\end{lemma}

\begin{proof}
《導来圏》補題 \ref{derived-lemma-subcategory-right-acyclics} を用いて証明する。
任意の入射的加群の層はフラスクなので，任意の $\mathcal{O}_X$-加群の層を
フラスク加群の層へ埋め込める。《入射的対象》補題
\ref{injectives-lemma-abelian-sheaves-space} を参照せよ。したがって，短完全列
$$
0 \to \mathcal{F} \to \mathcal{G} \to \mathcal{H} \to 0
$$
が与えられ，$\mathcal{F},\mathcal{G}$ がフラスクならば，
$\mathcal{H}$ もフラスクであり，かつ $X$ の任意の開集合上で切断をとった後も
列が短完全であることを示せば十分である。実際，第2の主張から第1の主張が
従う。そこで $U \subset X$ を開部分空間，$s \in \mathcal{H}(U)$ とする。
$s$ を $U$ 上の $\mathcal{G}$ の切断へ持ち上げられることを示す。
そのため，$V \subset U$ が開で，$t \in \mathcal{G}(V)$ が
$\mathcal{H}$ において $s|_V$ に移る切断であるような組 $(V,t)$ の集合を
$T$ とする。$V \subset V'$ かつ $t'|_V=t$ であるとき，またそのときに限り
$(V,t)\leq(V',t')$ と定め，$T$ に半順序を入れる。
$(V_\alpha,t_\alpha)$，$\alpha\in A$ が $T$ の全順序部分集合ならば，
$V=\bigcup V_\alpha$ は開であり，$\mathcal{G}$ の層条件により，各
$V_\alpha$ 上で $t_\alpha$ に制限される一意な切断
$t\in\mathcal{G}(V)$ が存在する。したがって Zorn の補題により，$T$ には
極大元 $(V,t)$ が存在する。$V=U$ を示せば証明が完了する。
任意の $x\in U$ をとる。$x$ の十分小さい開近傍 $W\subset U$ と，
$\mathcal{H}$ において $s|_W$ に移る $t'\in\mathcal{G}(W)$ を選べる。
このとき $t'|_{W\cap V}-t|_{W\cap V}$ は $\mathcal{H}$ において零に移るので，
ある切断 $r'\in\mathcal{F}(W\cap V)$ から来る。$\mathcal{F}$ はフラスクなので，
$W\cap V$ 上で $r'$ に制限される切断 $r\in\mathcal{F}(W)$ が存在する。
$t'$ を $r$ の像で修正すれば，$t$ と $t'$ の $W\cap V$ 上への制限が
同じであると仮定してよい。$\mathcal{G}$ の層条件により，$W\cup V$ 上の
$\mathcal{G}$ の切断 $\tilde t$ で，$t$ と $t'$ に制限されるものが存在する。
$(V,t)$ の極大性から $V\cup W=V$ を得る。したがって $x\in V$ であり，
証明が完了する。
\end{proof}

\noindent
次の補題はフラスク前層に対しては成り立たない。

\begin{lemma}
\label{cohomology-lemma-flasque-acyclic-cech}
$(X, \mathcal{O}_X)$ を環付き空間，$\mathcal{F}$ を
$\mathcal{O}_X$-加群の層，$\mathcal{U} : U = \bigcup U_i$ を開被覆とする。
$\mathcal{F}$ がフラスクならば，$p > 0$ に対して
$\check{H}^p(\mathcal{U}, \mathcal{F}) = 0$ である。
\end{lemma}

\begin{proof}
補題 \ref{cohomology-lemma-cech-spectral-sequence} の主張で用いた前層
$\underline{H}^q(\mathcal{F})$ は，補題 \ref{cohomology-lemma-flasque-acyclic} により
零である。したがって，再び補題 \ref{cohomology-lemma-flasque-acyclic} により
$\check{H}^p(U, \mathcal{F}) = H^p(U, \mathcal{F}) = 0$ である。
\end{proof}

\begin{lemma}
\label{cohomology-lemma-flasque-acyclic-pushforward}
$f : (X, \mathcal{O}_X) \to (Y, \mathcal{O}_Y)$ を環付き空間の射とし，
$\mathcal{F}$ を $\mathcal{O}_X$-加群の層とする。$\mathcal{F}$ がフラスク
ならば，$p > 0$ に対して $R^pf_*\mathcal{F} = 0$ である。
\end{lemma}

\begin{proof}
補題 \ref{cohomology-lemma-describe-higher-direct-images} と補題
\ref{cohomology-lemma-flasque-acyclic} から直ちに従う。
\end{proof}

\noindent
次の補題は有限被覆についての初等的な帰納法でも証明できる。
\cite{FOAG} の {\v C}ech コホモロジーの議論と比較せよ。

\begin{lemma}
\label{cohomology-lemma-vanishing-ravi}
$X$ を位相空間，$\mathcal{F}$ を $X$ 上のアーベル群の層，
$\mathcal{U} : U = \bigcup_{i \in I} U_i$ を開被覆とする。
$U'$ が $U_{i_0 \ldots i_p}$ の形の開集合の任意の和であるとき，制限写像
$\mathcal{F}(U) \to \mathcal{F}(U')$ は全射であると仮定する。
このとき，$p > 0$ に対して
$\check{H}^p(\mathcal{U}, \mathcal{F})$ は消滅する。
\end{lemma}

\begin{proof}
$Y$ を $I$ の空でない部分集合全体の集合とする。文字
$A,B,C,\ldots$ で $Y$ の元，すなわち $I$ の空でない部分集合を表す。
有限で空でない部分集合 $J \subset I$ に対し，
$$
V_J = \{A \in Y \mid J \subset A\}
$$
と置く。これは $V_{\{i\}} = \{A \in Y \mid i \in A\}$ であり，
$V_J = \bigcap_{j \in J} V_{\{j\}}$.
また $V_J \subset V_K$ であるための必要十分条件は $J \supset K$ である。
部分集合 $V_J$ の集まりが $Y$ の位相の基となるような $Y$ 上の位相が一意に
存在する。任意の開集合は，ある有限部分集合の族 $J_t$ に対して
$$
V = \bigcup\nolimits_{t \in T} V_{J_t}
$$
の形である。$J_t \subset J_{t'}$ ならば，$V$ を変えずに族から
$J_{t'}$ を除ける。したがって $J_t$ の間には包含関係がないと仮定してよい。
この場合，$V$ の極小元は集合 $A=J_t$ である。したがって開集合 $V$ から族
$(J_t)_{t \in T}$ を読み取れる。

\medskip\noindent
$Y$ の開被覆を完全に記述できる。まず，元 $A\in Y$ は $I$ の空でない
部分集合なので，
$$
Y = \bigcup\nolimits_{i \in I} V_{\{i\}}
$$
である。他の被覆を理解するため，$V$ を上のとおりとし，
$V_s\subset Y$ を族 $(J_{s,t})_{t\in T_s}$ に対応する開集合とする。このとき
$$
V = \bigcup\nolimits_{s \in S} V_s
$$
であるための必要十分条件は，各 $t\in T$ に対して，ある $s\in S$ と
$t_s\in T_s$ が存在して $J_t=J_{s,t_s}$ となることである。実際，族
$(J_t)_{t\in T}$ は極小なので，極小元 $A=J_t$ はある $s$ に対して
$V_s$ に属さなければならず，したがってある $t_s\in T_s$ に対して
$A\in V_{J_{t_s}}$ である。しかし $A$ は $V_s$ においても極小なので，
$J_{t_s}=J_t$ を得る。

\medskip\noindent
次に，$Y$ の開集合全体から $X$ の開集合全体への写像を作る。すなわち，
$Y$ を $U$ に移し，開集合 $V_J$ には規則
$$
V_J \mapsto U_J = \bigcap\nolimits_{i \in J} U_i
$$
を用い，任意の開集合 $V$ へは規則
$$
V = \bigcup\nolimits_{t \in T} V_{J_t}
\mapsto
\bigcup\nolimits_{t \in T} U_{J_t}
$$
で拡張する。上で与えた $Y$ の開被覆の分類から，この規則は開被覆を開被覆へ
移すことが分かる。したがって，$\mathcal{G}(Y)=\mathcal{F}(U)$ と置き，
$V=\bigcup\nolimits_{t\in T}V_{J_t}$ に対して
$$
\mathcal{G}(V) = \mathcal{F}\left(\bigcup\nolimits_{t \in T} U_{J_t}\right)
$$
と置き，$\mathcal{F}$ の制限写像を用いることにより，$Y$ 上の
アーベル群の層 $\mathcal{G}$ を得る。

\medskip\noindent
以上の準備のもとで補題を証明する。$Y$ の開被覆
$\mathcal{V} : Y = \bigcup_{i \in I} V_{\{i\}}$ がある。構成により，
{\v C}ech 複体の等式
$$
\check{C}^\bullet(\mathcal{V}, \mathcal{G}) =
\check{C}^\bullet(\mathcal{U}, \mathcal{F})
$$
がある。層 $\mathcal{G}$ は $Y$ 上フラスクなので（補題の主張における
$\mathcal{F}$ についての仮定による），消滅は補題
\ref{cohomology-lemma-flasque-acyclic-cech} から従う。
\end{proof}




\section{Leray スペクトル系列}
\label{cohomology-section-Leray}

\begin{lemma}
\label{cohomology-lemma-before-Leray}
$f : X \to Y$ を環付き空間の射とする。このとき可換図式
$$
\xymatrix{
D^{+}(X) \ar[rr]_-{R\Gamma(X, -)} \ar[d]_{Rf_*} & &
D^{+}(\mathcal{O}_X(X)) \ar[d]^{\text{スカラー制限}} \\
D^{+}(Y) \ar[rr]^-{R\Gamma(Y, -)} & &
D^{+}(\mathcal{O}_Y(Y))
}
$$
がある。より一般に，任意の開集合 $V \subset Y$ と
$U = f^{-1}(V)$ に対して可換図式
$$
\xymatrix{
D^{+}(X) \ar[rr]_-{R\Gamma(U, -)} \ar[d]_{Rf_*} & &
D^{+}(\mathcal{O}_X(U)) \ar[d]^{\text{スカラー制限}} \\
D^{+}(Y) \ar[rr]^-{R\Gamma(V, -)} & &
D^{+}(\mathcal{O}_Y(V))
}
$$
がある。さらに詳しい説明については注意 \ref{cohomology-remark-elucidate-lemma} も参照せよ。
\end{lemma}

\begin{proof}
$\mathcal{O}_X$-加群 $\mathcal{F}$ に，写像
$f^\sharp : \mathcal{O}_Y(Y) \to \mathcal{O}_X(X)$ を介して
$\mathcal{O}_Y(Y)$-加群とみなした $\mathcal{F}$ の大域切断を対応させる関手を
$\Gamma_{res} : \textit{Mod}(\mathcal{O}_X) \to \text{Mod}_{\mathcal{O}_Y(Y)}$
とする。また，
$restriction : \text{Mod}_{\mathcal{O}_X(X)} \to \text{Mod}_{\mathcal{O}_Y(Y)}$
を $f^\sharp : \mathcal{O}_Y(Y) \to \mathcal{O}_X(X)$ によって誘導される
スカラー制限関手とする。$restriction$ は完全であるから，その右導来関手は
スカラー制限関手をそのまま適用することで計算されることに注意せよ
（《導来圏》補題 \ref{derived-lemma-right-derived-exact-functor} を参照）。明らかに
$$
\Gamma_{res}
=
restriction \circ \Gamma(X, -)
=
\Gamma(Y, -) \circ f_*
$$
である。これら二つの合成のいずれにも《導来圏》補題
\ref{derived-lemma-compose-derived-functors} が適用できると主張する。
第一の合成については上の考察から明らかである。第二の合成については，
入射的 $\mathcal{O}_X$-加群が $Y$ 上の $\Gamma(Y, -)$-非輪状層に写ることを
意味する補題 \ref{cohomology-lemma-pushforward-injective} から従う。
\end{proof}

\begin{remark}
\label{cohomology-remark-elucidate-lemma}
補題 \ref{cohomology-lemma-before-Leray} の意味を素朴に説明しよう。この補題は，
$f : X \to Y$，$\mathcal{F} \in \textit{Mod}(\mathcal{O}_X)$，および
入射分解 $\mathcal{F} \to \mathcal{I}^\bullet$ が与えられたとき，
$$
\begin{matrix}
R\Gamma(X, \mathcal{F}) & \text{は次で表される} &
\Gamma(X, \mathcal{I}^\bullet) \\
Rf_*\mathcal{F} & \text{は次で表される} & f_*\mathcal{I}^\bullet \\
R\Gamma(Y, Rf_*\mathcal{F}) & \text{は次で表される} &
\Gamma(Y, f_*\mathcal{I}^\bullet)
\end{matrix}
$$
となることを述べている。最後の事実は Leray の非輪状性補題
（《導来圏》補題 \ref{derived-lemma-leray-acyclicity}）と
補題 \ref{cohomology-lemma-pushforward-injective} から従う。最後に，これと自明な等式
$$
\Gamma(X, \mathcal{I}^\bullet)
=
\Gamma(Y, f_*\mathcal{I}^\bullet).
$$
を合わせれば，補題の図式の可換性が得られる。
\end{remark}

\begin{lemma}
\label{cohomology-lemma-modules-abelian}
$X$ を環付き空間，$\mathcal{F}$ を $\mathcal{O}_X$-加群とする。
\begin{enumerate}
\item 開集合 $U \subset X$ に対する $\mathcal{F}$ のコホモロジー群
$H^i(U, \mathcal{F})$ は，$\mathcal{F}$ を $\mathcal{O}_X$-加群として
計算しても，アーベル群の層として計算しても同一である。
\item $f : X \to Y$ を環付き空間の射とする。$\mathcal{F}$ の高次順像
$R^if_*\mathcal{F}$ は，$\mathcal{F}$ を $\mathcal{O}_X$-加群として
計算しても，アーベル群の層として計算しても同一である。
\end{enumerate}
下に有界な $\mathcal{O}_X$-加群の複体についても同様の主張が成り立つ。
\end{lemma}

\begin{proof}
台となる位相空間上では恒等写像であり，環の層の一意な写像
$\underline{\mathbf{Z}}_X \to \mathcal{O}_X$ により与えられる
環付き空間の射
$(X, \mathcal{O}_X) \to (X, \underline{\mathbf{Z}}_X)$ を考える。
$\mathcal{F}$ を $\mathcal{O}_X$-加群とする。同じ層を
$\underline{\mathbf{Z}}_X$-加群，すなわちアーベル群の層とみなしたものを
$\mathcal{F}_{ab}$ と書く。$\mathcal{F} \to \mathcal{I}^\bullet$ を
入射分解とする。注意 \ref{cohomology-remark-elucidate-lemma} より，
$\Gamma(X, \mathcal{I}^\bullet)$ は $R\Gamma(X, \mathcal{F})$ と
$R\Gamma(X, \mathcal{F}_{ab})$ の両方を計算する。これで (1) が示された。

\medskip\noindent
(2) を示すには (1) と補題 \ref{cohomology-lemma-describe-higher-direct-images} を用いればよい。
結果は直ちに従う。
\end{proof}

\begin{lemma}[Leray スペクトル系列]
\label{cohomology-lemma-Leray}
$f : X \to Y$ を環付き空間の射とし，$\mathcal{F}^\bullet$ を
下に有界な $\mathcal{O}_X$-加群の複体とする。このときスペクトル系列
$$
E_2^{p, q} = H^p(Y, R^qf_*(\mathcal{F}^\bullet))
$$
があり，$H^{p + q}(X, \mathcal{F}^\bullet)$ に収束する。
\end{lemma}

\begin{proof}
これは関手の合成 $\Gamma_{res} = \Gamma(Y, -) \circ f_*$ から得られる
Grothendieck スペクトル系列（《導来圏》補題
\ref{derived-lemma-grothendieck-spectral-sequence}）にほかならない。ここで
$\Gamma_{res}$ は補題 \ref{cohomology-lemma-before-Leray} の証明で定めた関手である。
《導来圏》補題 \ref{derived-lemma-grothendieck-spectral-sequence} の仮定が
満たされることについては，補題 \ref{cohomology-lemma-before-Leray} の証明または
注意 \ref{cohomology-remark-elucidate-lemma} を参照せよ。
\end{proof}

\begin{remark}
\label{cohomology-remark-Leray-ss-more-structure}
補題 \ref{cohomology-lemma-Leray} で証明した形の Leray スペクトル系列は
$\Gamma(Y, \mathcal{O}_Y)$-加群のスペクトル系列である。しかし実際には，
これが $\Gamma(X, \mathcal{O}_X)$-加群のスペクトル系列であることも
容易に分かる。例えば $f$ から環付き空間の射
$f' :  (X, \mathcal{O}_X) \to (Y, f_*\mathcal{O}_X)$
が生じる。補題 \ref{cohomology-lemma-modules-abelian} により，
$\mathcal{O}_X$-加群 $\mathcal{F}$ と $f$ に対する Leray スペクトル系列の項
$E_r^{p, q}$ は，少なくとも $r \geq 2$ に対して，$\mathcal{F}$ と $f'$ に
対するものと同一である。すなわち，いずれも $\mathcal{F}$ をアーベル群の層と
みなしたときの Leray スペクトル系列の項と一致する。そして
$(f_*\mathcal{O}_X)(Y) = \mathcal{O}_X(X)$ だから結果が従う。
Leray スペクトル系列がさらに付加的な構造をもつこともしばしばある。
\end{remark}

\begin{lemma}
\label{cohomology-lemma-apply-Leray}
$f : X \to Y$ を環付き空間の射とし，$\mathcal{F}$ を
$\mathcal{O}_X$-加群とする。
\begin{enumerate}
\item $q > 0$ に対して $R^qf_*\mathcal{F} = 0$ ならば，すべての $p$ に対し
$H^p(X, \mathcal{F}) = H^p(Y, f_*\mathcal{F})$ である。
\item すべての $q$ と $p > 0$ に対して
$H^p(Y, R^qf_*\mathcal{F}) = 0$ ならば，すべての $q$ に対し
$H^q(X, \mathcal{F}) = H^0(Y, R^qf_*\mathcal{F})$ である。
\end{enumerate}
\end{lemma}

\begin{proof}
これらは Leray スペクトル系列を $E_2$ で退化させる二つの簡単な条件である。
これらの事実は（スペクトル系列を使わずに）直接証明することもでき，
層のコホモロジーのよい演習となる。
\end{proof}

\begin{lemma}
\label{cohomology-lemma-higher-direct-images-compose}
\begin{slogan}
合成の全導来関手は，全導来関手の合成である。
\end{slogan}
$f : X \to Y$ と $g : Y \to Z$ を環付き空間の射とする。このとき
$D^{+}(X) \to D^{+}(Z)$ の関手として
$Rg_* \circ Rf_* = R(g \circ f)_*$ である。
\end{lemma}

\begin{proof}
《導来圏》補題 \ref{derived-lemma-compose-derived-functors} を適用する。
$g_* \circ f_* = (g \circ f)_*$ であることは明らかである
（《層》補題 \ref{sheaves-lemma-pushforward-composition} を参照）。
$f_*\mathcal{I}$ が $g_*$-非輪状であることを示せばよい。これは補題
\ref{cohomology-lemma-pushforward-injective} と，補題
\ref{cohomology-lemma-describe-higher-direct-images} における高次順像 $R^ig_*$ の記述から従う。
\end{proof}

\begin{lemma}[相対 Leray スペクトル系列]
\label{cohomology-lemma-relative-Leray}
$f : X \to Y$ と $g : Y \to Z$ を環付き空間の射とし，
$\mathcal{F}$ を $\mathcal{O}_X$-加群とする。このとき
$$
E_2^{p, q} = R^pg_*(R^qf_*\mathcal{F})
$$
をもつスペクトル系列があり，$R^{p + q}(g \circ f)_*\mathcal{F}$ に収束する。
このスペクトル系列は $\mathcal{F}$ に関して関手的であり，下に有界な
$\mathcal{O}_X$-加群の複体に対する版もある。
\end{lemma}

\begin{proof}
これは関手の合成に対する Grothendieck スペクトル系列であり，補題
\ref{cohomology-lemma-higher-direct-images-compose} と《導来圏》補題
\ref{derived-lemma-grothendieck-spectral-sequence} から従う。
\end{proof}














\section{コホモロジーの関手性}
\label{cohomology-section-functoriality}

\begin{lemma}
\label{cohomology-lemma-functoriality}
$f : X \to Y$ を環付き空間の射とする。$\mathcal{G}^\bullet$ と
$\mathcal{F}^\bullet$ を，それぞれ下に有界な $\mathcal{O}_Y$-加群の複体と
$\mathcal{O}_X$-加群の複体とする。
$\varphi : \mathcal{G}^\bullet \to f_*\mathcal{F}^\bullet$
を複体の射とする。このとき $D^{+}(Y)$ に標準射
$$
\mathcal{G}^\bullet
\longrightarrow
Rf_*(\mathcal{F}^\bullet)
$$
がある。さらに，この構成は三つ組
$(\mathcal{G}^\bullet, \mathcal{F}^\bullet, \varphi)$ に関して関手的である。
\end{lemma}

\begin{proof}
入射分解 $\mathcal{F}^\bullet \to \mathcal{I}^\bullet$ を選ぶ。
定義により，$Rf_*(\mathcal{F}^\bullet)$ は
$K^{+}(\mathcal{O}_Y)$ において $f_*\mathcal{I}^\bullet$ で表される。合成
$$
\mathcal{G}^\bullet \to f_*\mathcal{F}^\bullet \to f_*\mathcal{I}^\bullet
$$
は $K^{+}(Y)$ における射であり，局所化関手
$j_Y : K^{+}(Y) \to D^{+}(Y)$ を適用すると補題の射になる。
\end{proof}

\noindent
$f : X \to Y$ を環付き空間の射とする。$\mathcal{G}$ を
$\mathcal{O}_Y$-加群，$\mathcal{F}$ を $\mathcal{O}_X$-加群とする。
$\mathcal{G}$ から $\mathcal{F}$ への $f$-写像 $\varphi$ とは，写像
$\varphi : \mathcal{G} \to f_*\mathcal{F}$，または同じことだが，写像
$\varphi : f^*\mathcal{G} \to \mathcal{F}$ のことであった
（《層》定義 \ref{sheaves-definition-f-map} を参照）。このような $f$-写像から，
$D^{+}(\mathcal{O}_Y(Y))$ における複体の射
\begin{equation}
\label{cohomology-equation-functorial-derived}
\varphi :
R\Gamma(Y, \mathcal{G})
\longrightarrow
R\Gamma(X, \mathcal{F})
\end{equation}
が生じる。実際，補題 \ref{cohomology-lemma-functoriality} の
$D^{+}(Y)$ における射 $\mathcal{G} \to Rf_*\mathcal{F}$ に
$R\Gamma(Y, -)$ を適用する。補題 \ref{cohomology-lemma-before-Leray} より
$R\Gamma(X, \mathcal{F}) = R\Gamma(Y, Rf_*\mathcal{F})$
であるから，表示した矢印を得る。以下の注意 \ref{cohomology-remark-explain-arrow} で
これを完全に詳述する。特に，コホモロジー上の写像
\begin{equation}
\label{cohomology-equation-functorial}
\varphi : H^i(Y, \mathcal{G}) \longrightarrow H^i(X, \mathcal{F}).
\end{equation}
が得られる。

\begin{remark}
\label{cohomology-remark-explain-arrow}
$f : X \to Y$ を環付き空間の射とする。$\mathcal{G}$ を
$\mathcal{O}_Y$-加群，$\mathcal{F}$ を $\mathcal{O}_X$-加群とし，
$\varphi$ を $\mathcal{G}$ から $\mathcal{F}$ への $f$-写像とする。
$\mathcal{F} \to \mathcal{I}^\bullet$ を入射的 $\mathcal{O}_X$-加群の
複体による分解とする。また，$\mathcal{G} \to \mathcal{J}^\bullet$ と
$f_*\mathcal{I}^\bullet \to (\mathcal{J}')^\bullet$ を入射的
$\mathcal{O}_Y$-加群の複体による分解とする。《導来圏》補題
\ref{derived-lemma-morphisms-lift} により，図式
\begin{equation}
\label{cohomology-equation-choice}
\xymatrix{
\mathcal{G} \ar[d] \ar[r] &
f_*\mathcal{F} \ar[r] &
f_*\mathcal{I}^\bullet \ar[d] \\
\mathcal{J}^\bullet \ar[rr]^\beta & &
(\mathcal{J}')^\bullet
}
\end{equation}
を可換にする複体の写像 $\beta$ が存在する。大域切断関手を適用すると図式
$$
\xymatrix{
 & & \Gamma(Y, f_*\mathcal{I}^\bullet) \ar[d]_{qis} \ar@{=}[r] &
\Gamma(X, \mathcal{I}^\bullet) \\
\Gamma(Y, \mathcal{J}^\bullet) \ar[rr]^\beta & &
\Gamma(Y, (\mathcal{J}')^\bullet) &
}
$$
を得る。左下の複体は $R\Gamma(Y, \mathcal{G})$ を表し，右上の複体は
$R\Gamma(X, \mathcal{F})$ を表す。補題 \ref{cohomology-lemma-before-Leray} により
垂直矢印は擬同型であり，局所化関手
$K^{+}(\mathcal{O}_Y(Y)) \to D^{+}(\mathcal{O}_Y(Y))$ を適用すると可逆になる。
矢印 (\ref{cohomology-equation-functorial-derived}) は，水平写像と垂直写像の逆との合成で
与えられる。
\end{remark}





\section{細分と {\v C}ech コホモロジー}
\label{cohomology-section-refinements-cech}

\noindent
$(X, \mathcal{O}_X)$ を環付き空間とする。
$\mathcal{U} : X = \bigcup_{i \in I} U_i$ と
$\mathcal{V} : X = \bigcup_{j \in J} V_j$ を開被覆とする。
$\mathcal{U}$ が $\mathcal{V}$ の細分であると仮定する。すべての $i \in I$ に
対して $U_i \subset V_{c(i)}$ となる写像 $c : I \to J$ を選ぶ。
これは $\mathcal{O}_X$-加群の層 $\mathcal{F}$ に関して関手的な
{\v C}ech 複体の写像
$$
\gamma :
\check{\mathcal{C}}^\bullet(\mathcal{V}, \mathcal{F})
\longrightarrow
\check{\mathcal{C}}^\bullet(\mathcal{U}, \mathcal{F}),
\quad
(\xi_{j_0 \ldots j_p})
\longmapsto
(\xi_{c(i_0) \ldots c(i_p)}|_{U_{i_0 \ldots i_p}})
$$
を誘導する。$c' : I \to J$ が，すべての $i \in I$ に対して
$U_i \subset V_{c'(i)}$ となるもう一つの写像であるとする。このとき，
対応する写像 $\gamma$ と $\gamma'$ はホモトピックである。実際，
$\gamma - \gamma' = \text{d} \circ h + h \circ \text{d}$
であり，ここで
$h : \check{\mathcal{C}}^{p + 1}(\mathcal{V}, \mathcal{F}) \to
\check{\mathcal{C}}^p(\mathcal{U}, \mathcal{F})$
は公式
$$
h(\alpha)_{i_0 \ldots i_p} =
\sum\nolimits_{a = 0}^{p}
(-1)^a
\alpha_{c(i_0)\ldots c(i_a) c'(i_a) \ldots c'(i_p)}
$$
で与えられる。これが成り立つことを示す計算は省略する。より一般の場合の
証明については，(\ref{cohomology-equation-transformation}) に続く議論を参照せよ。
特に，{\v C}ech コホモロジー群上の写像は $c$ の選択に依存しない。
さらに，$\mathcal{W} : X = \bigcup_{k \in K} W_k$ が第三の開被覆であり，
$\mathcal{V}$ が $\mathcal{W}$ の細分ならば，写像 $I \to J$ と $J \to K$ に
付随する写像の合成
$$
\check{\mathcal{C}}^\bullet(\mathcal{W}, \mathcal{F})
\longrightarrow
\check{\mathcal{C}}^\bullet(\mathcal{V}, \mathcal{F})
\longrightarrow
\check{\mathcal{C}}^\bullet(\mathcal{U}, \mathcal{F})
$$
は，合成 $I \to K$ に付随する写像であることが明らかである。
したがって，{\v C}ech コホモロジー群を
$$
\check{H}^p(X, \mathcal{F}) =
\colim_\mathcal{U} \check{H}^p(\mathcal{U}, \mathcal{F})
$$
と定義できる。ここで余極限は，細分によって前順序づけられた $X$ の
すべての開被覆をわたる。

\medskip\noindent
上で定義した写像 $\gamma$ はコホモロジーへの写像と整合する。すなわち，合成
$$
\check{H}^p(\mathcal{V}, \mathcal{F}) \to
\check{H}^p(\mathcal{U}, \mathcal{F})
\xrightarrow{\text{補題 \ref{cohomology-lemma-cech-cohomology}}}
H^p(X, \mathcal{F})
$$
は，補題 \ref{cohomology-lemma-cech-cohomology} の第一の群からコホモロジーへの標準写像で
ある。以下の補題で，これを少し一般的な状況において証明する。その帰結として，
{\v C}ech コホモロジーからコホモロジーへの良定義な写像
\begin{equation}
\label{cohomology-equation-cech-to-cohomology}
\check{H}^p(X, \mathcal{F}) =
\colim_\mathcal{U} \check{H}^p(\mathcal{U}, \mathcal{F})
\longrightarrow
H^p(X, \mathcal{F})
\end{equation}
を得る。

\begin{lemma}
\label{cohomology-lemma-functoriality-cech}
$f : X \to Y$ を環付き空間の射とする。
$\varphi : f^*\mathcal{G} \to \mathcal{F}$ を，
$\mathcal{O}_Y$-加群 $\mathcal{G}$ から $\mathcal{O}_X$-加群
$\mathcal{F}$ への $f$-写像とする。
$\mathcal{U} : X = \bigcup_{i \in I} U_i$ と
$\mathcal{V} : Y = \bigcup_{j \in J} V_j$ を開被覆とする。
$\mathcal{U}$ が
$f^{-1}\mathcal{V} : X = \bigcup_{j \in J} f^{-1}(V_j)$ の細分であると仮定する。
このとき，$D^{+}(\mathcal{O}_X(X))$ に可換図式
$$
\xymatrix{
\check{\mathcal{C}}^\bullet(\mathcal{U}, \mathcal{F}) \ar[r] &
R\Gamma(X, \mathcal{F}) \\
\check{\mathcal{C}}^\bullet(\mathcal{V}, \mathcal{G}) \ar[r]
\ar[u]^\gamma &
R\Gamma(Y, \mathcal{G}) \ar[u]
}
$$
が存在する。ここで水平矢印は補題 \ref{cohomology-lemma-cech-cohomology} によって，
右の垂直矢印は (\ref{cohomology-equation-functorial-derived}) によって与えられる。
特に，コホモロジー群の可換図式
$$
\xymatrix{
\check{H}^p(\mathcal{U}, \mathcal{F}) \ar[r] &
H^p(X, \mathcal{F}) \\
\check{H}^p(\mathcal{V}, \mathcal{G}) \ar[r]
\ar[u]^\gamma &
H^p(Y, \mathcal{G}) \ar[u]
}
$$
を得る。ここで右の垂直矢印は (\ref{cohomology-equation-functorial}) である。
\end{lemma}

\begin{proof}
まず左の垂直矢印を定義する。すなわち，すべての $i \in I$ に対して
$U_i \subset f^{-1}(V_{c(i)})$ となる写像 $c : I \to J$ を選ぶ。
次数 $p$ における写像を公式
$$
\gamma(s)_{i_0 \ldots i_p} = \varphi(s)_{c(i_0) \ldots c(i_p)}
$$
で定める。仮定により，$\varphi$ は実際に写像
$\mathcal{G}(V_{c(i_0) \ldots c(i_p)}) \to \mathcal{F}(U_{i_0 \ldots i_p})$
を誘導するので，これは意味をもつ。また，これが複体の射を定めることも明らかで
ある。$X$ 上の入射分解 $\mathcal{F} \to \mathcal{I}^\bullet$ と
$Y$ 上の入射分解 $\mathcal{G} \to J^\bullet$ を選ぶ。補題
\ref{cohomology-lemma-cech-cohomology} の証明に従い，項
$$
B^{p, q} = \check{\mathcal{C}}^p(\mathcal{V}, \mathcal{J}^q)
\quad
\text{および}
\quad
A^{p, q} = \check{\mathcal{C}}^p(\mathcal{U}, \mathcal{I}^q).
$$
をもつ二重複体 $A^{\bullet, \bullet}$ と $B^{\bullet, \bullet}$ を導入する。
また，上の注意 \ref{cohomology-remark-explain-arrow} と同様に，$Y$ 上の入射分解
$f_*\mathcal{I} \to (\mathcal{J}')^\bullet$ と，
(\ref{cohomology-equation-choice}) を可換にする複体の射
$\beta : \mathcal{J} \to (\mathcal{J}')^\bullet$ を選ぶ。さらに，項
$$
(B')^{p, q} = \check{\mathcal{C}}^p(\mathcal{V}, (\mathcal{J}')^q)
\quad
\text{および}
\quad
(B'')^{p, q} = \check{\mathcal{C}}^p(\mathcal{V}, f_*\mathcal{I}^q).
$$
をもつ二重複体 $(B')^{\bullet, \bullet}$ と
$(B'')^{\bullet, \bullet}$ を導入する。$f_*\mathcal{I}^\bullet$ から
$\mathcal{I}^\bullet$ への複体の $f$-写像があることに注意せよ。したがって，
上と同じ公式が二重複体の射
$$
\gamma : (B'')^{\bullet, \bullet} \longrightarrow A^{\bullet, \bullet}.
$$
を定めることが明らかである。複体の図式
$$
\xymatrix{
\check{\mathcal{C}}^\bullet(\mathcal{U}, \mathcal{F})
\ar[r] &
\text{Tot}(A^{\bullet, \bullet}) & & &
\Gamma(X, \mathcal{I}^\bullet) \ar[lll]^{qis}
\ar@{=}[ddl]\\
\check{\mathcal{C}}^\bullet(\mathcal{V}, \mathcal{G})
\ar[r] \ar[u]^\gamma &
\text{Tot}(B^{\bullet, \bullet}) \ar[r]^\beta &
\text{Tot}((B')^{\bullet, \bullet}) &
\text{Tot}((B'')^{\bullet, \bullet}) \ar[l] \ar[llu]_{s\gamma} \\
& \Gamma(Y, \mathcal{J}^\bullet) \ar[u]^{qis} \ar[r]^\beta &
\Gamma(Y, (\mathcal{J}')^\bullet) \ar[u] &
\Gamma(Y, f_*\mathcal{I}^\bullet) \ar[u] \ar[l]_{qis}
}
$$
を考える。$\text{Tot}(A^{\bullet, \bullet})$ と
$\text{Tot}(B^{\bullet, \bullet})$ を終点とする二つの水平矢印は，補題
\ref{cohomology-lemma-cech-cohomology} で説明したものである。
左上の図形（五角形）は，(\ref{cohomology-equation-choice}) が可換であることだけから
可換である。下の二つの正方形は自明に可換である。また，右上の図形
（正方形）が可換であることも定義から直ちに分かる。これで，定義と，
図式の外側を $\check{\mathcal{C}}^\bullet(\mathcal{V}, \mathcal{G})$ から
$\Gamma(X, \mathcal{I}^\bullet)$ まで上側または下側に沿って進んだ結果が
同じであるという事実から，補題の結果が従う
（途中に現れる擬同型はすべて逆にしなければならない）。
\end{proof}





\section{Hausdorff 準コンパクト空間上のコホモロジー}
\label{cohomology-section-cohomology-LC}

\noindent
このような空間では {\v C}ech コホモロジーとコホモロジーは一致する。

\begin{lemma}
\label{cohomology-lemma-cech-always}
$X$ を位相空間，$\mathcal{F}$ をアーベル群の層とする。このとき，
(\ref{cohomology-equation-cech-to-cohomology}) で定義した写像
$\check{H}^1(X, \mathcal{F}) \to H^1(X, \mathcal{F})$ は同型である。
\end{lemma}

\begin{proof}
$\mathcal{U}$ を $X$ の開被覆とする。補題
\ref{cohomology-lemma-cech-spectral-sequence} により完全列
$$
0 \to \check{H}^1(\mathcal{U}, \mathcal{F}) \to H^1(X, \mathcal{F})
\to \check{H}^0(\mathcal{U}, \underline{H}^1(\mathcal{F}))
$$
がある。したがって写像は単射である。全射性を示すには，
$\check{H}^0(\mathcal{U}, \underline{H}^1(\mathcal{F}))$ の任意の元が，
$\mathcal{U}$ を細分で置き換えた後に零へ写ることを示せば十分である。
これは定義と，$\underline{H}^1(\mathcal{F})$ がアーベル群の前層であって，
コホモロジーの局所性によりその層化が零であることから直ちに従う
（補題 \ref{cohomology-lemma-kill-cohomology-class-on-covering} を参照）。
\end{proof}

\begin{lemma}
\label{cohomology-lemma-cech-Hausdorff-quasi-compact}
$X$ を Hausdorff 準コンパクト位相空間とし，$\mathcal{F}$ を $X$ 上の
アーベル群の層とする。このとき，(\ref{cohomology-equation-cech-to-cohomology}) で
定義した写像 $\check{H}^n(X, \mathcal{F}) \to H^n(X, \mathcal{F})$ は，
すべての $n$ に対して同型である。
\end{lemma}

\begin{proof}
$n = 0, 1$ に対して $\check{H}^n(X, -) \to H^n(X, -)$ が関手の同型で
あることはすでに分かっている（補題 \ref{cohomology-lemma-cech-always} を参照）。
関手 $H^n(X, -)$ は普遍 $\delta$-関手をなす（《導来圏》補題
\ref{derived-lemma-higher-derived-functors} を参照）。
$\check{H}^n(X, -)$ が普遍 $\delta$-関手をなし，かつ
$\check{H}^n(X, -) \to H^n(X, -)$ が境界写像と整合することを示せば，
普遍 $\delta$-関手の一意性によりこの写像は自動的に同型となる
（《ホモロジー》補題
\ref{homology-lemma-uniqueness-universal-delta-functor} を参照）。

\medskip\noindent
$0 \to \mathcal{F} \to \mathcal{G} \to \mathcal{H} \to 0$ を
$X$ 上のアーベル群の層の短完全列とし，
$\mathcal{U} : X = \bigcup_{i \in I} U_i$ を開被覆とする。これから複体の列
$$
0 \to \check{\mathcal{C}}^\bullet(\mathcal{U}, \mathcal{F}) \to
\check{\mathcal{C}}^\bullet(\mathcal{U}, \mathcal{G}) \to
\check{\mathcal{C}}^\bullet(\mathcal{U}, \mathcal{H}) \to 0
$$
を得るが，これは一般には右端で完全でない。この列は写像
$$
\check{H}^n(\mathcal{U}, \mathcal{F}) \to
\check{H}^n(\mathcal{U}, \mathcal{G}) \to
\check{H}^n(\mathcal{U}, \mathcal{H})
$$
を定めるが，境界作用素
$\delta : \check{H}^n(\mathcal{U}, \mathcal{H}) \to
\check{H}^{n + 1}(\mathcal{U}, \mathcal{F})$ を定めるには不十分である。
実際，このようなものは一般には存在しない。しかし，コサイクル
$h = (h_{i_0 \ldots i_n})$ のコホモロジー類である元
$\overline{h} \in \check{H}^n(\mathcal{U}, \mathcal{H})$ が与えられたとする。
このとき，開被覆
$$
U_{i_0 \ldots i_n} = \bigcup W_{i_0 \ldots i_n, k}
$$
を，$h_{i_0 \ldots i_n}|_{W_{i_0 \ldots i_n, k}}$ が
$W_{i_0 \ldots i_n, k}$ 上の $\mathcal{G}$ の切断へ持ち上がるように選べる。
《位相空間》補題 \ref{topology-lemma-refine-covering} により
（ここで $X$ が Hausdorff かつ準コンパクトであるという仮定を用いる），
開被覆 $\mathcal{V} : X = \bigcup_{j \in J} V_j$ と写像
$\alpha : J \to I$ を，$V_j \subset U_{\alpha(j)}$（すなわち細分）となり，
かつすべての $j_0, \ldots, j_n \in J$ に対して，ある $k$ が存在して
$V_{j_0 \ldots j_n} \subset
W_{\alpha(j_0) \ldots \alpha(j_n), k}$ となるように選べる。複体の写像
$$
\xymatrix{
0 \ar[r] &
\check{\mathcal{C}}^\bullet(\mathcal{U}, \mathcal{F}) \ar[d] \ar[r] &
\check{\mathcal{C}}^\bullet(\mathcal{U}, \mathcal{G}) \ar[d] \ar[r] &
\check{\mathcal{C}}^\bullet(\mathcal{U}, \mathcal{H}) \ar[d] \ar[r] &
0 \\
0 \ar[r] &
\check{\mathcal{C}}^\bullet(\mathcal{V}, \mathcal{F}) \ar[r] &
\check{\mathcal{C}}^\bullet(\mathcal{V}, \mathcal{G}) \ar[r] &
\check{\mathcal{C}}^\bullet(\mathcal{V}, \mathcal{H}) \ar[r] &
0
}
$$
を得る。実際，垂直矢印は {\v C}ech コホモロジー群の間の遷移写像を
定義するために用いた複体の写像である。細分の選び方から，
$$
g_{j_0 \ldots j_n} \in
\mathcal{G}(V_{j_0 \ldots j_n}),\quad
g_{j_0 \ldots j_n} \longmapsto
h_{\alpha(j_0) \ldots \alpha(j_n)}|_{V_{j_0 \ldots j_n}}
$$
を選べる。余鎖 $g = (g_{j_0 \ldots j_n})$ は一般にはコサイクルではないが，
その {\v C}ech 余境界 $\text{d}(g)$ は
$\check{\mathcal{C}}^{n + 1}(\mathcal{V}, \mathcal{H})$ において零へ写る
（上の可換図式と $h$ がコサイクルであることによる）。したがって
$\text{d}(g)$ は
$\check{\mathcal{C}}^\bullet(\mathcal{V}, \mathcal{F})$ のコサイクルである。
これにより
$$
\delta(\overline{h}) = \text{次における }\text{d}(g)\text{ の類 }
\check{H}^{n + 1}(\mathcal{V}, \mathcal{F})
$$
を定義できる。さて，元 $\xi \in \check{H}^n(X, \mathcal{G})$ が与えられたと
する。開被覆 $\mathcal{U}$ と，{\v C}ech コホモロジーを定義する余極限において
$\xi$ へ写る元 $\overline{h} \in \check{H}^n(\mathcal{U}, \mathcal{G})$ を選ぶ。
次に上のように $\mathcal{V}$ と $g$ を選び，$\delta(\xi)$ を
$\check{H}^n(X, \mathcal{F})$ における $\delta(\overline{h})$ の像と定める。
ここで多くの性質を確認しなければならないが，いずれも直接的である。
例えば，この構成が
$\mathcal{U}, \overline{h}, \mathcal{V}, \alpha : J \to I, g$ の選択に
依存しないことを確認する必要がある。$\text{d}(g)$ の類は持ち上げ
$g_{i_0 \ldots i_n}$ の選択に依存しない。なぜなら，その差は余境界となるからで
ある。$\alpha$ からの独立性も成り立つ\footnote{これは重要な確認である。
$\alpha$ の非一意性だけが，$X$ のすべての開被覆をわたる {\v C}ech 複体の
余極限を取って，{\v C}ech コホモロジーを計算する複体の短完全列を得ることを
妨げているからである。}。実際，$\alpha$ の別の選択は，上の図式において
ホモトピックな複体の垂直写像を定める（第
\ref{cohomology-section-refinements-cech} 節を参照）。他の選択については，$X$ の開被覆の
有限個の族が与えられれば，それらすべてを細分する開被覆を常に見つけられることを
用いる。加法性も確認する必要があるが，同じ方法で示される。
最後に，写像 $\check{H}^n(X, -) \to H^n(X, -)$ が境界写像と整合することを
確認する必要がある。そのために，入射分解
$$
\xymatrix{
0 \ar[r] &
\mathcal{F} \ar[r] \ar[d] &
\mathcal{G} \ar[r] \ar[d] &
\mathcal{H} \ar[r] \ar[d] &
0 \\
0 \ar[r] &
\mathcal{I}_1^\bullet \ar[r] &
\mathcal{I}_2^\bullet \ar[r] &
\mathcal{I}_3^\bullet \ar[r] &
0
}
$$
を《導来圏》補題 \ref{derived-lemma-injective-resolution-ses} のように選ぶ。
これから可換図式
$$
\xymatrix{
0 \ar[r] &
\check{\mathcal{C}}^\bullet(\mathcal{U}, \mathcal{F}) \ar[r] \ar[d] &
\check{\mathcal{C}}^\bullet(\mathcal{U}, \mathcal{F}) \ar[r] \ar[d] &
\check{\mathcal{C}}^\bullet(\mathcal{U}, \mathcal{F}) \ar[r] \ar[d] &
0 \\
0 \ar[r] &
\text{Tot}(\check{\mathcal{C}}^\bullet(\mathcal{U}, \mathcal{I}_1^\bullet))
\ar[r] &
\text{Tot}(\check{\mathcal{C}}^\bullet(\mathcal{U}, \mathcal{I}_2^\bullet))
\ar[r] &
\text{Tot}(\check{\mathcal{C}}^\bullet(\mathcal{U}, \mathcal{I}_3^\bullet))
\ar[r] &
0
}
$$
を得る。ここで $\mathcal{U}$ は上のような開被覆であり，垂直写像は
$\check{H}^n(\mathcal{U}, -) \to H^n(X, -)$ を定義するために用いた写像である
（補題 \ref{cohomology-lemma-cech-cohomology} を参照）。入射対象の複体の列は各項で
分裂完全であるから，下段の複体は完全である。したがってコホモロジーの境界写像は，
この下段の完全列に対する通常の手順で計算される
（《ホモロジー》補題 \ref{homology-lemma-long-exact-sequence-cochain} を参照）。
{\v C}ech コホモロジーの境界写像を定義した細分 $\mathcal{V}$ に移っても
同じことが成り立つ。したがって，両者は同じ構成を用いるので境界写像は一致する
（第一の境界写像が，与えられた開被覆上の {\v C}ech コホモロジーの元に対して
定義されている場合）。これで，{\v C}ech コホモロジー上の
$\delta$-関手構造の構成，およびこの構造が通常のコホモロジーに与えられた
$\delta$-関手構造と整合する理由についての議論を終える。

\medskip\noindent
最後に補題 \ref{cohomology-lemma-injective-trivial-cech} を適用すると，入射的な層上では
高次 {\v C}ech コホモロジーが零であることが分かる。したがって《ホモロジー》
補題 \ref{homology-lemma-efface-implies-universal} により，
{\v C}ech コホモロジーは普遍 $\delta$-関手である。
\end{proof}

\begin{lemma}
\label{cohomology-lemma-cohomology-of-closed}
\begin{reference}
\cite[Expose V bis, 4.1.3]{SGA4}
\end{reference}
$X$ を位相空間とする。$Z \subset X$ を準コンパクト部分集合とし，
$Z$ の任意の相異なる二点が $X$ において互いに素な開近傍をもつとする。
$X$ 上の任意のアーベル群の層 $\mathcal{F}$ に対して，標準写像
$$
\colim H^p(U, \mathcal{F})
\longrightarrow
H^p(Z, \mathcal{F}|_Z)
$$
は同型である。ここで余極限は，$X$ における $Z$ の開近傍 $U$ をわたる。
\end{lemma}

\begin{proof}
まず $p = 0$ の場合を証明する。単射性は $\mathcal{F}|_Z$ の定義から従い，
一般に成り立つ（任意の位相空間 $X$ の任意の部分集合について成り立つ）。
次に $s \in H^0(Z, \mathcal{F}|_Z)$ とする。このとき
$Z \subset \bigcup U_i$ で，かつ $s|_{Z \cap U_i}$ が
$s_i \in \mathcal{F}(U_i)$ から来るような開集合 $U_i \subset X$ を見つけられる。
したがって，$W_{ij} \cap Z = U_i \cap U_j \cap Z$ かつ
$s_i|_{W_{ij}} = s_j|_{W_{ij}}$ となる開集合
$W_{ij} \subset U_i \cap U_j$ が存在する。《位相空間》補題
\ref{topology-lemma-lift-covering-of-quasi-compact-hausdorff-subset} を適用すると，
$V_i \subset U_i$ かつ $V_i \cap V_j \subset W_{ij}$ となる $X$ の開集合
$V_i$ が得られる。ゆえに，$s_i|_{V_i}$ は貼り合わさって，$Z$ の開近傍
$\bigcup V_i$ 上の $\mathcal{F}$ の切断になる。

\medskip\noindent
証明を終えるには，$\mathcal{I}$ が $X$ 上の入射的なアーベル群の層ならば，
$p > 0$ に対して $H^p(Z, \mathcal{I}|_Z) = 0$ であることを示せば十分である。
これは短完全列と次元ずらしを用いて従う。詳細は省略する。そこで，ある
$p > 0$ に対して $\overline{\xi}$ が $H^p(Z, \mathcal{I}|_Z)$ の元であると
する。補題 \ref{cohomology-lemma-cech-Hausdorff-quasi-compact} により，
$\overline{\xi}$ は，ある $Z$ の開被覆
$\mathcal{V} : Z = \bigcup V_i$ に対する
$\check{H}^p(\mathcal{V}, \mathcal{I}|_Z)$
から来る。$\overline{\xi}$ は，
$\check{\mathcal{C}}^p(\mathcal{V}, \mathcal{I}|_Z)$ のコサイクル
$\xi = (\xi_{i_0 \ldots i_p})$ の類の像であるとする。

\medskip\noindent
規則
$$
\mathcal{I}'(V) =
\{s \in \mathcal{I}|_Z(V) \mid
\exists (U, t),\ U \subset X\text{ は開集合},
\ t \in \mathcal{I}(U),\ V = Z \cap U,\ s = t|_{Z \cap U} \}
$$
で定義される部分前層 $\mathcal{I}' \subset \mathcal{I}|_Z$ を考える。
このとき $\mathcal{I}|_Z$ は $\mathcal{I}'$ の層化である。したがって，
任意の $(p + 1)$-組 $i_0 \ldots i_p$ に対して，
$\xi_{i_0 \ldots i_p}|_{W_{i_0 \ldots i_p, k}}$ が $\mathcal{I}'$ の切断となる
ような開被覆
$V_{i_0 \ldots i_p} = \bigcup W_{i_0 \ldots i_p, k}$ を見つけられる。
《位相空間》補題 \ref{topology-lemma-refine-covering} を適用すれば，
$\mathcal{V}$ を細分した後，各 $\xi_{i_0 \ldots i_p}$ が前層
$\mathcal{I}'$ の切断であると仮定してよい。

\medskip\noindent
ある開集合 $U_i \subset X$ によって $V_i = Z \cap U_i$ と書く。
$\mathcal{I}$ はフラスクであり（補題 \ref{cohomology-lemma-injective-flasque}），
任意の $(p + 1)$-組 $i_0 \ldots i_p$ に対して
$\xi_{i_0 \ldots i_p}$ は $\mathcal{I}'$ の切断であるから，
$V_{i_0 \ldots i_p} = Z \cap U_{i_0 \ldots i_p}$ 上で
$\xi_{i_0 \ldots i_p}$ に制限される切断
$s_{i_0 \ldots i_p} \in \mathcal{I}(U_{i_0 \ldots i_p})$ を選べる。
（入射対象がフラスクであることへのこの依存は，《位相空間》補題
\ref{topology-lemma-lift-covering-of-quasi-compact-hausdorff-subset} を
もう一度適用することで避けられる。）
$s = (s_{i_0 \ldots i_p})$ を，開被覆 $U = \bigcup U_i$ に対する対応する余鎖と
する。$\text{d}(\xi) = 0$ だから，切断
$\text{d}(s)_{i_0 \ldots i_{p + 1}}$ は
$Z \cap U_{i_0 \ldots i_{p + 1}}$ 上で零に制限される。したがって，証明冒頭の
考察により，開部分集合
$W_{i_0 \ldots i_{p + 1}} \subset U_{i_0 \ldots i_{p + 1}}$
で，$Z \cap W_{i_0 \ldots i_{p + 1}} =
Z \cap U_{i_0 \ldots i_{p + 1}}$ かつ
$\text{d}(s)_{i_0 \ldots i_{p + 1}}|_{W_{i_0 \ldots i_{p + 1}}} = 0$
となるものが存在する。《位相空間》補題
\ref{topology-lemma-lift-covering-of-quasi-compact-hausdorff-subset} により，
$Z \subset \bigcup U'_i$ かつ
$U'_{i_0 \ldots i_{p + 1}} \subset W_{i_0 \ldots i_{p + 1}}$ となる
$U'_i \subset U_i$ を見つけられる。このとき
$s' = (s'_{i_0 \ldots i_p})$，ただし
$s'_{i_0 \ldots i_p} = s_{i_0 \ldots i_p}|_{U'_{i_0 \ldots i_p}}$ は，
$Z$ の開近傍の開被覆 $U' = \bigcup U'_i$ に対する $\mathcal{I}$ の
コサイクルである。$\mathcal{I}$ の高次 {\v C}ech コホモロジー群は零である
（補題 \ref{cohomology-lemma-injective-trivial-cech}）から，$s'$ は余境界である。
したがって，開被覆 $Z = \bigcup Z \cap U'_i$ に対する {\v C}ech 複体における
$\xi$ の像は余境界であり，証明が完了する。
\end{proof}






\section{基底変換写像}
\label{cohomology-section-base-change-map}

\noindent
いくつかの場合に基底変換写像を構成する方法が必要となる。導来引き戻しについて
まだ論じていないので，ここでは環付き空間の平坦射による基底変換の場合だけを
扱う。結果を述べる前に，導来圏上の平坦引き戻しについて説明しよう。
$g : X \to Y$ を環付き空間の平坦射とする。《加群》補題
\ref{modules-lemma-pullback-flat} により，関手
$g^* : \textit{Mod}(\mathcal{O}_Y) \to
\textit{Mod}(\mathcal{O}_X)$ は完全である。したがって，導来関手
$$
g^* : D^{+}(Y) \to D^{+}(X)
$$
をもち，これは $D^{+}(Y)$ の与えられた対象の代表をそのまま引き戻すことで
計算される（《導来圏》補題
\ref{derived-lemma-right-derived-exact-functor} を参照）。したがって，上で
示したように，この関手を $Lg^*$ ではなく $g^*$ と記す。

\begin{lemma}
\label{cohomology-lemma-base-change-map-flat-case}
環付き空間の可換図式
$$
\xymatrix{
X' \ar[r]_{g'} \ar[d]_{f'} &
X \ar[d]^f \\
S' \ar[r]^g &
S
}
$$
と，下に有界な $\mathcal{O}_X$-加群の複体
$\mathcal{F}^\bullet$ が与えられているとする。$g$ と $g'$ がともに
平坦であると仮定する。このとき $D^{+}(S')$ に標準基底変換写像
$$
g^*Rf_*\mathcal{F}^\bullet
\longrightarrow
R(f')_*(g')^*\mathcal{F}^\bullet
$$
が存在する。
\end{lemma}

\begin{proof}
入射分解 $\mathcal{F}^\bullet \to \mathcal{I}^\bullet$ と
$(g')^*\mathcal{F}^\bullet \to \mathcal{J}^\bullet$ を選ぶ。
補題 \ref{cohomology-lemma-pushforward-injective-flat} により，
$(g')_*\mathcal{J}^\bullet$ は $R(g')_*(g')^*\mathcal{F}^\bullet$ を表す
入射対象の複体である。したがって，《導来圏》補題
\ref{derived-lemma-morphisms-lift} および
\ref{derived-lemma-morphisms-equal-up-to-homotopy} により，図式
$$
\xymatrix{
(g')_*(g')^*\mathcal{F}^\bullet \ar[r] &
(g')_*\mathcal{J}^\bullet \\
\mathcal{F}^\bullet \ar[u]^{adjunction} \ar[r] &
\mathcal{I}^\bullet \ar[u]_\beta
}
$$
における矢印 $\beta$ が存在し，ホモトピーを除いて一意である。
$S$ への順像を取ると
$$
f_*\beta :
f_*\mathcal{I}^\bullet
\longrightarrow
f_*(g')_*\mathcal{J}^\bullet
=
g_*(f')_*\mathcal{J}^\bullet
$$
を得る。$g^*$ と $g_*$ の随伴により，複体の写像
$g^*f_*\mathcal{I}^\bullet \to (f')_*\mathcal{J}^\bullet$ を得る。
全過程での唯一の選択は写像 $\beta$ の選択であり，すべて複体のレベルで
行ったので，この写像はホモトピーを除いて一意であることに注意せよ。
\end{proof}

\begin{remark}
\label{cohomology-remark-correct-version-base-change-map}
基底変換写像の「正しい」版は写像
$$
Lg^* Rf_* \mathcal{F}^\bullet
\longrightarrow
R(f')_* L(g')^*\mathcal{F}^\bullet.
$$
である。この写像の構成には非有界複体が関わる。注意
\ref{cohomology-remark-base-change} を参照せよ。
\end{remark}





\section{位相空間における固有基底変換}
\label{cohomology-section-proper-base-change}

\noindent
本節では，位相空間における固有基底変換定理の非常に一般的な版を証明する。
この定理は，高次順像 $R^pf_*$ の茎をファイバー上で計算できることを述べる。

\begin{lemma}
\label{cohomology-lemma-proper-base-change}
$f : (X, \mathcal{O}_X) \to (Y, \mathcal{O}_Y)$ を環付き空間の射とし，
$y \in Y$ とする。次を仮定する。
\begin{enumerate}
\item $f$ は閉写像である。
\item $f$ は分離的である。
\item $f^{-1}(y)$ は準コンパクトである。
\end{enumerate}
このとき，$D^+(\mathcal{O}_X)$ の $E$ に対して，
$D^+(\mathcal{O}_{Y, y})$ において
$(Rf_*E)_y = R\Gamma(f^{-1}(y), E|_{f^{-1}(y)})$ である。
\end{lemma}

\begin{proof}
補題 \ref{cohomology-lemma-base-change-map-flat-case} の基底変換写像から標準写像
$(Rf_*E)_y \to R\Gamma(f^{-1}(y), E|_{f^{-1}(y)})$ を得る。
これが同型であることを示すため，$E$ を下に有界な入射対象の複体
$\mathcal{I}^\bullet$ で表す。$Z = f^{-1}(\{y\})$ と置く。
補題 \ref{cohomology-lemma-cohomology-of-closed} の仮定は満たされる
（《位相空間》補題 \ref{topology-lemma-separated} を参照）。
したがって制限 $\mathcal{I}^n|_Z$ は $\Gamma(Z, -)$ に関して非輪状である。
ゆえに $R\Gamma(Z, E|_Z)$ は複体
$\Gamma(Z, \mathcal{I}^\bullet|_Z)$ で表される
（《導来圏》補題 \ref{derived-lemma-leray-acyclicity} を参照）。
言い換えれば，写像
$$
\colim_V \mathcal{I}^\bullet(f^{-1}(V))
\longrightarrow
\Gamma(Z, \mathcal{I}^\bullet|_Z)
$$
が同型であることを示さなければならない。補題
\ref{cohomology-lemma-cohomology-of-closed} を用いると，$Z = f^{-1}(\{y\})$ の開近傍
$f^{-1}(V)$ の族が，すべての開近傍の系において共終であることを示せば十分で
ある。$f^{-1}(\{y\}) \subset U$ が開近傍ならば，$f$ は閉写像なので，
$V = Y \setminus f(X \setminus U)$ は $y$ の開近傍であり，
$f^{-1}(V) \subset U$ となる。これで補題が示された。
\end{proof}

\begin{theorem}[固有基底変換]
\label{cohomology-theorem-proper-base-change}
\begin{reference}
\cite[Expose V bis, 4.1.1]{SGA4}
\end{reference}
位相空間の Cartesian 図式
$$
\xymatrix{
X' = Y' \times_Y X \ar[d]_{f'} \ar[r]_-{g'} & X \ar[d]^f \\
Y' \ar[r]^g & Y
}
$$
を考える。
$f$ が固有であると仮定し，$E$ を $D^+(X)$ の対象とする。このとき，補題
\ref{cohomology-lemma-base-change-map-flat-case} の基底変換写像
$$
g^{-1}Rf_*E \longrightarrow Rf'_*(g')^{-1}E
$$
は $D^+(Y')$ における同型である。
\end{theorem}

\begin{proof}
$y' \in Y'$ を像が $y \in Y$ である点とする。基底変換写像が $y'$ における
茎上で同型を誘導することを示せば十分である。$f$ が固有なので，$f'$ も固有で
あり，$f$ と $f'$ のファイバーは準コンパクトで，$f$ と $f'$ は閉写像である
（《位相空間》定理 \ref{topology-theorem-characterize-proper} および補題
\ref{topology-lemma-base-change-separated} を参照）。したがって補題
\ref{cohomology-lemma-proper-base-change} を二度適用すると
$$
(Rf'_*(g')^{-1}E)_{y'} = R\Gamma((f')^{-1}(y'), (g')^{-1}E|_{(f')^{-1}(y')})
$$
および
$$
(Rf_*E)_y = R\Gamma(f^{-1}(y), E|_{f^{-1}(y)})
$$
を得る。誘導されるファイバーの写像
$(f')^{-1}(y') \to f^{-1}(y)$ は位相空間の同相写像であり，
$E|_{f^{-1}(y)}$ の引き戻しは $(g')^{-1}E|_{(f')^{-1}(y')}$ である。
望む結果が従う。
\end{proof}

\begin{lemma}[集合の層に対する固有基底変換]
\label{cohomology-lemma-proper-base-change-sheaves-of-sets}
位相空間の Cartesian 図式
$$
\xymatrix{
X' \ar[d]_{f'} \ar[r]_-{g'} & X \ar[d]^f \\
Y' \ar[r]^g & Y
}
$$
を考える。
$f$ が固有であると仮定する。このとき，$X$ 上の任意の集合の層
$\mathcal{F}$ に対して
$g^{-1}f_*\mathcal{F} = f'_*(g')^{-1}\mathcal{F}$
である。
\end{lemma}

\begin{proof}
定理 \ref{cohomology-theorem-proper-base-change} の証明とまったく同様に議論すると，
$(f_*\mathcal{F})_y = \Gamma(X_y, \mathcal{F}|_{X_y})$ を示せば十分であると
分かる。次に補題 \ref{cohomology-lemma-proper-base-change} と同様に議論して，これを
集合の層に対する補題 \ref{cohomology-lemma-cohomology-of-closed} の $p = 0$ の場合に
帰着させる。補題 \ref{cohomology-lemma-cohomology-of-closed} の証明の前半は集合の層にも
通用し，これで証明が完了する。いくつかの詳細は省略する。
\end{proof}



\section{コホモロジーと余極限}
\label{cohomology-section-limits}

\noindent
$X$ を環付き空間とし，$(\mathcal{F}_i, \varphi_{ii'})$ を有向集合 $I$ 上の
$\mathcal{O}_X$-加群の層の系とする。《圏論》第
\ref{categories-section-posets-limits} 節を参照せよ。各 $i$ に対して標準写像
$\mathcal{F}_i \to \colim_i \mathcal{F}_i$ があるので，すべての
$p \geq 0$ に対して標準写像
$$
\colim_i H^p(X, \mathcal{F}_i)
\longrightarrow
H^p(X, \colim_i \mathcal{F}_i)
$$
を得る。もちろん，任意の開集合 $U \subset X$ に対しても同様の写像がある。
これらの写像は，$p = 0$ の場合でさえ一般には同型ではない。本節では
《層》補題 \ref{sheaves-lemma-directed-colimits-sections} の結果を一般化する。
《加群》補題 \ref{modules-lemma-finite-presentation-quasi-compact-colimit}
（特別な場合 $\mathcal{G} = \mathcal{O}_X$）も参照せよ。

\begin{lemma}
\label{cohomology-lemma-quasi-separated-cohomology-colimit}
$X$ を環付き空間とする。$X$ の台となる位相空間が次の性質をもつと仮定する。
\begin{enumerate}
\item 準コンパクト開部分集合からなる基底が存在する。
\item 任意の二つの準コンパクト開集合の共通部分は準コンパクトである。
\end{enumerate}
このとき，$\mathcal{O}_X$-加群の層の任意の有向系
$(\mathcal{F}_i, \varphi_{ii'})$ と任意の準コンパクト開集合
$U \subset X$ に対して，標準写像
$$
\colim_i H^q(U, \mathcal{F}_i)
\longrightarrow
H^q(U, \colim_i \mathcal{F}_i)
$$
はすべての $q \geq 0$ に対して同型である。
\end{lemma}

\begin{proof}
この証明では，すべての準コンパクト開集合 $U \subset X$ に対して同時に
議論することが重要である。《層》補題
\ref{sheaves-lemma-directed-colimits-sections} と《位相空間》補題
\ref{topology-lemma-topology-quasi-separated-scheme} を合わせると，
$q = 0$ と任意の準コンパクト開集合 $U \subset X$ に対して結果が成り立つ。
すべての $q \leq q_0$ に対して結果を証明済みであると仮定し，
$q = q_0 + 1$ の場合を証明しよう。

\medskip\noindent
有向系に関する規約により，添字集合 $I$ は有向であり，$I$ 上の任意の
$\mathcal{O}_X$-加群の系 $(\mathcal{F}_i, \varphi_{ii'})$ は有向である。
《入射対象》補題 \ref{injectives-lemma-sheaves-modules-space} により，
$\mathcal{O}_X$-加群の圏は関手的な入射埋め込みをもつ。したがって，任意の系
$(\mathcal{F}_i, \varphi_{ii'})$ に対し，各 $\mathcal{I}_i$ が入射的
$\mathcal{O}_X$-加群である系 $(\mathcal{I}_i, \varphi_{ii'})$ と，単射な
$\mathcal{O}_X$-加群写像 $\mathcal{F}_i \to \mathcal{I}_i$ で与えられる
系の射が存在する。その余核を $\mathcal{Q}_i$ と書けば，短完全列
$$
0 \to
\mathcal{F}_i \to
\mathcal{I}_i \to
\mathcal{Q}_i \to 0.
$$
を得る。列
$$
0 \to
\colim_i \mathcal{F}_i \to
\colim_i \mathcal{I}_i \to
\colim_i \mathcal{Q}_i \to 0.
$$
も $\mathcal{O}_X$-加群の短完全列であると主張する。これは茎上で確認してよい。
《層》第 \ref{sheaves-section-limits-presheaves} 節および第
\ref{sheaves-section-limits-sheaves} 節により，茎を取る操作は余極限と可換する。
アーベル群の短完全列の有向余極限は短完全であるから
（《代数》補題 \ref{algebra-lemma-directed-colimit-exact} を参照），結果が従う。
さらに，すべての準コンパクト開集合 $U \subset X$ とすべての $q \geq 1$ に
対して $H^q(U, \colim_i \mathcal{I}_i) = 0$
であると主張する。この主張をひとまず認めて，図式
$$
\xymatrix{
\colim_i H^{q_0}(U, \mathcal{I}_i) \ar[d] \ar[r] &
\colim_i H^{q_0}(U, \mathcal{Q}_i) \ar[d] \ar[r] &
\colim_i H^{q_0 + 1}(U, \mathcal{F}_i) \ar[d] \ar[r] &
0 \ar[d] \\
H^{q_0}(U, \colim_i \mathcal{I}_i) \ar[r] &
H^{q_0}(U, \colim_i \mathcal{Q}_i) \ar[r] &
H^{q_0 + 1}(U, \colim_i \mathcal{F}_i) \ar[r] &
0
}
$$
を考える。右下隅の零は上の主張から，右上隅の零は層
$\mathcal{I}_i$ が入射的であることから得られる。《代数》補題
\ref{algebra-lemma-directed-colimit-exact} を適用すると上段は完全である。
したがって蛇の補題により，$q = q_0 + 1$ に対する結果が従う。

\medskip\noindent
残るのは主張を示すことである。補題 \ref{cohomology-lemma-cech-vanish-basis} を用いる。
$\mathcal{B}$ を $X$ のすべての準コンパクト開部分集合からなる族とする。
仮定により，これは $X$ の位相の基底である。$\text{Cov}$ を，
$U$ と各 $U_j$ が $X$ の準コンパクト開集合であるような有限開被覆
$\mathcal{U} : U = \bigcup_{j = 1, \ldots, m} U_j$ の族とする。
$q = 0$ に対する結果により，$\mathcal{U} \in \text{Cov}$ に対して
$$
\check{\mathcal{C}}^\bullet(\mathcal{U}, \colim_i \mathcal{I}_i)
=
\colim_i \check{\mathcal{C}}^\bullet(\mathcal{U}, \mathcal{I}_i)
$$
である。実際，すべての多重共通部分 $U_{j_0 \ldots j_p}$ は
準コンパクトである。補題 \ref{cohomology-lemma-injective-trivial-cech} により，
{\v C}ech 複体の余極限に現れる各複体は次数 $\geq 1$ で非輪状である。
したがって《代数》補題 \ref{algebra-lemma-directed-colimit-exact} により，
{\v C}ech 複体
$\check{\mathcal{C}}^\bullet(\mathcal{U}, \colim_i \mathcal{I}_i)$ も
次数 $\geq 1$ で非輪状である。言い換えれば，
$\check{H}^p(\mathcal{U},  \colim_i \mathcal{I}_i) = 0$
がすべての $p \geq 1$ に対して成り立つ。ゆえに補題
\ref{cohomology-lemma-cech-vanish-basis} の仮定が満たされ，主張が従う。
\end{proof}

\begin{lemma}
\label{cohomology-lemma-higher-direct-image-colimit}
$f : X \to Y$ を位相空間の連続写像とし，
$(\mathcal{F}_i, \varphi_{ii'})$ を $X$ 上のアーベル群の層の系とする。
$\mathcal{F} = \colim \mathcal{F}_i$ と置き，$p \geq 0$ を整数とする。
$H^p(f^{-1}(V), \mathcal{F}) = \colim H^p(f^{-1}(V), \mathcal{F}_i)$
となる開集合 $V \subset Y$ の族が $Y$ の位相の基底であると仮定する。
このとき $R^pf_*\mathcal{F} = \colim R^pf_*\mathcal{F}_i$ である。
\end{lemma}

\begin{proof}
$R^pf_*\mathcal{F}$ は，$V$ を $H^p(f^{-1}(V), \mathcal{F})$ へ送る前層
$\mathcal{G}$ の層化であった（補題
\ref{cohomology-lemma-describe-higher-direct-images} を参照）。同様に，
$R^pf_*\mathcal{F}_i$ は，$V$ を
$H^p(f^{-1}(V), \mathcal{F}_i)$ へ送る前層 $\mathcal{G}_i$ の層化である。
層化は層から前層への包含の左随伴であった
（《層》第 \ref{sheaves-section-sheafification} 節を参照）。したがって層化は
余極限と可換する（《圏論》補題 \ref{categories-lemma-adjoint-exact} を参照）。
ゆえに，前層の圏で余極限を取った前層の写像
$$
\colim \mathcal{G}_i \longrightarrow \mathcal{G}
$$
が層化上で同型を誘導することを示せば十分である。そのためには，前層
$\mathcal{G}$ と $\colim \mathcal{G}_i$ が $Y$ の位相のある基底上で一致することを
示せばよい。実際，この場合，基底の元上の前層の値から直接計算できる両者の層化の
茎が一致する。必要な一致はまさに補題の仮定である。
\end{proof}

\noindent
次に，《層》補題 \ref{sheaves-lemma-descend-opens} のコホモロジーに対する
類似を定式化する。《位相空間》補題
\ref{topology-lemma-directed-inverse-limit-spectral-spaces} のように，
$X$ をスペクトル遷移射をもつ図式におけるスペクトル空間 $X_i$ の
余フィルター極限として書かれたスペクトル空間とする。次が与えられていると仮定する。
\begin{enumerate}
\item 各 $i \in \Ob(\mathcal{I})$ に対する $X_i$ 上のアーベル群の層
$\mathcal{F}_i$。
\item $a : j \to i$ に対するアーベル群の層の $f_a$-写像
$\varphi_a : \mathcal{F}_i \to \mathcal{F}_j$
（《層》定義 \ref{sheaves-definition-f-map} を参照）。
\end{enumerate}
$c = a \circ b$ のときは常に
$\varphi_c = \varphi_b \circ \varphi_a$ であるとする。$X$ 上で
$\mathcal{F} = \colim p_i^{-1}\mathcal{F}_i$ と置く。

\begin{lemma}
\label{cohomology-lemma-colimit}
上で論じた状況において，$i \in \Ob(\mathcal{I})$ とし，
$U_i \subset X_i$ を準コンパクト開集合とする。このとき
$$
\colim_{a : j \to i} H^p(f_a^{-1}(U_i), \mathcal{F}_j) =
H^p(p_i^{-1}(U_i), \mathcal{F})
$$
がすべての $p \geq 0$ に対して成り立つ。特に
$H^p(X, \mathcal{F}) = \colim H^p(X_i, \mathcal{F}_i)$ である。
\end{lemma}

\begin{proof}
$p = 0$ の場合は《層》補題 \ref{sheaves-lemma-descend-opens} である。

\medskip\noindent
この段落では，$\mathcal{G}_i$ が入射的なアーベル群の層で
$\gamma_i$ が単射であるような系の写像
$(\gamma_i) : (\mathcal{F}_i, \varphi_a) \to (\mathcal{G}_i, \psi_a)$ を
見つけられることを示す。各 $i$ に対して単射
$\mathcal{F}_i \to \mathcal{I}_i$ を選ぶ。ここで $\mathcal{I}_i$ は
$X_i$ 上の入射的なアーベル群の層である。このとき写像の族
$$
\gamma_i :
\mathcal{F}_i
\longrightarrow
\prod\nolimits_{b : k \to i} f_{b, *}\mathcal{I}_k = \mathcal{G}_i
$$
を考えられる。ここで成分写像は
$f_b^{-1}\mathcal{F}_i \to \mathcal{F}_k \to \mathcal{I}_k$ に随伴する
写像である。$\mathcal{I}$ の $a : j \to i$ に対して標準写像
$$
\psi_a : f_a^{-1}\mathcal{G}_i \to \mathcal{G}_j
$$
があり，その $b : k \to j$ に対応する成分は標準写像
$f_b^{-1}f_{a \circ b, *}\mathcal{I}_k \to f_{b, *}\mathcal{I}_k$ である。
したがって，アーベル群の層の系の単射
$\{\gamma_i\} : \{\mathcal{F}_i, \varphi_a) \to (\mathcal{G}_i, \psi_a)$
を得る。$\mathcal{G}_i$ は $X_i$ 上の入射的なアーベル群の層であることに
注意せよ。補題 \ref{cohomology-lemma-pushforward-injective-flat} および《ホモロジー》
補題 \ref{homology-lemma-product-injectives} を参照せよ。これで構成が完了する。

\medskip\noindent
補題 \ref{cohomology-lemma-quasi-separated-cohomology-colimit} の証明とまったく同様に
議論すると，$p > 0$ に対して
$H^p(X, \colim f_i^{-1}\mathcal{G}_i) = 0$ を証明すれば十分であると分かる。

\medskip\noindent
$\mathcal{G} = \colim f_i^{-1}\mathcal{G}_i$ と置く。$X$ の任意の
準コンパクト開集合上で $\mathcal{G}$ のコホモロジーが消滅することを示すには，
$X$ の準コンパクト開集合を有限個の準コンパクト開集合で被覆する任意の開被覆
$\mathcal{U}$ に対して，$\mathcal{G}$ の {\v C}ech コホモロジーが零であることを
示せば十分である。補題 \ref{cohomology-lemma-cech-vanish-basis} を参照せよ。
《位相空間》補題 \ref{topology-lemma-descend-opens} により，このような被覆は，
ある $i$ に対する空間 $X_i$ 上の同様の被覆 $\mathcal{U}_i$ の
$p_i$ による逆像である。$p = 0$ の場合により
$$
\check{\mathcal{C}}^\bullet(\mathcal{U}, \mathcal{G}) =
\colim_{a : j \to i}
\check{\mathcal{C}}^\bullet(f_a^{-1}(\mathcal{U}_i), \mathcal{G}_j)
$$
である。右辺は複体のフィルター余極限であり，補題
\ref{cohomology-lemma-injective-trivial-cech} により，その各複体は正次数で非輪状である。
したがって《代数》補題 \ref{algebra-lemma-directed-colimit-exact} により結論を得る。
\end{proof}









\section{ネーター位相空間上の消滅}
\label{cohomology-section-vanishing-Noetherian}

\noindent
目標は Grothendieck の定理，すなわち命題
\ref{cohomology-proposition-vanishing-Noetherian} を証明することである。\cite{Tohoku} を参照せよ。

\begin{lemma}
\label{cohomology-lemma-cohomology-and-closed-immersions}
$i : Z \to X$ を位相空間の閉埋め込みとする。$Z$ 上の任意のアーベル群の層
$\mathcal{F}$ に対して
$H^p(Z, \mathcal{F}) = H^p(X, i_*\mathcal{F})$ である。
\end{lemma}

\begin{proof}
$i_*$ は完全であり（《加群》補題
\ref{modules-lemma-i-star-exact} を参照），したがって関手として
$R^pi_* = 0$ であるから成り立つ（《導来圏》補題
\ref{derived-lemma-right-derived-exact-functor}）。ゆえに補題
\ref{cohomology-lemma-apply-Leray} を適用できる。
\end{proof}

\begin{lemma}
\label{cohomology-lemma-irreducible-constant-cohomology-zero}
$X$ を既約位相空間とする。このとき，任意のアーベル群 $A$ とすべての
$p > 0$ に対して $H^p(X, \underline{A}) = 0$ である。
\end{lemma}

\begin{proof}
$\underline{A}$ は《層》定義 \ref{sheaves-definition-constant-sheaf} で
定義した定数層であった。$X$ は既約なので，任意の空でない開集合 $U$ は既約で，
特に連結である。したがって，空でない開集合 $U \subset X$ に対して
$\underline{A}(U) = A$ である。また $\underline{A}(\emptyset) = 0$ である。
ゆえに $\underline{A}$ は $X$ 上のフラスクなアーベル群の層である。
消滅は補題 \ref{cohomology-lemma-flasque-acyclic} から従う。
\end{proof}

\begin{lemma}
\label{cohomology-lemma-subsheaf-of-constant-sheaf}
\begin{reference}
\cite[Page 168]{Tohoku}.
\end{reference}
任意の二つの準コンパクト開集合の共通部分が準コンパクトであるような位相空間
$X$ を考える。
$\mathcal{F} \subset \underline{\mathbf{Z}}$
を，準コンパクト開集合上の有限個の切断で生成される部分層とする。
このときアーベル部分層による有限フィルトレーション
$$
(0) = \mathcal{F}_0 \subset \mathcal{F}_1 \subset \ldots \subset
\mathcal{F}_n = \mathcal{F}
$$
で，任意の $0 < i \leq n$ に対して短完全列
$$
0 \to j'_!\underline{\mathbf{Z}}_V \to j_!\underline{\mathbf{Z}}_U \to
\mathcal{F}_i/\mathcal{F}_{i - 1} \to 0
$$
が存在するものがある。ここで $j : U \to X$ と $j' : V \to X$ は
準コンパクト開集合の $X$ への包含である。
\end{lemma}

\begin{proof}
$\mathcal{F}$ が準コンパクト開集合 $U_1, \ldots, U_t$ 上の切断
$s_1, \ldots, s_t$ で生成されるとする。$U_i$ は準コンパクトで，
$s_i$ は $\mathbf{Z}$ への局所定数関数だから，必要なら $U_i$ を開かつ閉な
部分集合への有限分解の各部分で置き換えることにより，$s_i$ が定数切断であると
仮定してよい。$n_i \in \mathbf{Z}$ に対して $s_i = n_i$ とする。
$n_i = 0$ ならば，もちろん $(U_i, n_i)$ をリストから除ける。必要なら符号を
反転して $n_i > 0$ と仮定してもよい。次に，任意の部分集合
$I \subset \{1, \ldots, t\}$ に対して
$\bigcap_{i \in I} U_i$ と $\gcd(n_i, i \in I)$ をリストに加えてよい。
こうすると，リスト $(U_1, n_1), \ldots, (U_t, n_t)$ は次の性質を満たす。
$x \in X$ に対して
$I_x = \{i \in \{1, \ldots, t\} \mid x \in U_i\}$ と置く。このとき，
$\gcd(n_i, i \in I_x)$ はある $i \in I_x$ に対する $n_i$ として達成される。

\medskip\noindent
フィルトレーションとして $\mathcal{F}_0 = (0)$ を取り，
$\mathcal{F}_n$ を，$n_i \leq n$ となる $i$ に対する $U_i$ 上の切断 $n_i$ で
生成されるものとする。$n \gg 0$ に対して
$\mathcal{F}_n = \mathcal{F}$ であることは明らかである。さらに，商
$\mathcal{F}_n/\mathcal{F}_{n - 1}$ は
$U = \bigcup_{n_i \leq n} U_i$ 上の切断 $n$ で生成され，写像
$j_!\underline{\mathbf{Z}}_U \to \mathcal{F}_n/\mathcal{F}_{n - 1}$ の核は
$V = \bigcup_{n_i \leq n - 1} U_i$ 上の切断 $n$ で生成される。したがって，
補題の主張にある短完全列を得る。
\end{proof}

\begin{lemma}
\label{cohomology-lemma-vanishing-generated-one-section}
\begin{reference}
これは \cite[Proposition 3.6.1]{Tohoku} の特別な場合である。
\end{reference}
$X$ を位相空間，$d \geq 0$ を整数とする。次を仮定する。
\begin{enumerate}
\item $X$ は準コンパクトである。
\item 準コンパクト開集合は $X$ の基底をなす。
\item 二つの準コンパクト開集合の共通部分は準コンパクトである。
\item 任意の準コンパクト開集合 $j : U \to X$ とすべての $p > d$ に対して
$H^p(X, j_!\underline{\mathbf{Z}}_U) = 0$ である。
\end{enumerate}
このとき，$X$ 上の任意のアーベル群の層 $\mathcal{F}$ とすべての
$p > d$ に対して $H^p(X, \mathcal{F}) = 0$ である。
\end{lemma}

\begin{proof}
$U$ が $X$ の準コンパクト開集合をわたるものとして
$S = \coprod_{U \subset X} \mathcal{F}(U)$ と置く。任意の有限部分集合
$A = \{s_1, \ldots, s_n\} \subset S$ に対して，$\mathcal{F}_A$ をすべての
$s_i$ で生成される $\mathcal{F}$ の部分層とする
（《加群》定義 \ref{modules-definition-generated-by-local-sections} を参照）。
$A \subset A'$ ならば
$\mathcal{F}_A \subset \mathcal{F}_{A'}$ であることに注意せよ。したがって
$\{\mathcal{F}_A\}$ は，$S$ の有限部分集合からなる有向半順序集合上の系をなす。
《加群》補題 \ref{modules-lemma-generated-by-local-sections-stalk} により，茎を
見れば
$$
\colim_A \mathcal{F}_A = \mathcal{F}
$$
であることが明らかである。補題
\ref{cohomology-lemma-quasi-separated-cohomology-colimit} により
$$
H^p(X, \mathcal{F}) =
\colim_A H^p(X, \mathcal{F}_A)
$$
である。したがって，アーベル群の層 $\mathcal{F}_A$ に対して消滅を証明すれば
十分である。言い換えれば，$\mathcal{F}$ が $X$ の準コンパクト開集合上の
有限個の局所切断で生成される場合に結果を証明すれば十分である。

\medskip\noindent
$\mathcal{F}$ が局所切断 $s_1, \ldots, s_n$ で生成されるとする。
$\mathcal{F}' \subset \mathcal{F}$ を $s_1, \ldots, s_{n - 1}$ で生成される
部分層とする。このとき短完全列
$$
0 \to \mathcal{F}' \to \mathcal{F} \to \mathcal{F}/\mathcal{F}' \to 0
$$
がある。コホモロジー長完全列から，$n$ 個未満の局所切断で生成される
アーベル群の層 $\mathcal{F}'$ と $\mathcal{F}/\mathcal{F}'$ に対して消滅を
証明すれば十分であると分かる。したがって，高々一つの局所切断で生成される層に
対して消滅を証明すれば十分である。これらの層は，$U$ が $X$ の準コンパクト
開集合であるような層 $j_!\underline{\mathbf{Z}}_U$ の商にほかならない。

\medskip\noindent
ここで，$U$ が $X$ の準コンパクト開集合であるような短完全列
$$
0 \to \mathcal{K} \to j_!\underline{\mathbf{Z}}_U \to \mathcal{F} \to 0
$$
があると仮定する。$q \geq d + 1$ に対して
$H^q(X, \mathcal{K})$ が零であることを示せば十分である。上と同様に，
$\mathcal{K}$ を，準コンパクト開集合上の有限個の切断で生成される部分層
$\mathcal{K}'$ のフィルター余極限として書ける。このとき $\mathcal{F}$ は層
$j_!\underline{\mathbf{Z}}_U/\mathcal{K}'$ のフィルター余極限である。このように
して，$\mathcal{K}$ が準コンパクト開集合上の有限個の切断で生成される場合に
帰着する。$\mathcal{K}$ は $\underline{\mathbf{Z}}_X$ の部分層であることに
注意せよ。したがって補題 \ref{cohomology-lemma-subsheaf-of-constant-sheaf} により，
$\mathcal{K}$ には有限フィルトレーションが存在し，その逐次商
$\mathcal{Q}$ は短完全列
$$
0 \to j''_!\underline{\mathbf{Z}}_W \to
j'_!\underline{\mathbf{Z}}_V \to \mathcal{Q} \to 0
$$
に入る。ここで $j'' : W \to X$ と $j' : V \to X$ は準コンパクト開集合の
包含である。したがって，$p > d$ に対する
$H^p(X, \mathcal{Q})$ の消滅は，$j''_!\underline{\mathbf{Z}}_W$ と
$j'_!\underline{\mathbf{Z}}_V$ のコホモロジー群の消滅に関する補題の仮定から
従う。$\mathcal{K}$ に戻ると，コホモロジー長完全列を用いる帰納法により，
望む消滅が $\mathcal{K}$ に対しても従う。
\end{proof}

\begin{example}
\label{cohomology-example-datta}
$X = \mathbf{N}$ に，開集合が $\emptyset$，$X$，および $n \geq 1$ に対する
$U_n = \{i \mid i \leq n\}$ であるような位相を入れる。$X$ 上のアーベル群の層
$\mathcal{F}$ は，アーベル群の逆系 $A_n = \mathcal{F}(U_n)$ と同じものであり，
$\Gamma(X, \mathcal{F}) = \lim A_n$ である。逆極限関手は逆系の圏上の完全関手では
ないので，$H^1$ が零でないアーベル群の層が存在することが分かる。最後に，
$j : U = U_n \to X$ が包含ならば，$p \geq 1$ に対して
$H^p(X, j_!\mathbf{Z}_U) = 0$ であることを読者は確認できる。したがって $X$ は，
$d = 0$ に対して補題 \ref{cohomology-lemma-vanishing-generated-one-section} の条件
(2)，(3)，(4) を満たすが結論を満たさない空間の例である。
\end{example}

\begin{lemma}
\label{cohomology-lemma-subsheaf-irreducible}
$X$ を既約位相空間とし，
$\mathcal{H} \subset \underline{\mathbf{Z}}$ を定数層のアーベル部分層とする。
このとき，ある $d \in \mathbf{Z}$ に対して
$\mathcal{H}|_U = \underline{d\mathbf{Z}}_U$ となる空でない開集合
$U \subset X$ が存在する。
\end{lemma}

\begin{proof}
$X$ の任意の空でない開集合 $V$ に対して
$\underline{\mathbf{Z}}(V) = \mathbf{Z}$ であった
（補題 \ref{cohomology-lemma-irreducible-constant-cohomology-zero} の証明を参照）。
$\mathcal{H} = 0$ ならば，補題は $d = 0$ として成り立つ。
$\mathcal{H} \not = 0$ ならば，
$\mathcal{H}(U) \not = 0$ となる空でない開集合 $U \subset X$ が存在する。
ある $n \geq 1$ に対して $\mathcal{H}(U) = n\mathbf{Z}$ とする。このとき
$\underline{n\mathbf{Z}}_U
\subset \mathcal{H}|_U \subset
\underline{\mathbf{Z}}_U$ である。第一の包含が真の包含ならば，
空でない $U' \subset U$ と整数 $1 \leq n' < n$ で，
$\underline{n'\mathbf{Z}}_{U'}
\subset \mathcal{H}|_{U'} \subset
\underline{\mathbf{Z}}_{U'}$
となるものを見つけられる。
この過程は有限回で停止しなければならない。したがって補題を得る。
\end{proof}

\begin{proposition}[Grothendieck]
\label{cohomology-proposition-vanishing-Noetherian}
\begin{reference}
\cite[Theorem 3.6.5]{Tohoku}.
\end{reference}
$X$ をネーター位相空間とする。$\dim(X) \leq d$ ならば，$X$ 上の任意の
アーベル群の層 $\mathcal{F}$ とすべての $p > d$ に対して
$H^p(X, \mathcal{F}) = 0$ である。
\end{proposition}

\begin{proof}
この補題を $d$ に関する帰納法で証明する。そこで $d$ を固定し，次元 $< d$ の
すべてのネーター位相空間に対して補題が成り立つと仮定する。

\medskip\noindent
$\mathcal{F}$ を $X$ 上のアーベル群の層とする。$U \subset X$ を開集合とし，
閉補集合を $Z \subset X$ と書く。包含写像を
$j : U \to X$ および $i : Z \to X$ と書く。このとき短完全列
$$
0 \to j_{!}j^*\mathcal{F} \to \mathcal{F} \to i_*i^*\mathcal{F} \to 0
$$
がある。《加群》補題 \ref{modules-lemma-canonical-exact-sequence} を参照せよ。
$j_!j^*\mathcal{F}$ は $U$ の位相的閉包 $Z'$ 上に台をもつことに注意せよ。
すなわち，これは $Z'$ 上のあるアーベル群の層 $\mathcal{F}'$ に対する
$i'_*\mathcal{F}'$ の形である。ここで $i' : Z' \to X$ は包含である。

\medskip\noindent
これを用いて $X$ が既約である場合に帰着できる。実際，《位相空間》補題
\ref{topology-lemma-Noetherian} により，$X$ は有限個の既約成分をもつ。
$X$ が複数の既約成分をもつならば，$Z \subset X$ を $X$ の一つの既約成分とし，
$U = X \setminus Z$ と置く。上の議論とコホモロジー長完全列により，
$p > d$ に対して $H^p(X, i_*i^*\mathcal{F})$ と
$H^p(X, i'_*\mathcal{F}')$ が消滅することを証明すれば十分である。
補題 \ref{cohomology-lemma-cohomology-and-closed-immersions} により，$p > d$ に対して
$H^p(Z, i^*\mathcal{F})$ と $H^p(Z', \mathcal{F}')$ が消滅することを
証明すれば十分である。$Z'$ と $Z$ の既約成分の個数はより少ないので，
実際に $X$ が既約である場合に帰着する。

\medskip\noindent
$d = 0$ かつ $X$ が既約ならば，$X$ は $X$ の唯一の空でない開部分集合である。
したがってすべての層は定数層であり，高次コホモロジー群は消滅する
（例えば補題 \ref{cohomology-lemma-irreducible-constant-cohomology-zero} による）。

\medskip\noindent
$X$ が次元 $d > 0$ の既約空間であるとする。補題
\ref{cohomology-lemma-vanishing-generated-one-section} により，ある開集合
$U \subset X$ に対して
$\mathcal{F} = j_!\underline{\mathbf{Z}}_U$ である場合に帰着する。
この場合，短完全列
$$
0 \to j_!(\underline{\mathbf{Z}}_U) \to
\underline{\mathbf{Z}}_X \to i_*\underline{\mathbf{Z}}_Z \to 0
$$
を考える。ここで $Z = X \setminus U$ である。補題
\ref{cohomology-lemma-irreducible-constant-cohomology-zero} により，すべての
$p \geq 1$ に対して $H^p(X, \underline{\mathbf{Z}}_X)$ は消滅する。
帰納法により
$H^p(X, i_*\underline{\mathbf{Z}}_Z) = H^p(Z, \underline{\mathbf{Z}}_Z) = 0$
が $p \geq d$ に対して成り立つ。したがってコホモロジー長完全列により結論を得る。
\end{proof}





\section{閉部分集合に台をもつコホモロジー}
\label{cohomology-section-cohomology-support}

\noindent
本節では必要最小限だけを論じる。議論は第
\ref{cohomology-section-cohomology-support-bis} 節で続ける。

\medskip\noindent
$X$ を位相空間，$Z \subset X$ を閉部分集合とし，$\mathcal{F}$ を
$X$ 上のアーベル群の層とする。
$$
\Gamma_Z(X, \mathcal{F}) =
\{s \in \mathcal{F}(X) \mid \text{Supp}(s) \subset Z\}
$$
を，台が $Z$ に含まれる切断の部分集合とする。切断の台は《加群》定義
\ref{modules-definition-support} で定義されている。《加群》補題
\ref{modules-lemma-support-section-closed} により，
$\Gamma_Z(X, \mathcal{F})$ は $\Gamma(X, \mathcal{F})$ の部分群である。
同じ補題により，対応 $\mathcal{F} \mapsto \Gamma_Z(X, \mathcal{F})$ は
$\mathcal{F}$ の関手となる。この関手は左完全だが，一般には完全でない。

\medskip\noindent
アーベル群の層の圏には十分多くの入射対象があるので（《入射対象》補題
\ref{injectives-lemma-abelian-sheaves-space}），《導来圏》補題
\ref{derived-lemma-enough-injectives-right-derived} により右導来関手
$$
R\Gamma_Z(X, -) : D^+(X) \longrightarrow D^+(\textit{Ab})
$$
を得る。対象 $K$ における $R\Gamma_Z(X, -)$ の値は，$K$ を入射的な
アーベル群の層の下に有界な複体 $\mathcal{I}^\bullet$ で表し，
$\Gamma_Z(X, \mathcal{I}^\bullet)$ を取ることで計算される
（《導来圏》補題 \ref{derived-lemma-injective-acyclic} を参照）。
アーベル群の層 $\mathcal{F}$ の $Z$ に台をもつコホモロジー群を
$H^q_Z(X, \mathcal{F}) = R^q\Gamma_Z(X, \mathcal{F})$ と定義する。

\medskip\noindent
$\mathcal{I}$ を $X$ 上の入射的なアーベル群の層とし，
$U = X \setminus Z$ とする。このとき制限写像
$\mathcal{I}(X) \to \mathcal{I}(U)$ は全射であり
（補題 \ref{cohomology-lemma-injective-restriction-surjective}），その核は
$\Gamma_Z(X, \mathcal{I})$ である。直ちに，$K \in D^+(X)$ に対して
$D^+(\textit{Ab})$ に識別三角形
$$
R\Gamma_Z(X, K) \to R\Gamma(X, K) \to R\Gamma(U, K) \to R\Gamma_Z(X, K)[1]
$$
があることが従う。その帰結として，任意の $D^+(X)$ の $K$ に対して
コホモロジー長完全列
$$
\ldots \to H^i_Z(X, K) \to H^i(X, K) \to H^i(U, K) \to
H^{i + 1}_Z(X, K) \to \ldots
$$
を得る。

\medskip\noindent
$X$ 上のアーベル群の層 $\mathcal{F}$ に対して，
$\mathcal{H}_Z(\mathcal{F})$ と書かれる
{\it $Z$ に台をもつ切断の部分層}を考えられる。これは規則
$$
\mathcal{H}_Z(\mathcal{F})(U) =
\{s \in \mathcal{F}(U) \mid \text{Supp}(s) \subset U \cap Z\} =
\Gamma_{Z \cap U}(U, \mathcal{F}|_U)
$$
で定義される。《加群》補題 \ref{modules-lemma-i-star-exact} の同値を用いると，
$\mathcal{H}_Z(\mathcal{F})$ を $Z$ 上のアーベル群の層とみなせる
（《加群》注意 \ref{modules-remark-sections-support-in-closed} を参照）。
したがって関手
$$
\textit{Ab}(X) \longrightarrow \textit{Ab}(Z),\quad
\mathcal{F} \longmapsto
\mathcal{H}_Z(\mathcal{F})\text{ を }Z\text{ 上の層とみなしたもの}
$$
を得る。この関手は左完全だが，一般には完全でない。上とまったく同様に
右導来関手
$$
R\mathcal{H}_Z : D^+(X) \longrightarrow D^+(Z)
$$
を得る。
$\mathcal{H}^q_Z(\mathcal{F}) = R^q\mathcal{H}_Z(\mathcal{F})$ と置く。このとき
$\mathcal{H}^0_Z(\mathcal{F}) = \mathcal{H}_Z(\mathcal{F})$ である。

\medskip\noindent
任意のアーベル群の層 $\mathcal{F}$ に対して
$\Gamma_Z(X, \mathcal{F}) = \Gamma(Z, \mathcal{H}_Z(\mathcal{F}))$ であることに
注意せよ。以下の補題 \ref{cohomology-lemma-sections-with-support-acyclic} により，
関手 $\mathcal{H}_Z$ は入射的なアーベル群の層を
$\Gamma(Z, -)$ に関して右非輪状な層へ写す。したがって《導来圏》補題
\ref{derived-lemma-grothendieck-spectral-sequence} により，収束する
Grothendieck スペクトル系列
$$
E_2^{p, q} = H^p(Z, \mathcal{H}^q_Z(K)) \Rightarrow H^{p + q}_Z(X, K)
$$
を得る。これは $D^+(X)$ の $K$ に関して関手的である。

\begin{lemma}
\label{cohomology-lemma-sections-with-support-acyclic}
$i : Z \to X$ を閉部分集合の包含とし，$\mathcal{I}$ を $X$ 上の入射的な
アーベル群の層とする。このとき $\mathcal{H}_Z(\mathcal{I})$ は
$Z$ 上の入射的なアーベル群の層である。
\end{lemma}

\begin{proof}
$\mathcal{H}_Z(-)$ は完全関手 $i_*$ の右随伴なので，《ホモロジー》補題
\ref{homology-lemma-adjoint-preserve-injectives} から従う。《加群》補題
\ref{modules-lemma-i-star-exact} および
\ref{modules-lemma-i-star-right-adjoint} を参照せよ。
\end{proof}




\section{スペクトル空間上のコホモロジー}
\label{cohomology-section-spectral}

\noindent
スペクトル空間のコホモロジーに関する鍵となる結果は補題 \ref{cohomology-lemma-colimit} で
あり，大まかに言えば，《位相空間》定義
\ref{topology-definition-spectral-space} で定義されるスペクトル空間の圏において，
コホモロジーが余フィルター極限と可換することを述べる。これを適用すると，
補題 \ref{cohomology-lemma-cohomology-of-closed} と補題
\ref{cohomology-lemma-proper-base-change} の類似を次のように得られる。

\begin{lemma}
\label{cohomology-lemma-cohomology-of-neighbourhoods-of-closed}
$X$ をスペクトル空間，$\mathcal{F}$ を $X$ 上のアーベル群の層とする。
$E \subset X$ を準コンパクト部分集合とし，$W \subset X$ を $E$ のある点へ
特殊化する $X$ の点全体の集合とする。
\begin{enumerate}
\item $H^p(W, \mathcal{F}|_W) = \colim H^p(U, \mathcal{F})$
である。ここで余極限は $E$ の準コンパクト開近傍をわたる。
\item $H^p(W \setminus E, \mathcal{F}|_{W \setminus E}) =
\colim H^p(U \setminus E, \mathcal{F}|_{U \setminus E})$
である。ただし $E$ は構成可能部分集合とする。
\end{enumerate}
\end{lemma}

\begin{proof}
《位相空間》補題 \ref{topology-lemma-make-spectral-space} により
$W = \lim U$ である。ここで極限は $E$ を含む準コンパクト開集合をわたる。
《位相空間》補題 \ref{topology-lemma-spectral-sub} により，各 $U$ は
スペクトル空間である。したがって補題 \ref{cohomology-lemma-colimit} を適用して
(1) を得る。(2) についても同じ証明が通用するが，《位相空間》補題
\ref{topology-lemma-make-spectral-space-minus} を用いる。
\end{proof}

\begin{lemma}
\label{cohomology-lemma-proper-base-change-spectral}
$f : X \to Y$ をスペクトル空間のスペクトル写像とし，$y \in Y$ とする。
$E \subset Y$ を $y$ へ特殊化する点全体の集合とし，$\mathcal{F}$ を
$X$ 上のアーベル群の層とする。このとき
$(R^pf_*\mathcal{F})_y = H^p(f^{-1}(E), \mathcal{F}|_{f^{-1}(E)})$ である。
\end{lemma}

\begin{proof}
$V$ が $Y$ における $y$ の準コンパクト開近傍をわたるとき
$E = \bigcap V$ であることに注意せよ。したがって
$f^{-1}(E) = \bigcap f^{-1}(V)$ であり，位相空間として
$f^{-1}(E) = \lim f^{-1}(V)$ である。$f$ はスペクトル写像なので，各
$f^{-1}(V)$ もスペクトル空間である
（《位相空間》補題 \ref{topology-lemma-spectral-sub}）。
ゆえに $f^{-1}(E)$ はスペクトル空間であり，補題 \ref{cohomology-lemma-colimit} により
$$
H^p(f^{-1}(E), \mathcal{F}|_{f^{-1}(E)}) =
\colim H^p(f^{-1}(V), \mathcal{F})
$$
である。一方，$y$ における $R^pf_*\mathcal{F}$ の茎は右辺の余極限で与えられる。
\end{proof}

\begin{lemma}
\label{cohomology-lemma-vanishing-for-profinite}
$X$ をプロ有限位相空間とする。このとき，すべての $q > 0$ とすべての
アーベル群の層 $\mathcal{F}$ に対して $H^q(X, \mathcal{F}) = 0$ である。
\end{lemma}

\begin{proof}
$X$ の任意の開被覆は，有限個の開部分への互いに素な分解によって細分できる
（《位相空間》補題
\ref{topology-lemma-profinite-refine-open-covering} を参照）。したがって
$\mathcal{F} \to \mathcal{G}$ が $X$ 上のアーベル群の層の全射ならば，
$\mathcal{F}(X) \to \mathcal{G}(X)$ は全射である。言い換えれば，大域切断関手は
完全関手である。ゆえにその高次導来関手は零である
（《導来圏》補題 \ref{derived-lemma-right-derived-exact-functor} を参照）。
\end{proof}

\noindent
次のコホモロジー消滅に関する結果は Grothendieck の結果
（命題 \ref{cohomology-proposition-vanishing-Noetherian}）を改良するものであり，
\cite{Scheiderer} に見いだせる。

\begin{proposition}
\label{cohomology-proposition-cohomological-dimension-spectral}
\begin{reference}
(1) は \cite{Scheiderer} の主定理である。
\end{reference}
$X$ を Krull 次元 $d$ のスペクトル空間，$\mathcal{F}$ を
$X$ 上のアーベル群の層とする。
\begin{enumerate}
\item $q > d$ に対して $H^q(X, \mathcal{F}) = 0$ である。
\item 任意の準コンパクト開集合 $U \subset X$ に対して
$H^d(X, \mathcal{F}) \to H^d(U, \mathcal{F})$ は全射である。
\item 任意の構成可能閉部分集合 $Z \subset X$ と $q > d$ に対して
$H^q_Z(X, \mathcal{F}) = 0$ である。
\end{enumerate}
\end{proposition}

\begin{proof}
この結果を $d$ に関する帰納法で証明する。

\medskip\noindent
$d = 0$ ならば $X$ はプロ有限空間である
（《位相空間》補題 \ref{topology-lemma-characterize-profinite-spectral} を参照）。
したがって補題 \ref{cohomology-lemma-vanishing-for-profinite} により (1) が成り立つ。
$U \subset X$ が準コンパクト開集合ならば，Hausdorff 空間の準コンパクト
部分集合として $U$ は閉でもある。したがって位相空間として
$X = U \amalg (X \setminus U)$ であり，(2) が成り立つことが分かる。
(3) のような $Z$ が与えられたとき，長完全列
$$
H^{q - 1}(X, \mathcal{F}) \to
H^{q - 1}(X \setminus Z, \mathcal{F}) \to
H^q_Z(X, \mathcal{F}) \to H^q(X, \mathcal{F})
$$
を考える。$X$ と $U = X \setminus Z$ はプロ有限であり（実際，$Z$ は
構成可能なので $U$ は準コンパクトである），また (2) と (1) が成り立つので，
$Z$ に台をもつコホモロジー群の望む消滅を得る。

\medskip\noindent
帰納段階に移る。$d \geq 1$ とし，次元 $< d$ のすべてのスペクトル空間に対して
命題が成り立つと仮定する。まず $X$ に対して (2) を証明する。
$U$ を準コンパクト開集合，$\xi \in H^d(U, \mathcal{F})$ とする。
$Z = X \setminus U$ と置き，$W \subset X$ を $Z$ へ特殊化する点全体の集合とする。
補題 \ref{cohomology-lemma-cohomology-of-neighbourhoods-of-closed} により
$$
H^d(W \setminus Z, \mathcal{F}|_{W \setminus Z}) =
\colim_{Z \subset V} H^d(V \setminus Z, \mathcal{F})
$$
である。ここで余極限は，$X$ における $Z$ の準コンパクト開近傍 $V$ をわたる。
《位相空間》補題 \ref{topology-lemma-make-spectral-space} により，
$W \setminus Z$ はスペクトル空間である。$W$ の各点は $Z$ のある点へ
特殊化するので，$W \setminus Z$ は Krull 次元 $< d$ のスペクトル空間である。
帰納法の仮定により，$H^d(W \setminus Z,
\mathcal{F}|_{W \setminus Z})$ における $\xi$ の像は零である。
表示した公式により，$\xi|_{V \setminus Z} = 0$ となる準コンパクト開集合
$Z \subset V \subset X$ が存在する。$V \setminus Z = V \cap U$ なので，被覆
$X = U \cup V$ に対する Mayer-Vietoris（補題
\ref{cohomology-lemma-mayer-vietoris}）により，$U$ 上では $\xi$ に，$V$ 上では零に制限
される $\tilde \xi \in H^d(X, \mathcal{F})$ が存在する。言い換えれば (2) が成り立つ。

\medskip\noindent
(2) を仮定して (1) を証明する。入射分解
$\mathcal{F} \to \mathcal{I}^\bullet$ を選び，
$$
\mathcal{G} = \Im(\mathcal{I}^{d - 1} \to \mathcal{I}^d) =
\Ker(\mathcal{I}^d \to \mathcal{I}^{d + 1})
$$
と置く。準コンパクト開集合 $U \subset X$ に対して，次の完全列の写像がある。
$$
\xymatrix{
\mathcal{I}^{d - 1}(X) \ar[r] \ar[d] &
\mathcal{G}(X) \ar[r] \ar[d] &
H^d(X, \mathcal{F}) \ar[d] \ar[r] & 0 \\
\mathcal{I}^{d - 1}(U) \ar[r] &
\mathcal{G}(U) \ar[r] &
H^d(U, \mathcal{F}) \ar[r] & 0
}
$$
補題 \ref{cohomology-lemma-injective-flasque} と $d \geq 1$ であることにより，層
$\mathcal{I}^{d - 1}$ はフラスクである。(2) により右の垂直矢印は全射である。
図式追跡から，写像 $\mathcal{G}(X) \to \mathcal{G}(U)$ は全射であると結論する。
補題 \ref{cohomology-lemma-vanishing-ravi} により，$q > 0$ であり，$\mathcal{U} :
U = U_1 \cup \ldots \cup U_n$ は準コンパクト開集合を準コンパクト開集合で
覆う任意の有限被覆であるとき，
$\check{H}^q(\mathcal{U}, \mathcal{G}) = 0$ である。補題
\ref{cohomology-lemma-cech-vanish-basis} を適用すると，
$X$ のすべての準コンパクト開集合 $U$ とすべての $q > 0$ に対して
$H^q(U, \mathcal{G}) = 0$ であることが分かる。Leray の非輪状性補題
（《導来圏》補題 \ref{derived-lemma-leray-acyclicity}）により
$$
H^q(X, \mathcal{F}) =
H^q\left(
\Gamma(X, \mathcal{I}^0) \to \ldots \to
\Gamma(X, \mathcal{I}^{d - 1}) \to \Gamma(X, \mathcal{G})
\right)
$$
である。特に $q > d$ ならばコホモロジー群は消滅する。

\medskip\noindent
(3) を証明する。(3) のような $Z$ が与えられたとき，長完全列
$$
H^{q - 1}(X, \mathcal{F}) \to
H^{q - 1}(X \setminus Z, \mathcal{F}) \to
H^q_Z(X, \mathcal{F}) \to H^q(X, \mathcal{F})
$$
を考える。$X$ と $U = X \setminus Z$ は次元 $\leq d$ のスペクトル空間であり
（《位相空間》補題 \ref{topology-lemma-spectral-sub}），また (2) と (1) が
成り立つので，望む消滅を得る。
\end{proof}













\section{交代 {\v C}ech 複体}
\label{cohomology-section-alternating-cech}

\noindent
本節では，{\v C}ech 複体を交代 {\v C}ech 複体およびいくつかの
関連する複体と比較する。

\medskip\noindent
$X$ を位相空間とする。$\mathcal{U} : U = \bigcup_{i \in I} U_i$ を
開被覆とする。$p \geq 0$ に対して
$$
\check{\mathcal{C}}_{alt}^p(\mathcal{U}, \mathcal{F})
=
\left\{
\begin{matrix}
s \in  \check{\mathcal{C}}^p(\mathcal{U}, \mathcal{F})
\text{ であって }
s_{i_0 \ldots i_p} = 0 \text{（もし } i_n = i_m
\text{ となるある } n \not = m \text{ が存在するならば），}\\
\text{ かつ }
s_{i_0\ldots i_n \ldots i_m \ldots i_p}
=
-s_{i_0\ldots i_m \ldots i_n \ldots i_p}
\text{ が常に成り立つ。}
\end{matrix}
\right\}
$$
と置く。式 (\ref{cohomology-equation-d-cech}) の微分 $d$ が
$\check{\mathcal{C}}^p_{alt}(\mathcal{U}, \mathcal{F})$ を
$\check{\mathcal{C}}^{p + 1}_{alt}(\mathcal{U}, \mathcal{F})$ へ写すことの
確認は省略する。

\begin{definition}
\label{cohomology-definition-alternating-cech-complex}
$X$ を位相空間とし，$\mathcal{U} : U = \bigcup_{i \in I} U_i$ を
開被覆とする。$\mathcal{F}$ を $X$ 上のアーベル前層とする。
複体 $\check{\mathcal{C}}_{alt}^\bullet(\mathcal{U}, \mathcal{F})$ を，
$\mathcal{F}$ と開被覆 $\mathcal{U}$ に付随する
{\it 交代 {\v C}ech 複体}という。
\end{definition}

\noindent
したがって，複体の標準的射
$$
\check{\mathcal{C}}_{alt}^\bullet(\mathcal{U}, \mathcal{F})
\longrightarrow
\check{\mathcal{C}}^\bullet(\mathcal{U}, \mathcal{F})
$$
がある。これは交代 {\v C}ech 複体から通常の {\v C}ech 複体への
包含にほかならない。

\medskip\noindent
被覆 $\mathcal{U} : U = \bigcup_{i \in I} U_i$ に，$I$ 上の全順序
$<$ が与えられていると仮定する。このとき
$$
\check{\mathcal{C}}_{ord}^p(\mathcal{U}, \mathcal{F})
=
\prod\nolimits_{(i_0, \ldots, i_p) \in I^{p + 1}, i_0 < \ldots < i_p}
\mathcal{F}(U_{i_0\ldots i_p}).
$$
と置く。これはアーベル群である。
$s \in \check{\mathcal{C}}_{ord}^p(\mathcal{U}, \mathcal{F})$ に対し，
$\mathcal{F}(U_{i_0\ldots i_p})$ におけるその値を
$s_{i_0\ldots i_p}$ と書く。
$$
d : \check{\mathcal{C}}_{ord}^p(\mathcal{U}, \mathcal{F})
\longrightarrow
\check{\mathcal{C}}_{ord}^{p + 1}(\mathcal{U}, \mathcal{F})
$$
を，任意の $i_0 < \ldots < i_{p + 1}$ に対する公式
$$
d(s)_{i_0\ldots i_{p + 1}}
=
\sum\nolimits_{j = 0}^{p + 1}
(-1)^j
s_{i_0\ldots \hat i_j \ldots i_{p + 1}}|_{U_{i_0\ldots i_{p + 1}}}
$$
によって定める。この公式は式 (\ref{cohomology-equation-d-cech}) と同一であることに
注意する。$d \circ d = 0$ は直ちに分かる。言い換えれば，
$\check{\mathcal{C}}_{ord}^\bullet(\mathcal{U}, \mathcal{F})$ は複体である。

\begin{definition}
\label{cohomology-definition-ordered-cech-complex}
$X$ を位相空間とする。
$\mathcal{U} : U = \bigcup_{i \in I} U_i$ を開被覆とする。
$I$ 上に全順序が与えられていると仮定する。
$\mathcal{F}$ を $X$ 上のアーベル前層とする。
複体 $\check{\mathcal{C}}_{ord}^\bullet(\mathcal{U}, \mathcal{F})$ を，
$\mathcal{F}$，開被覆 $\mathcal{U}$，および $I$ 上の所与の全順序に
付随する {\it 順序付き {\v C}ech 複体}という。
\end{definition}

\noindent
この複体を交代 {\v C}ech 複体と呼ぶこともある。その理由は，
順序付き {\v C}ech 複体と交代 {\v C}ech 複体の間に明らかな比較写像が
あるからである。すなわち，規則
$$
c :
\check{\mathcal{C}}_{ord}^\bullet(\mathcal{U}, \mathcal{F})
\longrightarrow
\check{\mathcal{C}}^\bullet(\mathcal{U}, \mathcal{F})
$$
を
$$
c(s)_{i_0\ldots i_p} =
\left\{
\begin{matrix}
0 &
\text{もし} &
i_n = i_m \text{ となるある } n \not = m \text{ が存在するならば}\\
\text{sgn}(\sigma) s_{i_{\sigma(0)}\ldots i_{\sigma(p)}} &
\text{もし} &
i_{\sigma(0)} < i_{\sigma(1)} < \ldots < i_{\sigma(p)}
\text{ のとき}
\end{matrix}
\right.
$$
によって定める。ここで $\sigma$ は $\{0, \ldots, p\}$ の置換を表し，
$\text{sgn}(\sigma)$ はその符号を表す。文献では，交代 {\v C}ech 複体と
順序付き {\v C}ech 複体を写像 $c$ によって同一視することが多い。
実際，次の容易な補題が成り立つ。

\begin{lemma}
\label{cohomology-lemma-ordered-alternating}
$X$ を位相空間とする。
$\mathcal{U} : U = \bigcup_{i \in I} U_i$ を開被覆とする。
$I$ に全順序が与えられていると仮定する。
写像 $c$ は複体の射である。実際，これは複体の同型
$$
c : \check{\mathcal{C}}_{ord}^\bullet(\mathcal{U}, \mathcal{F})
\to \check{\mathcal{C}}_{alt}^\bullet(\mathcal{U}, \mathcal{F})
$$
を誘導する。
\end{lemma}

\begin{proof}
省略する。
\end{proof}

\noindent
さらに，$i_0 < i_1 < \ldots < i_p$ のときの規則
$$
\pi :
\check{\mathcal{C}}^\bullet(\mathcal{U}, \mathcal{F})
\longrightarrow
\check{\mathcal{C}}_{ord}^\bullet(\mathcal{U}, \mathcal{F})
$$
を
$$
\pi(s)_{i_0\ldots i_p} = s_{i_0\ldots i_p}
$$
によって記述される写像 $\pi$ もある。

\begin{lemma}
\label{cohomology-lemma-project-to-ordered}
$X$ を位相空間とする。
$\mathcal{U} : U = \bigcup_{i \in I} U_i$ を開被覆とする。
$I$ に全順序が与えられていると仮定する。
写像 $\pi : \check{\mathcal{C}}^\bullet(\mathcal{U}, \mathcal{F})
\to \check{\mathcal{C}}_{ord}^\bullet(\mathcal{U}, \mathcal{F})$ は
複体の射である。これは複体の同型
$$
\pi : \check{\mathcal{C}}_{alt}^\bullet(\mathcal{U}, \mathcal{F})
\to \check{\mathcal{C}}_{ord}^\bullet(\mathcal{U}, \mathcal{F})
$$
を誘導し，この同型は射 $c$ の左逆である。
\end{lemma}

\begin{proof}
省略する。
\end{proof}

\begin{remark}
\label{cohomology-remark-compared-ordered-complexes}
これは，添字集合 $I$ 上に二つの全順序 $<_1$ と $<_2$ があるとき，
複体の同型
$\tau = \pi_2 \circ c_1 :
\check{\mathcal{C}}_{ord\text{-}1}(\mathcal{U}, \mathcal{F}) \to
\check{\mathcal{C}}_{ord\text{-}2}(\mathcal{U}, \mathcal{F})$
を得ることを意味する。明らかに
$$
\tau(s)_{i_0 \ldots i_p} =
\text{sign}(\sigma) s_{i_{\sigma(0)} \ldots i_{\sigma(p)}}
$$
である。ここで $i_0 <_1 i_1 <_1 \ldots <_1 i_p$ かつ
$i_{\sigma(0)} <_2 i_{\sigma(1)} <_2 \ldots <_2 i_{\sigma(p)}$ である。
この意味で，順序付き {\v C}ech 複体は選んだ全順序に依存しない。
\end{remark}

\begin{lemma}
\label{cohomology-lemma-alternating-usual}
$X$ を位相空間とする。
$\mathcal{U} : U = \bigcup_{i \in I} U_i$ を開被覆とする。
$I$ に全順序が与えられていると仮定する。
写像 $c \circ \pi$ は
$\check{\mathcal{C}}^\bullet(\mathcal{U}, \mathcal{F})$ 上の恒等写像と
ホモトピックである。特に包含写像
$\check{\mathcal{C}}_{alt}^\bullet(\mathcal{U}, \mathcal{F}) \to
\check{\mathcal{C}}^\bullet(\mathcal{U}, \mathcal{F})$
はホモトピー同値である。
\end{lemma}

\begin{proof}
任意の多重添字 $(i_0, \ldots, i_p) \in I^{p + 1}$ に対して，
$$
i_{\sigma(0)} \leq i_{\sigma(1)} \leq \ldots \leq i_{\sigma(p)}
\quad
\text{かつ}
\quad
\sigma(j) < \sigma(j + 1)
\quad
\text{（もし}
\quad
i_{\sigma(j)} = i_{\sigma(j + 1)}\text{ならば）}
$$
を満たす一意な置換
$\sigma : \{0, \ldots, p\} \to \{0, \ldots, p\}$ が存在する。
この置換を $\sigma = \sigma^{i_0 \ldots i_p}$ と書く。

\medskip\noindent
任意の置換 $\sigma : \{0, \ldots, p\} \to \{0, \ldots, p\}$ と
任意の $0 \leq a \leq p$ に対し，$0 \leq j < a$ では
$\sigma_a(j) = \sigma(j)$ であり，かつ
$\sigma_a(a) < \sigma_a(a + 1) < \ldots < \sigma_a(p)$ であるような
$\{0, \ldots, p\}$ の一意な置換を $\sigma_a$ と書く。
したがって $p = 3$ で，$\sigma$ と $\tau$ が
$$
\begin{matrix}
\text{id} & 0 & 1 & 2 & 3 \\
\sigma & 3 & 2 & 1 & 0
\end{matrix}
\quad \text{かつ} \quad
\begin{matrix}
\text{id} & 0 & 1 & 2 & 3 \\
\tau & 3 & 0 & 2 & 1
\end{matrix}
$$
で与えられるならば，
$$
\begin{matrix}
\text{id} & 0 & 1 & 2 & 3 \\
\sigma_0 & 0 & 1 & 2 & 3 \\
\sigma_1 & 3 & 0 & 1 & 2 \\
\sigma_2 & 3 & 2 & 0 & 1 \\
\sigma_3 & 3 & 2 & 1 & 0 \\
\end{matrix}
\quad \text{かつ} \quad
\begin{matrix}
\text{id} & 0 & 1 & 2 & 3 \\
\tau_0 & 0 & 1 & 2 & 3 \\
\tau_1 & 3 & 0 & 1 & 2 \\
\tau_2 & 3 & 0 & 1 & 2 \\
\tau_3 & 3 & 0 & 2 & 1 \\
\end{matrix}
$$
を得る。常に $\sigma_0 = \text{id}$ かつ $\sigma_p = \sigma$ であることは
明らかである。

\medskip\noindent
この記法のもとで，
$s \in \check{\mathcal{C}}^{p + 1}(\mathcal{U}, \mathcal{F})$ に対する元
$h(s) \in \check{\mathcal{C}}^p(\mathcal{U}, \mathcal{F})$ を，成分
\begin{equation}
\label{cohomology-equation-first-homotopy}
h(s)_{i_0\ldots i_p} =
\sum\nolimits_{0 \leq a \leq p}
(-1)^a \text{sign}(\sigma_a)
s_{i_{\sigma(0)} \ldots i_{\sigma(a)} i_{\sigma_a(a)} \ldots i_{\sigma_a(p)}}
\end{equation}
によって定める。ここで $\sigma = \sigma^{i_0 \ldots i_p}$ である。添字
$i_{\sigma(a)}$ は
$i_{\sigma(0)} \ldots i_{\sigma(a)} i_{\sigma_a(a)} \ldots i_{\sigma_a(p)}$
の中に二度現れる。実際，$\sigma_a$ の定義により，ある $j \geq a$ に
ついて $\sigma_a(j) = \sigma(a)$ なので，最初の $a + 1$ 個の添字の群に
一度，次の $p - a + 1$ 個の添字の群に一度現れる。したがって，各元
$s_{i_{\sigma(0)} \ldots i_{\sigma(a)} i_{\sigma_a(a)} \ldots i_{\sigma_a(p)}}$
が開集合 $U_{i_0 \ldots i_p}$ 上で定義されているから，この和には意味が
ある。また $a = 0$ のときは $s_{i_0 \ldots i_p}$ を得て，
$a = p$ のときは
$(-1)^p \text{sign}(\sigma) s_{i_{\sigma(0)} \ldots i_{\sigma(p)}}$
を得ることにも注意する。

\medskip\noindent
次を主張する：
$$
(dh + hd)(s)_{i_0 \ldots i_p} =
s_{i_0 \ldots i_p} -
\text{sign}(\sigma) s_{i_{\sigma(0)} \ldots i_{\sigma(p)}}
$$
ここで $\sigma = \sigma^{i_0 \ldots i_p}$ である。この主張の確認は
省略する。（stacks-project のサブディレクトリ scripts には
first-homotopy.gp という PARI/gp スクリプトがあり，この主張の有限個の
場合を確認するために用いることができる。符号が正しいことを確かめる
ために，このスクリプトを書いた。）規則
$$
\kappa :
\check{\mathcal{C}}^\bullet(\mathcal{U}, \mathcal{F})
\longrightarrow
\check{\mathcal{C}}^\bullet(\mathcal{U}, \mathcal{F})
$$
で与えられる作用素を $\kappa$ と書く：
$$
\kappa(s)_{i_0 \ldots i_p} =
\text{sign}(\sigma^{i_0 \ldots i_p}) s_{i_{\sigma(0)} \ldots i_{\sigma(p)}}.
$$
上の主張から，$\kappa$ は複体の射であり，{\v C}ech 複体の恒等写像と
ホモトピックであることが従う。しかし作用素 $\kappa$ の像は交代部分複体
ではないので，これだけから補題が直ちに従うわけではない。実際，
$\kappa$ の像は「半交代」複体
$\check{\mathcal{C}}_{semi\text{-}alt}^p(\mathcal{U}, \mathcal{F})$
であり，$s$ がこの複体の $p$ 次余鎖であるための必要十分条件は，
$\sigma = \sigma^{i_0 \ldots i_p}$ となる任意の
$(i_0, \ldots, i_p) \in I^{p + 1}$ に対して
$$
s_{i_0 \ldots i_p} = \text{sign}(\sigma) s_{i_{\sigma(0)} \ldots i_{\sigma(p)}}
$$
が成り立つことである。さらにもう一つの変種の {\v C}ech 複体，すなわち
次で定義される半順序付き {\v C}ech 複体を導入する：
$$
\check{\mathcal{C}}_{semi\text{-}ord}^p(\mathcal{U}, \mathcal{F})
=
\prod\nolimits_{i_0 \leq i_1 \leq \ldots \leq i_p}
\mathcal{F}(U_{i_0 \ldots i_p})
$$
式 (\ref{cohomology-equation-d-cech}) はここでも微分を定め，したがって複体を得る
ことは容易に分かる。また（補題 \ref{cohomology-lemma-project-to-ordered} と同様に），
射影写像
$$
\check{\mathcal{C}}_{semi\text{-}alt}^\bullet(\mathcal{U}, \mathcal{F})
\longrightarrow
\check{\mathcal{C}}_{semi\text{-}ord}^\bullet(\mathcal{U}, \mathcal{F})
$$
が複体の同型であることも明らかである。

\medskip\noindent
したがって，明らかな包含写像
$$
\check{\mathcal{C}}_{ord}^p(\mathcal{U}, \mathcal{F})
\longrightarrow
\check{\mathcal{C}}_{semi\text{-}ord}^p(\mathcal{U}, \mathcal{F})
$$
がホモトピー同値であることを示せば補題が従う。このためにホモトピー
\begin{equation}
\label{cohomology-equation-second-homotopy}
h(s)_{i_0 \ldots i_p} =
\left\{
\begin{matrix}
0 & \text{もし} & i_0 < i_1 < \ldots < i_p \\
(-1)^a s_{i_0 \ldots i_{a - 1} i_a i_a i_{a + 1} \ldots i_p}
& \text{もし} & i_0 < i_1 < \ldots < i_{a - 1} < i_a = i_{a + 1}
\end{matrix}
\right.
\end{equation}
を用いる。次を主張する：
$$
(dh + hd)(s)_{i_0 \ldots i_p} =
\left\{
\begin{matrix}
0 & \text{もし} & i_0 < i_1 < \ldots < i_p \\
s_{i_0 \ldots i_p}
& \text{それ以外} &
\end{matrix}
\right.
$$
確認は省略する。（stacks-project のサブディレクトリ scripts には
second-homotopy.gp という PARI/gp スクリプトがあり，この主張の有限個の
場合を確認するために用いることができる。符号が正しいことを確かめる
ために，このスクリプトを書いた。）この主張から，射影と自然な包含の合成
$$
\check{\mathcal{C}}_{semi\text{-}ord}^\bullet(\mathcal{U}, \mathcal{F})
\longrightarrow
\check{\mathcal{C}}_{ord}^\bullet(\mathcal{U}, \mathcal{F})
\longrightarrow
\check{\mathcal{C}}_{semi\text{-}ord}^\bullet(\mathcal{U}, \mathcal{F})
$$
が，望みどおり恒等写像とホモトピックであることが明らかに分かる。
\end{proof}

\begin{lemma}
\label{cohomology-lemma-alternating-cech-trivial}
$X$ を位相空間とし，$\mathcal{F}$ を $X$ 上のアーベル前層とする。
$\mathcal{U} : U = \bigcup_{i \in I} U_i$ を開被覆とする。ある
$i \in I$ に対して $U_i = U$ ならば，$\mathcal{F}(U)$ を次数 $-1$ に置き，
$\mathcal{F}(U)$ から $\check{\mathcal{C}}^0(\mathcal{U}, \mathcal{F})$
への標準写像を微分として得られる拡大交代 {\v C}ech 複体
$$
\mathcal{F}(U) \to \check{\mathcal{C}}_{alt}^\bullet(\mathcal{U}, \mathcal{F})
$$
は $0$ とホモトピー同値である。同様に，$I$ 上の任意の全順序に対して，
拡大順序付き {\v C}ech 複体
$$
\mathcal{F}(U) \to
\check{\mathcal{C}}_{ord}^\bullet(\mathcal{U}, \mathcal{F})
$$
も $0$ とホモトピー同値である。
\end{lemma}

\begin{proof}[第1の証明]
補題 \ref{cohomology-lemma-cech-trivial} と補題 \ref{cohomology-lemma-alternating-usual} を
組み合わせればよい。
\end{proof}

\begin{proof}[第2の証明]
交代 {\v C}ech 複体と順序付き {\v C}ech 複体は同型なので，順序付きの
場合を証明すれば十分である。標準的な記法を用いる。拡大順序付き
{\v C}ech 複体の次数 $p$ の余鎖 $s$ は
$s = (s_{i_0 \ldots i_p})$ という形をもち，ここで
$s_{i_0 \ldots i_p}$ は $\mathcal{F}(U_{i_0 \ldots i_p})$ の元であり，
$i_0 < \ldots < i_p$ である。この記法では
$$
d(x)_{i_0 \ldots i_{p + 1}} =
\sum\nolimits_j (-1)^j x_{i_0 \ldots \hat i_j \ldots i_p}
$$
となる。$U = U_i$ となる添字 $i \in I$ を固定する。ホモトピーとして，
規則
$$
h : \text{次数 }p + 1\text{ の余鎖} \to \text{次数 }p\text{ の余鎖}
$$
で与えられる写像を用いる：
$$
h(s)_{i_0 \ldots i_p} = 0 \text{（もし } i \in \{i_0, \ldots, i_p\}
\text{ ならば），かつそうでなければ }
h(s)_{i_0 \ldots i_p} =
(-1)^j s_{i_0 \ldots i_j i i_{j + 1} \ldots i_p}.
$$
第2の場合，ここで $j$ は $i_j < i < i_{j + 1}$ となる一意な添字である。
また $U = U_i$ なので，等式
$$
\mathcal{F}(U_{i_0 \ldots i_p}) =
\mathcal{F}(U_{i_0 \ldots i_j i i_{j + 1} \ldots i_p})
$$
があり，これにより
$(-1)^j s_{i_0 \ldots i_j i i_{j + 1} \ldots i_p}$
を $\mathcal{F}(U_{i_0 \ldots i_p})$ の元とみなすことに意味がある。
計算により $d h + h d$ が恒等写像の負に等しいことを示せば，証明は
完了する。そのため，次数 $p$ の余鎖 $s$ を固定し，
$i_0 < \ldots < i_p$ を $I$ の元とする。

\medskip\noindent
場合 I：$i \in \{i_0, \ldots, i_p\}$ とする。$i = i_t$ としよう。
このとき $h(d(s))_{i_0 \ldots i_p} = 0$ である。一方，
$$
d(h(s))_{i_0 \ldots i_p} =
\sum (-1)^j h(s)_{i_0 \ldots \hat i_j \ldots i_p} =
(-1)^t h(s)_{i_0 \ldots \hat i \ldots i_p} =
(-1)^t (-1)^{t - 1} s_{i_0 \ldots i_p}
$$
である。したがって，望みどおり
$(dh + hd)(s)_{i_0 \ldots i_p} = -s_{i_0 \ldots i_p}$ である。

\medskip\noindent
場合 II：$i \not \in \{i_0, \ldots, i_p\}$ とする。
$i_j < i < i_{j + 1}$ となる $j$ を取る。このとき
\begin{align*}
h(d(s))_{i_0 \ldots i_p}
& =
(-1)^j d(s)_{i_0 \ldots i_j i i_{j + 1} \ldots i_p} \\
& =
\sum\nolimits_{j' \leq j} (-1)^{j + j'}
s_{i_0 \ldots \hat i_{j'} \ldots i_j i i_{j + 1} \ldots i_p} -
s_{i_0 \ldots i_p} \\
&
+ \sum\nolimits_{j' > j} (-1)^{j + j' + 1}
s_{i_0 \ldots i_j i i_{j + 1} \ldots \hat i_{j'} \ldots i_p}
\end{align*}
であることが分かる。一方，
\begin{align*}
d(h(s))_{i_0 \ldots i_p}
& =
\sum\nolimits_{j'} (-1)^{j'} h(s)_{i_0 \ldots \hat i_{j'} \ldots i_p} \\
& =
\sum\nolimits_{j' \leq j} (-1)^{j' + j - 1}
s_{i_0 \ldots \hat i_{j'} \ldots i_j i i_{j + 1} \ldots i_p} \\
& +
\sum\nolimits_{j' > j} (-1)^{j' + j}
s_{i_0 \ldots i_j i i_{j + 1} \ldots \hat i_{j'} \ldots i_p}
\end{align*}
である。これらを加えると，望みどおり
$(dh + hd)(s)_{i_0 \ldots i_p} = - s_{i_0 \ldots i_p}$
を得る。
\end{proof}





\section{{\v C}ech 複体の別の見方}
\label{cohomology-section-locally-finite-cech}

\noindent
本節では，{\v C}ech 複体とコホモロジーの関係を確立する別の方法を論じる。

\begin{lemma}
\label{cohomology-lemma-covering-resolution}
$X$ を環付き空間とする。
$\mathcal{U} : X = \bigcup_{i \in I} U_i$ を $X$ の開被覆とする。
$\mathcal{F}$ を $\mathcal{O}_X$ 加群とする。$\mathcal{F}$ の
$U_{i_0 \ldots i_p}$ への制限を $\mathcal{F}_{i_0 \ldots i_p}$ と書く。
次を満たす $\mathcal{O}_X$ 加群の複体
${\mathfrak C}^\bullet(\mathcal{U}, \mathcal{F})$ が存在する：
$$
{\mathfrak C}^p(\mathcal{U}, \mathcal{F}) =
\prod\nolimits_{i_0 \ldots i_p}
(j_{i_0 \ldots i_p})_* \mathcal{F}_{i_0 \ldots i_p}
$$
であり，微分
$d : {\mathfrak C}^p(\mathcal{U}, \mathcal{F})
\to {\mathfrak C}^{p + 1}(\mathcal{U}, \mathcal{F})$
は式 (\ref{cohomology-equation-d-cech}) のとおりである。さらに，標準写像
$$
\mathcal{F} \to {\mathfrak C}^\bullet(\mathcal{U}, \mathcal{F})
$$
が存在し，これは擬同型である。すなわち，
${\mathfrak C}^\bullet(\mathcal{U}, \mathcal{F})$ は
$\mathcal{F}$ の分解である。
\end{lemma}

\begin{proof}
列
$$
0 \to \mathcal{F} \to \mathfrak{C}^0(\mathcal{U}, \mathcal{F}) \to
\mathfrak{C}^1(\mathcal{U}, \mathcal{F}) \to  \ldots
$$
が茎上で完全であることを確認する。$x \in X$ とし，
$x \in U_{i_{\text{fix}}}$ となる $i_{\text{fix}} \in I$ を選ぶ。
このとき，写像
$$
h : \mathfrak{C}^p(\mathcal{U}, \mathcal{F})_x
\to \mathfrak{C}^{p - 1}(\mathcal{U}, \mathcal{F})_x
$$
を次のように定める。$s \in \mathfrak{C}^p(\mathcal{U}, \mathcal{F})_x$
ならば，$x$ のある近傍 $V$ 上で定義された代表元
$$
\widetilde{s} \in
\mathfrak{C}^p(\mathcal{U}, \mathcal{F})(V) =
\prod\nolimits_{i_0 \ldots i_p}
\mathcal{F}(V \cap U_{i_0} \cap \ldots \cap U_{i_p})
$$
を取り，
$$
h(s)_{i_0 \ldots i_{p - 1}} =
\widetilde{s}_{i_{\text{fix}} i_0 \ldots i_{p - 1}, x}.
$$
と置く。同じ公式（$p = 0$ の場合）により，写像
$\mathfrak{C}^{0}(\mathcal{U},\mathcal{F})_x \to \mathcal{F}_x$ を得る。
形式的に計算すると次のようになる：
\begin{align*}
(dh + hd)(s)_{i_0 \ldots i_p}
& =
\sum\nolimits_{j = 0}^p
(-1)^j
h(s)_{i_0 \ldots \hat i_j \ldots i_p}
+
d(s)_{i_{\text{fix}} i_0 \ldots i_p}\\
& =
\sum\nolimits_{j = 0}^p
(-1)^j
s_{i_{\text{fix}} i_0 \ldots \hat i_j \ldots i_p}
+
s_{i_0 \ldots i_p}
+
\sum\nolimits_{j = 0}^p
(-1)^{j + 1}
s_{i_{\text{fix}} i_0 \ldots \hat i_j \ldots i_p} \\
& =
s_{i_0 \ldots i_p}
\end{align*}
これは，$h$ が拡大複体
$$
0 \to \mathcal{F}_x \to \mathfrak{C}^0(\mathcal{U}, \mathcal{F})_x
\to \mathfrak{C}^1(\mathcal{U}, \mathcal{F})_x \to \ldots
$$
の恒等写像から零写像へのホモトピーであることを示すので，結論を得る。
\end{proof}

\noindent
この補題を用いると，補題 \ref{cohomology-lemma-cech-spectral-sequence} の
{\v C}ech からコホモロジーへのスペクトル系列を容易に再証明できる。
実際，$X$，$\mathcal{U}$，$\mathcal{F}$ を補題
\ref{cohomology-lemma-covering-resolution} のとおりとし，
$\mathcal{F} \to \mathcal{I}^\bullet$ を入射分解とする。このとき二重複体
$$
A^{\bullet, \bullet} =
\Gamma(X, {\mathfrak C}^\bullet(\mathcal{U}, \mathcal{I}^\bullet)).
$$
を考えることができる。構成により
$$
A^{p, q} = \prod\nolimits_{i_0 \ldots i_p} \mathcal{I}^q(U_{i_0 \ldots i_p})
$$
である。この二重複体に付随する《ホモロジー》節
\ref{homology-section-double-complex} の二つのスペクトル系列を考える。
特に《ホモロジー》補題 \ref{homology-lemma-ss-double-complex} を参照せよ。
積を取ることは完全なので（《ホモロジー》補題
\ref{homology-lemma-product-abelian-groups-exact}），スペクトル系列
$({}'E_r, {}'d_r)_{r \geq 0}$ に対して
${}'E_2^{p, q} = \check{H}^p(\mathcal{U}, \underline{H}^q(\mathcal{F}))$
を得る。スペクトル系列 $({}''E_r, {}''d_r)_{r \geq 0}$ に対しては，
$p > 0$ のとき ${}''E_2^{p, q} = 0$ であり，かつ
${}''E_2^{0, q} = H^q(X, \mathcal{F})$ である。実際，固定した $q$ に対し，
層の複体
${\mathfrak C}^\bullet(\mathcal{U}, \mathcal{I}^q)$
は，入射層 $\mathcal{I}^q$ の入射層による分解である（補題
\ref{cohomology-lemma-covering-resolution}，補題 \ref{cohomology-lemma-cohomology-of-open}，
補題 \ref{cohomology-lemma-pushforward-injective-flat}，および《ホモロジー》補題
\ref{homology-lemma-product-injectives}）。したがって，
$\Gamma(X, {\mathfrak C}^\bullet(\mathcal{U}, \mathcal{I}^q))$
のコホモロジーは正次数で零であり，次数 $0$ では
$\Gamma(X, \mathcal{I}^q)$ に等しい。次の微分についてコホモロジーを
取ることで，スペクトル系列 $({}''E_r, {}''d_r)_{r \geq 0}$ に関する
上の主張を得る。両スペクトル系列は $A^{\bullet, \bullet}$ に付随する
全複体のコホモロジーへ収束するので，結果が従う。

\begin{definition}
\label{cohomology-definition-covering-locally-finite}
$X$ を位相空間とする。開被覆 $X = \bigcup_{i \in I} U_i$ が
{\it 局所有限}であるとは，任意の $x \in X$ に対し，
$\{i \in I \mid W \cap U_i \not = \emptyset\}$ が有限となるような
$x$ の開近傍 $W$ が存在することをいう。
\end{definition}

\begin{remark}
\label{cohomology-remark-locally-finite-sections}
$X = \bigcup_{i \in I} U_i$ を局所有限開被覆とする。
包含写像を $j_i : U_i \to X$ と書く。各 $i$ に対し，$U_i$ 上の
アーベル層 $\mathcal{F}_i$ が与えられていると仮定する。アーベル層
$\mathcal{G} = \bigoplus_{i \in I} (j_i)_*\mathcal{F}_i$ を考える。
すると，開集合 $V \subset X$ に対して実際に
$$
\Gamma(V, \mathcal{G}) = \prod\nolimits_{i \in I} \mathcal{F}_i(V \cap U_i).
$$
となる。言い換えれば，
$$
\bigoplus\nolimits_{i \in I} (j_i)_*\mathcal{F}_i =
\prod\nolimits_{i \in I} (j_i)_*\mathcal{F}_i
$$
層の族の直和は，直和になると素朴に予想される前層の層化であることに
気づけば，これは奇妙ではない。《加群》節 \ref{modules-section-kernels}
の議論を参照せよ。したがってこの場合，補題
\ref{cohomology-lemma-covering-resolution} の複体の項は
$$
{\mathfrak C}^p(\mathcal{U}, \mathcal{F}) =
\bigoplus\nolimits_{i_0 \ldots i_p}
(j_{i_0 \ldots i_p})_* \mathcal{F}_{i_0 \ldots i_p}
$$
となる。これは有用なことがある。
\end{remark}











\section{複体の {\v C}ech コホモロジー}
\label{cohomology-section-cech-cohomology-of-complexes}

\noindent
一般に，$X$ 上のアーベル群の層 ${\mathcal F}$ と ${\mathcal G}$ に対し，
カップ積写像
$$
H^i(X, {\mathcal F}) \times H^j(X, {\mathcal G})
\longrightarrow
H^{i + j}(X, {\mathcal F} \otimes_{\mathbf Z} {\mathcal G}).
$$
がある。本節では，カップ積の明示的な公式を用い，{\v C}ech コサイクルに
よってこれを定義する。コホモロジーが {\v C}ech コホモロジーに等しく
ないかもしれないことが気になるならば，超被覆を用いても同じコサイクル
記法を使える。この方法には，任意のトポス上の超コホモロジーのカップ積を
定義するためにも使えるという利点もある（ここに将来の参照を挿入）。

\medskip\noindent
$X$ 上のアーベル群の前層からなる下に有界な複体
${\mathcal F}^\bullet$ を取る。しばしば {\v C}ech コサイクルを用いて
$H^n(X, {\mathcal F}^\bullet)$ を計算できる。実際，
${\mathcal U} : X = \bigcup_{i \in I} U_i$
を $X$ の開被覆とする。{\v C}ech 複体
$\check{\mathcal{C}}^\bullet(\mathcal{U}, \mathcal{F})$
（定義 \ref{cohomology-definition-cech-complex}）は前層 $\mathcal{F}$ に関して
関手的なので，二重複体
$\check{\mathcal{C}}^\bullet(\mathcal{U}, \mathcal{F}^\bullet)$
を得る。$\check{\mathcal{C}}^\bullet({\mathcal U}, {\mathcal F}^\bullet)$
に付随する全複体は，次数 $n$ の項が
$$
\text{Tot}^n(\check{\mathcal{C}}^\bullet({\mathcal U}, {\mathcal F}^\bullet))
=
\bigoplus\nolimits_{p + q = n}
\prod\nolimits_{i_0\ldots i_p} {\mathcal F}^q(U_{i_0\ldots i_p})
$$
である複体である。《ホモロジー》定義
\ref{homology-definition-associated-simple-complex} を参照せよ。
$\text{Tot}^n$ の典型的な元を
$\alpha = \{\alpha_{i_0\ldots i_p}\}$ と書く。ここで
$\alpha_{i_0 \ldots i_p} \in \mathcal{F}^q(U_{i_0\ldots i_p})$ である。
言い換えれば，$\alpha_{i_0\ldots i_p}$ の $\mathcal{F}$-次数は
$q = n - p$ である。この記法を用いるには，$\alpha$ が属する次数を常に
意識する必要がある。この状況を公式
$\deg_{\mathcal F}(\alpha_{i_0\ldots i_p}) = q$
で表す。《ホモロジー》定義
\ref{homology-definition-associated-simple-complex} の規約によれば，
次数 $n$ の元 $\alpha$ の微分は
$$
d(\alpha)_{i_0\ldots i_{p + 1}}
=
\sum\nolimits_{j = 0}^{p + 1}
(-1)^j \alpha_{i_0 \ldots \hat i_j \ldots i_{p + 1}} + 
(-1)^{p + 1}d_{{\mathcal F}}(\alpha_{i_0 \ldots i_{p + 1}})
$$
で与えられる。ここで $d_\mathcal{F}$ は複体
$\mathcal{F}^\bullet$ の微分を表す。式
$\alpha_{i_0 \ldots \hat i_j \ldots i_{p + 1}}$ は，
$\alpha_{i_0 \ldots \hat i_j \ldots i_{p + 1}}
\in {\mathcal F}(U_{i_0\ldots\hat i_j\ldots i_{p + 1}})$ の
$U_{i_0 \ldots i_{p + 1}}$ への制限を意味する。

\medskip\noindent
$\text{Tot}(\check{\mathcal{C}}^\bullet({\mathcal U},
{\mathcal F}^\bullet))$ の構成は ${\mathcal F}^\bullet$ に関して関手的で
ある。また，複体の関手的変換
\begin{equation}
\label{cohomology-equation-global-sections-to-cech}
\Gamma(X, {\mathcal F}^\bullet)
\longrightarrow
\text{Tot}(\check{\mathcal{C}}^\bullet({\mathcal U}, {\mathcal F}^\bullet))
\end{equation}
もあり，次の規則で定められる。切断
$s\in \Gamma(X, {\mathcal F}^n)$ は，
$\alpha_{i_0} = s|_{U_{i_0}}$ かつ $p > 0$ のとき
$\alpha_{i_0\ldots i_p} = 0$ となる元
$\alpha = \{\alpha_{i_0\ldots i_p}\}$ へ写される。

\medskip\noindent
細分を考える。${\mathcal V} = \{ V_j \}_{j\in J}$ を
${\mathcal U}$ の細分とする。これは，すべての $j\in J$ に対して
$V_j \subset U_{t(j)}$ となる写像 $t: J \to I$ があることを意味する。
これにより関手的変換
\begin{equation}
\label{cohomology-equation-transformation}
T_t :
\text{Tot}(\check{\mathcal{C}}^\bullet({\mathcal U}, {\mathcal F}^\bullet))
\longrightarrow
\text{Tot}(\check{\mathcal{C}}^\bullet({\mathcal V}, {\mathcal F}^\bullet)).
\end{equation}
が得られ，規則
$$
T_t(\alpha)_{j_0\ldots j_p}
=
\alpha_{t(j_0)\ldots t(j_p)}|_{V_{j_0\ldots j_p}}.
$$
で定められる。上のような二つの写像 $t, t' : J \to I$ が与えられた
とき，上で構成した写像 $T_t$ と $T_{t'}$ はホモトピックである。
ホモトピーは，次数 $n$ の元 $\alpha$ に対して
$$
h(\alpha)_{j_0\ldots j_p}
=
\sum\nolimits_{a = 0}^{p}
(-1)^a
\alpha_{t(j_0)\ldots t(j_a) t'(j_a) \ldots t'(j_p)}
$$
で与えられる。これが成り立つのは，ここでも $\alpha$ を次数 $n$ の元と
した次の計算による（したがって $d(\alpha)$ の次数は $n + 1$，
$h(\alpha)$ の次数は $n - 1$ である）：
\begin{align*}
(
d(h(\alpha)) + h(d(\alpha))
)_{j_0\ldots j_p}
= &
\sum\nolimits_{k = 0}^p
(-1)^k
h(\alpha)_{j_0 \ldots \hat j_k \ldots j_p}
+ \\
&
(-1)^p
d_{\mathcal F}(h(\alpha)_{j_0 \ldots j_p})
+ \\
&
\sum\nolimits_{a = 0}^p
(-1)^a
d(\alpha)_{t(j_0) \ldots t(j_a) t'(j_a) \ldots t'(j_p)}
\\
= &
\sum\nolimits_{k = 0}^p
\sum\nolimits_{a = 0}^{k - 1}
(-1)^{k + a}
\alpha_{t(j_0)\ldots t(j_a)t'(j_a)\ldots \hat{t'(j_k)}\ldots t'(j_p)}
+ \\
&
\sum\nolimits_{k = 0}^p
\sum\nolimits_{a = k + 1}^p
(-1)^{k + a - 1}
\alpha_{t(j_0)\ldots \hat{t(j_k)}\ldots t(j_a)t'(j_a)\ldots t'(j_p)}
+ \\
&
\sum\nolimits_{a = 0}^p
(-1)^{p + a}
d_{\mathcal F}(\alpha_{t(j_0)\ldots t(j_a) t'(j_a) \ldots t'(j_p)})
+ \\
&
\sum\nolimits_{a = 0}^p
\sum\nolimits_{k = 0}^a
(-1)^{a + k}
\alpha_{t(j_0)\ldots\hat{t(j_k)}\ldots t(j_a)t'(j_a)\ldots t'(j_p)}
+ \\
&
\sum\nolimits_{a = 0}^p
\sum\nolimits_{k = a}^p
(-1)^{a + k + 1}
\alpha_{t(j_0) \ldots t(j_a) t'(j_a) \ldots \hat{t'(j_k)} \ldots t'(j_p)}
+ \\
&
\sum\nolimits_{a = 0}^p
(-1)^{a + p + 1}
d_{\mathcal F}(\alpha_{t(j_0)\ldots t(j_a) t'(j_a) \ldots t'(j_p)})
\\
= &
\alpha_{t'(j_0)\ldots t'(j_p)} +
(-1)^{2p + 1}\alpha_{t(j_0)\ldots t(j_p)}
\\
= &
T_{t'}(\alpha)_{j_0\ldots j_p} - T_t(\alpha)_{j_0\ldots j_p}
\end{align*}
相殺の確認は読者に委ねる。（第1項と第2項，および第4項と第5項における
$k$ と $a$ の両方を含む項は相殺する。ただし $a = k$ のものは第4項と
第5項にのみ現れ，これらも望む二項を除いて互いに相殺する。）
したがって誘導写像
$$
H^n(T_t) :
H^n(
\text{Tot}(\check{\mathcal{C}}^\bullet({\mathcal U}, {\mathcal F}^\bullet))
)
\to
H^n(
\text{Tot}(\check{\mathcal{C}}^\bullet({\mathcal V}, {\mathcal F}^\bullet))
)
$$
は $t$ の選択に依存しない。写像 $H^\bullet(T_t)$ によるすべての細分に
わたる {\v C}ech コホモロジー群の余極限を
{\it {\v C}ech 超コホモロジー}と定義する。

\medskip\noindent
（$X$ のすべての開被覆にわたる）余極限において，次の補題は
{\v C}ech 超コホモロジーからコホモロジーへの写像を与える。この写像は
しばしば同型であり，超被覆を用いれば常に同型である。

\begin{lemma}
\label{cohomology-lemma-cech-complex-complex}
$(X, \mathcal{O}_X)$ を環付き空間とする。
$\mathcal{U} : X = \bigcup_{i \in I} U_i$ を開被覆とする。
$\mathcal{O}_X$ 加群の下に有界な複体 $\mathcal{F}^\bullet$ に対して，
標準写像
$$
\text{Tot}(\check{\mathcal{C}}^\bullet(\mathcal{U}, \mathcal{F}^\bullet))
\longrightarrow
R\Gamma(X, \mathcal{F}^\bullet)
$$
があり，$\mathcal{F}^\bullet$ に関して関手的で，
(\ref{cohomology-equation-global-sections-to-cech}) および
(\ref{cohomology-equation-transformation}) と両立する。また，
$$
E_2^{p, q} =
H^p(\text{Tot}(\check{\mathcal{C}}^\bullet(\mathcal{U},
\underline{H}^q(\mathcal{F}^\bullet)))
$$
をもち，$H^{p + q}(X, \mathcal{F}^\bullet)$ へ収束するスペクトル系列
$(E_r, d_r)_{r \geq 0}$ がある。
\end{lemma}

\begin{proof}
${\mathcal I}^\bullet$ を入射対象からなる下に有界な複体とする。
$\mathcal{I}^\bullet$ に対する写像
(\ref{cohomology-equation-global-sections-to-cech}) は
$\Gamma(X, {\mathcal I}^\bullet) \to
\text{Tot}(\check{\mathcal{C}}^\bullet({\mathcal U}, {\mathcal I}^\bullet))$.
という写像である。《ホモロジー》補題
\ref{homology-lemma-double-complex-gives-resolution} を二重複体
$\check{\mathcal{C}}^\bullet({\mathcal U}, {\mathcal I}^\bullet)$ に適用すれば，
これはアーベル群の複体の擬同型であることが従う。ここでは補題
\ref{cohomology-lemma-injective-trivial-cech} を用いる。
${\mathcal F}^\bullet \to {\mathcal I}^\bullet$ を，
${\mathcal F}^\bullet$ から入射対象からなる下に有界な複体への擬同型とする。
$R\Gamma(X, {\mathcal F}^\bullet)$ は複体
$\Gamma(X, {\mathcal I}^\bullet)$ によって表されるので，
$$
\text{Tot}(\check{\mathcal{C}}^\bullet({\mathcal U}, {\mathcal F}^\bullet))
\longrightarrow
\text{Tot}(\check{\mathcal{C}}^\bullet({\mathcal U}, {\mathcal I}^\bullet)).
$$
を用いて補題の写像を得る。関手性と両立性の確認は省略する。
補題のスペクトル系列を構成するため，Cartan--Eilenberg 分解
$\mathcal{F}^\bullet \to \mathcal{I}^{\bullet, \bullet}$ を選ぶ。
《導来圏》補題 \ref{derived-lemma-cartan-eilenberg} を参照せよ。この場合，
$\mathcal{F}^\bullet \to \text{Tot}(\mathcal{I}^{\bullet, \bullet})$ は
入射分解であり，したがって
$$
\text{Tot}(\check{\mathcal{C}}^\bullet({\mathcal U},
\text{Tot}({\mathcal I}^{\bullet, \bullet})))
$$
は，上で見たように $R\Gamma(X, \mathcal{F}^\bullet)$ を計算する。
《ホモロジー》注意 \ref{homology-remark-triple-complex} により，これを
三重複体
$\check{\mathcal{C}}^\bullet({\mathcal U}, {\mathcal I}^{\bullet, \bullet})$
に付随する全複体とみなせる。したがって同じ注意を用いれば，これを項
$$
A^{n, m} =
\bigoplus\nolimits_{p + q = n}
\check{\mathcal{C}}^p({\mathcal U}, \mathcal{I}^{q, m})
$$
をもつ二重複体 $A^{\bullet, \bullet}$ に付随する全複体とみなせる。
$\mathcal{I}^{q, \bullet}$ は $\mathcal{F}^q$ の入射分解なので，
$A^{\bullet, \bullet}$ に付随する第1スペクトル系列
（《ホモロジー》補題 \ref{homology-lemma-ss-double-complex}）を適用して，
$$
E_1^{n, m} =
\bigoplus\nolimits_{p + q = n}
\check{\mathcal{C}}^p(\mathcal{U}, \underline{H}^m(\mathcal{F}^q))
$$
をもつスペクトル系列を得る。これは複体
$\text{Tot}(\check{\mathcal{C}}^\bullet(\mathcal{U},
\underline{H}^m(\mathcal{F}^\bullet))$ の第 $n$ 項である。したがって，
補題に記した $E_2$ 項を得る。収束性は《ホモロジー》補題
\ref{homology-lemma-first-quadrant-ss} による。
\end{proof}

\begin{lemma}
\label{cohomology-lemma-cech-complex-complex-computes}
$(X, \mathcal{O}_X)$ を環付き空間とする。
$\mathcal{U} : X = \bigcup_{i \in I} U_i$ を開被覆とする。
$\mathcal{F}^\bullet$ を $\mathcal{O}_X$ 加群の下に有界な複体とする。
すべての $i > 0$ およびすべての $p, i_0, \ldots, i_p, q$ に対して
$H^i(U_{i_0 \ldots i_p}, \mathcal{F}^q) = 0$ ならば，補題
\ref{cohomology-lemma-cech-complex-complex} の写像
$
\text{Tot}(\check{\mathcal{C}}^\bullet(\mathcal{U}, \mathcal{F}^\bullet))
\to
R\Gamma(X, \mathcal{F}^\bullet)
$
は同型である。
\end{lemma}

\begin{proof}
補題 \ref{cohomology-lemma-cech-complex-complex} のスペクトル系列から直ちに従う。
\end{proof}

\begin{remark}
\label{cohomology-remark-shift-complex-cech-complex}
$(X, \mathcal{O}_X)$ を環付き空間とする。
$\mathcal{U} : X = \bigcup_{i \in I} U_i$ を開被覆とする。
$\mathcal{F}^\bullet$ を $\mathcal{O}_X$ 加群の下に有界な複体とする。
$b$ を整数とする。導来圏において可換図式
$$
\xymatrix{
\text{Tot}(\check{\mathcal{C}}^\bullet(\mathcal{U}, \mathcal{F}^\bullet))[b]
\ar[r] \ar[d]_\gamma &
R\Gamma(X, \mathcal{F}^\bullet)[b] \ar[d] \\
\text{Tot}(\check{\mathcal{C}}^\bullet(\mathcal{U}, \mathcal{F}^\bullet[b]))
\ar[r] &
R\Gamma(X, \mathcal{F}^\bullet[b])
}
$$
があると主張する。ここで写像 $\gamma$ は，《ホモロジー》注意
\ref{homology-remark-shift-double-complex} で構成された複体上の写像である。
これは意味をもつ。なぜなら，二重複体
$\check{\mathcal{C}}^\bullet(\mathcal{U}, \mathcal{F}^\bullet[b])$
は，《ホモロジー》注意 \ref{homology-remark-shift-double-complex} で導入した
二重複体
$\check{\mathcal{C}}^\bullet(\mathcal{U}, \mathcal{F}^\bullet)[0, b]$
と明らかに同じだからである。図式が可換であることを確認するため，補題
\ref{cohomology-lemma-cech-complex-complex} の証明と同様に入射分解
$\mathcal{F}^\bullet \to \mathcal{I}^\bullet$ を選べる。図式追跡により，
$\mathcal{F}^\bullet$ を $\mathcal{I}^\bullet$ で置き換えた場合に図式が
可換であることを確認すれば十分だと分かる。そこで拡大図式
$$
\xymatrix{
\Gamma(X, \mathcal{I}^\bullet)[b] \ar[r] \ar[d] &
\text{Tot}(\check{\mathcal{C}}^\bullet(\mathcal{U}, \mathcal{I}^\bullet))[b]
\ar[r] \ar[d]_\gamma &
R\Gamma(X, \mathcal{I}^\bullet)[b] \ar[d] \\
\Gamma(X, \mathcal{I}^\bullet[b]) \ar[r] &
\text{Tot}(\check{\mathcal{C}}^\bullet(\mathcal{U}, \mathcal{I}^\bullet[b]))
\ar[r] &
R\Gamma(X, \mathcal{I}^\bullet[b])
}
$$
を考える。ここで左の水平矢印は
(\ref{cohomology-equation-global-sections-to-cech}) である。この場合，水平矢印は
導来圏における同型なので（補題 \ref{cohomology-lemma-cech-complex-complex} の証明を
参照），左の正方形が可換であることを示せば十分である。写像 $\gamma$ は
直和因子 $\check{\mathcal{C}}^0(\mathcal{U}, \mathcal{I}^{q + b})$ 上で符号
$1$ を用いるので，これは成り立つ。《ホモロジー》注意
\ref{homology-remark-shift-double-complex} の公式を参照せよ。
\end{remark}

\noindent
$X$ を位相空間とし，$\mathcal{U} : X = \bigcup_{i \in I} U_i$ を
開被覆とし，$\mathcal{F}^\bullet$ をアーベル群の前層からなる下に有界な
複体とする。次で定められる写像
$\tau :
\text{Tot}(\check{\mathcal{C}}^\bullet({\mathcal U}, {\mathcal F}^\bullet))
\to
\text{Tot}(\check{\mathcal{C}}^\bullet({\mathcal U}, {\mathcal F}^\bullet))$
を考える：
$$
\tau(\alpha)_{i_0 \ldots i_p} = (-1)^{p(p + 1)/2} \alpha_{i_p \ldots i_0}.
$$
このとき次数 $n$ の元 $\alpha$ に対して
\begin{align*}
& d(\tau(\alpha))_{i_0 \ldots i_{p + 1}} \\
& =
\sum\nolimits_{j = 0}^{p + 1}
(-1)^j
\tau(\alpha)_{i_0 \ldots \hat i_j \ldots i_{p + 1}}
+
(-1)^{p + 1}
d_{\mathcal F}(\tau(\alpha)_{i_0 \ldots i_{p + 1}})
\\
& =
\sum\nolimits_{j = 0}^{p + 1}
(-1)^{j + \frac{p(p + 1)}{2}}
\alpha_{i_{p + 1} \ldots \hat i_j \ldots i_0}
+
(-1)^{p + 1 + \frac{(p + 1)(p + 2)}{2}}
d_{\mathcal F}(\alpha_{i_{p + 1} \ldots i_0})
\end{align*}
である。一方，
\begin{align*}
& \tau(d(\alpha))_{i_0\ldots i_{p + 1}} \\
& =
(-1)^{\frac{(p + 1)(p + 2)}{2}} d(\alpha)_{i_{p + 1} \ldots i_0}
\\
& =
(-1)^{\frac{(p + 1)(p + 2)}{2}}
\left(
\sum\nolimits_{j = 0}^{p + 1}
(-1)^j
\alpha_{i_{p + 1}\ldots \hat i_{p + 1 - j} \ldots i_0}
+
(-1)^{p + 1}
d_{\mathcal F}(\alpha_{i_{p + 1}\ldots i_0})
\right)
\end{align*}
である。$p(p + 1)/2 \equiv (p + 1)(p + 2)/2 + p + 1 \bmod 2$ なので，
$d(\tau(\alpha)) = \tau(d(\alpha))$ と結論できる。言い換えれば，
$\tau$ は複体
$\text{Tot}(\check{\mathcal{C}}^\bullet({\mathcal U}, {\mathcal F}^\bullet))$
の自己準同型である。図式
$$
\begin{matrix}
\Gamma(X, {\mathcal F}^\bullet) &
\longrightarrow &
\text{Tot}(\check{\mathcal{C}}^\bullet({\mathcal U}, {\mathcal F}^\bullet)) \\
\downarrow \text{id} & & \downarrow \tau \\
\Gamma(X, {\mathcal F}^\bullet) &
\longrightarrow &
\text{Tot}(\check{\mathcal{C}}^\bullet({\mathcal U}, {\mathcal F}^\bullet))
\end{matrix}
$$
は可換であることに注意する。さらに $\tau$ は明らかに細分と両立する。
これは，$\tau$ が {\v C}ech コホモロジー上で恒等写像として作用することを
示唆する（すなわち余極限において。ただし {\v C}ech 超コホモロジーが
超コホモロジーと一致すると仮定する。超被覆を用いればこれは常に成り立つ）。
実際，$\tau$ は全 {\v C}ech 複体
$\text{Tot}(\check{\mathcal{C}}^\bullet({\mathcal U}, {\mathcal F}^\bullet))$
上で恒等写像とホモトピックであると主張する。これを証明するため，
ホモトピーとして，次数 $n$ の $\alpha$ に対する
$$
h(\alpha)_{i_0 \ldots i_p}
=
\sum\nolimits_{a = 0}^p
\epsilon_p(a)
\alpha_{i_0 \ldots i_a i_p \ldots i_a}
\quad\text{with}\quad
\epsilon_p(a) = (-1)^{\frac{(p - a)(p - a - 1)}{2} + p}
$$
を用いる。通常どおり $|_{U_{i_0 \ldots i_p}}$ は書かない。
これが成り立つのは，ここでも $\alpha$ を次数 $n$ の元とした次の計算に
よる：
\begin{align*}
(d(h(\alpha)) + h(d(\alpha)))_{i_0 \ldots i_p}
= &
\sum\nolimits_{k = 0}^p
(-1)^k
h(\alpha)_{i_0 \ldots \hat i_k \ldots i_p}
+ \\
&
(-1)^p
d_{\mathcal F}(h(\alpha)_{i_0 \ldots i_p})
+ \\
&
\sum\nolimits_{a = 0}^p
\epsilon_p(a)
d(\alpha)_{i_0 \ldots i_a i_p \ldots i_a}
\\
= &
\sum\nolimits_{k = 0}^p
\sum\nolimits_{a = 0}^{k - 1}
(-1)^k \epsilon_{p - 1}(a)
\alpha_{i_0 \ldots i_a i_p \ldots \hat{i_k} \ldots i_a}
+ \\
&
\sum\nolimits_{k = 0}^p
\sum\nolimits_{a = k + 1}^p
(-1)^k \epsilon_{p - 1}(a - 1)
\alpha_{i_0 \ldots \hat{i_k} \ldots i_a i_p \ldots i_a}
+ \\
&
\sum\nolimits_{a = 0}^p
(-1)^p \epsilon_p(a)
d_{\mathcal F}(\alpha_{i_0 \ldots i_a i_p \ldots i_a})
+ \\
&
\sum\nolimits_{a = 0}^p
\sum\nolimits_{k = 0}^a
\epsilon_p(a) (-1)^k
\alpha_{i_0 \ldots \hat{i_k} \ldots i_a i_p \ldots i_a}
+ \\
&
\sum\nolimits_{a = 0}^p
\sum\nolimits_{k = a}^p
\epsilon_p(a) (-1)^{p + a + 1 - k}
\alpha_{i_0 \ldots i_a i_p \ldots \hat{i_k} \ldots i_a}
+ \\
&
\sum\nolimits_{a = 0}^p
\epsilon_p(a) (-1)^{p + 1}
d_{\mathcal F}(\alpha_{i_0 \ldots i_a i_p \ldots i_a})
\\
= &
\epsilon_p(0) \alpha_{i_p \ldots i_0} +
\epsilon_p(p) (-1)^{p + 1} \alpha_{i_0 \ldots i_p} \\
= &
(-1)^{\frac{p(p + 1)}{2}}\alpha_{i_p \ldots i_0}
- \alpha_{i_0 \ldots i_p}
\end{align*}
相殺は
$$
(-1)^k \epsilon_{p - 1}(a) + \epsilon_p(a)(-1)^{p + a + 1 - k} = 0
\quad\text{and}\quad
(-1)^k\epsilon_{p - 1}(a - 1) + \epsilon_p(a) (-1)^k = 0
$$
から従う。相殺の確認は読者に委ねる。

\medskip\noindent
アーベル層からなる二つの下に有界な複体
${\mathcal F}^\bullet$ と ${\mathcal G}^\bullet$ があると仮定する。複体
$\text{Tot}({\mathcal F}^\bullet\otimes_{\mathbf Z} {\mathcal G}^\bullet)$
を，項が $\bigoplus_{p + q = n} {\mathcal F}^p \otimes {\mathcal G}^q$ で，
微分が斉次な $\alpha$ と $\beta$ に対する規則
\begin{equation}
\label{cohomology-equation-differential-tensor-product-complexes}
d(\alpha \otimes \beta) =
d(\alpha)\otimes \beta + (-1)^{\deg(\alpha)} \alpha \otimes d(\beta)
\end{equation}
に従う複体として定義する。《ホモロジー》定義
\ref{homology-definition-associated-simple-complex} を参照せよ。

\medskip\noindent
$M^\bullet$ と $N^\bullet$ をアーベル群の下に有界な二つの複体とする。
$m$ と $n$ がそれぞれ $M^\bullet$ と $N^\bullet$ のコサイクルならば，
$m \otimes n$ が $\text{Tot}(M^\bullet\otimes N^\bullet)$ のコサイクルで
あることは直ちに分かる。したがってカップ積
$$
H^i(M^\bullet) \times H^j(N^\bullet)
\longrightarrow
H^{i + j}(Tot(M^\bullet\otimes N^\bullet)).
$$
を得る。これは《代数続論》節 \ref{more-algebra-section-products-tor} でも
論じられている。

\medskip\noindent
したがって，複体の超コホモロジーにおけるカップ積の構成は，複体の写像
\begin{equation}
\label{cohomology-equation-needs-signs}
\text{Tot}\left(
\text{Tot}(\check{\mathcal{C}}^\bullet({\mathcal U}, {\mathcal F}^\bullet))
\otimes_{\mathbf Z}
\text{Tot}(\check{\mathcal{C}}^\bullet({\mathcal U}, {\mathcal G}^\bullet))
\right)
\longrightarrow
\text{Tot}(
\check{\mathcal{C}}^\bullet({\mathcal U},
\text{Tot}({\mathcal F}^\bullet\otimes {\mathcal G}^\bullet)
))
\end{equation}
の構成に帰着する。この写像を $\cup$ と書き，規則
$$
(\alpha \cup \beta)_{i_0 \ldots i_p}
=
\sum\nolimits_{r = 0}^p
\epsilon(n, m, p, r)
\alpha_{i_0 \ldots i_r} \otimes \beta_{i_r \ldots i_p}.
$$
で与える。ここで $\alpha$ の次数は $n$，$\beta$ の次数は $m$ であり，
$$
\epsilon(n, m, p, r) = (-1)^{(p + r)n + rp + r}.
$$
とする。$\epsilon(n, m, p, n) = 1$ に注意する。したがって，
$\mathcal{F}^\bullet = \mathcal{F}[0]$ が，次数 $0$ に置かれた一つの
アーベル層 $\mathcal{F}$ からなる複体ならば，$\cup$ の公式に符号は
現れない（この場合，$r = n$ でなければ
$\alpha_{i_0 \ldots i_r} = 0$ だからである）。符号が必要な理由とその
計算法については，Deligne による \cite[Exposee XVII]{SGA4} を参照せよ。
(\ref{cohomology-equation-needs-signs}) が複体の写像であることを確認するには，
$$
d(\alpha \cup \beta) =
d(\alpha) \cup \beta +
(-1)^{\deg(\alpha)} \alpha \cup d(\beta)
$$
を示さなければならない。これは，《ホモロジー》定義
\ref{homology-definition-associated-simple-complex} で与えられる
$\text{Tot}(
\text{Tot}(\check{\mathcal{C}}^\bullet({\mathcal U}, {\mathcal F}^\bullet))
\otimes_{\mathbf Z}
\text{Tot}(\check{\mathcal{C}}^\bullet({\mathcal U}, {\mathcal G}^\bullet))
)$ 上の微分の定義による。まず
\begin{align*}
d(\alpha \cup \beta)_{i_0 \ldots i_{p + 1}}
= &
\sum\nolimits_{j = 0}^{p + 1}
(-1)^j
(\alpha \cup \beta)_{i_0 \ldots \hat i_j \ldots i_{p + 1}}
+
(-1)^{p + 1}
d_{{\mathcal F} \otimes {\mathcal G}}
((\alpha \cup \beta)_{i_0 \ldots i_{p + 1}})
\\
= &
\sum\nolimits_{j = 0}^{p + 1}
\sum\nolimits_{r = 0}^{j - 1}
(-1)^j \epsilon(n, m, p, r)
\alpha_{i_0 \ldots i_r} \otimes \beta_{i_r \ldots \hat i_j \ldots i_{p + 1}}
+ \\
&
\sum\nolimits_{j = 0}^{p + 1}
\sum\nolimits_{r = j + 1}^{p + 1}
(-1)^j \epsilon(n, m, p, r - 1)
\alpha_{i_0 \ldots \hat i_j \ldots i_r} \otimes \beta_{i_r \ldots i_{p + 1}}
+ \\
&
\sum\nolimits_{r = 0}^{p + 1}
(-1)^{p + 1} \epsilon(n, m, p + 1, r)
d_{{\mathcal F} \otimes {\mathcal G}}
(\alpha_{i_0 \ldots i_r} \otimes \beta_{i_r \ldots i_{p + 1}})
\end{align*}
と計算する。最後の項の各加数は
$$
(-1)^{p + 1} \epsilon(n, m, p + 1, r)
\left(
d_{\mathcal F}(\alpha_{i_0 \ldots i_r}) \otimes 
\beta_{i_r \ldots i_{p + 1}} +
(-1)^{n - r}
\alpha_{i_0 \ldots i_r} \otimes d_{\mathcal G}(\beta_{i_r \ldots i_{p + 1}})
\right).
$$
に等しい。実際，
$\deg_\mathcal{F}(\alpha_{i_0 \ldots i_r}) = n - r$ である。一方，
\begin{align*}
(d(\alpha) \cup \beta)_{i_0\ldots i_{p + 1}}
= &
\sum\nolimits_{r = 0}^{p + 1}
\epsilon(n + 1, m, p + 1, r)
d(\alpha)_{i_0\ldots i_r} \otimes \beta_{i_r\ldots i_{p + 1}}
\\
= &
\sum\nolimits_{r = 0}^{p + 1}
\sum\nolimits_{j = 0}^{r}
\epsilon(n + 1, m, p + 1, r) (-1)^j
\alpha_{i_0\ldots\hat{i_j}\ldots i_r} \otimes \beta_{i_r\ldots i_{p + 1}}
+ \\
&
\sum\nolimits_{r = 0}^{p + 1}
\epsilon(n + 1, m, p + 1, r) (-1)^r
d_{\mathcal F}(\alpha_{i_0 \ldots i_r}) \otimes \beta_{i_r\ldots i_{p + 1}}
\end{align*}
かつ
\begin{align*}
(\alpha \cup d(\beta))_{i_0\ldots i_{p + 1}}
= &
\sum\nolimits_{r = 0}^{p + 1}
\epsilon(n, m + 1, p + 1, r)
\alpha_{i_0 \ldots i_r} \otimes d(\beta)_{i_r \ldots i_{p + 1}}
\\
= &
\sum\nolimits_{r = 0}^{p + 1}
\sum\nolimits_{j = r}^{p + 1}
\epsilon(n, m + 1, p + 1, r) (-1)^{j - r}
\alpha_{i_0 \ldots i_r} \otimes \beta_{i_r \ldots \hat{i_j}\ldots i_{p + 1}}
+ \\
&
\sum\nolimits_{r = 0}^{p + 1}
\epsilon(n, m + 1, p + 1, r) (-1)^{p + 1 - r}
\alpha_{i_0 \ldots i_r} \otimes d_{\mathcal G}(\beta_{i_r \ldots i_{p + 1}})
\end{align*}
である。次の等式が成り立てば，望む等式が得られる：
\begin{align*}
(-1)^{p + 1} \epsilon(n, m, p + 1, r)
& =
\epsilon(n + 1, m, p + 1, r) (-1)^r \\
(-1)^{p + 1} \epsilon(n, m, p + 1, r) (-1)^{n - r}
& =
(-1)^n \epsilon(n, m + 1, p + 1, r) (-1)^{p + 1 - r} \\
\epsilon(n + 1, m, p + 1, r) (-1)^r
& =
(-1)^{1 + n} \epsilon(n, m + 1, p + 1, r - 1) \\
(-1)^j \epsilon(n, m, p, r)
& =
(-1)^n \epsilon(n, m + 1, p + 1, r) (-1)^{j - r} \\
(-1)^j \epsilon(n, m, p, r - 1)
& =
\epsilon(n + 1, m, p + 1, r) (-1)^j
\end{align*}
（第3の等式は，$d(\alpha) \cup \beta$ と
$(-1)^n \alpha \cup d(\beta)$ から生じる $r = j$ の項を互いに相殺させる
ために必要である。）確認は読者に委ねる。（代わりに，Stacks project の
サブディレクトリ scripts にあるスクリプト signs.gp を確認してもよい。）

\medskip\noindent
カップ積の結合性を考える。${\mathcal F}^\bullet$，
${\mathcal G}^\bullet$，${\mathcal H}^\bullet$ を $X$ 上のアーベル群からなる
下に有界な複体とする。（符号を介在させない）明らかな写像
$$
\text{Tot}(
\text{Tot}({\mathcal F}^\bullet \otimes_{\mathbf Z} {\mathcal G}^\bullet)
\otimes_{\mathbf Z} {\mathcal H}^\bullet
)
\longrightarrow
\text{Tot}(
{\mathcal F}^\bullet \otimes_{\mathbf Z}
\text{Tot}({\mathcal G}^\bullet \otimes_{\mathbf Z} {\mathcal H}^\bullet)
).
$$
は複体の同型である。別の言い方をすれば，三重複体
${\mathcal F}^\bullet \otimes_{\mathbf Z} {\mathcal G}^\bullet
\otimes_{\mathbf Z} {\mathcal H}^\bullet$ は，斉次元に対して
$$
d(\alpha \otimes \beta \otimes \gamma) =
d(\alpha) \otimes \beta \otimes \gamma +
(-1)^{\deg(\alpha)} \alpha \otimes d(\beta) \otimes \gamma +
(-1)^{\deg(\alpha) + \deg(\beta)} \alpha \otimes \beta \otimes d(\gamma)
$$
を満たす微分をもつ，よく定義された全複体を与える。この写像を用いると，
$$
(\alpha \cup \beta) \cup \gamma = \alpha \cup ( \beta \cup \gamma)
$$
を確認するのは容易である。実際，$\alpha$ の次数を $a$，$\beta$ の次数を
$b$，$\gamma$ の次数を $c$ とすると，
\begin{align*}
((\alpha \cup \beta) \cup \gamma)_{i_0 \ldots i_p}
= &
\sum\nolimits_{r = 0}^p
\epsilon(a + b, c, p, r)
(\alpha \cup \beta)_{i_0 \ldots i_r} \otimes \gamma_{i_r \ldots i_p}
\\
= &
\sum\nolimits_{r = 0}^p
\sum\nolimits_{s = 0}^r
\epsilon(a + b, c, p, r) \epsilon(a, b, r, s)
\alpha_{i_0 \ldots i_s} \otimes
\beta_{i_s \ldots i_r} \otimes
\gamma_{i_r \ldots i_p}
\end{align*}
であり，また
\begin{align*}
(\alpha \cup (\beta \cup \gamma)_{i_0\ldots i_p}
= &
\sum\nolimits_{s = 0}^p
\epsilon(a, b + c, p, s)
\alpha_{i_0 \ldots i_s} \otimes (\beta \cup \gamma)_{i_s \ldots i_p}
\\
= &
\sum\nolimits_{s = 0}^p
\sum\nolimits_{r = s}^p
\epsilon(a, b + c, p, s) \epsilon(b, c, p - s, r - s)
\alpha_{i_0 \ldots i_s} \otimes \beta_{i_s \ldots i_r} \otimes
\gamma_{i_r \ldots i_p}
\end{align*}
である。法 $2$ の自明な計算により符号が一致することが分かる。
（代わりに，Stacks project のサブディレクトリ scripts にあるスクリプト
signs.gp を確認してもよい。）

\medskip\noindent
最後に，少なくともコホモロジーのレベルで，カップ積が次数付き可換構造を
保つ理由を示す。そのために，上で導入した作用素 $\tau$ を用いる。
${\mathcal F}^\bullet$ をアーベル群からなる下に有界な複体とし，
次数付き可換な乗法
$$
\wedge^\bullet :
\text{Tot}({\mathcal F}^\bullet\otimes {\mathcal F}^\bullet)
\longrightarrow
{\mathcal F}^\bullet.
$$
が与えられていると仮定する。これは次を意味する：$s$ が
${\mathcal F}^a$ の局所切断で，$t$ が ${\mathcal F}^b$ の局所切断ならば，
$s \wedge t$ は ${\mathcal F}^{a + b}$ の局所切断である。次数付き可換とは
$s \wedge t = (-1)^{ab} t \wedge s$ が成り立つことをいう。
$\wedge$ は複体の写像なので，
$d(s\wedge t) = d(s) \wedge t + (-1)^a s \wedge d(t)$ である。合成
$$
\text{Tot}(
\text{Tot}(\check{\mathcal{C}}^\bullet({\mathcal U}, {\mathcal F}^\bullet))
\otimes
\text{Tot}(\check{\mathcal{C}}^\bullet({\mathcal U}, {\mathcal F}^\bullet))
)
\to
\text{Tot}(
\check{\mathcal{C}}^\bullet({\mathcal U},
\text{Tot}({\mathcal F}^\bullet\otimes_{\mathbf Z}{\mathcal F}^\bullet))
)
\to
\text{Tot}(\check{\mathcal{C}}^\bullet({\mathcal U}, {\mathcal F}^\bullet))
$$
はコホモロジー上のカップ積
$$
H^n(
\text{Tot}(\check{\mathcal{C}}^\bullet({\mathcal U}, {\mathcal F}^\bullet))
)
\times
H^m(
\text{Tot}(\check{\mathcal{C}}^\bullet({\mathcal U}, {\mathcal F}^\bullet))
)
\longrightarrow
H^{n + m}(
\text{Tot}(\check{\mathcal{C}}^\bullet({\mathcal U}, {\mathcal F}^\bullet))
)
$$
を誘導する。したがって余極限において {\v C}ech コホモロジー上の積も
誘導し，ゆえに（必要なら超被覆を用いて）${\mathcal F}^\bullet$ の
コホモロジー上の積を誘導する。この積も（コホモロジー上で）次数付き可換
であると主張する。これを証明するため，まず
$\text{Tot}(\check{\mathcal{C}}^\bullet({\mathcal U}, {\mathcal F}^\bullet))$
の次数 $n$ の元 $\alpha$ と，同じ複体
$\text{Tot}(\check{\mathcal{C}}^\bullet({\mathcal U}, {\mathcal F}^\bullet))$
の次数 $m$ の元 $\beta$ を取り，計算すると
\begin{align*}
\wedge^\bullet(\alpha \cup \beta)_{i_0 \ldots i_p}
= &
\sum\nolimits_{r = 0}^p
\epsilon(n, m, p, r)
\alpha_{i_0 \ldots i_r} \wedge \beta_{i_r \ldots i_p} \\
= &
\sum\nolimits_{r = 0}^p
\epsilon(n, m, p, r)
(-1)^{\deg(\alpha_{i_0 \ldots i_r})\deg(\beta_{i_r \ldots i_p})}
\beta_{i_r \ldots i_p} \wedge \alpha_{i_0 \ldots i_r}
\end{align*}
となる。ここで $\wedge$ の次数付き可換性を用いた。一方，
\begin{align*}
\tau(\wedge^\bullet(\tau(\beta) \cup \tau(\alpha)))_{i_0 \ldots i_p}
= &
\chi(p)
\sum\nolimits_{r = 0}^p
\epsilon(m, n, p, r)
\tau(\beta)_{i_p \ldots i_{p - r}} \wedge \tau(\alpha)_{i_{p - r} \ldots i_0}
\\
= &
\chi(p)
\sum\nolimits_{r = 0}^p
\epsilon(m, n, p, r) \chi(r) \chi(p - r)
\beta_{i_{p - r} \ldots i_p} \wedge \alpha_{i_0 \ldots i_{p - r}}
\\
= &
\chi(p)
\sum\nolimits_{r = 0}^p
\epsilon(m, n, p, p - r) \chi(r) \chi(p - r)
\beta_{i_r \ldots i_p} \wedge \alpha_{i_0 \ldots i_r}
\end{align*}
である。ここで $\chi(t) = (-1)^{\frac{t(t + 1)}{2}}$ である。
$\tau$ がコホモロジー上で恒等写像として作用することは既に証明したので，
$$
\epsilon(n, m, p, r)
(-1)^{(n - r)(m - (p - r))}
=
(-1)^{nm} \chi(p)\epsilon(m, n, p, p - r) \chi(r) \chi(p - r)
$$
を確認すればよい。法 $2$ の自明な計算により，これらの符号が一致する
ことが分かる。（代わりに，Stacks project のサブディレクトリ scripts に
あるスクリプト signs.gp を確認してもよい。）

\medskip\noindent
最後に，カップ積と境界写像の両立性を調べる。次が
$$
0
\to
{\mathcal F}_1^\bullet
\to
{\mathcal F}_2^\bullet
\to
{\mathcal F}_3^\bullet
\to
0
\quad\text{and}\quad
0
\leftarrow
{\mathcal G}_1^\bullet
\leftarrow
{\mathcal G}_2^\bullet
\leftarrow
{\mathcal G}_3^\bullet
\leftarrow
0
$$
$X$ 上のアーベル層からなる下に有界な複体の短完全系列であると仮定する。
${\mathcal H}^\bullet$ をアーベル層からなるもう一つの下に有界な複体とし，
複体の写像
$$
\gamma_i :
\text{Tot}({\mathcal F}_i^\bullet \otimes_{\mathbf Z} {\mathcal G}_i^\bullet)
\longrightarrow
{\mathcal H}^\bullet
$$
が複体間の写像と両立すると仮定する。すなわち，図式
$$
\xymatrix{
\text{Tot}({\mathcal F}_1^\bullet \otimes_{\mathbf Z} {\mathcal G}_1^\bullet)
\ar[d]_{\gamma_1}
&
\text{Tot}({\mathcal F}_1^\bullet \otimes_{\mathbf Z} {\mathcal G}_2^\bullet)
\ar[l] \ar[d]
\\
\mathcal{H}^\bullet &
\text{Tot}({\mathcal F}_2^\bullet \otimes_{\mathbf Z} {\mathcal G}_2^\bullet)
\ar[l]_-{\gamma_2}
}
$$
および
$$
\xymatrix{
\text{Tot}({\mathcal F}_2^\bullet \otimes_{\mathbf Z} {\mathcal G}_2^\bullet)
\ar[d]_{\gamma_2}
&
\text{Tot}({\mathcal F}_2^\bullet \otimes_{\mathbf Z} {\mathcal G}_3^\bullet)
\ar[l] \ar[d]
\\
\mathcal{H}^\bullet &
\text{Tot}({\mathcal F}_3^\bullet \otimes_{\mathbf Z} {\mathcal G}_3^\bullet)
\ar[l]_-{\gamma_3}
}
$$
が可換であると仮定する。

\begin{lemma}
\label{cohomology-lemma-compute-sign-cup-product-boundaries}
上の状況で，層 $\mathcal{F}_i^p$ と $\mathcal{G}_j^q$ に対して
{\v C}ech コホモロジーがコホモロジーと一致すると仮定する。
$a_3 \in H^n(X, \mathcal{F}_3^\bullet)$ および
$b_1 \in H^m(X, \mathcal{G}_1^\bullet)$ とする。このとき
$$
\gamma_1( \partial a_3 \cup b_1) =
(-1)^{n + 1} \gamma_3( a_3 \cup \partial b_1)
$$
が $H^{n + m + 1}(X, \mathcal{H}^\bullet)$ において成り立つ。ここで
$\partial$ は，上の複体の短完全系列に付随するコホモロジー上の境界写像を
表す。
\end{lemma}

\begin{proof}
次の規約と記法を用いる。${\mathcal F}_1^p$ を
${\mathcal F}_2^p$ の部分層とみなし，${\mathcal G}_3^q$ を
${\mathcal G}_2^q$ の部分層とみなす。したがって，$s$ が
${\mathcal F}_1^p$ の局所切断ならば，${\mathcal F}_2^p$ の対応する切断も
$s$ と書く。${\mathcal G}_3^q$ の局所切断についても同様である。さらに，
$s$ が ${\mathcal F}_2^p$ の局所切断ならば，その ${\mathcal F}_3^p$ に
おける像を $\bar s$ と書く。写像
${\mathcal G}_2^q \to {\mathcal G}^q_1$ についても同様である。特に，
$s$ が ${\mathcal F}_2^p$ の局所切断で $\bar s = 0$ ならば，$s$ は
${\mathcal F}_1^p$ の局所切断である。上の図式の可換性から，
${\mathcal F}_2^p$ の局所切断 $s$ と ${\mathcal G}_3^q$ の局所切断 $t$
に対して，${\mathcal H}^{p + q}$ の切断として
$\gamma_2(s \otimes t) = \gamma_3(\bar s \otimes t)$ が成り立つ。

\medskip\noindent
${\mathcal U} : X =  \bigcup_{i \in I} U_i$ を $X$ の開被覆とする。
$\alpha_3$ と $\beta_1$ を，それぞれ
$\text{Tot}(
\check{\mathcal{C}}^\bullet({\mathcal U}, {\mathcal F}_3^\bullet))$,
および $\text{Tot}(
\check{\mathcal{C}}^\bullet({\mathcal U}, {\mathcal G}_1^\bullet))$
の次数 $n$ と次数 $m$ のコサイクルで，それぞれ $a_3$ と $b_1$ を表す
ものとする。必要なら $\mathcal{U}$ を細分した後，それぞれ
$\text{Tot}(
\check{\mathcal{C}}^\bullet({\mathcal U}, {\mathcal F}_2^\bullet))$,
および $\text{Tot}(
\check{\mathcal{C}}^\bullet({\mathcal U}, {\mathcal G}_2^\bullet))$
に属する次数 $n$ と次数 $m$ の余鎖 $\alpha_2$ と $\beta_2$ で，それぞれ
$\alpha_3$ と $\beta_1$ へ写るものを見つけられる。このとき
$$
\overline{d(\alpha_2)} = d(\bar \alpha_2) = 0
\quad\text{and}\quad
\overline{d(\beta_2)} = d(\bar \beta_2) = 0.
$$
である。これは，$\alpha_1 = d(\alpha_2)$ が
$\text{Tot}(\check{\mathcal{C}}^\bullet({\mathcal U}, {\mathcal F}_1^\bullet))$
における次数 $n + 1$ のコサイクルで，$\partial a_3$ を表すことを意味する。
同様に，$\beta_3 = d(\beta_2)$ は
$\text{Tot}(\check{\mathcal{C}}^\bullet({\mathcal U}, {\mathcal G}_3^\bullet))$
における次数 $m + 1$ のコサイクルで，$\partial b_1$ を表す。したがって
\begin{align*}
d(\gamma_2(\alpha_2 \cup \beta_2))
& =
\gamma_2(d(\alpha_2 \cup \beta_2))
\\
& =
\gamma_2(d(\alpha_2) \cup \beta_2 + (-1)^n \alpha_2 \cup d(\beta_2) )
\\
& =
\gamma_2( \alpha_1 \cup \beta_2)  + (-1)^n \gamma_2( \alpha_2 \cup \beta_3)
\\
& =
\gamma_1(\alpha_1 \cup \beta_1) + (-1)^n \gamma_3(\alpha_3 \cup \beta_3)
\end{align*}
と計算できる。これは，符号が補題に示したとおり $(-1)^{n + 1}$ である
ことまで教える\footnote{符号は，$D(X)$ の三角形に付随するコホモロジーの
長完全系列における符号の規約に依存する。Stacks project の規約では，
(a) 識別三角形は項ごとに分裂する完全系列に対応し，(b) 長完全系列の
境界写像は，符号を介在させずに蛇の補題の写像で与えられる。
《導来圏》節 \ref{derived-section-homotopy-triangulated} を参照せよ。}。
\end{proof}

\begin{lemma}
\label{cohomology-lemma-boundary-derivation-over-cup-product}
$X$ を位相空間とする。$\mathcal{O}' \to \mathcal{O}$ を環の層の全射で，
その核 $\mathcal{I} \subset \mathcal{O}'$ の平方が零であるものとする。
このとき $M = H^1(X, \mathcal{I})$ は
$R = H^0(X, \mathcal{O})$ 加群であり，短完全系列
$$
0 \to \mathcal{I} \to \mathcal{O}' \to \mathcal{O} \to 0
$$
に付随する境界写像 $\partial : R \to M$ は導分である
（《代数》定義 \ref{algebra-definition-derivation}）。
\end{lemma}

\begin{proof}
仮定により $\mathcal{I} \cdot \mathcal{I} = 0$ なので，写像
$\mathcal{O}' \to \SheafHom(\mathcal{I}, \mathcal{I})$ は
$\mathcal{O}$ を経由する。したがって $\mathcal{I}$ は
$\mathcal{O}$ 加群の層であり，これが $M$ 上の $R$ 加群構造を定める。
境界写像は加法的なので，Leibniz 則を証明すれば十分である。
$f \in R$ とする。$f|_{U_i} \in \mathcal{O}(U_i)$ の持ち上げ
$f_i \in \mathcal{O}'(U_i)$ が存在するような開被覆
$\mathcal{U} : X = \bigcup U_i$ を選ぶ。
$f_i - f_j$ は $\mathcal{I}(U_i \cap U_j)$ の元であることに注意する。
すると $\partial(f)$ は，
$\alpha_{i_0i_1} = f_{i_0} - f_{i_1}$ を満たす $1$-コサイクル $\alpha$ の
{\v C}ech コホモロジー類に対応する。（補題 \ref{cohomology-lemma-cech-h1} により，
$\mathcal{U}$ に関する第1 {\v C}ech コホモロジー群は $M$ の部分加群である
ことに注意する。）次に，$g \in R$ をもう一つの元とし，（必要なら開被覆を
細分した後）$g_i \in \mathcal{O}'(U_i)$ が
$g|_{U_i} \in \mathcal{O}(U_i)$ を持ち上げると仮定する。このとき
$\partial(g)$ は，$\beta_{i_0i_1} = g_{i_0} - g_{i_1}$ を満たすコサイクル
$\beta$ で与えられる。$f_ig_i \in \mathcal{O}'(U_i)$ は
$fg|_{U_i}$ を持ち上げるので，$\partial(fg)$ は，次を満たすコサイクル
$\gamma$ で与えられる：
$$
\gamma_{i_0i_1} = f_{i_0}g_{i_0} - f_{i_1}g_{i_1} =
(f_{i_0} - f_{i_1})g_{i_0} + f_{i_1}(g_{i_0} - g_{i_1}) =
\alpha_{i_0i_1}g + f\beta_{i_0i_1}
$$
ここでは $\mathcal{I}$ 上の $\mathcal{O}$ 加群構造の定義を用いた。
これで Leibniz 則が証明され，証明が完了する。
\end{proof}









\section{平坦分解}
\label{cohomology-section-flat}

\noindent
本節の内容の参考文献は \cite{Spaltenstein} である。
$(X, \mathcal{O}_X)$ を環付き空間とする。《加群》補題
\ref{modules-lemma-module-quotient-flat} により，任意の
$\mathcal{O}_X$ 加群は平坦な $\mathcal{O}_X$ 加群の商である。
《導来圏》補題 \ref{derived-lemma-subcategory-left-resolution} により，
$\mathcal{O}_X$ 加群の任意の上に有界な複体は，平坦な
$\mathcal{O}_X$ 加群からなる上に有界な複体による左分解をもつ。
しかし非有界複体に対しては，平坦分解だけでは十分でないことが分かる。

\begin{lemma}
\label{cohomology-lemma-derived-tor-exact}
$(X, \mathcal{O}_X)$ を環付き空間とする。
$\mathcal{G}^\bullet$ を $\mathcal{O}_X$ 加群の複体とする。関手
$$
K(\textit{Mod}(\mathcal{O}_X))
\longrightarrow
K(\textit{Mod}(\mathcal{O}_X)),
\quad
\mathcal{F}^\bullet \longmapsto
\text{Tot}(\mathcal{G}^\bullet \otimes_{\mathcal{O}_X} \mathcal{F}^\bullet)
$$
および
$$
K(\textit{Mod}(\mathcal{O}_X))
\longrightarrow
K(\textit{Mod}(\mathcal{O}_X)),
\quad
\mathcal{F}^\bullet \longmapsto
\text{Tot}(\mathcal{F}^\bullet \otimes_{\mathcal{O}_X} \mathcal{G}^\bullet)
$$
はいずれも三角圏の完全関手である。
\end{lemma}

\begin{proof}
《導来圏》注意 \ref{derived-remark-double-complex-as-tensor-product-of}
から従う。
\end{proof}

\begin{definition}
\label{cohomology-definition-K-flat}
$(X, \mathcal{O}_X)$ を環付き空間とする。$\mathcal{O}_X$ 加群の複体
$\mathcal{K}^\bullet$ が {\it K-平坦}であるとは，任意の
$\mathcal{O}_X$ 加群の非輪状複体 $\mathcal{F}^\bullet$ に対して，複体
$$
\text{Tot}(\mathcal{F}^\bullet \otimes_{\mathcal{O}_X} \mathcal{K}^\bullet)
$$
が非輪状であることをいう。
\end{definition}

\begin{lemma}
\label{cohomology-lemma-K-flat-quasi-isomorphism}
$(X, \mathcal{O}_X)$ を環付き空間とする。
$\mathcal{K}^\bullet$ を K-平坦複体とする。このとき関手
$$
K(\textit{Mod}(\mathcal{O}_X))
\longrightarrow
K(\textit{Mod}(\mathcal{O}_X)), \quad
\mathcal{F}^\bullet
\longmapsto
\text{Tot}(\mathcal{F}^\bullet \otimes_{\mathcal{O}_X} \mathcal{K}^\bullet)
$$
は擬同型を擬同型へ写す。
\end{lemma}

\begin{proof}
補題 \ref{cohomology-lemma-derived-tor-exact} と，擬同型がその錐の非輪状性によって
特徴づけられることから従う。
\end{proof}

\begin{lemma}
\label{cohomology-lemma-check-K-flat-stalks}
$(X, \mathcal{O}_X)$ を環付き空間とし，$\mathcal{K}^\bullet$ を
$\mathcal{O}_X$ 加群の複体とする。このとき $\mathcal{K}^\bullet$ が
K-平坦であるための必要十分条件は，すべての $x \in X$ に対して
$\mathcal{O}_{X, x}$ 加群の複体 $\mathcal{K}_x^\bullet$ が K-平坦である
ことである（《代数続論》定義 \ref{more-algebra-definition-K-flat}）。
\end{lemma}

\begin{proof}
すべての $x \in X$ に対して $\mathcal{K}_x^\bullet$ が K-平坦ならば，
$\otimes$ と直和は茎を取ることと可換であり，完全性は茎で確認できるので，
$\mathcal{K}^\bullet$ は K-平坦である。《加群》補題
\ref{modules-lemma-abelian} を参照せよ。逆に，$\mathcal{K}^\bullet$ が
K-平坦であると仮定する。$x \in X$ を取り，$M^\bullet$ を
$\mathcal{O}_{X, x}$ 加群の非輪状複体とする。このとき
$i_{x, *}M^\bullet$ は $\mathcal{O}_X$ 加群の非輪状複体である。したがって
$\text{Tot}(i_{x, *}M^\bullet \otimes_{\mathcal{O}_X} \mathcal{K}^\bullet)$
は非輪状である。$x$ における茎を取れば，
$\text{Tot}(M^\bullet \otimes_{\mathcal{O}_{X, x}} \mathcal{K}_x^\bullet)$
が非輪状であることが分かる。
\end{proof}

\begin{lemma}
\label{cohomology-lemma-tensor-product-K-flat}
$(X, \mathcal{O}_X)$ を環付き空間とする。$\mathcal{K}^\bullet$ と
$\mathcal{L}^\bullet$ が $\mathcal{O}_X$ 加群の K-平坦複体ならば，
$\text{Tot}(\mathcal{K}^\bullet \otimes_{\mathcal{O}_X} \mathcal{L}^\bullet)$
は $\mathcal{O}_X$ 加群の K-平坦複体である。
\end{lemma}

\begin{proof}
同型
$$
\text{Tot}(\mathcal{M}^\bullet \otimes_{\mathcal{O}_X}
\text{Tot}(\mathcal{K}^\bullet \otimes_{\mathcal{O}_X} \mathcal{L}^\bullet))
=
\text{Tot}(\text{Tot}(\mathcal{M}^\bullet \otimes_{\mathcal{O}_X}
\mathcal{K}^\bullet) \otimes_{\mathcal{O}_X} \mathcal{L}^\bullet)
$$
と定義から従う。
\end{proof}

\begin{lemma}
\label{cohomology-lemma-K-flat-two-out-of-three}
$(X, \mathcal{O}_X)$ を環付き空間とする。
$(\mathcal{K}_1^\bullet, \mathcal{K}_2^\bullet, \mathcal{K}_3^\bullet)$ を
$K(\textit{Mod}(\mathcal{O}_X))$ の識別三角形とする。
$\mathcal{K}_i^\bullet$ の三つのうち二つが K-平坦ならば，残る一つも
K-平坦である。
\end{lemma}

\begin{proof}
補題 \ref{cohomology-lemma-derived-tor-exact} と，
$K(\textit{Mod}(\mathcal{O}_X))$ の識別三角形において三つのうち二つが
非輪状ならば残る一つも非輪状であることから従う。
\end{proof}

\begin{lemma}
\label{cohomology-lemma-K-flat-two-out-of-three-ses}
$(X, \mathcal{O}_X)$ を環付き空間とする。
$0 \to \mathcal{K}_1^\bullet \to \mathcal{K}_2^\bullet \to
\mathcal{K}_3^\bullet \to 0$ を複体の短完全系列で，
$\mathcal{K}_3^\bullet$ の各項が平坦な $\mathcal{O}_X$ 加群であるものとする。
$\mathcal{K}_i^\bullet$ の三つのうち二つが K-平坦ならば，残る一つも
K-平坦である。
\end{lemma}

\begin{proof}
《加群》補題 \ref{modules-lemma-flat-tor-zero} により，任意の複体
$\mathcal{L}^\bullet$ に対して複体の短完全系列
$$
0 \to
\text{Tot}(\mathcal{L}^\bullet \otimes_{\mathcal{O}_X} \mathcal{K}_1^\bullet)
\to
\text{Tot}(\mathcal{L}^\bullet \otimes_{\mathcal{O}_X} \mathcal{K}_1^\bullet)
\to
\text{Tot}(\mathcal{L}^\bullet \otimes_{\mathcal{O}_X} \mathcal{K}_1^\bullet)
\to 0
$$
を得る。したがって，コホモロジー層の長完全系列と K-平坦複体の定義から
補題が従う。
\end{proof}

\begin{lemma}
\label{cohomology-lemma-pullback-K-flat}
$f : (X, \mathcal{O}_X) \to (Y, \mathcal{O}_Y)$ を環付き空間の射とする。
$\mathcal{O}_Y$ 加群の K-平坦複体の引き戻しは，$\mathcal{O}_X$ 加群の
K-平坦複体である。
\end{lemma}

\begin{proof}
これは茎で確認できる。補題 \ref{cohomology-lemma-check-K-flat-stalks} を参照せよ。
したがって，《層》補題 \ref{sheaves-lemma-stalk-pullback-modules} と
《代数続論》補題 \ref{more-algebra-lemma-base-change-K-flat} から従う。
\end{proof}

\begin{lemma}
\label{cohomology-lemma-bounded-flat-K-flat}
$(X, \mathcal{O}_X)$ を環付き空間とする。平坦な $\mathcal{O}_X$ 加群から
なる上に有界な複体は K-平坦である。
\end{lemma}

\begin{proof}
これは茎で確認できる。補題 \ref{cohomology-lemma-check-K-flat-stalks} を参照せよ。
したがってこの補題は，《加群》補題 \ref{modules-lemma-flat-stalks-flat}
と《代数続論》補題
\ref{more-algebra-lemma-derived-tor-quasi-isomorphism} から従う。
\end{proof}

\noindent
次の補題では，複体の系の余極限とは項ごとの余極限を意味する。

\begin{lemma}
\label{cohomology-lemma-colimit-K-flat}
$(X, \mathcal{O}_X)$ を環付き空間とする。
$\mathcal{K}_1^\bullet \to \mathcal{K}_2^\bullet \to \ldots$ を
K-平坦複体の系とする。このとき $\colim_i \mathcal{K}_i^\bullet$ は
K-平坦である。
\end{lemma}

\begin{proof}
項ごとの余極限を取っているので，明らかに
$$
\colim_i \text{Tot}(
\mathcal{F}^\bullet \otimes_{\mathcal{O}_X} \mathcal{K}_i^\bullet)
=
\text{Tot}(\mathcal{F}^\bullet \otimes_{\mathcal{O}_X}
\colim_i \mathcal{K}_i^\bullet)
$$
である。したがって，フィルター余極限が完全であることから補題が従う。
\end{proof}

\begin{lemma}
\label{cohomology-lemma-resolution-by-direct-sums-extensions-by-zero}
$(X, \mathcal{O}_X)$ を環付き空間とする。任意の $\mathcal{O}_X$ 加群の複体
$\mathcal{G}^\bullet$ に対して，$\mathcal{O}_X$ 加群の複体の可換図式
$$
\xymatrix{
\mathcal{K}_1^\bullet \ar[d] \ar[r] &
\mathcal{K}_2^\bullet \ar[d] \ar[r] & \ldots \\
\tau_{\leq 1}\mathcal{G}^\bullet \ar[r] &
\tau_{\leq 2}\mathcal{G}^\bullet \ar[r] & \ldots
}
$$
で，次の性質をもつものが存在する：(1) 垂直矢印は擬同型であり，項ごとに
全射である。(2) 各 $\mathcal{K}_n^\bullet$ は上に有界な複体で，その各項は
$j_{U!}\mathcal{O}_U$ という形の $\mathcal{O}_X$ 加群の直和である。
(3) 写像 $\mathcal{K}_n^\bullet \to \mathcal{K}_{n + 1}^\bullet$ は項ごとに
分裂する単射で，その余核は $j_{U!}\mathcal{O}_U$ という形の
$\mathcal{O}_X$ 加群の直和である。さらに，写像
$\colim \mathcal{K}_n^\bullet \to \mathcal{G}^\bullet$ は擬同型である。
\end{lemma}

\begin{proof}
図式の存在と性質 (1)，(2)，(3) は，《加群》補題
\ref{modules-lemma-module-quotient-flat} と《導来圏》補題
\ref{derived-lemma-special-direct-system} から直ちに従う。フィルター余極限は
完全なので，誘導写像
$\colim \mathcal{K}_n^\bullet \to \mathcal{G}^\bullet$ は擬同型である。
\end{proof}

\begin{lemma}
\label{cohomology-lemma-K-flat-resolution}
$(X, \mathcal{O}_X)$ を環付き空間とする。任意の複体
$\mathcal{G}^\bullet$ に対して，各項が平坦な $\mathcal{O}_X$ 加群である
K-平坦複体 $\mathcal{K}^\bullet$ と，項ごとに全射な擬同型
$\mathcal{K}^\bullet \to \mathcal{G}^\bullet$ が存在する。
\end{lemma}

\begin{proof}
補題 \ref{cohomology-lemma-resolution-by-direct-sums-extensions-by-zero} のような図式を
選ぶ。各複体 $\mathcal{K}_n^\bullet$ は平坦加群からなる上に有界な複体で
ある。《加群》補題 \ref{modules-lemma-j-shriek-flat} を参照せよ。したがって
補題 \ref{cohomology-lemma-bounded-flat-K-flat} により $\mathcal{K}_n^\bullet$ は
K-平坦である。ゆえに補題 \ref{cohomology-lemma-colimit-K-flat} により
$\colim \mathcal{K}_n^\bullet$ は K-平坦である。構成により，誘導写像
$\colim \mathcal{K}_n^\bullet \to \mathcal{G}^\bullet$ は擬同型であり，
項ごとに全射である。補題
\ref{cohomology-lemma-resolution-by-direct-sums-extensions-by-zero} の性質 (3) から，
$\colim \mathcal{K}_n^m$ は平坦加群の直和であり，したがって平坦である。
これで最後の主張も証明された。
\end{proof}

\begin{lemma}
\label{cohomology-lemma-derived-tor-quasi-isomorphism-other-side}
$(X, \mathcal{O}_X)$ を環付き空間とする。
$\alpha : \mathcal{P}^\bullet \to \mathcal{Q}^\bullet$ を，
$\mathcal{O}_X$ 加群の K-平坦複体の擬同型とする。任意の
$\mathcal{O}_X$ 加群の複体 $\mathcal{F}^\bullet$ に対して，誘導写像
$$
\text{Tot}(\text{id}_{\mathcal{F}^\bullet} \otimes \alpha) :
\text{Tot}(\mathcal{F}^\bullet \otimes_{\mathcal{O}_X} \mathcal{P}^\bullet)
\longrightarrow
\text{Tot}(\mathcal{F}^\bullet \otimes_{\mathcal{O}_X} \mathcal{Q}^\bullet)
$$
は擬同型である。
\end{lemma}

\begin{proof}
$\mathcal{K}^\bullet$ が K-平坦複体であるような擬同型
$\mathcal{K}^\bullet \to \mathcal{F}^\bullet$ を選ぶ。補題
\ref{cohomology-lemma-K-flat-resolution} を参照せよ。可換図式
$$
\xymatrix{
\text{Tot}(\mathcal{K}^\bullet
\otimes_{\mathcal{O}_X} \mathcal{P}^\bullet) \ar[r] \ar[d] &
\text{Tot}(\mathcal{K}^\bullet
\otimes_{\mathcal{O}_X} \mathcal{Q}^\bullet) \ar[d] \\
\text{Tot}(\mathcal{F}^\bullet
\otimes_{\mathcal{O}_X} \mathcal{P}^\bullet) \ar[r] &
\text{Tot}(\mathcal{F}^\bullet
\otimes_{\mathcal{O}_X} \mathcal{Q}^\bullet)
}
$$
を考える。補題 \ref{cohomology-lemma-K-flat-quasi-isomorphism} により，垂直矢印と上の
水平矢印は擬同型なので，結果が従う。
\end{proof}

\noindent
$(X, \mathcal{O}_X)$ を環付き空間とする。
$\mathcal{F}^\bullet$ を $D(\mathcal{O}_X)$ の対象とする。
K-平坦分解 $\mathcal{K}^\bullet \to \mathcal{F}^\bullet$ を選ぶ。補題
\ref{cohomology-lemma-K-flat-resolution} を参照せよ。補題
\ref{cohomology-lemma-derived-tor-exact} により，三角圏の完全関手
$$
K(\mathcal{O}_X)
\longrightarrow
K(\mathcal{O}_X),
\quad
\mathcal{G}^\bullet
\longmapsto
\text{Tot}(\mathcal{G}^\bullet \otimes_{\mathcal{O}_X} \mathcal{K}^\bullet)
$$
を得る。補題 \ref{cohomology-lemma-K-flat-quasi-isomorphism} により，この関手は関手
$D(\mathcal{O}_X) \to D(\mathcal{O}_X)$ を誘導する。実際，
$D(\mathcal{O}_X)$ は $K(\mathcal{O}_X)$ を擬同型で局所化したものだからで
ある。$\mathcal{F}^\bullet$ の K-平坦分解の圏は余フィルター圏であり，
補題 \ref{cohomology-lemma-derived-tor-quasi-isomorphism-other-side} があるので，得られる
関手は（同型を除いて）$\mathcal{K}^\bullet$ の選択に依存しない。

\begin{definition}
\label{cohomology-definition-derived-tor}
$(X, \mathcal{O}_X)$ を環付き空間とする。
$\mathcal{F}^\bullet$ を $D(\mathcal{O}_X)$ の対象とする。上で記述した
三角圏の完全関手
$$
- \otimes_{\mathcal{O}_X}^{\mathbf{L}} \mathcal{F}^\bullet :
D(\mathcal{O}_X)
\longrightarrow
D(\mathcal{O}_X)
$$
を {\it 導来テンソル積}という。
\end{definition}

\noindent
明示的な構成から，$D(\mathcal{O}_X)$ の
$\mathcal{G}^\bullet$ と $\mathcal{F}^\bullet$ に対して標準同型
$$
\mathcal{F}^\bullet \otimes_{\mathcal{O}_X}^{\mathbf{L}} \mathcal{G}^\bullet
\cong
\mathcal{G}^\bullet \otimes_{\mathcal{O}_X}^{\mathbf{L}} \mathcal{F}^\bullet
$$
があることは明らかである。ここでは《代数続論》節
\ref{more-algebra-section-sign-rules} で与えられる符号規則を用いる。
したがって
$\mathcal{F}^\bullet \otimes_{\mathcal{O}_X}^{\mathbf{L}} \mathcal{G}^\bullet$
と書くとき，通常はどちらの変数を用いて導来テンソル積を定義するかを
区別しない。

\begin{definition}
\label{cohomology-definition-tor}
$(X, \mathcal{O}_X)$ を環付き空間とする。
$\mathcal{F}$，$\mathcal{G}$ を $\mathcal{O}_X$ 加群とする。
$\mathcal{F}$ と $\mathcal{G}$ の {\it Tor} を，公式
$$
\text{Tor}_p^{\mathcal{O}_X}(\mathcal{F}, \mathcal{G}) =
H^{-p}(\mathcal{F} \otimes_{\mathcal{O}_X}^\mathbf{L} \mathcal{G})
$$
によって定義する。導来テンソル積は上で定義したものを用いる。
\end{definition}

\noindent
この定義から，任意の $\mathcal{O}_X$ 加群の短完全系列
$0 \to \mathcal{F}_1 \to \mathcal{F}_2 \to \mathcal{F}_3 \to 0$
と任意の $\mathcal{O}_X$ 加群 $\mathcal{G}$ に対して，コホモロジーの
長完全系列
$$
\xymatrix{
\mathcal{F}_1 \otimes_{\mathcal{O}_X} \mathcal{G} \ar[r] &
\mathcal{F}_2 \otimes_{\mathcal{O}_X} \mathcal{G} \ar[r] &
\mathcal{F}_3 \otimes_{\mathcal{O}_X} \mathcal{G} \ar[r] & 0 \\
\text{Tor}_1^{\mathcal{O}_X}(\mathcal{F}_1, \mathcal{G}) \ar[r] &
\text{Tor}_1^{\mathcal{O}_X}(\mathcal{F}_2, \mathcal{G}) \ar[r] &
\text{Tor}_1^{\mathcal{O}_X}(\mathcal{F}_3, \mathcal{G}) \ar[ull]
}
$$
が得られる。これをこの状況に付随する $\text{Tor}$ の長完全系列という。

\begin{lemma}
\label{cohomology-lemma-flat-tor-zero}
\begin{slogan}
Tor は平坦性からのずれを測る。
\end{slogan}
$(X, \mathcal{O}_X)$ を環付き空間とする。
$\mathcal{F}$ を $\mathcal{O}_X$ 加群とする。次は同値である：
\begin{enumerate}
\item $\mathcal{F}$ は平坦な $\mathcal{O}_X$ 加群である。
\item $\text{Tor}_1^{\mathcal{O}_X}(\mathcal{F}, \mathcal{G}) = 0$
が任意の $\mathcal{O}_X$ 加群 $\mathcal{G}$ に対して成り立つ。
\end{enumerate}
\end{lemma}

\begin{proof}
もし $\mathcal{F}$ が平坦ならば，$\mathcal{F} \otimes_{\mathcal{O}_X} -$
は完全関手であり，衛星関手は消滅する。逆に (2) が成り立つと仮定する。
$\mathcal{G} \to \mathcal{H}$ が単射で，余核が $\mathcal{Q}$ ならば，
$\text{Tor}$ の長完全系列から，
$\mathcal{F} \otimes_{\mathcal{O}_X} \mathcal{G} \to
\mathcal{F} \otimes_{\mathcal{O}_X} \mathcal{H}$
の核は
$\text{Tor}_1^{\mathcal{O}_X}(\mathcal{F}, \mathcal{Q})$
の商であることが分かる。これは仮定により零である。したがって
$\mathcal{F}$ は平坦である。
\end{proof}

\begin{lemma}
\label{cohomology-lemma-factor-through-K-flat}
$(X, \mathcal{O}_X)$ を環付き空間とする。
$a : \mathcal{K}^\bullet \to \mathcal{L}^\bullet$ を
$\mathcal{O}_X$ 加群の複体の写像とする。$\mathcal{K}^\bullet$ が K-平坦
ならば，複体 $\mathcal{N}^\bullet$ と複体の写像
$b : \mathcal{K}^\bullet \to \mathcal{N}^\bullet$
および $c : \mathcal{N}^\bullet \to \mathcal{L}^\bullet$ で，次を満たす
ものが存在する：
\begin{enumerate}
\item $\mathcal{N}^\bullet$ は K-平坦である。
\item $c$ は擬同型である。
\item $a$ は $c \circ b$ とホモトピックである。
\end{enumerate}
$\mathcal{K}^\bullet$ の各項が平坦ならば，$\mathcal{N}^\bullet$ についても
同じことが成り立つように $\mathcal{N}^\bullet$，$b$，$c$ を選べる。
\end{lemma}

\begin{proof}
$K(\textit{Mod}(\mathcal{O}_X))$ が三角圏であることを用いる。
《導来圏》命題 \ref{derived-proposition-homotopy-category-triangulated} を
参照せよ。識別三角形
$\mathcal{K}^\bullet \to \mathcal{L}^\bullet \to
\mathcal{C}^\bullet \to \mathcal{K}^\bullet[1]$.
を選ぶ。$\mathcal{M}^\bullet$ が平坦な項をもつ K-平坦複体であるような
擬同型 $\mathcal{M}^\bullet \to \mathcal{C}^\bullet$ を選ぶ。補題
\ref{cohomology-lemma-K-flat-resolution} を参照せよ。三角圏の公理により，合成
$\mathcal{M}^\bullet \to \mathcal{C}^\bullet \to \mathcal{K}^\bullet[1]$
を識別三角形
$\mathcal{K}^\bullet \to \mathcal{N}^\bullet \to
\mathcal{M}^\bullet \to \mathcal{K}^\bullet[1]$.
に組み込める。補題 \ref{cohomology-lemma-K-flat-two-out-of-three} により，
$\mathcal{N}^\bullet$ は K-平坦である。再び三角圏の公理を用いると，
次の識別三角形の射に組み込まれる写像
$\mathcal{N}^\bullet \to \mathcal{L}^\bullet$ を選べる：
$$
\xymatrix{
\mathcal{K}^\bullet \ar[r] \ar[d] &
\mathcal{N}^\bullet \ar[r] \ar[d] &
\mathcal{M}^\bullet \ar[r] \ar[d] &
\mathcal{K}^\bullet[1] \ar[d] \\
\mathcal{K}^\bullet \ar[r] &
\mathcal{L}^\bullet \ar[r] &
\mathcal{C}^\bullet \ar[r] &
\mathcal{K}^\bullet[1]
}
$$
三つの矢印のうち二つが擬同型なので，これらの識別三角形に付随する
コホモロジーの長完全系列により，第3の矢印
$\mathcal{N}^\bullet \to \mathcal{L}^\bullet$ も擬同型である
（好みに応じて，この図式の $D(\mathcal{O}_X)$ における像を見て，
《導来圏》補題 \ref{derived-lemma-third-isomorphism-triangle} を用いてもよい）。
これで (1)，(2)，(3) の証明が完了する。最後の主張を証明するため，
$\mathcal{N}^n \cong \mathcal{M}^n \oplus \mathcal{K}^n$ となるように
$\mathcal{N}^\bullet$ を選べる。《導来圏》補題
\ref{derived-lemma-improve-distinguished-triangle-homotopy} を参照せよ。
したがって，$\mathcal{K}^\bullet$ の各項が平坦ならば，望む平坦性を得る。
\end{proof}











\section{導来引き戻し}
\label{cohomology-section-derived-pullback}

\noindent
$f : (X, \mathcal{O}_X) \to (Y, \mathcal{O}_Y)$ を環付き空間の射とする。
K-平坦分解を用いて導来引き戻し関手
$$
Lf^* : D(\mathcal{O}_Y) \to D(\mathcal{O}_X)
$$
を定義できる。実際，任意の $\mathcal{O}_Y$ 加群の複体
$\mathcal{G}^\bullet$ に対して K-平坦分解
$\mathcal{K}^\bullet \to \mathcal{G}^\bullet$ を選び，
$Lf^*\mathcal{G}^\bullet = f^*\mathcal{K}^\bullet$ と置ける。
補題 \ref{cohomology-lemma-pullback-K-flat}，補題 \ref{cohomology-lemma-K-flat-resolution}，および
補題 \ref{cohomology-lemma-derived-tor-quasi-isomorphism-other-side} を用いれば，これが
よく定義されることが分かる。ただし細部を完全に処理するには，いくらかの
一般論を用いる方が便利かもしれない。

\begin{lemma}
\label{cohomology-lemma-derived-base-change}
上の構成は選択に依存せず，三角圏の完全関手
$Lf^* : D(\mathcal{O}_Y) \to D(\mathcal{O}_X)$ を定める。
\end{lemma}

\begin{proof}
これを見るため，《導来圏》節 \ref{derived-section-derived-functors} で
展開された一般論を用いる。
$\mathcal{D} = K(\mathcal{O}_Y)$ および
$\mathcal{D}' = D(\mathcal{O}_X)$ と置く。規則
$F(\mathcal{G}^\bullet) = f^*\mathcal{G}^\bullet$ で定まる三角圏の完全関手を
$F : \mathcal{D} \to \mathcal{D}'$ と書く。
$S$ を $\mathcal{D} = K(\mathcal{O}_Y)$ における擬同型の集合とする。
これにより《導来圏》状況 \ref{derived-situation-derived-functor} の状況を得る
ので，《導来圏》定義
\ref{derived-definition-right-derived-functor-defined} が適用できる。
$LF$ はいたるところ定義されると主張する。これは《導来圏》補題
\ref{derived-lemma-find-existence-computes} から従う。ここで
$\mathcal{P} \subset \Ob(\mathcal{D})$ を K-平坦複体の族とする。
(1) は補題 \ref{cohomology-lemma-K-flat-resolution} から従う。(2) を見るには，
$\mathcal{O}_Y$ 加群の K-平坦複体の間の擬同型
$\mathcal{K}_1^\bullet  \to \mathcal{K}_2^\bullet$ に対し，写像
$f^*\mathcal{K}_1^\bullet  \to f^*\mathcal{K}_2^\bullet$ が擬同型であることを
示さなければならない。この写像を
$$
f^{-1}\mathcal{K}_1^\bullet \otimes_{f^{-1}\mathcal{O}_Y} \mathcal{O}_X
\longrightarrow
f^{-1}\mathcal{K}_2^\bullet \otimes_{f^{-1}\mathcal{O}_Y} \mathcal{O}_X
$$
と書く。関手 $f^{-1}$ は完全なので，写像
$f^{-1}\mathcal{K}_1^\bullet  \to f^{-1}\mathcal{K}_2^\bullet$ は擬同型である。
射 $(X, f^{-1}\mathcal{O}_Y) \to (Y, \mathcal{O}_Y)$ に補題
\ref{cohomology-lemma-pullback-K-flat} を適用すると，複体
$f^{-1}\mathcal{K}_1^\bullet$ と $f^{-1}\mathcal{K}_2^\bullet$ は
$f^{-1}\mathcal{O}_Y$ 加群の K-平坦複体である。したがって，補題
\ref{cohomology-lemma-derived-tor-quasi-isomorphism-other-side} により，表示した写像は
擬同型である。ゆえに導来関手
$$
LF :
D(\mathcal{O}_Y) = S^{-1}\mathcal{D}
\longrightarrow
\mathcal{D}' = D(\mathcal{O}_X)
$$
を得る。《導来圏》式 (\ref{derived-equation-everywhere}) を参照せよ。
最後に，《導来圏》補題 \ref{derived-lemma-find-existence-computes} は，
$\mathcal{K}^\bullet$ が K-平坦なとき
$LF(\mathcal{K}^\bullet) = F(\mathcal{K}^\bullet) = f^*\mathcal{K}^\bullet$
であることも保証する。すなわち，$Lf^* = LF$ は確かに上で記述した方法で
計算される。
\end{proof}

\begin{lemma}
\label{cohomology-lemma-derived-pullback-composition}
$f : X \to Y$ および $g : Y \to Z$ を環付き空間の射とする。このとき関手
$D(\mathcal{O}_Z) \to D(\mathcal{O}_X)$ として
$Lf^* \circ Lg^* = L(g \circ f)^*$ である。
\end{lemma}

\begin{proof}
$E$ を $D(\mathcal{O}_Z)$ の対象とする。構成により，$Z$ 上で $E$ を表す
K-平坦複体 $\mathcal{K}^\bullet$ を選び，
$Lg^*E = g^*\mathcal{K}^\bullet$ と置くことで $Lg^*E$ が計算される。
補題 \ref{cohomology-lemma-pullback-K-flat} により，$g^*\mathcal{K}^\bullet$ は
$Y$ 上で K-平坦である。したがって $Lf^*Lg^*E$ は
$f^*g^*\mathcal{K}^\bullet = (g \circ f)^*\mathcal{K}^\bullet$
で与えられ，これは $L(g \circ f)^*E$ も表す。
\end{proof}

\begin{lemma}
\label{cohomology-lemma-pullback-tensor-product}
$f : (X, \mathcal{O}_X) \to (Y, \mathcal{O}_Y)$ を環付き空間の射とする。
$\mathcal{F}^\bullet, \mathcal{G}^\bullet \in \Ob(D(\mathcal{O}_Y))$ に対し，
二変数について関手的な標準同型
$$
Lf^*(
\mathcal{F}^\bullet \otimes_{\mathcal{O}_Y}^{\mathbf{L}} \mathcal{G}^\bullet
) =
Lf^*\mathcal{F}^\bullet 
\otimes_{\mathcal{O}_X}^{\mathbf{L}}
Lf^*\mathcal{G}^\bullet 
$$
がある。
\end{lemma}

\begin{proof}
${\mathcal F}^\bullet$ と ${\mathcal G}^\bullet$ が K-平坦複体であると
仮定してよい。この場合，
$\mathcal{F}^\bullet \otimes_{\mathcal{O}_Y}^{\mathbf{L}} \mathcal{G}^\bullet$
は二重複体
$\mathcal{F}^\bullet \otimes_{\mathcal{O}_Y} \mathcal{G}^\bullet$ に付随する
全複体にほかならない。補題 \ref{cohomology-lemma-tensor-product-K-flat} により，
$\text{Tot}(\mathcal{F}^\bullet \otimes_{\mathcal{O}_Y} \mathcal{G}^\bullet)$
も K-平坦である。したがって補題の同型は，同型
$$
\text{Tot}(f^*\mathcal{F}^\bullet \otimes_{\mathcal{O}_X}
f^*\mathcal{G}^\bullet)
\longrightarrow
f^*\text{Tot}(\mathcal{F}^\bullet \otimes_{\mathcal{O}_Y} \mathcal{G}^\bullet)
$$
から生じる。その成分は，《加群》補題
\ref{modules-lemma-tensor-product-pullback} の同型
$f^*\mathcal{F}^p \otimes_{\mathcal{O}_X} f^*\mathcal{G}^q \to
f^*(\mathcal{F}^p \otimes_{\mathcal{O}_Y} \mathcal{G}^q)$ である。
\end{proof}

\begin{lemma}
\label{cohomology-lemma-variant-derived-pullback}
$f : (X, \mathcal{O}_X) \to (Y, \mathcal{O}_Y)$ を環付き空間の射とする。
$D(\mathcal{O}_X)$ の $\mathcal{F}^\bullet$ と $D(\mathcal{O}_Y)$ の
$\mathcal{G}^\bullet$ に対し，二変数について関手的な標準同型
$$
\mathcal{F}^\bullet
\otimes_{\mathcal{O}_X}^{\mathbf{L}}
Lf^*\mathcal{G}^\bullet
=
\mathcal{F}^\bullet 
\otimes_{f^{-1}\mathcal{O}_Y}^{\mathbf{L}}
f^{-1}\mathcal{G}^\bullet 
$$
がある。
\end{lemma}

\begin{proof}
$\mathcal{F}$ を $\mathcal{O}_X$ 加群とし，$\mathcal{G}$ を
$\mathcal{O}_Y$ 加群とする。このとき
$\mathcal{F} \otimes_{\mathcal{O}_X} f^*\mathcal{G} =
\mathcal{F} \otimes_{f^{-1}\mathcal{O}_Y} f^{-1}\mathcal{G}$
である。実際，
$f^*\mathcal{G} =
\mathcal{O}_X \otimes_{f^{-1}\mathcal{O}_Y} f^{-1}\mathcal{G}$ である。
補題はこれと定義から従う。
\end{proof}

\begin{lemma}
\label{cohomology-lemma-tensor-pull-compatibility}
$f : (X, \mathcal{O}_X) \to (Y, \mathcal{O}_Y)$ を環付き空間の射とする。
$\mathcal{K}^\bullet$ と $\mathcal{M}^\bullet$ を $\mathcal{O}_Y$ 加群の
複体とする。図式
$$
\xymatrix{
Lf^*(\mathcal{K}^\bullet
\otimes_{\mathcal{O}_Y}^\mathbf{L}
\mathcal{M}^\bullet) \ar[r] \ar[d] &
Lf^*\text{Tot}(\mathcal{K}^\bullet
\otimes_{\mathcal{O}_Y}
\mathcal{M}^\bullet) \ar[d] \\
Lf^*\mathcal{K}^\bullet \otimes_{\mathcal{O}_X}^\mathbf{L}
Lf^*\mathcal{M}^\bullet \ar[d] &
f^*\text{Tot}(\mathcal{K}^\bullet
\otimes_{\mathcal{O}_Y}
\mathcal{M}^\bullet) \ar[d] \\
f^*\mathcal{K}^\bullet \otimes_{\mathcal{O}_X}^\mathbf{L}
f^*\mathcal{M}^\bullet \ar[r] &
\text{Tot}(f^*\mathcal{K}^\bullet \otimes_{\mathcal{O}_X}
f^*\mathcal{M}^\bullet)
}
$$
は可換である。
\end{lemma}

\begin{proof}
補題 \ref{cohomology-lemma-pullback-K-flat} のような K-平坦分解の存在を用いる。
そのような分解 $\mathcal{P}^\bullet \to \mathcal{K}^\bullet$ および
$\mathcal{Q}^\bullet \to \mathcal{M}^\bullet$ を選ぶと，
$$
\xymatrix{
Lf^*\text{Tot}(\mathcal{P}^\bullet
\otimes_{\mathcal{O}_Y}
\mathcal{Q}^\bullet) \ar[r] \ar[d] &
Lf^*\text{Tot}(\mathcal{K}^\bullet
\otimes_{\mathcal{O}_Y}
\mathcal{M}^\bullet) \ar[d] \\
f^*\text{Tot}(\mathcal{P}^\bullet
\otimes_{\mathcal{O}_Y}
\mathcal{Q}^\bullet) \ar[d] \ar[r] &
f^*\text{Tot}(\mathcal{K}^\bullet
\otimes_{\mathcal{O}_Y}
\mathcal{M}^\bullet) \ar[d] \\
\text{Tot}(f^*\mathcal{P}^\bullet \otimes_{\mathcal{O}_X}
f^*\mathcal{Q}^\bullet) \ar[r] &
\text{Tot}(f^*\mathcal{K}^\bullet \otimes_{\mathcal{O}_X}
f^*\mathcal{M}^\bullet)
}
$$
が可換であることが分かる。しかし今や，$\mathcal{P}^\bullet$ と
$\mathcal{Q}^\bullet$ の選択および補題 \ref{cohomology-lemma-tensor-product-K-flat}
により，この図式の左辺は図式の左辺である。
\end{proof}









\section{非有界複体のコホモロジー}
\label{cohomology-section-unbounded}

\noindent
$(X, \mathcal{O}_X)$ を環付き空間とする。圏
$\textit{Mod}(\mathcal{O}_X)$ は Grothendieck アーベル圏である。すなわち，
すべての余極限をもち，フィルター余極限は完全であり，さらに次の生成元をもつ：
$$
\bigoplus\nolimits_{U \subset X\text{ open}} j_{U!}\mathcal{O}_U,
$$
《加群》第 \ref{modules-section-kernels} 節および補題
\ref{modules-lemma-j-shriek-flat} と
\ref{modules-lemma-module-quotient-flat} を参照せよ。
《入射対象》定理
\ref{injectives-theorem-K-injective-embedding-grothendieck}
により，$\mathcal{O}_X$-加群の任意の複体 $\mathcal{F}^\bullet$ に対し，
単射な擬同型 $\mathcal{F}^\bullet \to \mathcal{I}^\bullet$ が存在する。
ここで $\mathcal{I}^\bullet$ は，すべての項が入射的
$\mathcal{O}_X$-加群である K-入射的な $\mathcal{O}_X$-加群複体であり，
しかもこの埋め込みは複体 $\mathcal{F}^\bullet$ に関して関手的に選べる。
《導来圏》補題 \ref{derived-lemma-enough-K-injectives-implies}
から次が従う。
\begin{enumerate}
\item 三角圏 $\mathcal{D}$ への任意の完全関手
$F : K(\textit{Mod}(\mathcal{O}_X)) \to \mathcal{D}$ は右導来関手
$RF : D(\mathcal{O}_X) \to \mathcal{D}$,
をもつ。
\item アーベル圏 $\mathcal{A}$ への任意の加法関手
$F : \textit{Mod}(\mathcal{O}_X) \to \mathcal{A}$
に対し，$F$ が誘導する完全関手
$F : K(\textit{Mod}(\mathcal{O}_X)) \to K(\mathcal{A})$ を考えると，
右導来関手
$RF : D(\mathcal{O}_X) \to D(\mathcal{A})$.
が得られる。
\end{enumerate}
構成により $RF(\mathcal{F}^\bullet) = F(\mathcal{I}^\bullet)$ である。
ここで $\mathcal{F}^\bullet \to \mathcal{I}^\bullet$ は上のものである。

\medskip\noindent
以上の例をいくつか挙げる。
\begin{enumerate}
\item 関手 $\Gamma(X, -) : \textit{Mod}(\mathcal{O}_X) \to
\text{Mod}_{\Gamma(X, \mathcal{O}_X)}$ から
$$
R\Gamma(X, -) : D(\mathcal{O}_X) \to D(\Gamma(X, \mathcal{O}_X))
$$
が得られる。コホモロジーには記法
$H^i(X, K) = H^i(R\Gamma(X, K))$ を用いる。
\item 開集合 $U \subset X$ に対し，関手
$\Gamma(U, -) : \textit{Mod}(\mathcal{O}_X) \to
\text{Mod}_{\Gamma(U, \mathcal{O}_X)}$ を考える。これから
$$
R\Gamma(U, -) : D(\mathcal{O}_X) \to D(\Gamma(U, \mathcal{O}_X))
$$
が得られる。コホモロジーには記法
$H^i(U, K) = H^i(R\Gamma(U, K))$ を用いる。
\item 環付き空間の射
$f : (X, \mathcal{O}_X) \to (Y, \mathcal{O}_Y)$ に対し，関手
$f_* : \textit{Mod}(\mathcal{O}_X) \to \textit{Mod}(\mathcal{O}_Y)$
を考える。これから非有界導来圏上の全導来順像
$$
Rf_* : D(\mathcal{O}_X) \longrightarrow D(\mathcal{O}_Y)
$$
が得られる。
\end{enumerate}

\begin{lemma}
\label{cohomology-lemma-adjoint}
$f : (X, \mathcal{O}_X) \to (Y, \mathcal{O}_Y)$ を環付き空間の射とする。
上で定義した関手 $Rf_*$ と，補題 \ref{cohomology-lemma-derived-base-change}
で定義した関手 $Lf^*$ は随伴である：
$$
\Hom_{D(\mathcal{O}_X)}(Lf^*\mathcal{G}^\bullet, \mathcal{F}^\bullet)
=
\Hom_{D(\mathcal{O}_Y)}(\mathcal{G}^\bullet, Rf_*\mathcal{F}^\bullet)
$$
この同型は $\mathcal{F}^\bullet \in \Ob(D(\mathcal{O}_X))$ と
$\mathcal{G}^\bullet \in \Ob(D(\mathcal{O}_Y))$ に関して双関手的である。
\end{lemma}

\begin{proof}
これは $Rf_*$ と $Lf^*$ が存在することから形式的に従う。
《導来圏》補題 \ref{derived-lemma-derived-adjoint-functors} を参照せよ。
\end{proof}

\begin{lemma}
\label{cohomology-lemma-derived-pushforward-composition}
$f : X \to Y$ と $g : Y \to Z$ を環付き空間の射とする。このとき
$D(\mathcal{O}_X) \to D(\mathcal{O}_Z)$ の関手として
$Rg_* \circ Rf_* = R(g \circ f)_*$ である。
\end{lemma}

\begin{proof}
補題 \ref{cohomology-lemma-adjoint} により，$Rg_* \circ Rf_*$ は
$Lf^* \circ Lg^*$ の随伴である。また補題
\ref{cohomology-lemma-derived-pullback-composition} により
$Lf^* \circ Lg^* = L(g \circ f)^*$ であるから，
随伴関手の一意性により $Rg_* \circ Rf_* = R(g \circ f)_*$ を得る。
\end{proof}

\begin{remark}
\label{cohomology-remark-base-change}
非有界導来関手 $Lf^*$ と $Rf_*$ の構成により，
完全な一般性の下で基底変換写像を構成できる。すなわち，
$$
\xymatrix{
X' \ar[r]_{g'} \ar[d]_{f'} &
X \ar[d]^f \\
S' \ar[r]^g &
S
}
$$
が環付き空間の可換図式であると仮定する。$K$ を
$D(\mathcal{O}_X)$ の対象とする。このとき標準的な基底変換写像
$$
Lg^*Rf_*K \longrightarrow R(f')_*L(g')^*K
$$
が $D(\mathcal{O}_{S'})$ に存在する。実際，この写像は写像
$L(f')^*Lg^*Rf_*K \to L(g')^*K$
に随伴する。$L(f')^*Lg^* = L(g')^*Lf^*$ であるから，これは写像
$L(g')^*Lf^*Rf_*K \to L(g')^*K$
と同じであり，後者として随伴写像 $Lf^*Rf_*K \to K$ に
$L(g')^*$ を施したものを取れる。
\end{remark}

\begin{remark}
\label{cohomology-remark-compose-base-change}
環付き空間の可換図式
$$
\xymatrix{
X' \ar[r]_k \ar[d]_{f'} & X \ar[d]^f \\
Y' \ar[r]^l \ar[d]_{g'} & Y \ar[d]^g \\
Z' \ar[r]^m & Z
}
$$
を考える。このとき二つの正方形に対する注意
\ref{cohomology-remark-base-change} の基底変換写像を合成すると，
外側の長方形に対する基底変換写像が得られる。より正確には，合成
\begin{align*}
Lm^* \circ R(g \circ f)_*
& =
Lm^* \circ Rg_* \circ Rf_* \\
& \to Rg'_* \circ Ll^* \circ Rf_* \\
& \to Rg'_* \circ Rf'_* \circ Lk^* \\
& = R(g' \circ f')_* \circ Lk^*
\end{align*}
がこの長方形に対する基底変換写像である。検証は省略する。
\end{remark}

\begin{remark}
\label{cohomology-remark-compose-base-change-horizontal}
環付き空間の可換図式
$$
\xymatrix{
X'' \ar[r]_{g'} \ar[d]_{f''} & X' \ar[r]_g \ar[d]_{f'} & X \ar[d]^f \\
Y'' \ar[r]^{h'} & Y' \ar[r]^h & Y
}
$$
を考える。このとき二つの正方形に対する注意
\ref{cohomology-remark-base-change} の基底変換写像を合成すると，
外側の長方形に対する基底変換写像が得られる。より正確には，合成
\begin{align*}
L(h \circ h')^* \circ Rf_*
& =
L(h')^* \circ Lh_* \circ Rf_* \\
& \to L(h')^* \circ Rf'_* \circ Lg^* \\
& \to Rf''_* \circ L(g')^* \circ Lg^* \\
& = Rf''_* \circ L(g \circ g')^*
\end{align*}
がこの長方形に対する基底変換写像である。検証は省略する。
\end{remark}

\begin{lemma}
\label{cohomology-lemma-adjoints-push-pull-compatibility}
$f : (X, \mathcal{O}_X) \to (Y, \mathcal{O}_Y)$ を環付き空間の射とし，
$\mathcal{K}^\bullet$ を $\mathcal{O}_X$-加群の複体とする。
複体上の $Lf^* \to f^*$，複体上の $f_* \to Rf_*$，および随伴
$Lf^* \circ Rf_* \to \text{id}$ から得られる図式
$$
\xymatrix{
Lf^*f_*\mathcal{K}^\bullet \ar[r] \ar[d] &
f^*f_*\mathcal{K}^\bullet \ar[d] \\
Lf^*Rf_*\mathcal{K}^\bullet \ar[r] &
\mathcal{K}^\bullet
}
$$
は $D(\mathcal{O}_X)$ で可換である。
\end{lemma}

\begin{proof}
K-平坦分解と K-入射分解の存在を用いる。補題
\ref{cohomology-lemma-pullback-K-flat} および上の議論を参照せよ。
$\mathcal{I}^\bullet$ が $\mathcal{O}_X$-加群の複体として K-入射的であるような
擬同型 $\mathcal{K}^\bullet \to \mathcal{I}^\bullet$ を選ぶ。
また，$\mathcal{Q}^\bullet$ が $\mathcal{O}_Y$-加群の複体として
K-平坦であるような擬同型
$\mathcal{Q}^\bullet \to f_*\mathcal{I}^\bullet$ を選ぶ。
$\mathcal{O}_Y$-加群の K-平坦複体 $\mathcal{P}^\bullet$ と，
複体の射の図式
$$
\xymatrix{
\mathcal{P}^\bullet \ar[r] \ar[d] &
f_*\mathcal{K}^\bullet \ar[d] \\
\mathcal{Q}^\bullet \ar[r] & f_*\mathcal{I}^\bullet
}
$$
を，ホモトピーを除いて可換で，かつ上の水平射が擬同型となるように選べる。
実際，複体のホモトピー圏において擬同型は乗法系をなすので，
まずある複体 $\mathcal{P}^\bullet$ に対してこのような図式を選び，
次に $\mathcal{P}^\bullet$ を K-平坦複体で置き換えればよい。
引き戻しを取ると，複体の射の図式
$$
\xymatrix{
f^*\mathcal{P}^\bullet \ar[r] \ar[d] &
f^*f_*\mathcal{K}^\bullet \ar[d] \ar[r] &
\mathcal{K}^\bullet \ar[d] \\
f^*\mathcal{Q}^\bullet \ar[r] &
f^*f_*\mathcal{I}^\bullet \ar[r] &
\mathcal{I}^\bullet
}
$$
を得る。これはホモトピーを除いて可換である。外側の長方形が，
補題の主張が成り立つことを示している。
\end{proof}

\begin{remark}
\label{cohomology-remark-cup-product}
$f : (X, \mathcal{O}_X) \to (Y, \mathcal{O}_Y)$ を環付き空間の射とする。
$Lf^*$ と $Rf_*$ の随伴性により，$D(\mathcal{O}_X)$ の任意の
$K, L$ に対して，$D(\mathcal{O}_Y)$ における相対カップ積
$$
Rf_*K \otimes_{\mathcal{O}_Y}^\mathbf{L} Rf_*L
\longrightarrow
Rf_*(K \otimes_{\mathcal{O}_X}^\mathbf{L} L)
$$
を構成できる。実際，この写像は写像
$Lf^*(Rf_*K \otimes_{\mathcal{O}_Y}^\mathbf{L} Rf_*L) \to
K \otimes_{\mathcal{O}_X}^\mathbf{L} L$ に随伴する。後者としては，同型
$Lf^*(Rf_*K \otimes_{\mathcal{O}_Y}^\mathbf{L} Rf_*L) =
Lf^*Rf_*K \otimes_{\mathcal{O}_X}^\mathbf{L} Lf^*Rf_*L$
（補題 \ref{cohomology-lemma-pullback-tensor-product}）と，余単位
$Lf^* \circ Rf_* \to \text{id}$ から得られる写像
$Lf^*Rf_*K \otimes_{\mathcal{O}_X}^\mathbf{L} Lf^*Rf_*L
\to K \otimes_{\mathcal{O}_X}^\mathbf{L} L$
との合成を取れる。
\end{remark}






\section{フィルトレーション付き複体のコホモロジー}
\label{cohomology-section-cohomology-filtered-object}

\noindent
代数多様体のコホモロジーを研究すると，層のフィルトレーション付き複体が
しばしば自然に現れる。例えば de Rham 複体には Hodge フィルトレーションが
備わっている。本節では，非常に一般的な《入射対象》補題
\ref{injectives-lemma-K-injective-embedding-filtration} を用いて
コホモロジー上のスペクトル系列を構成し，これを先に構成した
スペクトル系列と関係づける。

\begin{lemma}
\label{cohomology-lemma-spectral-sequence-filtered-object}
$(X, \mathcal{O}_X)$ を環付き空間とし，$\mathcal{F}^\bullet$ を
$\mathcal{O}_X$-加群のフィルトレーション付き複体とする。
$d_r$ が双次数 $(r, -r + 1)$ をもち，かつ
$$
E_1^{p, q} = H^{p + q}(X, \text{gr}^p\mathcal{F}^\bullet)
$$
を満たす，双次数付き $\Gamma(X, \mathcal{O}_X)$-加群の標準的な
スペクトル系列 $(E_r, \text{d}_r)_{r \geq 1}$ が存在する。
すべての $n$ に対して
$$
H^n(X, F^p\mathcal{F}^\bullet) = 0\text{ for }p \gg 0
\quad\text{and}\quad
H^n(X, F^p\mathcal{F}^\bullet) = H^n(X, \mathcal{F}^\bullet)\text{ for }p \ll 0
$$
であるならば，このスペクトル系列は有界で，
$H^*(X, \mathcal{F}^\bullet)$ に収束する。
\end{lemma}

\begin{proof}
（複体が有限フィルトレーションをもつ下に有界な加群複体である場合の証明は，
下の注意を参照せよ。）《入射対象》補題
\ref{injectives-lemma-K-injective-embedding-filtration}
のようなフィルトレーション付き複体の写像
$j : \mathcal{F}^\bullet \to \mathcal{J}^\bullet$ を選ぶ。
求めるスペクトル系列は，フィルトレーション付き複体
$$
\Gamma(X, \mathcal{J}^\bullet)
\quad\text{with}\quad
F^p\Gamma(X, \mathcal{J}^\bullet) = \Gamma(X, F^p\mathcal{J}^\bullet)
$$
に付随する《ホモロジー》第
\ref{homology-section-filtered-complex} 節のスペクトル系列である。
コホモロジーは K-入射的な代表に値を取ることで計算されるから，
$E_1$ 頁は補題で述べたとおりである。仮定した条件の下での収束性と
有界性は《ホモロジー》補題
\ref{homology-lemma-ss-converges-trivial} から従う。
\end{proof}

\begin{remark}
\label{cohomology-remark-spectral-sequence-filtered-object}
$(X, \mathcal{O}_X)$ を環付き空間とし，$\mathcal{F}^\bullet$ を
$\mathcal{O}_X$-加群のフィルトレーション付き複体とする。
$\mathcal{F}^\bullet$ が下に有界で，各 $n$ に対して
$\mathcal{F}^n$ 上のフィルトレーションが有限ならば，補題
\ref{cohomology-lemma-spectral-sequence-filtered-object} のスペクトル系列を，
《入射対象》補題
\ref{injectives-lemma-K-injective-embedding-filtration} を用いずに構成できる。
実際，《導来圏》補題
\ref{derived-lemma-right-resolution-by-filtered-injectives}
により，フィルトレーション付き複体のフィルトレーション付き擬同型
$i : \mathcal{F}^\bullet \to \mathcal{I}^\bullet$
が存在し，$\mathcal{I}^\bullet$ は下に有界，各 $n$ に対して
$\mathcal{I}^n$ 上のフィルトレーションは有限，かつ各
$\text{gr}^p\mathcal{I}^n$ は入射的 $\mathcal{O}_X$-加群である。
そこで，
$$
\Gamma(X, \mathcal{I}^\bullet)
\quad\text{with}\quad
F^p\Gamma(X, \mathcal{I}^\bullet) = \Gamma(X, F^p\mathcal{I}^\bullet)
$$
に付随するスペクトル系列を取る。コホモロジーは入射的対象からなる
下に有界な複体に値を取ることで計算できるから，$E_1$ 頁は補題で
述べたとおりである。仮定した条件の下での収束性と有界性は
《ホモロジー》補題 \ref{homology-lemma-biregular-ss-converges} から従う。
実際，これは《導来圏》補題
\ref{derived-lemma-ss-filtered-derived} のスペクトル系列の特別な場合である。
\end{remark}

\begin{example}
\label{cohomology-example-spectral-sequence}
$(X, \mathcal{O}_X)$ を環付き空間とし，$\mathcal{F}^\bullet$ を
$\mathcal{O}_X$-加群の複体とする。
$F^p\mathcal{F}^\bullet = \tau_{\leq -p}\mathcal{F}^\bullet$ として
補題 \ref{cohomology-lemma-spectral-sequence-filtered-object} を適用できる。
（$\mathcal{F}^\bullet$ が下に有界ならば，注意
\ref{cohomology-remark-spectral-sequence-filtered-object} を用いてもよい。）
するとスペクトル系列
$$
E_1^{p, q} = H^{p + q}(X, H^{-p}(\mathcal{F}^\bullet)[p]) =
H^{2p + q}(X, H^{-p}(\mathcal{F}^\bullet))
$$
を得る。$p = -j$ および $q = i + 2j$ と番号を付け替えると，任意の
$K \in D(\mathcal{O}_X)$ に対し，$d'_r$ が双次数 $(r, -r + 1)$ をもち，
$$
(E'_2)^{i, j} = H^i(X, H^j(K))
$$
を満たす双次数付き加群のスペクトル系列
$(E'_r, d'_r)_{r \geq 2}$ が存在することが分かる。
例えば $K$ が下に有界ならば，このスペクトル系列は有界で，
$H^{i + j}(X, K)$ に収束する。下に有界な場合，これは
《導来圏》補題 \ref{derived-lemma-two-ss-complex-functor} の
第二スペクトル系列（Cartan--Eilenberg 分解を用いて構成される）の一例である。
\end{example}

\begin{example}
\label{cohomology-example-spectral-sequence-bis}
$(X, \mathcal{O}_X)$ を環付き空間とし，$\mathcal{F}^\bullet$ を
$\mathcal{O}_X$-加群の複体とする。
$F^p\mathcal{F}^\bullet = \sigma_{\geq p}\mathcal{F}^\bullet$ として
補題 \ref{cohomology-lemma-spectral-sequence-filtered-object} を適用できる。
するとスペクトル系列
$$
E_1^{p, q} = H^{p + q}(X, \mathcal{F}^p[-p]) = H^q(X, \mathcal{F}^p)
$$
を得る。$\mathcal{F}^\bullet$ が下に有界ならば，
\begin{enumerate}
\item 注意 \ref{cohomology-remark-spectral-sequence-filtered-object} を用いて
このスペクトル系列を構成できる。
\item このスペクトル系列は有界で，
$H^{i + j}(X, \mathcal{F}^\bullet)$ に収束する。
\item このスペクトル系列は《導来圏》補題
\ref{derived-lemma-two-ss-complex-functor} の第一スペクトル系列
（Cartan--Eilenberg 分解を用いて構成される）に等しい。
\end{enumerate}
\end{example}

\begin{lemma}
\label{cohomology-lemma-relative-spectral-sequence-filtered-object}
$f : (X, \mathcal{O}_X) \to (Y, \mathcal{O}_Y)$ を環付き空間の射とし，
$\mathcal{F}^\bullet$ を $\mathcal{O}_X$-加群の
フィルトレーション付き複体とする。$d_r$ が双次数 $(r, -r + 1)$ をもち，かつ
$$
E_1^{p, q} = R^{p + q}f_*\text{gr}^p\mathcal{F}^\bullet
$$
を満たす，双次数付き $\mathcal{O}_Y$-加群の標準的なスペクトル系列
$(E_r, \text{d}_r)_{r \geq 1}$ が存在する。すべての $n$ に対して
$$
R^nf_*F^p\mathcal{F}^\bullet = 0 \text{ for }p \gg 0
\quad\text{and}\quad
R^nf_*F^p\mathcal{F}^\bullet = R^nf_*\mathcal{F}^\bullet \text{ for }p \ll 0
$$
であるならば，このスペクトル系列は有界で，
$Rf_*\mathcal{F}^\bullet$ に収束する。
\end{lemma}

\begin{proof}
証明は補題 \ref{cohomology-lemma-spectral-sequence-filtered-object} の証明とまったく同じである。
\end{proof}





\section{Godement 分解}
\label{cohomology-section-godement}

\noindent
参考文献は \cite{Godement} である。

\medskip\noindent
$(X, \mathcal{O}_X)$ を環付き空間とする。$X$ と同じ点をもつ離散位相空間を
$X_{disc}$ と書き，自明な連続写像を $f : X_{disc} \to X$ と書く。
$\mathcal{O}_{X_{disc}} = f^{-1}\mathcal{O}_X$ と置く。このとき
$f : (X_{disc}, \mathcal{O}_{X_{disc}}) \to (X, \mathcal{O}_X)$ は
環付き空間の平坦射である。加群の層上の随伴関手の対 $f^*, f_*$ に，
《単体的対象》第 \ref{simplicial-section-standard} 節の内容の
{\it 双対}を適用できる。こうして増大余単体的対象
$$
\xymatrix{
\text{id} \ar[r] &
f_*f^* \ar@<1ex>[r] \ar@<-1ex>[r] &
f_*f^*f_*f^* \ar@<0ex>[l] \ar@<-2ex>[r] \ar@<0ex>[r] \ar@<2ex>[r] &
f_*f^*f_*f^*f_*f^* \ar@<1ex>[l] \ar@<-1ex>[l]
}
$$
を，$\textit{Mod}(\mathcal{O}_X)$ からそれ自身への関手の圏で得る。
《単体的対象》補題 \ref{simplicial-lemma-standard-simplicial} を参照せよ。
さらに，増大
$$
\xymatrix{
f^* \ar[r] &
f^*f_*f^* \ar@<1ex>[r] \ar@<-1ex>[r] &
f^*f_*f^*f_*f^* \ar@<0ex>[l] \ar@<-2ex>[r] \ar@<0ex>[r] \ar@<2ex>[r] &
f^*f_*f^*f_*f^*f_*f^* \ar@<1ex>[l] \ar@<-1ex>[l]
}
$$
はホモトピー同値である。《単体的対象》補題
\ref{simplicial-lemma-standard-simplicial-homotopy} を参照せよ。

\begin{lemma}
\label{cohomology-lemma-godement-resolution}
$(X, \mathcal{O}_X)$ を環付き空間とする。任意の
$\mathcal{O}_X$-加群の層 $\mathcal{F}$ に対して分解
$$
0 \to
\mathcal{F} \to
f_*f^*\mathcal{F} \to
f_*f^*f_*f^*\mathcal{F} \to
f_*f^*f_*f^*f_*f^*\mathcal{F} \to \ldots
$$
が存在する。これは $\mathcal{F}$ に関して関手的で，各項
$f_*f^* \ldots f_*f^*\mathcal{F}$ はフラスク
$\mathcal{O}_X$-加群であり，かつすべての $x \in X$ に対して写像
$$
\mathcal{F}_x[0] \to \Big(
(f_*f^*\mathcal{F})_x \to
(f_*f^*f_*f^*\mathcal{F})_x \to
(f_*f^*f_*f^*f_*f^*\mathcal{F})_x \to \ldots
\Big)
$$
は $\mathcal{O}_{X, x}$-加群の複体の圏におけるホモトピー同値である。
\end{lemma}

\begin{proof}
複体 $f_*f^*\mathcal{F} \to  f_*f^*f_*f^*\mathcal{F} \to
f_*f^*f_*f^*f_*f^*\mathcal{F} \to \ldots$ は，上で述べた，項が\\
$f_*f^*\mathcal{F}, f_*f^*f_*f^*\mathcal{F},
f_*f^*f_*f^*f_*f^*\mathcal{F}, \ldots$ である余単体的対象に
付随する複体である。《単体的対象》第
\ref{simplicial-section-dold-kan-cosimplicial} 節を参照せよ。
増大から，表示した写像 $\mathcal{F} \to f_*f^*\mathcal{F}$ が得られる。
$X_{disc}$ 上の任意のアーベル層 $\mathcal{H}$ に対して，
順像 $f_*\mathcal{H}$ はフラスクである。実際，$X_{disc}$ は離散空間であり，
フラスク層の順像はフラスクである。したがって，複体の各項は
フラスク $\mathcal{O}_X$-加群である。

\medskip\noindent
点 $x \in X_{disc} = X$ と任意の $\mathcal{O}_X$-加群
$\mathcal{G}$ に対し，$(f^*\mathcal{G})_x = \mathcal{G}_x$ である。
したがって $f^*$ は完全関手であり，$\mathcal{O}_X$-加群の複体
$\mathcal{G}_1 \to \mathcal{G}_2 \to \mathcal{G}_3$
が完全であることと，
$f^*\mathcal{G}_1 \to f^*\mathcal{G}_2 \to f^*\mathcal{G}_3$
が完全であることは同値である（《加群》補題
\ref{modules-lemma-abelian} を参照）。本節の導入で述べた結果により，
$f^*$ を適用すると，定余単体的対象 $f^*\mathcal{F}$ から，項が
$f_*f^*\mathcal{F}, f_*f^*f_*f^*\mathcal{F},
f_*f^*f_*f^*f_*f^*\mathcal{F}, \ldots$ である余単体的対象への
ホモトピー同値が得られる。《単体的対象》補題
\ref{simplicial-lemma-homotopy-equivalence-s-Q} により，
$$
f^*\mathcal{F}[0] \to \Big(
f^*f_*f^*\mathcal{F} \to
f^*f_*f^*f_*f^*\mathcal{F} \to
f^*f_*f^*f_*f^*f_*f^*\mathcal{F} \to \ldots
\Big)
$$
はホモトピー同値である。これから補題の残る二つの主張が直ちに従う。
\end{proof}

\begin{lemma}
\label{cohomology-lemma-godement-resolution-bounded-below}
$(X, \mathcal{O}_X)$ を環付き空間とし，
$\mathcal{F}^\bullet$ を $\mathcal{O}_X$-加群の下に有界な複体とする。
擬同型
$\mathcal{F}^\bullet \to \mathcal{G}^\bullet$
が存在する。ここで $\mathcal{G}^\bullet$ はフラスク
$\mathcal{O}_X$-加群からなる下に有界な複体であり，すべての
$x \in X$ に対して写像
$\mathcal{F}^\bullet_x \to \mathcal{G}^\bullet_x$ は
$\mathcal{O}_{X, x}$-加群の複体の圏におけるホモトピー同値である。
\end{lemma}

\begin{proof}
$\mathcal{A}$ を $\mathcal{O}_X$-加群の複体の圏とし，
$\mathcal{B}$ を $\mathcal{O}_X$-加群の複体の圏とする。このとき
$\mathcal{A}$ と $\mathcal{B}$ の間の随伴関手 $f^*$ と $f_*$ に
上の議論を適用できる。補題 \ref{cohomology-lemma-godement-resolution} の証明と
まったく同様に論じると，アーベル圏 $\mathcal{A}$ における分解
$$
0 \to
\mathcal{F}^\bullet \to
f_*f^*\mathcal{F}^\bullet \to
f_*f^*f_*f^*\mathcal{F}^\bullet \to
f_*f^*f_*f^*f_*f^*\mathcal{F}^\bullet \to \ldots
$$
を得る。各 $f_*f^*\ldots f_*f^*\mathcal{F}^\bullet$ の各項は
フラスク $\mathcal{O}_X$-加群であり，すべての $x \in X$ に対して写像
$$
\mathcal{F}^\bullet_x[0] \to \Big(
(f_*f^*\mathcal{F}^\bullet)_x \to
(f_*f^*f_*f^*\mathcal{F}^\bullet)_x \to
(f_*f^*f_*f^*f_*f^*\mathcal{F}^\bullet)_x \to \ldots
\Big)
$$
は $\mathcal{O}_{X, x}$-加群の複体の複体の圏における
ホモトピー同値である。複体の複体は二重複体と同じものなので，
誘導される写像
$$
\mathcal{F}^\bullet \to
\mathcal{G}^\bullet =
\text{Tot}(
f_*f^*\mathcal{F}^\bullet \to
f_*f^*f_*f^*\mathcal{F}^\bullet \to
f_*f^*f_*f^*f_*f^*\mathcal{F}^\bullet \to \ldots
)
$$
を考えられる。複体 $\mathcal{F}^\bullet$ は下に有界なので，
$\mathcal{G}^\bullet$ も同様である。実際，$\mathcal{G}^\bullet$ の各項は，
複体 $f_*f^*\ldots f_*f^*\mathcal{F}^\bullet$ の項の有限直和であり，
したがってフラスクである。補題の最後の主張は
《ホモロジー》補題
\ref{homology-lemma-homotopy-complex-complexes} から従う。特にこれは
$\mathcal{F}^\bullet \to \mathcal{G}^\bullet$
が擬同型であることを示すので，証明は完了する。
\end{proof}














\section{カップ積}
\label{cohomology-section-cup-product}

\noindent
$(X, \mathcal{O}_X)$ を環付き空間とし，$K, M$ を
$D(\mathcal{O}_X)$ の対象とする。$A = \Gamma(X, \mathcal{O}_X)$ と置く。
この状況での（大域）カップ積とは，$D(A)$ における写像
$$
\mu :
R\Gamma(X, K) \otimes_A^\mathbf{L} R\Gamma(X, M)
\longrightarrow
R\Gamma(X, K \otimes_{\mathcal{O}_X}^\mathbf{L} M)
$$
である。これを，$D(pt, A) = D(A)$ を通じて，注意
\ref{cohomology-remark-cup-product} における環付き空間の射
$f : (X, \mathcal{O}_X) \to (pt, A)$ に対する相対カップ積として定義する。
特にこの写像はペアリング
$$
\cup :
H^i(X, K) \times H^j(X, M)
\longrightarrow
H^{i + j}(X, K \otimes_{\mathcal{O}_X}^\mathbf{L} M)
$$
を定める。実際，$\xi \in H^i(X, K) = H^i(R\Gamma(X, K))$ と
$\eta \in H^j(X, M) = H^j(R\Gamma(X, M))$ が与えられたとする。
まず両者を「テンソル積」して，
$H^{i + j}(R\Gamma(X, K) \otimes_A^\mathbf{L} R\Gamma(X, M))$ の元
$\xi \otimes \eta$ を得ることができる。《代数続論》第
\ref{more-algebra-section-products-tor} 節を参照せよ。
次に $\mu$ を適用すると，求める元
$\xi \cup \eta = \mu(\xi \otimes \eta)$
を $H^{i + j}(X, K \otimes_{\mathcal{O}_X}^\mathbf{L} M)$ に得る。

\medskip\noindent
$\xi$ と $\eta$ のカップ積を考える別の方法を述べよう。すなわち，
$$
H^i(R\Gamma(X, K)) = \Hom_{D(\mathcal{O}_X)}(\mathcal{O}_X[-i], K)
\quad\text{and}\quad
H^j(R\Gamma(X, M)) = \Hom_{D(\mathcal{O}_X)}(\mathcal{O}_X[-j], M)
$$
と書ける。これは $\Hom(\mathcal{O}_X, -) = \Gamma(X, -)$ だからである。
したがって $\xi$ と $\eta$ は，それぞれ写像
$$
\tilde \xi : \mathcal{O}_X[-i] \to K
\quad\text{and}\quad
\tilde \eta : \mathcal{O}_X[-j] \to M
$$
と「同じ」ものである。これと導来テンソル積の関手性を組み合わせると，
$$
\mathcal{O}_X[-i - j] =
\mathcal{O}_X[-i] \otimes_{\mathcal{O}_X}^\mathbf{L} \mathcal{O}_X[-j]
\xrightarrow{\tilde \xi \otimes \tilde \eta}
K \otimes_{\mathcal{O}_X}^\mathbf{L} M
$$
を得る。これは上と同じ理由により
$H^{i + j}(X, K \otimes_{\mathcal{O}_X}^\mathbf{L} M)$ の元である。

\begin{lemma}
\label{cohomology-lemma-second-cup-equals-first}
この構成はカップ積を与える。
\end{lemma}

\begin{proof}
上の $f : (X, \mathcal{O}_X) \to (pt, A)$ に対して
$Rf_*(-) = R\Gamma(X, -)$ であり，写像 $\mu$ は写像
$$
Lf^*(Rf_*K \otimes_A^\mathbf{L} Rf_*M) =
Lf^*Rf_*K \otimes_{\mathcal{O}_X}^\mathbf{L} Lf^*Rf_*M
\xrightarrow{\epsilon_K \otimes \epsilon_M}
K \otimes_{\mathcal{O}_X}^\mathbf{L} M
$$
に随伴する。ここで $\epsilon$ は $Lf^*$ と $Rf_*$ の随伴の余単位である。
$\xi$ と $\eta$ を写像 $\xi : A[-i] \to R\Gamma(X, K)$ および
$\eta : A[-j] \to R\Gamma(X, M)$ と考えると，テンソル積
$\xi \otimes \eta$ は写像\footnote{ここには符号が隠れている。
すなわち，この等式は合成
$$
A[-i - j] \to (A \otimes_A^\mathbf{L} A)[-i - j] \to
A[-i] \otimes_A^\mathbf{L} A[-j]
$$
によって定義される。第二段階では《代数続論》項
(\ref{more-algebra-item-shift-tensor}) の同一視を用いるが，
原理的にはここで符号が現れる。ただし，本書の規約ではこの場合の符号は
$+1$ である。また $+1$ でなかったとしても，同一視
$\mathcal{O}_X[-i - j] =
\mathcal{O}_X[-i] \otimes_{\mathcal{O}_X}^\mathbf{L} \mathcal{O}_X[-j]$
に同じ符号を用いたので問題はない。}
$$
A[-i - j] = A[-i] \otimes_A^\mathbf{L} A[-j]
\xrightarrow{\xi \otimes \eta}
R\Gamma(X, K) \otimes_A^\mathbf{L} R\Gamma(X, M)
$$
に対応する。定義により，カップ積 $\xi \cup \eta$ は写像
$A[-i - j] \to R\Gamma(X, K \otimes_{\mathcal{O}_X}^\mathbf{L} M)$
であり，これは
$$
(\epsilon_K \otimes \epsilon_M) \circ Lf^*(\xi \otimes \eta) =
(\epsilon_K \circ Lf^*\xi) \otimes (\epsilon_M \circ Lf^*\eta)
$$
に随伴する。一方，
$\epsilon_K \circ Lf^*\xi = \tilde \xi$ かつ
$\epsilon_M \circ Lf^*\eta = \tilde \eta$ であることは容易に分かる。
したがって $\widetilde{\xi \cup \eta} = \tilde \xi \otimes \tilde \eta$
であり，求める一致を得る。
\end{proof}

\begin{remark}
\label{cohomology-remark-cup-with-element-map-total-cohomology}
$(X, \mathcal{O}_X)$ を環付き空間とし，$K, M$ を
$D(\mathcal{O}_X)$ の対象とする。$A = \Gamma(X, \mathcal{O}_X)$ と置く。
$\xi \in H^i(X, K)$ が与えられると，付随する写像
$$
\xi = ``\xi \cup -'' :
R\Gamma(X, M)[-i]
\to
R\Gamma(X, K \otimes_{\mathcal{O}_X}^\mathbf{L} M)
$$
を得る。実際，補題 \ref{cohomology-lemma-second-cup-equals-first} の証明と同様に
$\xi$ を写像 $\xi : A[-i] \to R\Gamma(X, K)$ で表し，次の合成を用いる：
$$
R\Gamma(X, M)[-i] = A[-i] \otimes_A^\mathbf{L} R\Gamma(X, M)
\xrightarrow{\xi \otimes 1}
R\Gamma(X, K) \otimes_A^\mathbf{L} R\Gamma(X, M)
\to
R\Gamma(X, K \otimes_{\mathcal{O}_X}^\mathbf{L} M)
$$
ここで第二の矢印は上の大域カップ積 $\mu$ である。補題
\ref{cohomology-lemma-second-cup-equals-first} とその証明から明らかなように，
コホモロジー上ではこれは $\xi$ によるカップ積を再現する。
\end{remark}

\noindent
相対カップ積の自然な整合性を定式化して証明しよう。すなわち，
環付き空間の射
$f : (X, \mathcal{O}_X) \to (Y, \mathcal{O}_Y)$ があると仮定し，
$\mathcal{K}^\bullet$ と $\mathcal{M}^\bullet$ を
$\mathcal{O}_X$-加群の複体とする。素朴なカップ積
$$
\text{Tot}(
f_*\mathcal{K}^\bullet
\otimes_{\mathcal{O}_Y}
f_*\mathcal{M}^\bullet)
\longrightarrow
f_*\text{Tot}(\mathcal{K}^\bullet
\otimes_{\mathcal{O}_X}
\mathcal{M}^\bullet)
$$
が存在する。これが相対カップ積と関係することを主張する。

\begin{lemma}
\label{cohomology-lemma-cup-compatible-with-naive}
上の状況において，次の図式は可換である。
$$
\xymatrix{
f_*\mathcal{K}^\bullet
\otimes_{\mathcal{O}_Y}^\mathbf{L}
f_*\mathcal{M}^\bullet \ar[r] \ar[d]
&
Rf_*\mathcal{K}^\bullet
\otimes_{\mathcal{O}_Y}^\mathbf{L}
Rf_*\mathcal{M}^\bullet \ar[d]^{\text{Remark \ref{cohomology-remark-cup-product}}} \\
\text{Tot}(
f_*\mathcal{K}^\bullet
\otimes_{\mathcal{O}_Y}
f_*\mathcal{M}^\bullet) \ar[d]_{\text{naive cup product}} &
Rf_*(\mathcal{K}^\bullet
\otimes_{\mathcal{O}_X}^\mathbf{L}
\mathcal{M}^\bullet) \ar[d] \\
f_*\text{Tot}(\mathcal{K}^\bullet
\otimes_{\mathcal{O}_X}
\mathcal{M}^\bullet) \ar[r] &
Rf_*\text{Tot}(\mathcal{K}^\bullet
\otimes_{\mathcal{O}_X}
\mathcal{M}^\bullet)
}
$$
\end{lemma}

\begin{proof}
注意 \ref{cohomology-remark-cup-product} の構成により，図式を時計回りに進んで得る写像
$$
f_*\mathcal{K}^\bullet
\otimes_{\mathcal{O}_Y}^\mathbf{L}
f_*\mathcal{M}^\bullet 
\longrightarrow
Rf_*\text{Tot}(\mathcal{K}^\bullet
\otimes_{\mathcal{O}_X}
\mathcal{M}^\bullet)
$$
は，写像
\begin{align*}
Lf^*(f_*\mathcal{K}^\bullet
\otimes_{\mathcal{O}_Y}^\mathbf{L}
f_*\mathcal{M}^\bullet)
& =
Lf^*f_*\mathcal{K}^\bullet
\otimes_{\mathcal{O}_Y}^\mathbf{L}
Lf^*f_*\mathcal{M}^\bullet \\
& \to
Lf^*Rf_*\mathcal{K}^\bullet
\otimes_{\mathcal{O}_Y}^\mathbf{L}
Lf^*Rf_*\mathcal{M}^\bullet \\
& \to
\mathcal{K}^\bullet
\otimes_{\mathcal{O}_Y}^\mathbf{L}
\mathcal{M}^\bullet \\
& \to
\text{Tot}(\mathcal{K}^\bullet
\otimes_{\mathcal{O}_X}
\mathcal{M}^\bullet)
\end{align*}
に随伴する。補題 \ref{cohomology-lemma-adjoints-push-pull-compatibility} により，
これはまた
\begin{align*}
Lf^*(f_*\mathcal{K}^\bullet
\otimes_{\mathcal{O}_Y}^\mathbf{L}
f_*\mathcal{M}^\bullet)
& =
Lf^*f_*\mathcal{K}^\bullet
\otimes_{\mathcal{O}_Y}^\mathbf{L}
Lf^*f_*\mathcal{M}^\bullet \\
& \to
f^*f_*\mathcal{K}^\bullet
\otimes_{\mathcal{O}_Y}^\mathbf{L}
f^*f_*\mathcal{M}^\bullet \\
& \to
\mathcal{K}^\bullet
\otimes_{\mathcal{O}_Y}^\mathbf{L}
\mathcal{M}^\bullet \\
& \to
\text{Tot}(\mathcal{K}^\bullet
\otimes_{\mathcal{O}_X}
\mathcal{M}^\bullet)
\end{align*}
に等しい。反時計回りに進むと，写像
\begin{align*}
Lf^*(f_*\mathcal{K}^\bullet
\otimes_{\mathcal{O}_Y}^\mathbf{L}
f_*\mathcal{M}^\bullet)
& \to
Lf^*\text{Tot}(
f_*\mathcal{K}^\bullet
\otimes_{\mathcal{O}_Y}
f_*\mathcal{M}^\bullet) \\
& \to
Lf^*f_*\text{Tot}(\mathcal{K}^\bullet
\otimes_{\mathcal{O}_X}
\mathcal{M}^\bullet) \\
& \to
Lf^*Rf_*\text{Tot}(\mathcal{K}^\bullet
\otimes_{\mathcal{O}_X}
\mathcal{M}^\bullet) \\
& \to
\text{Tot}(\mathcal{K}^\bullet
\otimes_{\mathcal{O}_X}
\mathcal{M}^\bullet)
\end{align*}
に随伴する写像を得る。補題
\ref{cohomology-lemma-adjoints-push-pull-compatibility} により，これはまた
\begin{align*}
Lf^*(f_*\mathcal{K}^\bullet
\otimes_{\mathcal{O}_Y}^\mathbf{L}
f_*\mathcal{M}^\bullet)
& \to
Lf^*\text{Tot}(
f_*\mathcal{K}^\bullet
\otimes_{\mathcal{O}_Y}
f_*\mathcal{M}^\bullet) \\
& \to
Lf^*f_*\text{Tot}(\mathcal{K}^\bullet
\otimes_{\mathcal{O}_X}
\mathcal{M}^\bullet) \\
& \to
f^*f_*\text{Tot}(\mathcal{K}^\bullet
\otimes_{\mathcal{O}_X}
\mathcal{M}^\bullet) \\
& \to
\text{Tot}(\mathcal{K}^\bullet
\otimes_{\mathcal{O}_X}
\mathcal{M}^\bullet)
\end{align*}
に等しい。ここで次の図式を調べれば証明が完了する。
$$
\xymatrix{
Lf^*(f_*\mathcal{K}^\bullet
\otimes_{\mathcal{O}_Y}^\mathbf{L}
f_*\mathcal{M}^\bullet) \ar[d] \ar[rr] & &
Lf^*f_*\mathcal{K}^\bullet \otimes_{\mathcal{O}_X}^\mathbf{L}
Lf^*f_*\mathcal{M}^\bullet \ar[d] \\
Lf^*\text{Tot}(
f_*\mathcal{K}^\bullet
\otimes_{\mathcal{O}_Y}
f_*\mathcal{M}^\bullet) \ar[d]_{naive} \ar[r] &
f^*\text{Tot}(
f_*\mathcal{K}^\bullet
\otimes_{\mathcal{O}_Y}
f_*\mathcal{M}^\bullet) \ar[ldd]^{naive} \ar[dd] &
f^*f_*\mathcal{K}^\bullet \otimes_{\mathcal{O}_X}^\mathbf{L}
f^*f_*\mathcal{M}^\bullet \ar[dd] \ar[ldd] \\
Lf^*f_*\text{Tot}(\mathcal{K}^\bullet
\otimes_{\mathcal{O}_X}
\mathcal{M}^\bullet) \ar[d] \\
f^*f_*\text{Tot}(\mathcal{K}^\bullet \otimes_{\mathcal{O}_X}
\mathcal{M}^\bullet) \ar[rd] &
\text{Tot}(f^*f_*\mathcal{K}^\bullet \otimes_{\mathcal{O}_X}
f^*f_*\mathcal{M}^\bullet) \ar[d] &
\mathcal{K}^\bullet \otimes_{\mathcal{O}_X}^\mathbf{L}
\mathcal{M}^\bullet \ar[ld] \\
& \text{Tot}(\mathcal{K}^\bullet
\otimes_{\mathcal{O}_X}
\mathcal{M}^\bullet)
}
$$
この図式のすべての多角形は可換である。上の多角形が可換であることは
補題 \ref{cohomology-lemma-tensor-pull-compatibility} による。
二つの素朴なカップ積を含む正方形は，$Lf^* \to f^*$ が
加群の複体に関して関手的であるため可換である。二つの写像
$\mathcal{A}^\bullet \otimes^\mathbf{L} \mathcal{B}^\bullet \to
\text{Tot}(\mathcal{A}^\bullet \otimes \mathcal{B}^\bullet)$
を含む正方形についても同様である。最後に，残る正方形の可換性は
複体の段階で成り立ち，$f^*$ と $f_*$ の随伴性による
素朴なカップ積の定義とみなせる。図式の外周を二通りに進むと
上で与えた二つの写像になるので，証明は完了する。
\end{proof}

\noindent
$(X, \mathcal{O}_X)$ を環付き空間とし，$\mathcal{K}^\bullet$ と
$\mathcal{M}^\bullet$ を $\mathcal{O}_X$-加群の複体とする。
このとき「素朴な」カップ積
$$
\mu' :
\text{Tot}(
\Gamma(X, \mathcal{K}^\bullet) \otimes_A \Gamma(X, \mathcal{M}^\bullet))
\longrightarrow
\Gamma(X, \text{Tot}(
\mathcal{K}^\bullet \otimes_{\mathcal{O}_X} \mathcal{M}^\bullet))
$$
がある。補題 \ref{cohomology-lemma-cup-compatible-with-naive} を射
$(X, \mathcal{O}_X) \to (pt, A)$ に適用すると，この素朴なカップ積は，
本節の最初の段落で定義したカップ積 $\mu$ と次の可換図式で関係づけられる。
$$
\xymatrix{
\Gamma(X, \mathcal{K}^\bullet)
\otimes_A^\mathbf{L}
\Gamma(X, \mathcal{M}^\bullet) \ar[d] \ar[r] &
R\Gamma(X, \mathcal{K}^\bullet)
\otimes_A^\mathbf{L}
R\Gamma(X, \mathcal{M}^\bullet) \ar[d]^-\mu \\
\text{Tot}(\Gamma(X, \mathcal{K}^\bullet)
\otimes_A
\Gamma(X, \mathcal{M}^\bullet)) \ar[d]_-{\mu'} &
R\Gamma(X, \mathcal{K}^\bullet
\otimes_{\mathcal{O}_X}^\mathbf{L}
\mathcal{M}^\bullet) \ar[d] \\
\Gamma(X, \text{Tot}(\mathcal{K}^\bullet
\otimes_{\mathcal{O}_X}
\mathcal{M}^\bullet)) \ar[r] &
R\Gamma(X, \text{Tot}(\mathcal{K}^\bullet
\otimes_{\mathcal{O}_X}
\mathcal{M}^\bullet))
}
$$
これは $D(A)$ における図式である。コホモロジー上では可換図式
$$
\xymatrix{
H^i(\Gamma(X, \mathcal{K}^\bullet)) \times
H^j(\Gamma(X, \mathcal{M}^\bullet)) \ar[d] \ar[r] &
H^{i + j}(X,
\text{Tot}(\mathcal{K}^\bullet \otimes_{\mathcal{O}_X} \mathcal{M}^\bullet)) \\
H^i(X, \mathcal{K}^\bullet) \times
H^j(X, \mathcal{M}^\bullet) \ar[r]^\cup &
H^{i + j}(X, \mathcal{K}^\bullet \otimes_{\mathcal{O}_X}^\mathbf{L}
\mathcal{M}^\bullet) \ar[u]
}
$$
を得る。これは素朴なカップ積と真のカップ積を関係づける。

\begin{lemma}
\label{cohomology-lemma-diagrams-commute}
$(X, \mathcal{O}_X)$ を環付き空間とし，
$\mathcal{K}^\bullet$ と $\mathcal{M}^\bullet$ を
$\mathcal{O}_X$-加群の下に有界な複体とする。
$\mathcal{U} : X = \bigcup_{i \in I} U_i$ を開被覆とする。このとき，
$$
\xymatrix{
\text{Tot}(\check{\mathcal{C}}^\bullet(\mathcal{U}, \mathcal{K}^\bullet))
\otimes_A^\mathbf{L}
\text{Tot}(\check{\mathcal{C}}^\bullet(\mathcal{U}, \mathcal{M}^\bullet))
\ar[d] \ar[r] &
R\Gamma(X, \mathcal{K}^\bullet)
\otimes_A^\mathbf{L}
R\Gamma(X, \mathcal{M}^\bullet) \ar[d]^\mu \\
\text{Tot}(
\text{Tot}(\check{\mathcal{C}}^\bullet(\mathcal{U}, \mathcal{K}^\bullet))
\otimes_A
\text{Tot}(\check{\mathcal{C}}^\bullet(\mathcal{U}, \mathcal{M}^\bullet)))
\ar[d]^{(\ref{cohomology-equation-needs-signs})} &
R\Gamma(X,
\mathcal{K}^\bullet \otimes_{\mathcal{O}_X}^\mathbf{L} \mathcal{M}^\bullet)
\ar[d] \\
\text{Tot}(
\check{\mathcal{C}}^\bullet({\mathcal U},
\text{Tot}(\mathcal{K}^\bullet \otimes_{\mathcal{O}_X} \mathcal{M}^\bullet)
)) \ar[r] &
R\Gamma(X,
\text{Tot}(\mathcal{K}^\bullet \otimes_{\mathcal{O}_X} \mathcal{M}^\bullet))
}
$$
は $D(A)$ において可換である。ここで水平射は補題
\ref{cohomology-lemma-cech-complex-complex} のものである。
\end{lemma}

\begin{proof}
補題 \ref{cohomology-lemma-godement-resolution-bounded-below} のような複体の擬同型
$a : \mathcal{K}^\bullet \to \mathcal{K}_1^\bullet$ および
$b : \mathcal{M}^\bullet \to \mathcal{M}_1^\bullet$
を選ぶ。茎上の写像 $a$ と $b$ はホモトピー同値なので，誘導される写像
$$
\text{Tot}(\mathcal{K}^\bullet \otimes_{\mathcal{O}_X} \mathcal{M}^\bullet)
\to
\text{Tot}(\mathcal{K}_1^\bullet \otimes_{\mathcal{O}_X} \mathcal{M}_1^\bullet)
$$
も茎上のホモトピー同値であり（《代数続論》補題
\ref{more-algebra-lemma-derived-tor-homotopy}），したがって擬同型である。
ゆえに二つの図式の終域
$$
R\Gamma(X,
\text{Tot}(\mathcal{K}^\bullet
\otimes_{\mathcal{O}_X} \mathcal{M}^\bullet)) =
R\Gamma(X,
\text{Tot}(\mathcal{K}_1^\bullet
\otimes_{\mathcal{O}_X} \mathcal{M}_1^\bullet))
$$
は $D(A)$ において同じである。したがって，$\mathcal{K}$ と $\mathcal{M}$ を
$\mathcal{K}_1$ と $\mathcal{M}_1$ で置き換えた場合に図式が可換であることを
証明すれば十分である。こうして次の段落で扱う場合に帰着する。

\medskip\noindent
$\mathcal{K}^\bullet$ と $\mathcal{M}^\bullet$ がフラスク
$\mathcal{O}_X$-加群からなる下に有界な複体であると仮定し，
カップ積と {\v C}ech 複体上のカップ積
(\ref{cohomology-equation-needs-signs}) を関係づける図式を考える。
このとき次の可換図式を考えられる。
$$
\xymatrix{
\Gamma(X, \mathcal{K}^\bullet)
\otimes_A^\mathbf{L}
\Gamma(X, \mathcal{M}^\bullet) \ar[d] \ar[r] &
\text{Tot}(\check{\mathcal{C}}^\bullet(\mathcal{U}, \mathcal{K}^\bullet))
\otimes_A^\mathbf{L}
\text{Tot}(\check{\mathcal{C}}^\bullet(\mathcal{U}, \mathcal{M}^\bullet))
\ar[d] \\
\text{Tot}(\Gamma(X, \mathcal{K}^\bullet)
\otimes_A
\Gamma(X, \mathcal{M}^\bullet)) \ar[d] \ar[r] &
\text{Tot}(
\text{Tot}(\check{\mathcal{C}}^\bullet(\mathcal{U}, \mathcal{K}^\bullet))
\otimes_A
\text{Tot}(\check{\mathcal{C}}^\bullet(\mathcal{U}, \mathcal{M}^\bullet)))
\ar[d]^{(\ref{cohomology-equation-needs-signs})} \\
\Gamma(X, \text{Tot}(\mathcal{K}^\bullet
\otimes_{\mathcal{O}_X}
\mathcal{M}^\bullet)) \ar[r] &
\text{Tot}(
\check{\mathcal{C}}^\bullet({\mathcal U},
\text{Tot}(\mathcal{K}^\bullet \otimes_{\mathcal{O}_X} \mathcal{M}^\bullet)
))
}
$$
この図式の水平射は $D(A)$ における同型である。実際，
$\mathcal{K}^\bullet$ のようなフラスク加群からなる下に有界な複体に対し，
$$
\Gamma(X, \mathcal{K}^\bullet) =
\text{Tot}(\check{\mathcal{C}}^\bullet(\mathcal{U}, \mathcal{K}^\bullet)) =
R\Gamma(X, \mathcal{K}^\bullet)
$$
が $D(A)$ において成り立つからである。これは補題
\ref{cohomology-lemma-flasque-acyclic}，《導来圏》補題
\ref{derived-lemma-leray-acyclicity}，および補題
\ref{cohomology-lemma-cech-complex-complex-computes} から従う。
したがって，(\ref{cohomology-equation-needs-signs}) を含む補題の図式の可換性は，
$f$ が一点への射である場合にすでに証明した補題
\ref{cohomology-lemma-cup-compatible-with-naive} の可換性から従う
（補題 \ref{cohomology-lemma-cup-compatible-with-naive} の後の議論を参照）。
\end{proof}

\begin{lemma}
\label{cohomology-lemma-cup-product-associative}
$f : (X, \mathcal{O}_X) \to (Y, \mathcal{O}_Y)$ を環付き空間の射とする。
注意 \ref{cohomology-remark-cup-product} の相対カップ積は結合的である。すなわち，
$D(\mathcal{O}_X)$ のすべての $K, L, M$ に対して図式
$$
\xymatrix{
Rf_*K \otimes_{\mathcal{O}_Y}^\mathbf{L}
Rf_*L \otimes_{\mathcal{O}_Y}^\mathbf{L}
Rf_*M \ar[r] \ar[d] &
Rf_*(K \otimes_{\mathcal{O}_X}^\mathbf{L} L)
\otimes_{\mathcal{O}_Y}^\mathbf{L} Rf_*M \ar[d] \\
Rf_*K \otimes_{\mathcal{O}_Y}^\mathbf{L}
Rf_*(L \otimes_{\mathcal{O}_X}^\mathbf{L} M) \ar[r] &
Rf_*(K \otimes_{\mathcal{O}_X}^\mathbf{L} 
L \otimes_{\mathcal{O}_X}^\mathbf{L} M)
}
$$
は $D(\mathcal{O}_Y)$ において可換である。
\end{lemma}

\begin{proof}
どちらの側を進んでも，$D(\mathcal{O}_X)$ における自明な写像
\begin{align*}
Lf^*(Rf_*K \otimes_{\mathcal{O}_Y}^\mathbf{L}
Rf_*L \otimes_{\mathcal{O}_Y}^\mathbf{L}
Rf_*M) & =
Lf^*(Rf_*K) \otimes_{\mathcal{O}_X}^\mathbf{L}
Lf^*(Rf_*L) \otimes_{\mathcal{O}_X}^\mathbf{L}
Lf^*(Rf_*M) \\
& \to
K \otimes_{\mathcal{O}_X}^\mathbf{L} 
L \otimes_{\mathcal{O}_X}^\mathbf{L} M
\end{align*}
に随伴する写像が得られる。
\end{proof}

\begin{lemma}
\label{cohomology-lemma-cup-product-commutative}
$f : (X, \mathcal{O}_X) \to (Y, \mathcal{O}_Y)$ を環付き空間の射とする。
注意 \ref{cohomology-remark-cup-product} の相対カップ積は可換である。すなわち，
$D(\mathcal{O}_X)$ のすべての $K, L$ に対して図式
$$
\xymatrix{
Rf_*K \otimes_{\mathcal{O}_Y}^\mathbf{L} Rf_*L \ar[r] \ar[d]_\psi &
Rf_*(K \otimes_{\mathcal{O}_X}^\mathbf{L} L) \ar[d]^{Rf_*\psi} \\
Rf_*L \otimes_{\mathcal{O}_Y}^\mathbf{L} Rf_*K \ar[r] &
Rf_*(L \otimes_{\mathcal{O}_X}^\mathbf{L} K)
}
$$
は $D(\mathcal{O}_Y)$ において可換である。ここで $\psi$ は導来圏上の
可換性制約である（補題 \ref{cohomology-lemma-symmetric-monoidal-derived}）。
\end{lemma}

\begin{proof}
省略する。
\end{proof}

\begin{lemma}
\label{cohomology-lemma-compose-cup-product}
$f : (X, \mathcal{O}_X) \to (Y, \mathcal{O}_Y)$ および
$g : (Y, \mathcal{O}_Y) \to (Z, \mathcal{O}_Z)$ を環付き空間の射とする。
注意 \ref{cohomology-remark-cup-product} の相対カップ積は合成と整合する。すなわち，
$D(\mathcal{O}_X)$ のすべての $K, L$ に対して図式
$$
\xymatrix{
R(g \circ f)_*K \otimes_{\mathcal{O}_Z}^\mathbf{L} R(g \circ f)_*L
\ar@{=}[rr] \ar[d] & &
Rg_*Rf_*K \otimes_{\mathcal{O}_Z}^\mathbf{L} Rg_*Rf_*L \ar[d] \\
R(g \circ f)_*(K \otimes_{\mathcal{O}_X}^\mathbf{L} L) \ar@{=}[r] &
Rg_*Rf_*(K \otimes_{\mathcal{O}_X}^\mathbf{L} L) &
Rg_*(Rf_*K \otimes_{\mathcal{O}_Y}^\mathbf{L}  Rf_*L) \ar[l]
}
$$
は $D(\mathcal{O}_Z)$ において可換である。
\end{lemma}

\begin{proof}
図式をいずれの向きに進んでも，$D(\mathcal{O}_X)$ における写像
\begin{align*}
& L(g \circ f)^*\left(R(g \circ f)_*K
\otimes_{\mathcal{O}_Z}^\mathbf{L}
R(g \circ f)_*L\right) \\
& =
L(g \circ f)^*R(g \circ f)_*K
\otimes_{\mathcal{O}_X}^\mathbf{L}
L(g \circ f)^*R(g \circ f)_*L) \\
& \to
K \otimes_{\mathcal{O}_X}^\mathbf{L} L
\end{align*}
に随伴する写像が得られるからである。これを見るには，余単位の合成
$$
L(g \circ f)^*R(g \circ f)_* =
Lf^* Lg^* Rg_* Rf_*  \to
Lf^* Rf_* \to \text{id}
$$
が $L(g \circ f)^*$ と $R(g \circ f)_*$ に対する余単位であることを用いる。
《圏論》補題 \ref{categories-lemma-compose-counits} を参照せよ。
\end{proof}

\begin{lemma}
\label{cohomology-lemma-base-change-cup-product}
環付き空間の可換正方形
$$
\xymatrix{
(X', \mathcal{O}_{X'}) \ar[r]_{g'} \ar[d]_{f'} &
(X, \mathcal{O}_X) \ar[d]^f \\
(Y', \mathcal{O}_{Y'}) \ar[r]^g & (Y, \mathcal{O}_Y) 
}
$$
を考える。$K, L$ を $D(\mathcal{O}_X)$ の対象とする。
相対カップ積はこの正方形と整合する。すなわち，図式
$$
\xymatrix{
Lg^*(Rf_*K \otimes_{\mathcal{O}_Y}^\mathbf{L} Rf_*L)
\ar[r] \ar@{=}[d] &
Lg^*(Rf_*(K \otimes_{\mathcal{O}_X}^\mathbf{L} L)) \ar[d] \\
Lg^*Rf_*K \otimes_{\mathcal{O}_{Y'}}^\mathbf{L} Lg^*Rf_*L \ar[d] &
R(f')_*L(g')^*(K \otimes_{\mathcal{O}_X}^\mathbf{L} L) \ar@{=}[d] \\
R(f')_*(L(g')^*K \otimes_{\mathcal{O}_{Y'}} R(f')_*(L(g')^*L \ar[r] &
R(f')_*(L(g')^*K \otimes_{\mathcal{O}_{X'}}^\mathbf{L} L(g')^*L)
}
$$
は $D(\mathcal{O}_{Y'})$ において可換である。水平射は相対カップ積
（注意 \ref{cohomology-remark-cup-product}）で与えられ，垂直射は基底変換写像
（注意 \ref{cohomology-remark-base-change}）および補題
\ref{cohomology-lemma-pullback-tensor-product} によって与えられる。
\end{lemma}

\begin{proof}
省略する。
\end{proof}














\section{K-入射的複体のいくつかの性質}
\label{cohomology-section-properties-K-injective}

\noindent
$(X, \mathcal{O}_X)$ を環付き空間とし，$U \subset X$ を開部分集合とする。
対応する開埋め込みを
$j : (U, \mathcal{O}_U) \to (X, \mathcal{O}_X)$ と書く。
引き戻し関手 $j^*$ は単なる制限関手なので完全である。したがって，
導来引き戻し $Lj^*$ は任意の複体を単に制限することで計算される。
対応する関手をしばしば単に
$$
D(\mathcal{O}_X) \to D(\mathcal{O}_U),
\quad
E \mapsto j^*E = E|_U
$$
と書く。同様に，零による延長
$j_! : \textit{Mod}(\mathcal{O}_U) \to \textit{Mod}(\mathcal{O}_X)$
（《層》第 \ref{sheaves-section-open-immersions} 節を参照）は完全関手である
（《加群》補題 \ref{modules-lemma-j-shriek-exact}）。したがって関手
$$
j_! : D(\mathcal{O}_U) \to D(\mathcal{O}_X),\quad
F \mapsto j_!F
$$
を誘導する。これは対象 $F$ を表す任意の複体に $j_!$ を施すだけで得られる。

\begin{lemma}
\label{cohomology-lemma-restrict-K-injective-to-open}
$X$ を環付き空間とし，$U \subset X$ を開部分空間とする。
$\mathcal{O}_X$-加群の K-入射的複体を $U$ に制限したものは，
$\mathcal{O}_U$-加群の K-入射的複体である。
\end{lemma}

\begin{proof}
これは《導来圏》補題
\ref{derived-lemma-adjoint-preserve-K-injectives} と，制限関手が
完全な左随伴 $j_!$ をもつことから従う。$j_!$ の構成については
《層》第 \ref{sheaves-section-open-immersions} 節を，完全性については
《加群》補題 \ref{modules-lemma-j-shriek-exact} を参照せよ。
\end{proof}

\begin{lemma}
\label{cohomology-lemma-unbounded-cohomology-of-open}
$X$ を環付き空間とし，$U \subset X$ を開部分空間とする。
$D(\mathcal{O}_X)$ の $K$ に対して
$H^p(U, K) = H^p(U, K|_U)$ である。
\end{lemma}

\begin{proof}
$\mathcal{I}^\bullet$ を $K$ を表す $\mathcal{O}_X$-加群の
K-入射的複体とする。このとき
$$
H^q(U, K) = H^q(\Gamma(U, \mathcal{I}^\bullet)) =
H^q(\Gamma(U, \mathcal{I}^\bullet|_U))
$$
である。これはコホモロジーの構成による。補題
\ref{cohomology-lemma-restrict-K-injective-to-open} により，
複体 $\mathcal{I}^\bullet|_U$ は $K|_U$ を表す K-入射的複体であり，
補題が従う。
\end{proof}

\begin{lemma}
\label{cohomology-lemma-sheafification-cohomology}
$(X, \mathcal{O}_X)$ を環付き空間とし，$K$ を
$D(\mathcal{O}_X)$ の対象とする。
$$
U \mapsto H^q(U, K) = H^q(U, K|_U)
$$
の層化は $K$ の第 $q$ コホモロジー層 $H^q(K)$ である。
\end{lemma}

\begin{proof}
等式 $H^q(U, K) = H^q(U, K|_U)$ は補題
\ref{cohomology-lemma-unbounded-cohomology-of-open} による。
$K$ を表す K-入射的複体 $\mathcal{I}^\bullet$ を選ぶ。このとき次式が成り立つ。
$$
H^q(U, K) =
\frac{\Ker(\mathcal{I}^q(U) \to \mathcal{I}^{q + 1}(U))}
{\Im(\mathcal{I}^{q - 1}(U) \to \mathcal{I}^q(U))}.
$$
これはコホモロジーの構成による。また，
$H^q(K) = \Ker(\mathcal{I}^q \to \mathcal{I}^{q + 1})/
\Im(\mathcal{I}^{q - 1} \to \mathcal{I}^q)$ であるから，結果は明らかである。
\end{proof}

\begin{lemma}
\label{cohomology-lemma-restrict-direct-image-open}
$f : (X, \mathcal{O}_X) \to (Y, \mathcal{O}_Y)$ を環付き空間の射とする。
開部分空間 $V \subset Y$ が与えられたとき，$U = f^{-1}(V)$ と置き，
誘導される射を $g : U \to V$ と書く。このとき
$D(\mathcal{O}_X)$ の $E$ に対して
$(Rf_*E)|_V = Rg_*(E|_U)$ である。
\end{lemma}

\begin{proof}
$E$ を $\mathcal{O}_X$-加群の K-入射的複体
$\mathcal{I}^\bullet$ で表す。このとき
$Rf_*(E) = f_*\mathcal{I}^\bullet$ であり，補題
\ref{cohomology-lemma-restrict-K-injective-to-open} により
$Rg_*(E|_U) = g_*(\mathcal{I}^\bullet|_U)$ である。
$X$ 上の任意の層 $\mathcal{F}$ に対して
$(f_*\mathcal{F})|_V = g_*(\mathcal{F}|_U)$ であることは明らかなので，
結果が従う。
\end{proof}

\begin{lemma}
\label{cohomology-lemma-Leray-unbounded}
$f : X \to Y$ を環付き空間の射とする。このとき
$D(\mathcal{O}_X) \to D(\Gamma(Y, \mathcal{O}_Y))$ の関手として
$R\Gamma(Y, -) \circ Rf_* = R\Gamma(X, -)$ である。
より一般に，$V \subset Y$ が開で $U = f^{-1}(V)$ ならば
$R\Gamma(U, -) = R\Gamma(V, -) \circ Rf_*$ である。
\end{lemma}

\begin{proof}
$Z$ を，$\Gamma(Z, \mathcal{O}_Z) = \Gamma(Y, \mathcal{O}_Y)$ を満たす
一点空間からなる環付き空間とする。構造層の大域切断上でこの同一視を
誘導する環付き空間の標準射 $Y \to Z$ がある。このとき
$D(\mathcal{O}_Z) = D(\Gamma(Y, \mathcal{O}_Y))$ である。
したがって主張 $R\Gamma(Y, -) \circ Rf_* = R\Gamma(X, -)$ は，
補題 \ref{cohomology-lemma-derived-pushforward-composition} を
$X \to Y \to Z$ に適用して得られる。

\medskip\noindent
第二の（より一般的な）主張は，第一の主張に補題
\ref{cohomology-lemma-restrict-direct-image-open} を適用して従う。
\end{proof}

\begin{lemma}
\label{cohomology-lemma-unbounded-describe-higher-direct-images}
$f : (X, \mathcal{O}_X) \to (Y, \mathcal{O}_Y)$ を環付き空間の射とし，
$K$ を $D(\mathcal{O}_X)$ の対象とする。このとき $H^i(Rf_*K)$ は，
前層
$$
V \mapsto H^i(f^{-1}(V), K) = H^i(V, Rf_*K)
$$
に付随する層である。
\end{lemma}

\begin{proof}
等式 $H^i(f^{-1}(V), K) = H^i(V, Rf_*K)$ は，補題
\ref{cohomology-lemma-Leray-unbounded} の第二の主張でコホモロジーを取って得られる。
層化に関する主張は補題
\ref{cohomology-lemma-sheafification-cohomology} から従う。
\end{proof}

\begin{lemma}
\label{cohomology-lemma-modules-abelian-unbounded}
$X$ を環付き空間とし，$K$ を $D(\mathcal{O}_X)$ の対象とする。
$D(\underline{\mathbf{Z}}_X)$ におけるその像を $K_{ab}$ と書く。
\begin{enumerate}
\item 任意の開集合 $U \subset X$ に対して標準写像
$R\Gamma(U, K) \to R\Gamma(U, K_{ab})$
があり，これは $D(\textit{Ab})$ における同型である。
\item $f : X \to Y$ を環付き空間の射とする。標準写像
$Rf_*K \to Rf_*(K_{ab})$ があり，これは
$D(\underline{\mathbf{Z}}_Y)$ における同型である。
\end{enumerate}
\end{lemma}

\begin{proof}
写像を次のように構成する。$K$ を表す K-入射的複体
$\mathcal{I}^\bullet$ を選ぶ。また，$\mathcal{J}^\bullet$ が
アーベル群の層の K-入射的複体であるような擬同型
$\mathcal{I}^\bullet \to \mathcal{J}^\bullet$ を選ぶ。このとき (1) の写像は
$\Gamma(U, \mathcal{I}^\bullet) \to \Gamma(U, \mathcal{J}^\bullet)$ で与えられ，
(2) の写像は $f_*\mathcal{I}^\bullet \to f_*\mathcal{J}^\bullet$ で与えられる。
これらの写像が同型であることを示すには，コホモロジー群および
コホモロジー層上で同型を誘導することを証明すれば十分である。補題
\ref{cohomology-lemma-unbounded-cohomology-of-open} および
\ref{cohomology-lemma-unbounded-describe-higher-direct-images}
により，写像
$$
H^0(X, K) \longrightarrow H^0(X, K_{ab})
$$
が同型であることを示せば十分である。ここで
$$
H^0(X, K) = \Hom_{D(\mathcal{O}_X)}(\mathcal{O}_X, K)
$$
であり，もう一方の群についても同様である。$K$ を表す任意の
$\mathcal{O}_X$-加群の複体 $\mathcal{K}^\bullet$ を選ぶ。
導来圏を局所化として構成することにより，
$$
\Hom_{D(\mathcal{O}_X)}(\mathcal{O}_X, K) =
\colim_{s : \mathcal{F}^\bullet \to \mathcal{O}_X}
\Hom_{K(\mathcal{O}_X)}(\mathcal{F}^\bullet, \mathcal{K}^\bullet)
$$
を得る。ここで余極限は，$\mathcal{O}_X$-加群の複体の擬同型 $s$ 全体に
わたって取る。同様に，
$$
\Hom_{D(\underline{\mathbf{Z}}_X)}(\underline{\mathbf{Z}}_X, K) =
\colim_{s : \mathcal{G}^\bullet \to \underline{\mathbf{Z}}_X}
\Hom_{K(\underline{\mathbf{Z}}_X)}(\mathcal{G}^\bullet, \mathcal{K}^\bullet)
$$
を得る。次に，$\mathcal{G}^\bullet$ が平坦
$\underline{\mathbf{Z}}_X$-加群からなる上に有界な複体であるような擬同型
$s : \mathcal{G}^\bullet \to \underline{\mathbf{Z}}_X$
はこの系で共終であることに注意する（これは《加群》補題
\ref{modules-lemma-module-quotient-flat} および《導来圏》補題
\ref{derived-lemma-subcategory-left-resolution} から従う。第
\ref{cohomology-section-flat} 節の議論を参照せよ）。したがって写像
$H^0(X, K) \longrightarrow H^0(X, K_{ab})$
の逆を次のように構成できる。元 $\xi \in H^0(X, K_{ab})$ を対
$$
(s : \mathcal{G}^\bullet \to \underline{\mathbf{Z}}_X,
a : \mathcal{G}^\bullet \to \mathcal{K}^\bullet)
$$
で表す。ここで $\mathcal{G}^\bullet$ は平坦
$\underline{\mathbf{Z}}_X$-加群からなる上に有界な複体である。
これを
$$
(\mathcal{G}^\bullet \otimes_{\underline{\mathbf{Z}}_X} \mathcal{O}_X
\to \mathcal{O}_X,
\mathcal{G}^\bullet  \otimes_{\underline{\mathbf{Z}}_X} \mathcal{O}_X
\to \mathcal{K}^\bullet)
$$
に送る。ここで注意すべき唯一の点は，補題
\ref{cohomology-lemma-derived-tor-quasi-isomorphism-other-side} および
\ref{cohomology-lemma-bounded-flat-K-flat} により，第一の矢印が擬同型であることである。
この構成が実際に逆写像であることの詳細な検証は省略する。
\end{proof}

\begin{lemma}
\label{cohomology-lemma-adjoint-lower-shriek-restrict}
$(X, \mathcal{O}_X)$ を環付き空間とし，$U \subset X$ を開部分集合とする。
対応する開埋め込みを
$j : (U, \mathcal{O}_U) \to (X, \mathcal{O}_X)$ と書く。制限関手
$D(\mathcal{O}_X) \to D(\mathcal{O}_U)$ は，零による延長
$j_! : D(\mathcal{O}_U) \to D(\mathcal{O}_X)$ の右随伴である。
\end{lemma}

\begin{proof}
これは $j_!$ と $j^*$ が随伴かつ完全であること
（したがって $Lj_! = j_!$ と $Rj^* = j^*$ が存在する）から形式的に従う。
《導来圏》補題 \ref{derived-lemma-derived-adjoint-functors} を参照せよ。
\end{proof}

\begin{lemma}
\label{cohomology-lemma-K-injective-flat}
$f : X \to Y$ を環付き空間の平坦射とする。
$\mathcal{I}^\bullet$ が $\mathcal{O}_X$-加群の K-入射的複体ならば，
$f_*\mathcal{I}^\bullet$ は $\mathcal{O}_Y$-加群の複体として K-入射的である。
\end{lemma}

\begin{proof}
実際，
$$
\Hom_{K(\mathcal{O}_Y)}(\mathcal{F}^\bullet, f_*\mathcal{I}^\bullet)
=
\Hom_{K(\mathcal{O}_X)}(f^*\mathcal{F}^\bullet, \mathcal{I}^\bullet)
$$
である。これは《層》補題
\ref{sheaves-lemma-adjoint-pullback-pushforward-modules}
と，$f$ が平坦であるため $f^*$ が完全であることによる。
\end{proof}





\section{非有界 Mayer--Vietoris}
\label{cohomology-section-unbounded-mayer-vietoris}

\noindent
非有界コホモロジーに対しても Mayer--Vietoris 系列が存在する。

\begin{lemma}
\label{cohomology-lemma-exact-sequence-lower-shriek}
$(X, \mathcal{O}_X)$ を環付き空間とし，
$X = U \cup V$ を二つの開部分空間の和とする。
$D(\mathcal{O}_X)$ の任意の対象 $E$ に対して，完全三角
$$
j_{U \cap V!}E|_{U \cap V} \to
j_{U!}E|_U \oplus j_{V!}E|_V \to E \to 
j_{U \cap V!}E|_{U \cap V}[1]
$$
が $D(\mathcal{O}_X)$ に存在する。
\end{lemma}

\begin{proof}
第 \ref{cohomology-section-properties-K-injective} 節で，制限関手と
零による延長関手はいずれも，任意の複体にその関手を施すだけで
計算されることを見た。$\mathcal{E}^\bullet$ を $E$ を表す
$\mathcal{O}_X$-加群の複体とする。補題の完全三角は，
$\mathcal{O}_X$-加群の複体の短完全列
$$
0 \to j_{U \cap V!}\mathcal{E}^\bullet|_{U \cap V} \to
j_{U!}\mathcal{E}^\bullet|_U \oplus j_{V!}\mathcal{E}^\bullet|_V
\to \mathcal{E}^\bullet \to 0
$$
に付随する完全三角である（《導来圏》第
\ref{derived-section-canonical-delta-functor} 節，特に補題
\ref{derived-lemma-derived-canonical-delta-functor} による）。
この列が完全であることは，《層》補題
\ref{sheaves-lemma-j-shriek-modules} を用いて茎上で確認できる
（計算は省略する）。
\end{proof}

\begin{lemma}
\label{cohomology-lemma-exact-sequence-j-star}
$(X, \mathcal{O}_X)$ を環付き空間とし，
$X = U \cup V$ を二つの開部分空間の和とする。
$D(\mathcal{O}_X)$ の任意の対象 $E$ に対して，完全三角
$$
E \to 
Rj_{U, *}E|_U \oplus Rj_{V, *}E|_V \to
Rj_{U \cap V, *}E|_{U \cap V} \to
E[1]
$$
が $D(\mathcal{O}_X)$ に存在する。
\end{lemma}

\begin{proof}
各項 $\mathcal{I}^n$ が $\textit{Mod}(\mathcal{O}_X)$ の入射対象であるような，
$E$ を表す K-入射的複体 $\mathcal{I}^\bullet$ を選ぶ。
《入射対象》定理
\ref{injectives-theorem-K-injective-embedding-grothendieck}
を参照せよ。$\mathcal{I}^\bullet|U$ も K-入射的複体であることを
すでに見た（補題 \ref{cohomology-lemma-restrict-K-injective-to-open}）。したがって
$Rj_{U, *}E|_U$ は $j_{U, *}\mathcal{I}^\bullet|_U$ で表される。
$V$ と $U \cap V$ についても同様である。ゆえに補題の完全三角は，
複体の短完全列
$$
0 \to
\mathcal{I}^\bullet \to
j_{U, *}\mathcal{I}^\bullet|_U \oplus j_{V, *}\mathcal{I}^\bullet|_V \to
j_{U \cap V, *}\mathcal{I}^\bullet|_{U \cap V} \to
0.
$$
に付随する完全三角である（《導来圏》第
\ref{derived-section-canonical-delta-functor} 節，特に補題
\ref{derived-lemma-derived-canonical-delta-functor} による）。
この列が完全であるのは，任意の開集合 $W \subset X$ と任意の $n$ に対して列
$$
0 \to
\mathcal{I}^n(W) \to
\mathcal{I}^n(W \cap U) \oplus \mathcal{I}^n(W \cap V) \to
\mathcal{I}^n(W \cap U \cap V) \to
0
$$
が完全だからである（補題 \ref{cohomology-lemma-mayer-vietoris} の証明を参照）。
\end{proof}

\begin{lemma}
\label{cohomology-lemma-mayer-vietoris-hom}
$(X, \mathcal{O}_X)$ を環付き空間とし，
$X = U \cup V$ を $X$ の二つの開部分空間の和とする。
$D(\mathcal{O}_X)$ の対象 $E$, $F$ に対して Mayer--Vietoris 系列
$$
\xymatrix{
& \ldots \ar[r] & \Ext^{-1}(E_{U \cap V}, F_{U \cap V}) \ar[lld] \\
\Hom(E, F) \ar[r] &
\Hom(E_U, F_U) \oplus
\Hom(E_V, F_V) \ar[r] &
\Hom(E_{U \cap V}, F_{U \cap V})
}
$$
が存在する。ここで下付き添字は対応する開集合への制限を表し，
$\Hom$ と $\Ext$ は対応する導来圏で取る。
\end{lemma}

\begin{proof}
補題 \ref{cohomology-lemma-exact-sequence-lower-shriek} の完全三角を用いて
$\Hom$ の長完全列を得る（《導来圏》補題
\ref{derived-lemma-representable-homological} による）。さらに，
$$
\Hom_{D(\mathcal{O}_X)}(j_{U!}E|_U, F) =
\Hom_{D(\mathcal{O}_U)}(E|_U, F|_U)
$$
を補題 \ref{cohomology-lemma-adjoint-lower-shriek-restrict} により用いる。
\end{proof}

\begin{lemma}
\label{cohomology-lemma-unbounded-mayer-vietoris}
$(X, \mathcal{O}_X)$ を環付き空間とし，
$X = U \cup V$ が二つの開部分集合の和であると仮定する。
$D(\mathcal{O}_X)$ の対象 $E$ に対して完全三角
$$
R\Gamma(X, E) \to R\Gamma(U, E) \oplus R\Gamma(V, E) \to
R\Gamma(U \cap V, E) \to R\Gamma(X, E)[1]
$$
があり，特にコホモロジーの長完全列
$$
\ldots \to
H^n(X, E) \to
H^n(U, E) \oplus H^0(V, E) \to
H^n(U \cap V, E) \to
H^{n + 1}(X, E) \to \ldots
$$
が得られる。完全三角と長完全列の構成は $E$ に関して関手的である。
\end{lemma}

\begin{proof}
各 $n$ に対して $\mathcal{I}^n$ が
$\textit{Mod}(\mathcal{O}_X)$ の入射対象であるような，
$E$ を表す K-入射的複体 $\mathcal{I}^\bullet$ を選べる。
《入射対象》定理
\ref{injectives-theorem-K-injective-embedding-grothendieck}
を参照せよ。このとき $R\Gamma(X, E)$ は
$\Gamma(X, \mathcal{I}^\bullet)$ で計算される。補題
\ref{cohomology-lemma-restrict-K-injective-to-open} により，
$U$, $V$, $U \cap V$ についても同様である。したがって補題の完全三角は，
複体の短完全列
$$
0 \to
\mathcal{I}^\bullet(X) \to
\mathcal{I}^\bullet(U) \oplus \mathcal{I}^\bullet(V) \to
\mathcal{I}^\bullet(U \cap V) \to
0.
$$
に付随する完全三角である（《導来圏》第
\ref{derived-section-canonical-delta-functor} 節，特に補題
\ref{derived-lemma-derived-canonical-delta-functor} による）。
これが短完全列であることは補題
\ref{cohomology-lemma-mayer-vietoris} の証明で見た。最後の主張は
《入射対象》定理
\ref{injectives-theorem-K-injective-embedding-grothendieck}
における構成の関手性から従う。
\end{proof}

\begin{lemma}
\label{cohomology-lemma-unbounded-relative-mayer-vietoris}
$f : X \to Y$ を環付き空間の射とし，
$X = U \cup V$ が二つの開部分集合の和であると仮定する。
$a = f|_U : U \to Y$, $b = f|_V : V \to Y$，および
$c = f|_{U \cap V} : U \cap V \to Y$ と書く。
$D(\mathcal{O}_X)$ の任意の対象 $E$ に対して完全三角
$$
Rf_*E \to
Ra_*(E|_U) \oplus Rb_*(E|_V) \to
Rc_*(E|_{U \cap V}) \to
Rf_*E[1]
$$
が存在する。この三角は $E$ に関して関手的である。
\end{lemma}

\begin{proof}
各 $n$ に対して $\mathcal{I}^n$ が
$\textit{Mod}(\mathcal{O}_X)$ の入射対象であるような，
$E$ を表す K-入射的複体 $\mathcal{I}^\bullet$ を選べる。
《入射対象》定理
\ref{injectives-theorem-K-injective-embedding-grothendieck}
を参照せよ。このとき $Rf_*E$ は $f_*\mathcal{I}^\bullet$ で計算される。
補題 \ref{cohomology-lemma-restrict-K-injective-to-open} により，
$U$, $V$, $U \cap V$ についても同様である。したがって補題の完全三角は，
複体の短完全列
$$
0 \to
f_*\mathcal{I}^\bullet \to
a_*\mathcal{I}^\bullet|_U \oplus b_*\mathcal{I}^\bullet|_V \to
c_*\mathcal{I}^\bullet|_{U \cap V} \to
0.
$$
に付随する完全三角である（《導来圏》第
\ref{derived-section-canonical-delta-functor} 節，特に補題
\ref{derived-lemma-derived-canonical-delta-functor} による）。
これは補題 \ref{cohomology-lemma-relative-mayer-vietoris} と，
$\textit{Mod}(\mathcal{O}_X)$ の入射対象 $\mathcal{I}$ に対して
$R^1f_*\mathcal{I} = 0$ であることにより，複体の短完全列である。
最後の主張は《入射対象》定理
\ref{injectives-theorem-K-injective-embedding-grothendieck}
における構成の関手性から従う。
\end{proof}

\begin{lemma}
\label{cohomology-lemma-pushforward-restriction}
$(X, \mathcal{O}_X)$ を環付き空間とし，$j : U \to X$ を
開部分空間の包含とする。$T \subset X$ を $U$ に含まれる閉部分集合とする。
\begin{enumerate}
\item $E$ が $D(\mathcal{O}_X)$ の対象で，そのコホモロジー層の台が
$T$ に含まれるならば，$E \to Rj_*(E|_U)$ は同型である。
\item $F$ が $D(\mathcal{O}_U)$ の対象で，そのコホモロジー層の台が
$T$ に含まれるならば，$j_!F \to Rj_*F$ は同型である。
\end{enumerate}
\end{lemma}

\begin{proof}
$V = X \setminus T$ および $W = U \cap V$ と置く。
$X = U \cup V$ は $X$ の開被覆である。開埋め込みを
$j_W : W \to V$ と書く。$E$ を，コホモロジー層の台が $T$ に含まれる
$D(\mathcal{O}_X)$ の対象とする。補題
\ref{cohomology-lemma-restrict-direct-image-open} により
$(Rj_*E|_U)|_V = Rj_{W, *}(E|_W) = 0$ である。実際，仮定により
$E|_W = 0$ である。一方，$Rj_*(E|_U)|_U = E|_U$ である。
したがって (1) は明らかである。
$F$ を，コホモロジー層の台が $T$ に含まれる
$D(\mathcal{O}_U)$ の対象とする。補題
\ref{cohomology-lemma-restrict-direct-image-open} により
$(Rj_*F)|_V = Rj_{W, *}(F|_W) = 0$ である。実際，仮定により
$F|_W = 0$ である。また $(j_!F)|_V = j_{W!}(F|_W) = 0$ である
（第一の等式は零による延長の定義から直ちに従う）。
$(Rj_*F)|_U = F$ かつ $(j_!F)|_U = F$ なので，(2) が成り立つ。
\end{proof}

\begin{lemma}
\label{cohomology-lemma-mayer-vietoris-cup}
$(X, \mathcal{O}_X)$ を環付き空間とし，
$A = \Gamma(X, \mathcal{O}_X)$ と置く。
$X = U \cup V$ が二つの開部分集合の和であると仮定する。
$D(\mathcal{O}_X)$ の対象 $K$ と $M$ に対して，完全三角の写像
$$
\xymatrix{
R\Gamma(X, K) \otimes_A^\mathbf{L} R\Gamma(X, M) \ar[r] \ar[d] &
R\Gamma(X, K \otimes_{\mathcal{O}_X}^\mathbf{L} M) \ar[d] \\
R\Gamma(X, K) \otimes_A^\mathbf{L}
(R\Gamma(U, M) \oplus R\Gamma(V, M)) \ar[r] \ar[d] &
R\Gamma(U, K \otimes_{\mathcal{O}_X}^\mathbf{L} M)
\oplus R\Gamma(V, K \otimes_{\mathcal{O}_X}^\mathbf{L} M)) \ar[d] \\
R\Gamma(X, K) \otimes_A^\mathbf{L} R\Gamma(U \cap V, M) \ar[r] \ar[d] &
R\Gamma(U \cap V, K \otimes_{\mathcal{O}_X}^\mathbf{L} M) \ar[d] \\
R\Gamma(X, K) \otimes_A^\mathbf{L} R\Gamma(X, M)[1] \ar[r] &
R\Gamma(X, K \otimes_{\mathcal{O}_X}^\mathbf{L} M)[1]
}
$$
が存在する。ここで，
\begin{enumerate}
\item 水平射はカップ積で与えられる。
\item 右辺は，$K \otimes_{\mathcal{O}_X}^\mathbf{L} M$ に対する補題
\ref{cohomology-lemma-unbounded-mayer-vietoris} の完全三角である。
\item 左辺は，$M$ に対する補題
\ref{cohomology-lemma-unbounded-mayer-vietoris} の完全三角に完全関手
$R\Gamma(X, K) \otimes_A^\mathbf{L} - $ を適用したものである。
\end{enumerate}
\end{lemma}

\begin{proof}
$R\Gamma(X, K)$ を表す，平坦 $A$-加群からなる K-平坦複体
$T^\bullet$ を選ぶ。《代数続論》補題
\ref{more-algebra-lemma-K-flat-resolution} を参照せよ。
環付き空間の射 $(X, \mathcal{O}_X) \to (pt, A)$ による
$T^\bullet$ の引き戻しを $T^\bullet \otimes_A \mathcal{O}_X$ と書く。
$D(\mathcal{O}_X)$ に自然な随伴写像
$\epsilon : T^\bullet \otimes_A \mathcal{O}_X \to K$ がある。
$T^\bullet \otimes_A \mathcal{O}_X$ は，平坦な項をもつ
$\mathcal{O}_X$-加群の K-平坦複体である。補題
\ref{cohomology-lemma-pullback-K-flat} および《加群》補題
\ref{modules-lemma-pullback-flat} を参照せよ。
補題 \ref{cohomology-lemma-factor-through-K-flat} により，$\epsilon$ を表す
$\mathcal{O}_X$-加群の複体の射
$$
T^\bullet \otimes_A \mathcal{O}_X \longrightarrow \mathcal{K}^\bullet
$$
を，$\mathcal{K}^\bullet$ が平坦な項をもつ K-平坦複体となるように選べる。
実際，$D(\mathcal{O}_X)$ の構成により，まず $\epsilon$ を
$\mathcal{O}_X$-加群の複体のある写像
$e : T^\bullet \otimes_A \mathcal{O}_X \to \mathcal{L}^\bullet$
で表し，次に $e$ に補題を適用できる。各項が入射的
$\mathcal{O}_X$-加群である，$M$ を表す K-入射的複体
$\mathcal{I}^\bullet$ を選ぶ。最後に擬同型
$$
\text{Tot}(\mathcal{K}^\bullet \otimes_\mathcal{O} \mathcal{I}^\bullet)
\longrightarrow
\mathcal{J}^\bullet
$$
を，各項が入射的 $\mathcal{O}_X$-加群である K-入射的複体
$\mathcal{J}^\bullet$ へ取る。この矢印の始域と終域は，
$D(\mathcal{O}_X)$ において
$K \otimes_{\mathcal{O}_X}^\mathbf{L} M$ を表すことに注意する。
ここで任意の開集合 $W \subset X$ に対して，$A$-加群の複体の写像
$$
\text{Tot}(T^\bullet \otimes_A \mathcal{I}^\bullet(W))
\to
\text{Tot}(\mathcal{K}^\bullet(W) \otimes_A \mathcal{I}^\bullet(W))
\to
\mathcal{J}^\bullet(W)
$$
を得る。その合成は $D(A)$ における写像
$$
R\Gamma(X, K) \otimes_A^\mathbf{L} R\Gamma(W, M)
\longrightarrow
R\Gamma(W, K \otimes_{\mathcal{O}_X}^\mathbf{L} M)
$$
を表す。明らかに，これらの写像は制限写像と整合する。
そこで，次の $A$-加群の複体の可換図式を考える。
$$
\xymatrix{
0 \ar[d] & 0 \ar[d] \\
\text{Tot}(T^\bullet \otimes_A \mathcal{I}^\bullet(X)) \ar[d] \ar[r] &
\mathcal{J}^\bullet(X) \ar[d] \\
\text{Tot}(T^\bullet \otimes_A
(\mathcal{I}^\bullet(U) \oplus \mathcal{I}^\bullet(V)) \ar[d] \ar[r] &
\mathcal{J}^\bullet(U) \oplus \mathcal{J}^\bullet(V) \ar[d] \\
\text{Tot}(T^\bullet \otimes_A \mathcal{I}^\bullet(U \cap V)) \ar[r] \ar[d] &
\mathcal{J}^\bullet(U \cap V) \ar[d] \\
0 & 0
}
$$
補題 \ref{cohomology-lemma-mayer-vietoris} の証明により，各列は
$A$-加群の複体の完全列である（ここではさらに，$T^\bullet$ の各項が
平坦 $A$-加群なので，$\text{Tot}(T^\bullet \otimes_A -)$ が
$A$-加群の複体の短完全列を短完全列へ写すことを用いる）。
補題 \ref{cohomology-lemma-unbounded-mayer-vietoris} の完全三角は，これらの
複体の短完全列に付随する完全三角である。したがって求める結果は，
《導来圏》第 \ref{derived-section-canonical-delta-functor} 節で論じた
「付随する完全三角を取ること」の関手性から従う。
\end{proof}












\section{閉部分集合に台をもつコホモロジー II}
\label{cohomology-section-cohomology-support-bis}

\noindent
第 \ref{cohomology-section-cohomology-support} 節で始めた議論を続ける。

\medskip\noindent
$(X, \mathcal{O}_X)$ を環付き空間とし，$Z \subset X$ を閉部分集合とする。
この状況では，$\mathcal{F} \mapsto \Gamma_Z(X, \mathcal{F})$ で与えられる関手
$\textit{Mod}(\mathcal{O}_X) \to \textit{Mod}(\mathcal{O}_X(X))$
を考えられる。《加群》定義 \ref{modules-definition-support} および
《加群》補題 \ref{modules-lemma-support-section-closed} を参照せよ。
K-入射分解を用いると（第 \ref{cohomology-section-unbounded} 節を参照），右導来関手
$$
R\Gamma_Z(X, - ) : D(\mathcal{O}_X) \to D(\mathcal{O}_X(X))
$$
を得る。$D(\mathcal{O}_X)$ の対象 $K$ に対し，
$H^q_Z(X, K) = H^q(R\Gamma_Z(X, K))$ を $Z$ に台をもつ
コホモロジー加群と呼ぶ。後で補題
\ref{cohomology-lemma-sections-support-abelian-unbounded} において，これが第
\ref{cohomology-section-cohomology-support} 節の構成と一致することを見る。

\medskip\noindent
$\mathcal{O}_X$-加群 $\mathcal{F}$ に対し，規則
$$
\mathcal{H}_Z(\mathcal{F})(U) =
\{s \in \mathcal{F}(U) \mid \text{Supp}(s) \subset U \cap Z\} =
\Gamma_{Z \cap U}(U, \mathcal{F}|_U)
$$
で定義される {\it $Z$ に台をもつ切断の部分層}
$\mathcal{H}_Z(\mathcal{F})$ を考えられる。《加群》注意
\ref{modules-remark-sections-support-in-closed-modules} で論じたように，
$\mathcal{H}_Z(\mathcal{F})$ を $Z$ 上の
$\mathcal{O}_X|_Z$-加群とみなせるので，関手
$$
\textit{Mod}(\mathcal{O}_X) \longrightarrow \textit{Mod}(\mathcal{O}_X|_Z),
\quad
\mathcal{F} \longmapsto \mathcal{H}_Z(\mathcal{F})
\text{ viewed as an }\mathcal{O}_X|_Z\text{-module on }Z
$$
を得る。この関手は左完全だが，一般には完全でない。上とまったく同様に
右導来関手
$$
R\mathcal{H}_Z : D(\mathcal{O}_X) \longrightarrow D(\mathcal{O}_X|_Z)
$$
を得る。$\mathcal{H}^q_Z(K) = H^q(R\mathcal{H}_Z(K))$ と置く。
すると任意の $\mathcal{O}_X$-加群の層 $\mathcal{F}$ に対して
$\mathcal{H}^0_Z(\mathcal{F}) = \mathcal{H}_Z(\mathcal{F})$ である。

\begin{lemma}
\label{cohomology-lemma-cohomology-with-support-sheaf-on-support}
$(X, \mathcal{O}_X)$ を環付き空間とし，$i : Z \to X$ を
閉部分集合の包含とする。
\begin{enumerate}
\item $R\mathcal{H}_Z : D(\mathcal{O}_X) \to D(\mathcal{O}_X|_Z)$
は $i_* : D(\mathcal{O}_X|_Z) \to D(\mathcal{O}_X)$ の右随伴である。
\item $D(\mathcal{O}_X|_Z)$ の $K$ に対して
$R\mathcal{H}_Z(i_*K) = K$ である。
\item $\mathcal{G}$ を $Z$ 上の $\mathcal{O}_X|_Z$-加群の層とする。
このとき $p > 0$ に対して
$\mathcal{H}^p_Z(i_*\mathcal{G}) = 0$ である。
\end{enumerate}
\end{lemma}

\begin{proof}
$i_*$ は完全なので $i_* = Ri_* = Li_*$ である。したがって (1) は，
《加群》補題 \ref{modules-lemma-adjoint-section-with-support} および
《導来圏》補題 \ref{derived-lemma-derived-adjoint-functors} から従う。
$K$ を (2) のような対象とする。$K$ を
$\mathcal{O}_X|_Z$-加群の K-入射的複体
$\mathcal{I}^\bullet$ で表せる。補題
\ref{cohomology-lemma-K-injective-flat} により，$i_*K$ を表す複体
$i_*\mathcal{I}^\bullet$ は $\mathcal{O}_X$-加群の K-入射的複体である。
したがって $R\mathcal{H}_Z(i_*K)$ は
$\mathcal{H}_Z(i_*\mathcal{I}^\bullet) = \mathcal{I}^\bullet$
で計算され，(2) が証明される。(3) は (2) の特別な場合である。
\end{proof}

\noindent
$(X, \mathcal{O}_X)$ を環付き空間とし，$Z \subset X$ を閉部分集合とする。
台が $Z$ に含まれる $\mathcal{O}_X$-加群の圏は，すべての
$\mathcal{O}_X$-加群の圏の Serre 部分圏である。《ホモロジー》定義
\ref{homology-definition-serre-subcategory} および《加群》補題
\ref{modules-lemma-support-section-closed} を参照せよ。
コホモロジー層の台が $Z$ に含まれる複体からなる
$D(\mathcal{O}_X)$ の厳密充満飽和三角部分圏を
$D_Z(\mathcal{O}_X)$ と書く。《導来圏》第
\ref{derived-section-triangulated-sub} 節を参照せよ。

\begin{lemma}
\label{cohomology-lemma-complexes-with-support-on-closed}
$(X, \mathcal{O}_X)$ を環付き空間とし，$i : Z \to X$ を
閉部分集合の包含とする。
\begin{enumerate}
\item $D(\mathcal{O}_X|_Z)$ の $K$ に対して
$i_*K$ は $D_Z(\mathcal{O}_X)$ に属する。
\item 関手 $i_* : D(\mathcal{O}_X|_Z) \to D_Z(\mathcal{O}_X)$ は同値であり，
その擬逆は
$i^{-1}|_{D_Z(\mathcal{O}_X)} = R\mathcal{H}_Z|_{D_Z(\mathcal{O}_X)}$.
\item 関手
$i_* \circ R\mathcal{H}_Z : D(\mathcal{O}_X) \to D_Z(\mathcal{O}_X)$
は包含関手 $D_Z(\mathcal{O}_X) \to D(\mathcal{O}_X)$ の右随伴である。
\end{enumerate}
\end{lemma}

\begin{proof}
(1) は定義から直ちに従う。(3) は (2) と補題
\ref{cohomology-lemma-cohomology-with-support-sheaf-on-support} の形式的帰結である。
以下では (2) を証明する。

\medskip\noindent
$i$ を環付き空間の射
$i : (Z, \mathcal{O}_X|_Z) \to (X, \mathcal{O}_X)$
と考える。$i^*$ と $i_*$ が随伴関手の対であることを想起せよ。
$i$ は閉埋め込みなので $i_*$ は完全である。また
$i^{-1}\mathcal{O}_X = \mathcal{O}_X|_Z$ は
$(Z, \mathcal{O}_X|_Z)$ の構造層なので，$i^* = i^{-1}$ は完全であり，
$i^*i_* = i^{-1}i_*$ は恒等関手と同型である。《加群》補題
\ref{modules-lemma-exactness-pushforward-pullback} および
\ref{modules-lemma-i-star-exact} を参照せよ。したがって
$i_* : D(\mathcal{O}_X|_Z) \to D_Z(\mathcal{O}_X)$
は充満忠実であり，$i^{-1}$ は左逆を与える。一方，$K$ を
$D_Z(\mathcal{O}_X)$ の対象とし，随伴写像 $K \to i_*i^{-1}K$ を考える。
$i_*$ と $i^{-1}$ の完全性により，これはコホモロジー層上に随伴写像
$H^n(K) \to i_*i^{-1}H^n(K)$ を誘導する。これらのコホモロジー層の台は
$Z$ に含まれるので，この随伴写像は同型である。したがって
$i_* : D(\mathcal{O}_X|_Z) \to D_Z(\mathcal{O}_X)$ は同値である。

\medskip\noindent
証明を終えるには，$K$ が $D_Z(\mathcal{O}_X)$ の対象ならば
$R\mathcal{H}_Z(K) = i^{-1}K$ であることを示せば十分である。
今証明したように $K = i_*i^{-1}K$ なので，補題
\ref{cohomology-lemma-cohomology-with-support-sheaf-on-support} から望む結論を得る。
\end{proof}

\begin{lemma}
\label{cohomology-lemma-sections-with-support-K-injective}
$(X, \mathcal{O}_X)$ を環付き空間とし，$i : Z \to X$ を
閉部分集合の包含とする。$\mathcal{I}^\bullet$ が
$\mathcal{O}_X$-加群の K-入射的複体ならば，
$\mathcal{H}_Z(\mathcal{I}^\bullet)$ は
$\mathcal{O}_X|_Z$-加群の K-入射的複体である。
\end{lemma}

\begin{proof}
$i_* : \textit{Mod}(\mathcal{O}_X|_Z) \to \textit{Mod}(\mathcal{O}_X)$
は完全で，$\mathcal{H}_Z$ の左随伴である
（《加群》補題 \ref{modules-lemma-adjoint-section-with-support}）。
したがって《導来圏》補題
\ref{derived-lemma-adjoint-preserve-K-injectives} から従う。
\end{proof}

\begin{lemma}
\label{cohomology-lemma-local-to-global-sections-with-support}
$(X, \mathcal{O}_X)$ を環付き空間とし，$i : Z \to X$ を
閉部分集合の包含とする。このとき
$R\Gamma(Z, - ) \circ R\mathcal{H}_Z = R\Gamma_Z(X, - )$
は $D(\mathcal{O}_X) \to D(\mathcal{O}_X(X))$ の関手として成り立つ。
\end{lemma}

\begin{proof}
これは K-入射分解を用いた右導来関手の構成，補題
\ref{cohomology-lemma-sections-with-support-K-injective}，および等式
$\Gamma_Z(X, -) = \Gamma(Z, -) \circ \mathcal{H}_Z$ から従う。
\end{proof}

\begin{lemma}
\label{cohomology-lemma-triangle-sections-with-support}
$(X, \mathcal{O}_X)$ を環付き空間とし，$i : Z \to X$ を
閉部分集合の包含とする。$U = X \setminus Z$ と置く。
$D(\mathcal{O}_X(X))$ に完全三角
$$
R\Gamma_Z(X, K) \to R\Gamma(X, K) \to R\Gamma(U, K) \to
R\Gamma_Z(X, K)[1]
$$
があり，これは $D(\mathcal{O}_X)$ の $K$ に関して関手的である。
\end{lemma}

\begin{proof}
各項が入射的 $\mathcal{O}_X$-加群である，$K$ を表す
K-入射的複体 $\mathcal{I}^\bullet$ を選ぶ。第
\ref{cohomology-section-unbounded} 節を参照せよ。$\mathcal{I}^\bullet|_U$ は
$\mathcal{O}_U$-加群の K-入射的複体であることを想起せよ
（補題 \ref{cohomology-lemma-restrict-K-injective-to-open}）。
したがって完全三角の各導来関手は，$\mathcal{I}^\bullet$ に
元の関手を施すことで得られる。ゆえに，入射的
$\mathcal{O}_X$-加群 $\mathcal{I}$ に対して短完全列
$$
0 \to \Gamma_Z(X, \mathcal{I}) \to \Gamma(X, \mathcal{I})
\to \Gamma(U, \mathcal{I}) \to 0
$$
が存在することを証明すれば十分である。これは補題
\ref{cohomology-lemma-injective-restriction-surjective} と定義から従う。
\end{proof}

\begin{lemma}
\label{cohomology-lemma-triangle-sections-with-support-sheaves}
$(X, \mathcal{O}_X)$ を環付き空間とし，$i : Z \to X$ を
閉部分集合の包含とする。補集合の包含を
$j : U = X \setminus Z \to X$ と書く。$D(\mathcal{O}_X)$ に完全三角
$$
i_*R\mathcal{H}_Z(K) \to K \to Rj_*(K|_U) \to
i_*R\mathcal{H}_Z(K)[1]
$$
があり，これは $D(\mathcal{O}_X)$ の $K$ に関して関手的である。
\end{lemma}

\begin{proof}
各項が入射的 $\mathcal{O}_X$-加群である，$K$ を表す
K-入射的複体 $\mathcal{I}^\bullet$ を選ぶ。第
\ref{cohomology-section-unbounded} 節を参照せよ。$\mathcal{I}^\bullet|_U$ は
$\mathcal{O}_U$-加群の K-入射的複体であることを想起せよ
（補題 \ref{cohomology-lemma-restrict-K-injective-to-open}）。
したがって完全三角の各導来関手は，$\mathcal{I}^\bullet$ に
元の関手を施すことで得られる。ゆえに，入射的
$\mathcal{O}_X$-加群 $\mathcal{I}$ に対して短完全列
$$
0 \to i_*\mathcal{H}_Z(\mathcal{I}) \to \mathcal{I}
\to j_*(\mathcal{I}|_U) \to 0
$$
が存在することを証明すれば十分である。これは補題
\ref{cohomology-lemma-injective-restriction-surjective} と定義から従う。
\end{proof}

\begin{lemma}
\label{cohomology-lemma-sections-support-in-closed-disjoint-open}
$(X, \mathcal{O}_X)$ を環付き空間とし，$Z \subset X$ を閉部分集合とする。
$j : U \to X$ を $U \cap Z = \emptyset$ を満たす開部分集合の包含とする。
このとき $D(\mathcal{O}_U)$ のすべての $K$ に対して
$R\mathcal{H}_Z(Rj_*K) = 0$ である。
\end{lemma}

\begin{proof}
$K$ を表す $\mathcal{O}_U$-加群の K-入射的複体
$\mathcal{I}^\bullet$ を選ぶ。このとき $j_*\mathcal{I}^\bullet$ は
$Rj_*K$ を表す。補題 \ref{cohomology-lemma-K-injective-flat} により，
$j_*\mathcal{I}^\bullet$ は $\mathcal{O}_X$-加群の K-入射的複体である。
したがって $\mathcal{H}_Z(j_*\mathcal{I}^\bullet)$ は
$R\mathcal{H}_Z(Rj_*K)$ を表す。ゆえに，$U$ 上の任意のアーベル層
$\mathcal{G}$ に対して $\mathcal{H}_Z(j_*\mathcal{G}) = 0$ を
示せば十分である。すなわち，ある開集合 $W$ 上の
$j_*\mathcal{G}$ の切断 $s$ で台が $W \cap Z$ に含まれるものが
零であることを示せばよい。台の条件は
$s|_{W \setminus W \cap Z} = 0$ を意味する。また
$j_*\mathcal{G}(W) =
\mathcal{G}(U \cap W) = j_*\mathcal{G}(W \setminus W \cap Z)$
なので，望むとおり $s$ は零である。
\end{proof}

\begin{lemma}
\label{cohomology-lemma-sections-support-abelian-unbounded}
$(X, \mathcal{O}_X)$ を環付き空間とし，$Z \subset X$ を閉部分集合とする。
$K$ を $D(\mathcal{O}_X)$ の対象とし，
$D(\underline{\mathbf{Z}}_X)$ におけるその像を $K_{ab}$ と書く。
\begin{enumerate}
\item 標準写像 $R\Gamma_Z(X, K) \to R\Gamma_Z(X, K_{ab})$ があり，
これは $D(\textit{Ab})$ における同型である。
\item 標準写像
$R\mathcal{H}_Z(K) \to R\mathcal{H}_Z(K_{ab})$
があり，これは $D(\underline{\mathbf{Z}}_Z)$ における同型である。
\end{enumerate}
\end{lemma}

\begin{proof}
(1) を証明する。写像を次のように構成する。$K$ を表す
$\mathcal{O}_X$-加群の K-入射的複体 $\mathcal{I}^\bullet$ を選ぶ。
また，$\mathcal{J}^\bullet$ がアーベル群の K-入射的複体であるような
擬同型 $\mathcal{I}^\bullet \to \mathcal{J}^\bullet$ を選ぶ。
このとき (1) の写像は
$$
\Gamma_Z(X, \mathcal{I}^\bullet) \to \Gamma_Z(X, \mathcal{J}^\bullet)
$$
で与えられる。これは $\Gamma_Z$ がアーベル層上の関手であることから
定まる。容易な確認により，得られた写像と補題
\ref{cohomology-lemma-modules-abelian-unbounded} の標準写像は，補題
\ref{cohomology-lemma-triangle-sections-with-support} の完全三角の写像
$$
\xymatrix{
R\Gamma_Z(X, K) \ar[r] \ar[d] &
R\Gamma(X, K) \ar[r] \ar[d] &
R\Gamma(U, K) \ar[d] \\
R\Gamma_Z(X, K_{ab}) \ar[r] &
R\Gamma(X, K_{ab}) \ar[r] &
R\Gamma(U, K_{ab})
}
$$
をなす。引用した補題により三つの矢印のうち二つは同型なので，
《導来圏》補題
\ref{derived-lemma-third-isomorphism-triangle}
から結論を得る。

\medskip\noindent
(2) の証明は省略する。ヒント：完全三角には補題
\ref{cohomology-lemma-triangle-sections-with-support-sheaves} を用い，
同じ議論を適用せよ。
\end{proof}

\begin{remark}
\label{cohomology-remark-support-cup-product}
$(X, \mathcal{O}_X)$ を環付き空間とし，$i : Z \to X$ を
閉部分集合の包含とする。$D(\mathcal{O}_X)$ の $K$ と $M$ に対し，
標準写像
$$
K|_Z \otimes_{\mathcal{O}_X|_Z}^\mathbf{L} R\mathcal{H}_Z(M)
\longrightarrow
R\mathcal{H}_Z(K \otimes_{\mathcal{O}_X}^\mathbf{L} M)
$$
が $D(\mathcal{O}_X|_Z)$ に存在する。ここで $K|_Z = i^{-1}K$ は，
$K$ の $Z$ への制限を $D(\mathcal{O}_X|_Z)$ の対象とみなしたものである。
補題 \ref{cohomology-lemma-cohomology-with-support-sheaf-on-support} による
$i_*$ と $R\mathcal{H}_Z$ の随伴性から，この写像を構成するには
標準写像
$$
i_*\left(K|_Z \otimes_{\mathcal{O}_X|_Z}^\mathbf{L} R\mathcal{H}_Z(M)\right)
\longrightarrow
K \otimes_{\mathcal{O}_X}^\mathbf{L} M
$$
を作れば十分である。この写像を構成するため，$M$ を表す
$\mathcal{O}_X$-加群の K-入射的複体 $\mathcal{I}^\bullet$ と，
$K$ を表す $\mathcal{O}_X$-加群の K-平坦複体
$\mathcal{K}^\bullet$ を選ぶ。$\mathcal{K}^\bullet|_Z$ は
$K|_Z$ を表す $\mathcal{O}_X|_Z$-加群の K-平坦複体である。
補題 \ref{cohomology-lemma-pullback-K-flat} を参照せよ。したがって，
$\mathcal{O}_X$-加群の複体の写像
$$
i_*\text{Tot}\left(
\mathcal{K}^\bullet|_Z \otimes_{\mathcal{O}_X|_Z}
\mathcal{H}_Z(\mathcal{I}^\bullet)\right)
\longrightarrow
\text{Tot}(\mathcal{K}^\bullet \otimes_{\mathcal{O}_X} \mathcal{I}^\bullet)
$$
を作る必要がある。そのためには写像
$$
i_*(\mathcal{K}^a|_Z \otimes_{\mathcal{O}_X|_Z}
\mathcal{H}_Z(\mathcal{I}^b))
\longrightarrow
\mathcal{K}^a \otimes_{\mathcal{O}_X} \mathcal{I}^b
$$
を作れば十分である。例えば茎を調べると，この式の左辺は
$\mathcal{K}^a \otimes_{\mathcal{O}_X} i_*\mathcal{H}_Z(\mathcal{I}^b)$
に等しいことが分かる。そこで包含
$\mathcal{H}_Z(\mathcal{I}^b) \to \mathcal{I}^b$ を用いれば
求める写像を得る。
\end{remark}

\begin{remark}
\label{cohomology-remark-support-cup-product-global}
注意 \ref{cohomology-remark-support-cup-product} の記法の下で，標準的なカップ積
\begin{align*}
H^a(X, K) \times H^b_Z(X, M)
& =
H^a(X, K) \times H^b(Z, R\mathcal{H}_Z(M)) \\
& \to
H^a(Z, K|_Z) \times H^b(Z, R\mathcal{H}_Z(M)) \\
& \to
H^{a + b}(Z, K|_Z \otimes_{\mathcal{O}_X|_Z}^\mathbf{L} R\mathcal{H}_Z(M)) \\
& \to
H^{a + b}(Z, R\mathcal{H}_Z(K \otimes_{\mathcal{O}_X}^\mathbf{L} M)) \\
& =
H^{a + b}_Z(X, K \otimes_{\mathcal{O}_X}^\mathbf{L} M)
\end{align*}
を得る。ここで等号は補題
\ref{cohomology-lemma-local-to-global-sections-with-support} により与えられる。
第一の矢印は $Z$ への制限，第二の矢印はカップ積
（第 \ref{cohomology-section-cup-product} 節），第三の矢印は注意
\ref{cohomology-remark-support-cup-product} の写像である。
\end{remark}

\begin{lemma}
\label{cohomology-lemma-support-cup-product}
注意 \ref{cohomology-remark-support-cup-product} の記法の下で，図式
$$
\xymatrix{
H^i(X, K) \times H^j_Z(X, M) \ar[r] \ar[d] &
H^{i + j}_Z(X, K \otimes_{\mathcal{O}_X}^\mathbf{L} M) \ar[d] \\
H^i(X, K) \times H^j(X, M) \ar[r] &
H^{i + j}(X, K \otimes_{\mathcal{O}_X}^\mathbf{L} M)
}
$$
は可換である。ここで上の水平射は注意
\ref{cohomology-remark-support-cup-product-global} のカップ積である。
\end{lemma}

\begin{proof}
省略する。
\end{proof}

\begin{remark}
\label{cohomology-remark-support-functorial}
$f : (X', \mathcal{O}_{X'}) \to (X, \mathcal{O}_X)$ を環付き空間の射とする。
$Z \subset X$ を閉部分集合とし，$Z' = f^{-1}(Z)$ と置く。
誘導される環付き空間の射を
$f|_{Z'} : (Z', \mathcal{O}_{X'}|_{Z'}) \to (Z, \mathcal{O}_X|Z)$ と書く。
$D(\mathcal{O}_X)$ の任意の $K$ に対して標準写像
$$
L(f|_{Z'})^*R\mathcal{H}_Z(K) \longrightarrow R\mathcal{H}_{Z'}(Lf^*K)
$$
が $D(\mathcal{O}_{X'}|_{Z'})$ に存在する。包含写像を
$i : Z \to X$ および $i' : Z' \to X'$ と書く。補題
\ref{cohomology-lemma-complexes-with-support-on-closed} の (2) を $i'$ に適用すると，
この写像を与えることは，写像
$$
i'_* L(f|_{Z'})^* R\mathcal{H}_Z(K)
\longrightarrow
i'_*R\mathcal{H}_{Z'}(Lf^*K)
$$
を $D_{Z'}(\mathcal{O}_{X'})$ で与えることと同じである。
注意 \ref{cohomology-remark-base-change} の関手の写像
$Lf^* \circ i_* \to i'_* \circ L(f|_{Z'})^*$ はこの場合には同型である
（$i_*$ と $i'_*$ が完全であり，
$\mathcal{O}_{Z, z} = \mathcal{O}_{X, z}$ であること，および
$Z'$ についての同様の事実を用いて茎上で確認すれば従う）。
したがって次の図式の上の水平射を構成すれば十分である。
$$
\xymatrix{
Lf^* i_* R\mathcal{H}_Z(K) \ar[rr] \ar[rd] & &
i'_* R\mathcal{H}_{Z'}(Lf^*K) \ar[ld] \\
& Lf^*K
}
$$
複体 $Lf^* i_* R\mathcal{H}_Z(K)$ の台は $Z'$ に含まれる。
南東向きの矢印は随伴写像
$i_*R\mathcal{H}_Z(K) \to K$ から得られる
（補題 \ref{cohomology-lemma-cohomology-with-support-sheaf-on-support}）。
随伴写像 $i'_* R\mathcal{H}_{Z'}(Lf^*K) \to Lf^*K$ は補題
\ref{cohomology-lemma-complexes-with-support-on-closed} の (3) により普遍的なので，
南東向きの矢印は南西向きの矢印を通じて一意に分解し，求める矢印を得る。
\end{remark}

\begin{lemma}
\label{cohomology-lemma-support-functorial}
注意 \ref{cohomology-remark-support-functorial} の記法と仮定の下で，図式
$$
\xymatrix{
H^p_Z(X, K) \ar[r] \ar[d] & H^p_{Z'}(X', Lf^*K) \ar[d] \\
H^p(X, K) \ar[r] & H^p(X', Lf^*K)
}
$$
は可換である。ここで上の水平射は，同一視
$H^p_Z(X, K) = H^p(Z, R\mathcal{H}_Z(K))$ および
$H^p_{Z'}(X', Lf^*K) = H^p(Z', R\mathcal{H}_{Z'}(K'))$,
，引き戻し写像
$H^p(Z, R\mathcal{H}_Z(K)) \to H^p(Z', L(f|_{Z'})^*R\mathcal{H}_Z(K))$,
および注意 \ref{cohomology-remark-support-functorial} で構成した写像から得られる。
\end{lemma}

\begin{proof}
省略する。ヒント：
$H^p(Z, R\mathcal{H}_Z(K)) = H^p(X, i_*R\mathcal{H}_Z(K))$ および
$R\mathcal{H}_{Z'}(Lf^*K)$ に対する同様の等式を用いれば，
引き戻し写像の関手性と，注意
\ref{cohomology-remark-support-functorial} の写像を定義するために用いた可換図式から従う。
\end{proof}






\section{逆系とコホモロジー I}
\label{cohomology-section-inverse-systems}

\noindent
$A$ を環とし，$I \subset A$ をイデアルとする。
$A/I^n$-加群の層の逆系に関するいくつかの結果を証明する。

\begin{lemma}
\label{cohomology-lemma-ML-general}
$I$ を環 $A$ のイデアル，$X$ を位相空間とする。
$$
\ldots \to \mathcal{F}_3 \to \mathcal{F}_2 \to \mathcal{F}_1
$$
を，$\mathcal{F}_n = \mathcal{F}_{n + 1}/I^n\mathcal{F}_{n + 1}$
を満たす $X$ 上の $A$-加群の層の逆系とする。$p \geq 0$ とする。
$$
\bigoplus\nolimits_{n \geq 0} H^{p + 1}(X, I^n\mathcal{F}_{n + 1})
$$
が次数付き $\bigoplus_{n \geq 0} I^n/I^{n + 1}$-加群として
昇鎖条件を満たすと仮定する。このとき逆系
$M_n = H^p(X, \mathcal{F}_n)$ は Mittag--Leffler 条件を満たす。
\footnote{実際，ある $c \geq 0$ が存在して，すべての $n \geq c$ に対し
$\Im(M_n \to M_{n - c})$ が安定像となる。}
\end{lemma}

\begin{proof}
$N_n = H^{p + 1}(X, I^n\mathcal{F}_{n + 1})$ と置き，
$\delta_n : M_n \to N_n$ を短完全列
$0 \to I^n\mathcal{F}_{n + 1} \to \mathcal{F}_{n + 1} \to \mathcal{F}_n \to 0$
から得られるコホモロジー上の境界写像とする。このとき
$\bigoplus \Im(\delta_n) \subset \bigoplus N_n$ は次数付き部分加群である。
実際，$s \in M_n$ かつ $f \in I^m$ ならば，可換図式
$$
\xymatrix{
0 \ar[r] &
I^n\mathcal{F}_{n + 1} \ar[d]_f \ar[r] &
\mathcal{F}_{n + 1} \ar[d]_f \ar[r] &
\mathcal{F}_n \ar[d]_f \ar[r] & 0 \\
0 \ar[r] &
I^{n + m}\mathcal{F}_{n + m + 1} \ar[r] &
\mathcal{F}_{n + m + 1} \ar[r] &
\mathcal{F}_{n + m} \ar[r] & 0
}
$$
を得る。中央の垂直写像は，$\mathcal{F}_{n + 1}$ の局所切断を
$\mathcal{F}_{n + m + 1}$ の切断へ持ち上げてから $f$ を掛けることで
与えられる。他の垂直射についても同様である。したがって
$\delta_{n + m}(fs) = f \delta_n(s)$ である。仮定により，
$\delta_{n_j}(s_j)$ が $\bigoplus \Im(\delta_n)$ を次数付き加群として
生成するような $s_j \in M_{n_j}$, $j = 1, \ldots, N$ を取れる。
$n > c = \max(n_j)$ とし，$s \in M_n$ とする。このとき
$\delta_n(s) = \sum f_j \delta_{n_j}(s_j)$ を満たす
$f_j \in I^{n - n_j}$ が存在する。したがって
$\delta(s - \sum f_j s_j) = 0$ であり，すなわち
$M_n$ で $s - \sum f_js_j$ へ写る $s' \in M_{n + 1}$ が存在する。
ゆえに
$$
\Im(M_{n + 1} \to M_{n - c}) = \Im(M_n \to M_{n - c})
$$
である。実際，元 $f_js_j$ は $M_{n - c}$ で零に写る。
これで補題が証明された。
\end{proof}

\begin{lemma}
\label{cohomology-lemma-ML-general-better}
$I$ を環 $A$ のイデアル，$X$ を位相空間とする。
$$
\ldots \to \mathcal{F}_3 \to \mathcal{F}_2 \to \mathcal{F}_1
$$
を，$\mathcal{F}_n = \mathcal{F}_{n + 1}/I^n\mathcal{F}_{n + 1}$
を満たす $X$ 上の $A$-加群の逆系とする。$p \geq 0$ とする。
$n$ に対して
$$
N_n =
\bigcap\nolimits_{m \geq n}
\Im\left(
H^{p + 1}(X, I^n\mathcal{F}_{m + 1}) \to H^{p + 1}(X, I^n\mathcal{F}_{n + 1})
\right)
$$
と定める。$\bigoplus N_n$ が次数付き
$\bigoplus_{n \geq 0} I^n/I^{n + 1}$-加群として昇鎖条件を満たすならば，
逆系 $M_n = H^p(X, \mathcal{F}_n)$ は Mittag--Leffler 条件を満たす。
\footnote{実際，ある $c \geq 0$ が存在して，すべての $n \geq c$ に対し
$\Im(M_n \to M_{n - c})$ が安定像となる。}
\end{lemma}

\begin{proof}
証明は補題 \ref{cohomology-lemma-ML-general} の証明とまったく同じである。
実際，$\bigoplus N_n$ が
$\bigoplus H^{p + 1}(X, I^n\mathcal{F}_{n + 1})$ の次数付き
$\bigoplus_{n \geq 0} I^n/I^{n + 1}$-部分加群であること，および境界写像
$\delta_n : M_n \to H^{p + 1}(X, I^n\mathcal{F}_{n + 1})$
の像が $N_n$ に含まれることを示せば，上の議論から結果が従う。

\medskip\noindent
$\xi \in N_n$ および $f \in I^k$ と仮定する。
$m \gg n + k$ を選び，$\xi$ の持ち上げ
$\xi' \in H^{p + 1}(X, I^n\mathcal{F}_{m + 1})$ を選ぶ。
図式
$$
\xymatrix{
0 \ar[r] &
I^n\mathcal{F}_{m + 1} \ar[d]_f \ar[r] &
\mathcal{F}_{m + 1} \ar[d]_f \ar[r] &
\mathcal{F}_n \ar[d]_f \ar[r] & 0 \\
0 \ar[r] &
I^{n + k}\mathcal{F}_{m + 1} \ar[r] &
\mathcal{F}_{m + 1} \ar[r] &
\mathcal{F}_{n + k} \ar[r] & 0
}
$$
を補題 \ref{cohomology-lemma-ML-general} の証明と同様に構成する。
コホモロジー上に誘導される写像を得て，
$f \xi' \in H^{p + 1}(X, I^{n + k}\mathcal{F}_{m + 1})$ が
$f \xi$ に写ることが分かる。これはすべての $m \gg n + k$ に対して
成り立つので，望むとおり $f\xi$ は $N_{n + k}$ に属する。

\medskip\noindent
境界写像 $\delta_n$ の像が $N_n$ に含まれることを見るため，図式
$$
\xymatrix{
0 \ar[r] &
I^n\mathcal{F}_{m + 1} \ar[d] \ar[r] &
\mathcal{F}_{m + 1} \ar[d] \ar[r] &
\mathcal{F}_n \ar[d] \ar[r] & 0 \\
0 \ar[r] &
I^n\mathcal{F}_{n + 1} \ar[r] &
\mathcal{F}_{n + 1} \ar[r] &
\mathcal{F}_n \ar[r] & 0
}
$$
を $m \geq n$ に対して考える。コホモロジー上に誘導される写像を
調べれば結論を得る。
\end{proof}

\begin{lemma}
\label{cohomology-lemma-topology-I-adic-general}
$I$ を環 $A$ のイデアル，$X$ を位相空間とする。
$$
\ldots \to \mathcal{F}_3 \to \mathcal{F}_2 \to \mathcal{F}_1
$$
を，$\mathcal{F}_n = \mathcal{F}_{n + 1}/I^n\mathcal{F}_{n + 1}$
を満たす $X$ 上の $A$-加群の層の逆系とする。$p \geq 0$ とする。
$$
\bigoplus\nolimits_{n \geq 0} H^p(X, I^n\mathcal{F}_{n + 1})
$$
が次数付き $\bigoplus_{n \geq 0} I^n/I^{n + 1}$-加群として
昇鎖条件を満たすと仮定する。このとき
$M = \lim H^p(X, \mathcal{F}_n)$ 上の極限位相は $I$-進位相である。
\end{lemma}

\begin{proof}
$n \geq 1$ に対して
$F^n = \Ker(M \to H^p(X, \mathcal{F}_n))$ と置き，$F^0 = M$ とする。
$I F^n \subset F^{n + 1}$ であり，特に $I^n M \subset F^n$ である。
したがって $I$-進位相は極限位相より細かい。逆を示すため，任意の $n$ に対し
$F^m \subset I^nM$ を満たす $m \geq n$ が存在することを示す。
\footnote{実際，ある $c \geq 0$ が存在して，すべての $n$ に対し
$F^{n + c} \subset I^nM$ となる。}
単射
$$
F^n/F^{n + 1} \longrightarrow H^p(X, \mathcal{F}_{n + 1})
$$
があり，その像は
$H^p(X, I^n\mathcal{F}_{n + 1}) \to H^p(X, \mathcal{F}_{n + 1})$
の像に含まれる。この写像による $F^n/F^{n + 1}$ の逆像を
$$
E_n \subset H^p(X, I^n\mathcal{F}_{n + 1})
$$
と書く。このとき $\bigoplus E_n$ は
$\bigoplus H^p(X, I^n\mathcal{F}_{n + 1})$ の次数付き
$\bigoplus I^n/I^{n + 1}$-部分加群であり，
$\bigoplus E_n \to \bigoplus F^n/F^{n + 1}$ は次数付き加群の準同型である。
詳細は省略する。仮定により，$\bigoplus E_n$ は
$\bigoplus I^n/I^{n + 1}$ 上有限個の斉次な元で生成される。
$E_n \to F^n/F^{n + 1}$ は全射なので，
$\bigoplus F^n/F^{n + 1}$ についても同じことが成り立つ。
したがって，$r$，$c_1, \ldots, c_r \geq 0$，および
$\bigoplus F^n/F^{n + 1}$ における像がこれを生成する
$a_i \in F^{c_i}$ を取れる。$c = \max(c_i)$ と置く。

\medskip\noindent
$n \geq c$ に対して $I F^n = F^{n + 1}$ と主張する。この主張から，
望むように $F^{n + c} = I^nF^c \subset I^nM$ が従う。
主張を証明するため $a \in F^{n + 1}$ とする。
$F^{n + 1}/F^{n + 2}$ における $a$ の像は $a_i$ の線形結合である。
したがって，ある $f_i \in I^{n + 1 - c_i}$ に対して
$a  - \sum f_i a_i \in F^{n + 2}$ である。
$n \geq c_i$ なので
$I^{n + 1 - c_i} = I \cdot I^{n - c_i}$ であり，
$g_{i, j} \in I$ かつ $h_{i, j}a_i \in F^n$ となるように
$f_i = \sum g_{i, j} h_{i, j}$ と書ける。ゆえに
$F^{n + 1} = F^{n + 2} + IF^n$ である。
簡単な帰納法により，すべての $e > 0$ に対して
$F^{n + 1} = F^{n + e} + IF^n$ を得る。したがって $IF^n$ は
$F^{n + 1}$ で稠密である。$I$ の生成元
$k_1, \ldots, k_r$ を選び，連続写像
$$
u : (F^n)^{\oplus r} \longrightarrow F^{n + 1},\quad
(x_1, \ldots, x_r) \mapsto \sum k_i x_i
$$
を考える（極限位相に関して）。上の議論により，すべての $m \geq n$ に対し，
$u$ による $(F^m)^{\oplus r}$ の像は $F^{m + 1}$ で稠密である。
開写像補題（《代数続論》補題
\ref{more-algebra-lemma-open-mapping}）により，$u$ は開写像である。
したがって $u$ は全射である。ゆえに $n \geq c$ に対して
$IF^n = F^{n + 1}$ であり，証明は完了する。
\end{proof}

\begin{lemma}
\label{cohomology-lemma-topology-I-adic-general-better}
$I$ を環 $A$ のイデアル，$X$ を位相空間とする。
$$
\ldots \to \mathcal{F}_3 \to \mathcal{F}_2 \to \mathcal{F}_1
$$
を，$\mathcal{F}_n = \mathcal{F}_{n + 1}/I^n\mathcal{F}_{n + 1}$
を満たす $X$ 上の $A$-加群の層の逆系とする。$p \geq 0$ とする。
$n$ に対して
$$
N_n =
\bigcap\nolimits_{m \geq n}
\Im\left(
H^p(X, I^n\mathcal{F}_{m + 1}) \to H^p(X, I^n\mathcal{F}_{n + 1})
\right)
$$
と定める。$\bigoplus N_n$ が次数付き
$\bigoplus_{n \geq 0} I^n/I^{n + 1}$-加群として昇鎖条件を満たすならば，
$M = \lim H^p(X, \mathcal{F}_n)$ 上の極限位相は $I$-進位相である。
\end{lemma}

\begin{proof}
証明は補題 \ref{cohomology-lemma-topology-I-adic-general} の証明とまったく同じである。
実際，$\bigoplus N_n$ が
$\bigoplus H^{p + 1}(X, I^n\mathcal{F}_{n + 1})$ の次数付き
$\bigoplus_{n \geq 0} I^n/I^{n + 1}$-部分加群であること，および
$F^n/F^{n + 1} \subset H^p(X, \mathcal{F}_{n + 1})$ が
$N_n \to H^p(X, \mathcal{F}_{n + 1})$ の像に含まれることを示せば，
上の議論から結果が従う。加群構造に関する主張は補題
\ref{cohomology-lemma-ML-general-better} の証明で見た。

\medskip\noindent
$t \in F^n$ とする。$H^p(X, \mathcal{F}_{n + 1})$ における $t$ の像へ
写る元 $s \in H^p(X, I^n\mathcal{F}_{n + 1})$ を選ぶ。
$s \in N_n$ を示す必要がある。$F^n$ は写像
$M \to H^p(X, \mathcal{F}_n)$ の核なので，すべての $m \geq n$ に対し，
$t$ を $H^p(X, \mathcal{F}_n)$ で零に写る元
$t_m \in H^p(X, \mathcal{F}_{m + 1})$ へ写せる。短完全列
$0 \to I^n\mathcal{F}_{m + 1} \to \mathcal{F}_{m + 1} \to \mathcal{F}_n \to 0$
から得られるコホモロジー列
$$
H^{p - 1}(X, \mathcal{F}_n) \to
H^p(X, I^n\mathcal{F}_{m + 1}) \to
H^p(X, \mathcal{F}_{m + 1}) \to
H^p(X, \mathcal{F}_n)
$$
を考える。$t_m$ へ写る
$s_m \in H^p(X, I^n\mathcal{F}_{m + 1})$ を選べる。
上の列を $m = n$ の場合の列と比較すると，$s_m$ は
$H^{p - 1}(X, \mathcal{F}_n) \to H^p(X, I^n\mathcal{F}_{n + 1})$
の像に属する元を除いて $s$ へ写ることが分かる。しかしこの写像は
$H^p(X, I^n\mathcal{F}_{m + 1}) \to H^p(X, I^n\mathcal{F}_{n + 1})$
を経由するので，望むとおり $s$ はその像に属する。
\end{proof}





\section{逆系とコホモロジー II}
\label{cohomology-section-inverse-systems-bis}

\noindent
この節では，イデアルが単項である場合に
節 \ref{cohomology-section-inverse-systems} の議論を続ける。

\begin{lemma}
\label{cohomology-lemma-equivalent-f-good}
$(X, \mathcal{O}_X)$ を環付き空間とし，
$f \in \Gamma(X, \mathcal{O}_X)$ とする。
$$
\ldots \to \mathcal{F}_3 \to \mathcal{F}_2 \to \mathcal{F}_1
$$
を $\mathcal{O}_X$-加群の逆系とする。次の条件を考える。
\begin{enumerate}
\item すべての $n \geq 1$ に対し，写像
$f : \mathcal{F}_{n + 1} \to \mathcal{F}_{n + 1}$ は
$\mathcal{F}_{n + 1} \to \mathcal{F}_n$ を経由し，短完全列
$0 \to \mathcal{F}_n \to \mathcal{F}_{n + 1} \to \mathcal{F}_1 \to 0$,
\item すべての $n \geq 1$ に対し，写像
$f^n : \mathcal{F}_{n + 1} \to \mathcal{F}_{n + 1}$
$\mathcal{F}_{n + 1} \to \mathcal{F}_1$ を経由し，短完全列
$0 \to \mathcal{F}_1 \to \mathcal{F}_{n + 1} \to \mathcal{F}_n \to 0$
\item $\mathcal{F}_n = \mathcal{G}[f^n]$ を満たす
$f$-可除な $\mathcal{O}_X$-加群 $\mathcal{G}$ が存在する。
\item $\mathcal{F}_n = \mathcal{F}/f^n\mathcal{F}$ を満たす
$f$-ねじれなし $\mathcal{O}_X$-加群 $\mathcal{F}$ が存在する。
\end{enumerate}
このとき (4) $\Rightarrow$ (3) $\Leftrightarrow$ (2) $\Leftrightarrow$ (1)
が成り立つ。
\end{lemma}

\begin{proof}
条件 (1) と (2) の同値性の証明，および (3) が (1) を導くことの証明は省略する。
(1) のような $\mathcal{F}_n$ が与えられたとき，(3) を示すため
$\mathcal{G} = \colim \mathcal{F}_n$ とおく。ただし写像
$\mathcal{F}_1 \to \mathcal{F}_2 \to \mathcal{F}_3 \to \ldots$
は (1) によるものとする。短完全列 (1) により
$\mathcal{F}_{n + 1} \subset \mathcal{G}$ の像は
$\mathcal{F}_n \subset \mathcal{G}$ だから，写像
$f : \mathcal{G} \to \mathcal{G}$ は全射である。
したがって $\mathcal{G}$ は
$\mathcal{F}_n = \mathcal{G}[f^n]$ を満たす
$f$-可除な $\mathcal{O}_X$-加群である。

\medskip\noindent
(4) のような $\mathcal{F}$ が与えられたとする。写像
$\mathcal{F}/f^{n + 1}\mathcal{F} \to \mathcal{F}/f^n\mathcal{F}$
は常に全射であり，その核は $f^n$ 倍によって誘導される写像
$\mathcal{F}/f\mathcal{F} \to \mathcal{F}/f^{n + 1}\mathcal{F}$
の像である。(2) を確かめるには，
$f^n : \mathcal{F} \to \mathcal{F}/f^{n + 1}\mathcal{F}$
の核が $f\mathcal{F}$ であることを見れば十分である。このためには，
開集合 $U \subset X$ 上の $\mathcal{F}$ の切断 $s,t$ が
$f^ns = f^{n + 1}t$ を満たすとき $s = ft$ となることを示せばよい。
$\mathcal{F}$ は $f$-ねじれなしであり，
$f : \mathcal{F} \to \mathcal{F}$ は単射だから，これは明らかである。
\end{proof}

\begin{lemma}
\label{cohomology-lemma-topology-I-adic-f}
$X$, $f$, $(\mathcal{F}_n)$ が補題
\ref{cohomology-lemma-equivalent-f-good} の条件 (1) を満たすとする。
$p \geq 0$ とし，
$H^p = \lim H^p(X, \mathcal{F}_n)$.
とおく。このとき，すべての $c \geq 1$ に対し $f^cH^p$ は
$H^p \to H^p(X, \mathcal{F}_c)$ の核である。したがって $H^p$ 上の
極限位相は $f$-進位相である。
\end{lemma}

\begin{proof}
$c \geq 1$ とする。$f^c H^p$ が
$H^p(X, \mathcal{F}_c)$ で零に写ることは明らかである。
$\xi = (\xi_n) \in H^p$ が極限位相に関して小さいとすると
$\xi_c = 0$ であり，したがって $n \geq c$ に対して
$\xi_n$ は $H^p(X, \mathcal{F}_c)$ で零に写る。
短完全列の逆系
$$
0 \to \mathcal{F}_{n - c} \xrightarrow{f^c} \mathcal{F}_n \to
\mathcal{F}_c \to 0
$$
と，対応するコホモロジー長完全列の逆系
$$
H^{p - 1}(X, \mathcal{F}_c) \to
H^p(X, \mathcal{F}_{n - c}) \to
H^p(X, \mathcal{F}_n) \to
H^p(X, \mathcal{F}_c)
$$
を考える。項 $H^{p - 1}(X, \mathcal{F}_c)$ は $n$ に依存しないので，
$\xi_n$ を持ち上げる元の整合的な列
$\xi'_n \in H^p(X, \mathcal{F}_{n - c})$
を選べる。$\xi' = (\xi'_n)$ とおけば，
望むとおり $\xi = f^c \xi'$ となる。
\end{proof}

\begin{lemma}
\label{cohomology-lemma-limit-finite}
$A$ を単項イデアル $(f)$ に関して完備な Noether 環とし，
$X$ を位相空間とする。
$$
\ldots \to \mathcal{F}_3 \to \mathcal{F}_2 \to \mathcal{F}_1
$$
を $A$-加群の層の逆系とする。次を仮定する。
\begin{enumerate}
\item $\Gamma(X, \mathcal{F}_1)$ は有限 $A$-加群である。
\item $X$, $f$, $(\mathcal{F}_n)$ は補題
\ref{cohomology-lemma-equivalent-f-good} の条件 (1) を満たす。
\end{enumerate}
このとき
$$
M = \lim \Gamma(X, \mathcal{F}_n)
$$
は有限 $A$-加群であり，$f$ は $M$ 上の非零因子であり，
$M/fM$ は $\Gamma(X, \mathcal{F}_1)$ における $M$ の像である。
\end{lemma}

\begin{proof}
補題 \ref{cohomology-lemma-topology-I-adic-f} により
$M/fM \subset H^0(X, \mathcal{F}_1)$ である。(1) と $A$ の
Noether 性から，$M/fM$ は有限 $A$-加群である。
$f^nM$ は $H^0(X, \mathcal{F}_n)$ で零に写るので，
$\bigcap f^nM = 0$ であることに注意する。
Algebra, Lemma \ref{algebra-lemma-finite-over-complete-ring}
により，$M$ は $A$ 上有限であると結論する。
最後に，$s = (s_n) \in M = \lim H^0(X, \mathcal{F}_n)$ が
$fs = 0$ を満たすとする。補題 \ref{cohomology-lemma-equivalent-f-good}
の条件 (1) により，$s_{n + 1}$ は
$\mathcal{F}_{n + 1} \to \mathcal{F}_n$ の核に属する。
したがって $s_n = 0$ である。$n$ は任意だったので $s = 0$ となる。
ゆえに $f$ は $M$ 上の非零因子である。
\end{proof}

\begin{lemma}
\label{cohomology-lemma-ML}
$A$ を環，$f \in A$，$X$ を位相空間とする。
$$
\ldots \to \mathcal{F}_3 \to \mathcal{F}_2 \to \mathcal{F}_1
$$
を $A$-加群の層の逆系とし，$p \geq 0$ とする。次を仮定する。
\begin{enumerate}
\item $H^{p + 1}(X, \mathcal{F}_1)$ が有限長の $A$-加群であるか，
または $A$ が Noether 環で $H^{p + 1}(X, \mathcal{F}_1)$ が有限
$A$-加群である。
\item $X$, $f$, $(\mathcal{F}_n)$ は補題
\ref{cohomology-lemma-equivalent-f-good} の条件 (1) を満たす。
\end{enumerate}
このとき逆系 $M_n = H^p(X, \mathcal{F}_n)$ は
Mittag--Leffler 条件を満たす。
\end{lemma}

\begin{proof}
$I = (f)$ とおく。補題 \ref{cohomology-lemma-ML-general} の判定条件を用いる。
すべての $n \geq 0$ に対して
$f^n : \mathcal{F}_1 \to I^n\mathcal{F}_{n + 1}$
は同型であることに注意する。したがって
$$
\bigoplus\nolimits_{n \geq 1} H^{p + 1}(X, \mathcal{F}_1) \cdot f^{n + 1}
$$
が次数付き $S = \bigoplus_{n \geq 0} A/(f) \cdot f^n$-加群として
昇鎖条件を満たすことを示せば十分である。$A$ が Noether 環でなければ，
$H^{p + 1}(X, \mathcal{F}_1)$ は有限長なので結果が従う。
$A$ が Noether 環ならば $S$ は Noether 環であり，
$H^{p + 1}(X, \mathcal{F}_1)$ に仮定した有限性により上の加群は
$S$ 上有限だから，やはり結果が従う。いくつかの詳細は省略する。
\end{proof}

\begin{lemma}
\label{cohomology-lemma-ML-better}
$A$ を環，$f \in A$，$X$ を位相空間とする。
$$
\ldots \to \mathcal{F}_3 \to \mathcal{F}_2 \to \mathcal{F}_1
$$
を $A$-加群の層の逆系とし，$p \geq 0$ とする。次を仮定する。
\begin{enumerate}
\item ある $m \geq 1$ に対して
$H^{p + 1}(X, \mathcal{F}_m) \to H^{p + 1}(X, \mathcal{F}_1)$
の像が有限長の $A$-加群であるか，または $A$ が Noether 環で，
写像
$H^{p + 1}(X, \mathcal{F}_m) \to H^{p + 1}(X, \mathcal{F}_1)$
の像たちの共通部分が有限 $A$-加群である。
\item $X$, $f$, $(\mathcal{F}_n)$ は補題
\ref{cohomology-lemma-equivalent-f-good} の条件 (1) を満たす。
\end{enumerate}
このとき逆系 $M_n = H^p(X, \mathcal{F}_n)$ は
Mittag--Leffler 条件を満たす。
\end{lemma}

\begin{proof}
$I = (f)$ とおく。加群 $N_n$ を用いる補題
\ref{cohomology-lemma-ML-general-better} の判定条件を使う。
$m \geq n$ に対して
$I^n\mathcal{F}_{m + 1} = \mathcal{F}_{m + 1 - n}$ である。したがって
$$
N_n = \bigcap\nolimits_{m \geq 1} \Im\left(
H^{p + 1}(X, \mathcal{F}_m) \to H^{p + 1}(X, \mathcal{F}_1)
\right)
$$
は $n$ に依存せず，
$\bigoplus N_n = \bigoplus N_1 \cdot f^{n + 1}$ である。
ゆえに補題 \ref{cohomology-lemma-ML} の証明とまったく同様に結論を得る。
\end{proof}

\begin{remark}
\label{cohomology-remark-compare-ML}
$(X, \mathcal{O}_X)$ を環付き空間，
$f \in \Gamma(X, \mathcal{O}_X)$ とし，
$\mathcal{F}$ を $\mathcal{O}_X$-加群とする。
$\mathcal{F}$ が $f$-ねじれなしならば，すべての $p \geq 0$ に対し，
逆系の短完全列
$$
0 \to \{H^p(X, \mathcal{F})/f^nH^p(X, \mathcal{F})\}
\to \{H^p(X, \mathcal{F}/f^n\mathcal{F})\}
\to \{H^{p + 1}(X, \mathcal{F})[f^n]\} \to 0
$$
を得る。最初の逆系は Mittag--Leffler 条件 (ML) を満たすので，
ここから次の三点が分かる。
\begin{enumerate}
\item 短完全列
$$
0 \to \widehat{H^p(X, \mathcal{F})} \to
\lim H^p(X, \mathcal{F}/f^n\mathcal{F}) \to
T_f(H^{p + 1}(X, \mathcal{F})) \to 0
$$
が存在する。ここで $\widehat{\ }$ は通常の $f$-進完備化を表し，
$T_f( - )$ は More on Algebra, Example
\ref{more-algebra-example-spectral-sequence-principal}.
の $f$-進 Tate 加群を表す。
\item 次が成り立つ。
$R^1\lim H^p(X, \mathcal{F}/f^n\mathcal{F}) =
R^1\lim H^{p + 1}(X, \mathcal{F})[f^n]$.
\item 系 $\{H^{p + 1}(X, \mathcal{F})[f^n]\}$ が ML であるための
必要十分条件は，$\{H^p(X, \mathcal{F}/f^n\mathcal{F})\}$ が ML
であることである。
\end{enumerate}
Homology, Lemma \ref{homology-lemma-Mittag-Leffler} および
More on Algebra, Lemmas \ref{more-algebra-lemma-six-term-Rlim},
\ref{more-algebra-lemma-Mittag-Leffler} を参照せよ。
\end{remark}










\section{導来極限}
\label{cohomology-section-derived-limits}

\noindent
$(X, \mathcal{O}_X)$ を環付き空間とする。三角圏
$D(\mathcal{O}_X)$ は積をもつので
（Injectives, Lemma \ref{injectives-lemma-derived-products}），
$D(\mathcal{O}_X)$ は導来極限をもつ。Derived Categories, Definition
\ref{derived-definition-derived-limit} を参照せよ。
$(K_n)$ が $D(\mathcal{O}_X)$ の逆系ならば，その導来極限を
$R\lim K_n$ と表す。

\begin{lemma}
\label{cohomology-lemma-RGamma-commutes-with-Rlim}
$(X, \mathcal{O}_X)$ を環付き空間とする。開集合 $U \subset X$ に対し，
関手 $R\Gamma(U, -)$ は $R\lim$ と可換である。さらに，
$D(\mathcal{O}_X)$ の任意の逆系 $(K_n)$ と任意の
$m \in \mathbf{Z}$ に対して短完全列
$$
0 \to
R^1\lim H^{m - 1}(U, K_n) \to H^m(U, R\lim K_n) \to
\lim H^m(U, K_n) \to 0
$$
が存在する。
\end{lemma}

\begin{proof}
最初の主張は
Injectives, Lemma \ref{injectives-lemma-RF-commutes-with-Rlim}
から従う。次に
$R\lim R\Gamma(U, K_n) = R\Gamma(U, R\lim K_n)$ に
More on Algebra, Remark \ref{more-algebra-remark-compare-derived-limit}
を適用すれば，短完全列を得る。
\end{proof}

\begin{lemma}
\label{cohomology-lemma-Rf-commutes-with-Rlim}
$f : (X, \mathcal{O}_X) \to (Y, \mathcal{O}_Y)$ を環付き空間の射とする。
このとき $Rf_*$ は $R\lim$ と可換する。すなわち，$Rf_*$ は
導来極限と可換する。
\end{lemma}

\begin{proof}
$(K_n)$ を $D(\mathcal{O}_X)$ の逆系とする。$D(\mathcal{O}_X)$ における
定義完全三角
$$
R\lim K_n \to \prod K_n \to \prod K_n
$$
を考える。完全関手 $Rf_*$ を適用すると，$D(\mathcal{O}_Y)$ における
完全三角
$$
Rf_*(R\lim K_n) \to Rf_*\left(\prod K_n\right) \to Rf_*\left(\prod K_n\right)
$$
を得る。したがって，$Rf_*$ が導来圏の積と可換することを示せば十分である
（この積は単に複体の積で与えられるわけではない。Injectives, Lemma
\ref{injectives-lemma-derived-products} を参照せよ）。
一方，補題 \ref{cohomology-lemma-adjoint} により $Rf_*$ は右随伴なので，
これは形式的に従う（Categories, Lemma
\ref{categories-lemma-adjoint-exact} を参照せよ）。
注意：$R\lim K_n$ は $D(\mathcal{O}_X)$ における極限ではないので，
Categories, Lemma \ref{categories-lemma-adjoint-exact}
を直接適用することはできない。
\end{proof}

\begin{remark}
\label{cohomology-remark-discuss-derived-limit}
$(X, \mathcal{O}_X)$ を環付き空間とし，$(K_n)$ を
$D(\mathcal{O}_X)$ の逆系とする。$K = R\lim K_n$ とおく。
各 $n,m$ に対し，$\mathcal{H}^m_n = H^m(K_n)$ を $K_n$ の
第 $m$ コホモロジー層とし，同様に
$\mathcal{H}^m = H^m(K)$ とおく。前層
$$
U \longmapsto \underline{\mathcal{H}}^m_n(U) = H^m(U, K_n)
$$
を $\underline{\mathcal{H}}^m_n$ と表す。同様に
$\underline{\mathcal{H}}^m(U) = H^m(U, K)$ とおく。
補題 \ref{cohomology-lemma-sheafification-cohomology} により，
$\mathcal{H}^m_n$ は $\underline{\mathcal{H}}^m_n$ の層化であり，
$\mathcal{H}^m$ は $\underline{\mathcal{H}}^m$ の層化である。
次の図がある。
$$
\xymatrix{
K \ar@{=}[d] &
\underline{\mathcal{H}}^m \ar[d] \ar[r] & 
\mathcal{H}^m \ar[d] \\
R\lim K_n &
\lim \underline{\mathcal{H}}^m_n \ar[r] & 
\lim \mathcal{H}^m_n
}
$$
一般には，$\lim \mathcal{H}^m_n$ が
$\lim \underline{\mathcal{H}}^m_n$ の層化であるとは限らない。
$U \subset X$ が開集合ならば，補題
\ref{cohomology-lemma-RGamma-commutes-with-Rlim} により短完全列
\begin{equation}
\label{cohomology-equation-ses-Rlim-over-U}
0 \to
R^1\lim \underline{\mathcal{H}}^{m - 1}_n(U) \to
\underline{\mathcal{H}}^m(U) \to
\lim \underline{\mathcal{H}}^m_n(U) \to 0
\end{equation}
を得る。
\end{remark}

\noindent
次の補題は，スキーム上の準連接加群で，遷移写像が全射であるような
逆系に適用できる。

\begin{lemma}
\label{cohomology-lemma-inverse-limit-is-derived-limit}
$(X, \mathcal{O}_X)$ を環付き空間，$(\mathcal{F}_n)$ を
$\mathcal{O}_X$-加群の逆系，$\mathcal{B}$ を $X$ の開集合からなる
集合とする。次を仮定する。
\begin{enumerate}
\item $X$ の各開集合は，その要素が $\mathcal{B}$ に属する被覆をもつ。
\item $p > 0$ かつ $U \in \mathcal{B}$ ならば
$H^p(U, \mathcal{F}_n) = 0$ である。
\item $U \in \mathcal{B}$ に対し，逆系 $\mathcal{F}_n(U)$ の
$R^1\lim$ は零である。
\end{enumerate}
このとき $R\lim \mathcal{F}_n = \lim \mathcal{F}_n$ であり，
$p > 0$ かつ $U \in \mathcal{B}$ に対して
$H^p(U, \lim \mathcal{F}_n) = 0$ である。
\end{lemma}

\begin{proof}
$K_n = \mathcal{F}_n$，$K = R\lim \mathcal{F}_n$ とおく。
注記 \ref{cohomology-remark-discuss-derived-limit} の記法と仮定 (2) により，
$U \in \mathcal{B}$ に対し，$m \not = 0$ ならば
$\underline{\mathcal{H}}_n^m(U) = 0$ であり，
$\underline{\mathcal{H}}_n^0(U) = \mathcal{F}_n(U)$ である。
式 (\ref{cohomology-equation-ses-Rlim-over-U}) と仮定 (3) から，
$m \not = 0$ ならば $\underline{\mathcal{H}}^m(U) = 0$ であり，
$m = 0$ ならばこれは $\lim \mathcal{F}_n(U)$ に等しい。
(1) を用いて層化すると，$m \not = 0$ ならば
$\mathcal{H}^m = 0$ であり，$m = 0$ ならばこれは
$\lim \mathcal{F}_n$ に等しい。ゆえに $K = \lim \mathcal{F}_n$ である。
$m > 0$ に対し
$H^m(U, K) = \underline{\mathcal{H}}^m(U) = 0$ であるから
（上を参照），第二の主張も従う。
\end{proof}

\begin{lemma}
\label{cohomology-lemma-cohomology-derived-limit-injective}
$(X, \mathcal{O}_X)$ を環付き空間，$(K_n)$ を
$D(\mathcal{O}_X)$ の逆系とし，$x \in X$，$m \in \mathbf{Z}$ とする。
整数 $n(x)$ と $x$ の開近傍の基本系 $\mathfrak{U}_x$ が存在し，
$U \in \mathfrak{U}_x$ に対して次が成り立つと仮定する。
\begin{enumerate}
\item $R^1\lim H^{m - 1}(U, K_n) = 0$ である。
\item $n \geq n(x)$ に対し
$H^m(U, K_n) \to H^m(U, K_{n(x)})$ は単射である。
\end{enumerate}
このとき茎上の写像
$H^m(R\lim K_n)_x \to H^m(K_{n(x)})_x$ は単射である。
\end{lemma}

\begin{proof}
$\gamma$ を，$H^m(K_{n(x)})_x$ で零に写る
$H^m(R\lim K_n)_x$ の元とする。補題
\ref{cohomology-lemma-sheafification-cohomology} により
$H^m(R\lim K_n)$ は $U \mapsto H^m(U, R\lim K_n)$ の層化なので，
$U \in \mathfrak{U}_x$ と，$\gamma$ に写る元
$\tilde \gamma \in H^m(U, R\lim K_n)$ を選べる。
すると $\tilde\gamma$ は
$\tilde\gamma_{n(x)} \in H^m(U, K_{n(x)})$ に写る。
再び補題 \ref{cohomology-lemma-sheafification-cohomology} により
$H^m(K_{n(x)})$ は $U \mapsto H^m(U, K_{n(x)})$ の層化なので，
$U$ を縮小することにより
$\tilde\gamma_{n(x)} = 0$ と仮定できる。この $U$ に対し，補題
\ref{cohomology-lemma-RGamma-commutes-with-Rlim} の短完全列
$$
0 \to
R^1\lim H^{m - 1}(U, K_n) \to H^m(U, R\lim K_n) \to
\lim H^m(U, K_n) \to 0
$$
を考える。仮定 (1) により左の群は零であり，仮定 (2) により
右の群は $H^m(U, K_{n(x)})$ に単射で写る。
よって $\tilde\gamma = 0$，したがって望むとおり $\gamma = 0$ となる。
\end{proof}

\begin{lemma}
\label{cohomology-lemma-is-limit-per-point}
$(X, \mathcal{O}_X)$ を環付き空間とし，
$E \in D(\mathcal{O}_X)$ とする。すべての $x \in X$ に対して，
関数 $p(x, -) : \mathbf{Z} \to \mathbf{Z}$ と
$x$ の開近傍の基本系 $\mathfrak{U}_x$ が存在し，
$$
H^p(U, H^{m - p}(E)) = 0 \text{ for }
U \in \mathfrak{U}_x \text{ and } p > p(x, m)
$$
が成り立つと仮定する。このとき Derived Categories, Remark
\ref{derived-remark-map-into-derived-limit-truncations}
の写像 $E \to R\lim \tau_{\geq -n} E$ は
$D(\mathcal{O}_X)$ における同型である。
\end{lemma}

\begin{proof}
$K_n = \tau_{\geq -n}E$，$K = R\lim K_n$ とおく。
標準写像 $E \to K$ は標準写像
$E \to K_n = \tau_{\geq -n}E$ から生じる。
$E \to K$ がコホモロジー層の同型
$H^m(E) \to H^m(K)$ を誘導することを示さなければならない。
以下の証明では $m$ を固定する。$n \geq -m$ ならば，写像
$E \to \tau_{\geq -n}E = K_n$ は同型
$H^m(E) \to H^m(K_n)$ を誘導する。
証明を終えるには，各 $x \in X$ に対して，
$H^m(K)_x \to H^m(K_{n(x)})_x$ が単射となるような整数
$n(x) \geq -m$ が存在することを示せば十分である。実際，そのとき合成
$$
H^m(E)_x \to H^m(K)_x \to H^m(K_{n(x)})_x
$$
は全単射で第二の矢印は単射だから，第一の矢印は全単射である。ここで
$$
n(x) = 1 + \max\{-m, p(x, m - 1) - m, -1 + p(x, m) - m, -2 + p(x, m + 1) - m\}.
$$
とおけば，いずれの場合も $n(x) \geq -m$ である。主張：写像
$$
H^{m - 1}(U, K_{n + 1}) \to H^{m - 1}(U, K_n)
\quad\text{and}\quad
H^m(U, K_{n + 1}) \to H^m(U, K_n)
$$
は $n \geq n(x)$ かつ $U \in \mathfrak{U}_x$ に対して同型である。
この主張から補題 \ref{cohomology-lemma-cohomology-derived-limit-injective}
の条件 (1), (2) が満たされ，したがって望む単射性が従う。
完全三角
$$
H^{-n - 1}(E)[n + 1] \to
K_{n + 1} \to K_n \to H^{-n - 1}(E)[n + 2]
$$
が存在することを思い出そう（Derived Categories, Remark
\ref{derived-remark-truncation-distinguished-triangle}）。
対応するコホモロジー長完全列を見ると，$n \geq n(x)$ かつ
$U \in \mathfrak{U}_x$ に対して
$$
H^{m + n}(U, H^{-n - 1}(E)),\quad
H^{m + n + 1}(U, H^{-n - 1}(E)),\quad
H^{m + n + 2}(U, H^{-n - 1}(E))
$$
が零ならば主張が従う。これは $n(x)$ の選び方と補題の仮定から従う。
\end{proof}

\begin{lemma}
\label{cohomology-lemma-is-limit-spaltenstein}
\begin{reference}
\cite[Proposition 3.13]{Spaltenstein}
\end{reference}
$(X, \mathcal{O}_X)$ を環付き空間とし，
$E \in D(\mathcal{O}_X)$ とする。すべての $x \in X$ に対して，
整数 $d_x \geq 0$ と $x$ の開近傍の基本系 $\mathfrak{U}_x$ が存在し，
$$
H^p(U, H^q(E)) = 0 \text{ for }
U \in \mathfrak{U}_x,\ p > d_x, \text{ and }q < 0
$$
が成り立つと仮定する。このとき Derived Categories, Remark
\ref{derived-remark-map-into-derived-limit-truncations}
の写像 $E \to R\lim \tau_{\geq -n} E$ は
$D(\mathcal{O}_X)$ における同型である。
\end{lemma}

\begin{proof}
補題 \ref{cohomology-lemma-is-limit-per-point} で
$p(x, m) = d_x + \max(0, m)$ とすれば従う。
\end{proof}

\begin{lemma}
\label{cohomology-lemma-is-limit}
$(X, \mathcal{O}_X)$ を環付き空間とし，
$E \in D(\mathcal{O}_X)$ とする。関数
$p(-) : \mathbf{Z} \to \mathbf{Z}$ と $X$ の開集合からなる集合
$\mathcal{B}$ が存在し，次を満たすと仮定する。
\begin{enumerate}
\item $X$ の各開集合は，その要素が $\mathcal{B}$ に属する被覆をもつ。
\item $p > p(m)$ かつ $U \in \mathcal{B}$ ならば
$H^p(U, H^{m - p}(E)) = 0$ である。
\end{enumerate}
このとき Derived Categories, Remark
\ref{derived-remark-map-into-derived-limit-truncations}
の写像 $E \to R\lim \tau_{\geq -n} E$ は
$D(\mathcal{O}_X)$ における同型である。
\end{lemma}

\begin{proof}
補題 \ref{cohomology-lemma-is-limit-per-point} を
$p(x, m) = p(m)$ および
$\mathfrak{U}_x = \{U \in \mathcal{B} \mid x \in U\}$.
に適用する。
\end{proof}

\begin{lemma}
\label{cohomology-lemma-is-limit-dimension}
$(X, \mathcal{O}_X)$ を環付き空間とし，
$E \in D(\mathcal{O}_X)$ とする。整数 $d \geq 0$ と
$X$ の位相の基底 $\mathcal{B}$ が存在し，
$$
H^p(U, H^q(E)) = 0 \text{ for }
U \in \mathcal{B},\ p > d, \text{ and }q < 0
$$
が成り立つと仮定する。このとき Derived Categories, Remark
\ref{derived-remark-map-into-derived-limit-truncations}
の写像 $E \to R\lim \tau_{\geq -n} E$ は
$D(\mathcal{O}_X)$ における同型である。
\end{lemma}

\begin{proof}
補題 \ref{cohomology-lemma-is-limit-spaltenstein} を $d_x = d$ および
$\mathfrak{U}_x = \{U \in \mathcal{B} \mid x \in U\}$ に適用する。
\end{proof}

\noindent
上の補題を用いると，いくつかの場合にコホモロジーを計算できる。

\begin{lemma}
\label{cohomology-lemma-cohomology-over-U-trivial}
$(X, \mathcal{O}_X)$ を環付き空間とし，
$K$ を $D(\mathcal{O}_X)$ の対象とする。
$\mathcal{B}$ を $X$ の開集合からなる集合とする。次を仮定する。
\begin{enumerate}
\item $X$ の各開集合は，その要素が $\mathcal{B}$ に属する被覆をもつ。
\item すべての $p > 0$，$q \in \mathbf{Z}$，$U \in \mathcal{B}$ に対し
$H^p(U, H^q(K)) = 0$ である。
\end{enumerate}
このとき $q \in \mathbf{Z}$，$U \in \mathcal{B}$ に対し
$H^q(U, K) = H^0(U, H^q(K))$ である。
\end{lemma}

\begin{proof}
補題 \ref{cohomology-lemma-is-limit-dimension} で $d = 0$ とすれば，
$K = R\lim \tau_{\geq -n} K$ であることに注意する。
$U \in \mathcal{B}$ とする。式 (\ref{cohomology-equation-ses-Rlim-over-U})
により短完全列
$$
0 \to R^1\lim H^{q - 1}(U, \tau_{\geq -n}K) \to
H^q(U, K) \to \lim H^q(U, \tau_{\geq -n}K) \to 0
$$
を得る。例 \ref{cohomology-example-spectral-sequence} のスペクトル系列を用いると，
条件 (2) からすべての $q$ に対して
$H^q(U, \tau_{\geq -n} K) = H^0(U, H^q(\tau_{\geq -n} K))$
が従う。$\tau_{\geq -n}K$ は下に有界なので，このスペクトル系列は
収束する。$n > -q$ ならば
$H^q(\tau_{\geq -n}K) = H^q(K)$ である。
したがって，表示した短完全列の左辺と右辺に現れる系は，それぞれ
$H^0(U, H^{q - 1}(K))$ と $H^0(U, H^q(K))$ を値として最終的に定常となる。
これで補題が従う。
\end{proof}

\noindent
導来極限を記述できる別の場合を挙げる。

\begin{lemma}
\label{cohomology-lemma-derived-limit-suitable-system}
$(X, \mathcal{O}_X)$ を環付き空間，$(K_n)$ を
$D(\mathcal{O}_X)$ の対象の逆系とする。
$\mathcal{B}$ を $X$ の開集合からなる集合とする。次を仮定する。
\begin{enumerate}
\item $X$ の各開集合は，その要素が $\mathcal{B}$ に属する被覆をもつ。
\item すべての $U \in \mathcal{B}$ とすべての
$q \in \mathbf{Z}$ に対し，次が成り立つ。
\begin{enumerate}
\item $p > 0$ に対し $H^p(U, H^q(K_n)) = 0$ である。
\item 逆系 $H^0(U, H^q(K_n))$ の $R^1\lim$ は零である。
\end{enumerate}
\end{enumerate}
このとき $q \in \mathbf{Z}$ に対し
$H^q(R\lim K_n) = \lim H^q(K_n)$ である。
\end{lemma}

\begin{proof}
$K = R\lim K_n$ とおき，$U \in \mathcal{B}$ とする。
補題 \ref{cohomology-lemma-cohomology-over-U-trivial} と (2)(a) により
$H^q(U, K_n) = H^0(U, H^q(K_n))$ である。
補題 \ref{cohomology-lemma-RGamma-commutes-with-Rlim} と (2)(b) により
$H^q(U, K) = \lim H^0(U, H^q(K_n))$ である。したがって
$H^q(U, K)$ は，$U$ 上の層 $H^q(K_n)$ の切断の逆極限である。
$\lim H^q(K_n)$ は層なので，仮定 (1) を用いると，
前層 $U \mapsto H^q(U, K)$ の層化である $H^q(K)$ は
$\lim H^q(K_n)$ に等しい。これで補題が証明された。
\end{proof}








\section{K-入射的分解の構成}
\label{cohomology-section-K-injective}


\noindent
$(X, \mathcal{O}_X)$ を環付き空間とし，
$\mathcal{F}^\bullet$ を $\mathcal{O}_X$-加群の複体とする。
圏 $\textit{Mod}(\mathcal{O}_X)$ は十分な入射的対象をもつので，
Derived Categories, Lemma \ref{derived-lemma-special-inverse-system}
を用いて，$\mathcal{O}_X$-加群の複体の圏における図式
$$
\xymatrix{
\ldots \ar[r] &
\tau_{\geq -2}\mathcal{F}^\bullet \ar[r] \ar[d] &
\tau_{\geq -1}\mathcal{F}^\bullet \ar[d] \\
\ldots \ar[r] & \mathcal{I}_2^\bullet \ar[r] & \mathcal{I}_1^\bullet
}
$$
で，次を満たすものを構成できる。
\begin{enumerate}
\item 垂直の矢印は擬同型である。
\item $\mathcal{I}_n^\bullet$ は入射的対象からなる下に有界な複体である。
\item 矢印 $\mathcal{I}_{n + 1}^\bullet \to \mathcal{I}_n^\bullet$ は
各次数で分裂する全射である。
\end{enumerate}
$\mathcal{O}_X$-加群の圏は極限をもち（前層のレベルで計算される），
したがって各次数ごとの極限
$\mathcal{I}^\bullet = \lim_n \mathcal{I}_n^\bullet$ を作れる。
Derived Categories, Lemmas
\ref{derived-lemma-bounded-below-injectives-K-injective} および
\ref{derived-lemma-limit-K-injectives}
により，これは K-入射的複体である。一般には標準写像
\begin{equation}
\label{cohomology-equation-into-candidate-K-injective}
\mathcal{F}^\bullet \to \mathcal{I}^\bullet
\end{equation}
は擬同型とは限らない。次の補題で，これが擬同型となるための
いくつかの条件を述べる。

\begin{lemma}
\label{cohomology-lemma-K-injective}
上で述べた状況を考える。
$\mathcal{H}^m = H^m(\mathcal{F}^\bullet)$ を第 $m$ コホモロジー層とする。
$\mathcal{B}$ を $X$ の開部分集合からなる集合とし，
$d \in \mathbf{N}$ とする。次を仮定する。
\begin{enumerate}
\item $X$ の各開集合は，その要素が $\mathcal{B}$ に属する被覆をもつ。
\item すべての $U \in \mathcal{B}$ に対し，$p > d$ かつ $q < 0$
ならば $H^p(U, \mathcal{H}^q) = 0$ である\footnote{条件
$\forall m$, $\exists p(m)$, $H^p(U. \mathcal{H}^{m - p}) = 0$ が
$p > p(m)$ のとき成り立てば十分である。補題
\ref{cohomology-lemma-is-limit} を参照せよ。}。
\end{enumerate}
このとき (\ref{cohomology-equation-into-candidate-K-injective}) は擬同型である。
\end{lemma}

\begin{proof}
Derived Categories, Lemma \ref{derived-lemma-difficulty-K-injectives}
により，写像
$\mathcal{F}^\bullet \to R\lim \tau_{\geq -n} \mathcal{F}^\bullet$
が同型であることを示せば十分である。これは補題
\ref{cohomology-lemma-is-limit-dimension} である。
\end{proof}

\noindent
層の複体からなる系の逆極限のコホモロジー層に関する技術的な補題を挙げる。
ある意味では，実際の逆極限ではなく導来極限を取るべきなので，
この補題を証明しようとすること自体が適切ではない。

\begin{lemma}
\label{cohomology-lemma-inverse-limit-complexes}
$(X, \mathcal{O}_X)$ を環付き空間とし，
$(\mathcal{F}_n^\bullet)$ を $\mathcal{O}_X$-加群の複体の逆系とする。
$m \in \mathbf{Z}$ とする。$X$ の開部分集合からなる集合
$\mathcal{B}$ と整数 $n_0$ が存在し，次を満たすと仮定する。
\begin{enumerate}
\item $X$ の各開集合は，その要素が $\mathcal{B}$ に属する被覆をもつ。
\item すべての $U \in \mathcal{B}$ に対し，次が成り立つ。
\begin{enumerate}
\item アーベル群の系
$\mathcal{F}_n^{m - 2}(U)$ および $\mathcal{F}_n^{m - 1}(U)$
の $R^1\lim$ は零である（例えば，これらが Mittag--Leffler
条件を満たせばよい）。
\item アーベル群の系 $H^{m - 1}(\mathcal{F}_n^\bullet(U))$
の $R^1\lim$ は零である（例えば，この系が Mittag--Leffler
条件を満たせばよい）。
\item すべての $n \geq n_0$ に対し
$H^m(\mathcal{F}_n^\bullet(U)) = H^m(\mathcal{F}_{n_0}^\bullet(U))$
である。
\end{enumerate}
\end{enumerate}
このとき写像
$H^m(\mathcal{F}^\bullet) \to \lim H^m(\mathcal{F}_n^\bullet) \to
H^m(\mathcal{F}_{n_0}^\bullet)$
は層の同型である。ここで
$\mathcal{F}^\bullet = \lim \mathcal{F}_n^\bullet$ は各次数ごとの
逆極限である。
\end{lemma}

\begin{proof}
$U \in \mathcal{B}$ とする。$H^m(\mathcal{F}^\bullet(U))$ は
$$
\lim_n \mathcal{F}_n^{m - 2}(U) \to
\lim_n \mathcal{F}_n^{m - 1}(U) \to
\lim_n \mathcal{F}_n^m(U) \to
\lim_n \mathcal{F}_n^{m + 1}(U)
$$
の左から三番目におけるコホモロジーであることに注意する。
仮定 (2)(a), (2)(b) により
More on Algebra, Lemma \ref{more-algebra-lemma-apply-Mittag-Leffler-again}
を適用でき，
$$
H^m(\mathcal{F}^\bullet(U)) = \lim H^m(\mathcal{F}_n^\bullet(U))
$$
と結論する。仮定 (2)(c) により，すべての $n \geq n_0$ に対して
$$
H^m(\mathcal{F}^\bullet(U)) = H^m(\mathcal{F}_n^\bullet(U))
$$
となる。仮定 (1) により，すべての $n \geq n_0$ に対して，
$U \mapsto H^m(\mathcal{F}^\bullet(U))$ の層化は
$U \mapsto H^m(\mathcal{F}_n^\bullet(U))$ の層化に等しい。
したがって層の逆系 $H^m(\mathcal{F}_n^\bullet)$ は
$n \geq n_0$ において値 $H^m(\mathcal{F}^\bullet)$ をもつ定値系である。
これで補題が証明された。
\end{proof}












\section{逆系とコホモロジー III}
\label{cohomology-section-inverse-systems-tri}

\noindent
この節では，導来極限を用いて
節 \ref{cohomology-section-inverse-systems-bis} の議論を続ける。

\begin{lemma}
\label{cohomology-lemma-formal-functions-principal}
$(X, \mathcal{O}_X)$ を環付き空間とし，
$A \to \Gamma(X, \mathcal{O}_X)$ を環準同型，$f \in A$ とする。
$E$ を $D(\mathcal{O}_X)$ の対象とする。
$$
E_n = E \otimes_{\mathcal{O}_X} (\mathcal{O}_X \xrightarrow{f^n} \mathcal{O}_X)
$$
と表し，$E^\wedge = R\lim E_n$ とおく。$p \in \mathbf{Z}$ に対し，
行と列が完全な次の標準可換図式が存在する。
$$
\xymatrix{
& 0 & 0 \\
0 \ar[r] &
\widehat{H^p(X, E)} \ar[r] \ar[u] &
\lim H^p(X, E_n) \ar[r] \ar[u] &
T_f(H^{p + 1}(X, E)) \ar[r] &
0 \\
0 \ar[r] &
H^0(H^p(X, E)^\wedge) \ar[r] \ar[u] &
H^p(X, E^\wedge) \ar[r] \ar[u] &
T_f(H^{p + 1}(X, E)) \ar[r] \ar@{=}[u] &
0 \\
&
R^1\lim H^p(X, E)[f^n] \ar[u] \ar[r]^\cong &
R^1\lim H^{p - 1}(X, E_n) \ar[u] \\
& 0 \ar[u] & 0 \ar[u]
}
$$
ここで $\widehat{H^p(X, E)} = \lim H^p(X, E)/f^n H^p(X, E)$ は通常の
$f$-進完備化，$H^p(X, E)^\wedge$ は導来 $f$-進完備化，
$T_f(H^{p + 1}(X, E))$ は $f$-進 Tate 加群である。More on Algebra,
Example
\ref{more-algebra-example-spectral-sequence-principal}.
を参照せよ。最後に
$H^p(X, E^\wedge) = H^p(R\Gamma(X, E)^\wedge)$ が成り立つ。
\end{lemma}

\begin{proof}
補題 \ref{cohomology-lemma-Rf-commutes-with-Rlim} により
$R\Gamma(X, E^\wedge) = R\lim R\Gamma(X, E_n)$ であることに注意する。
他方，
$$
R\Gamma(X, E_n) =
R\Gamma(X, E) \otimes_A^\mathbf{L} (A \xrightarrow{f^n} A)
$$
である（詳細は省略する）。したがって $R\Gamma(X, E^\wedge)$ は
$R\Gamma(X, E)$ の導来 $f$-進完備化
$R\Gamma(X, E)^\wedge$ である。この図式は
More on Algebra, Lemma
\ref{more-algebra-lemma-compare-completion-derived-completion} から従う。
\end{proof}

\begin{lemma}
\label{cohomology-lemma-stabilizes}
$\mathcal{A}$ をアーベル圏とし，$f : M \to M$ を
$\mathcal{A}$ の射とする。
$M[f^n] = \Ker(f^n : M \to M)$ が安定するならば，逆系
$$
(M \xrightarrow{f^n} M)
\quad\text{and}\quad
\Coker(f^n : M \to M)
$$
は $D(\mathcal{A})$ においてプロ同型である。
\end{lemma}

\begin{proof}
第一の逆系から第二の逆系への写像が明らかに存在する。
$M[f^c] = M[f^{c + 1}] = M[f^{c + 2}] = \ldots$.
と仮定する。このとき，図式
$$
\xymatrix{
M/M[f^c] \ar[r]_-{f^{n + c}} \ar[d]_{f^c} &
M \ar[d]^1 \\
M \ar[r]^{f^n} & M
}
$$
によって，$D(\mathcal{A})$ における逆系の逆向きの矢印を定義できる。
上の水平矢印は単射なので，上段の複体は
$\Coker(f^{n + c} : M \to M)$ と擬同型である。いくつかの詳細は省略する。
\end{proof}

\begin{example}
\label{cohomology-example-formal-functions-principal}
$(X, \mathcal{O}_X)$ を環付き空間とし，
$A \to \Gamma(X, \mathcal{O}_X)$ を環準同型，$f \in A$ とする。
$\mathcal{F}$ を $\mathcal{O}_X$-加群とする。すべての $n \geq c$ に対し
$\mathcal{F}[f^c] = \mathcal{F}[f^n]$ となる $c$ が存在すると仮定する。
$E = \mathcal{F}$ として補題
\ref{cohomology-lemma-formal-functions-principal} を適用する。
補題 \ref{cohomology-lemma-stabilizes} により，逆系 $(E_n)$ は逆系
$(\mathcal{F}/f^n\mathcal{F})$ とプロ同型である。
したがって $p \in \mathbf{Z}$ に対し，行と列が完全な可換図式
$$
\xymatrix{
& 0 & 0 \\
0 \ar[r] &
\widehat{H^p(X, \mathcal{F})} \ar[r] \ar[u] &
\lim H^p(X, \mathcal{F}/f^n\mathcal{F}) \ar[r] \ar[u] &
T_f(H^{p + 1}(X, \mathcal{F})) \ar[r] &
0 \\
0 \ar[r] &
H^0(H^p(X, \mathcal{F})^\wedge) \ar[r] \ar[u] &
H^p(R\Gamma(X, \mathcal{F})^\wedge) \ar[r] \ar[u] &
T_f(H^{p + 1}(X, \mathcal{F})) \ar[r] \ar@{=}[u] &
0 \\
&
R^1\lim H^p(X, \mathcal{F})[f^n] \ar[u] \ar[r]^\cong &
R^1\lim H^{p - 1}(X, \mathcal{F}/f^n\mathcal{F}) \ar[u] \\
& 0 \ar[u] & 0 \ar[u]
}
$$
を得る。ここで
$\widehat{H^p(X, \mathcal{F})} =
\lim H^p(X, \mathcal{F})/f^n H^p(X, \mathcal{F})$
は通常の $f$-進完備化であり，$D(A)$ の $M$ に対して
$M^\wedge$ は導来 $f$-進完備化を表す。
\end{example}
















\section{非有界複体の {\v C}ech コホモロジー}
\label{cohomology-section-cech-cohomology-of-unbounded-complexes}

\noindent
節 \ref{cohomology-section-cech-cohomology-of-complexes} の構成は，
非有界複体に対して「正しい」ものではない。問題は，Stacks project では
二重複体の全複体化に直和を用いており，これを積に置き換える必要があることにある。
本節ではそうする代わりに，被覆が有限であると仮定し，交代 {\v C}ech 複体を用いる。

\medskip\noindent
$(X, \mathcal{O}_X)$ を環付き空間とする。
${\mathcal F}^\bullet$ を $\mathcal{O}_X$-加群の前層の複体とし，
${\mathcal U} : X = \bigcup_{i \in I} U_i$ を $X$ の
{\bf 有限}開被覆とする。交代 {\v C}ech 複体
$\check{\mathcal{C}}_{alt}^\bullet(\mathcal{U}, \mathcal{F})$
（節 \ref{cohomology-section-alternating-cech}）は前層 $\mathcal{F}$ に関して
関手的なので，二重複体
$\check{\mathcal{C}}^\bullet_{alt}(\mathcal{U}, \mathcal{F}^\bullet)$.
を得る。本節では，付随する全複体を用いる。
$\text{Tot}(\check{\mathcal{C}}^\bullet_{alt}({\mathcal U},
{\mathcal F}^\bullet))$
の構成は ${\mathcal F}^\bullet$ に関して関手的である。さらに，複体の
関手的変換
\begin{equation}
\label{cohomology-equation-global-sections-to-alternating-cech}
\Gamma(X, {\mathcal F}^\bullet)
\longrightarrow
\text{Tot}(\check{\mathcal{C}}^\bullet_{alt}({\mathcal U},
{\mathcal F}^\bullet))
\end{equation}
があり，次の規則で定義される。切断
$s\in \Gamma(X, {\mathcal F}^n)$ は，
$\alpha_{i_0} = s|_{U_{i_0}}$ かつ $p > 0$ に対して
$\alpha_{i_0\ldots i_p} = 0$ となる元
$\alpha = \{\alpha_{i_0\ldots i_p}\}$ に写る。

\begin{lemma}
\label{cohomology-lemma-alternating-cech-complex-complex}
$(X, \mathcal{O}_X)$ を環付き空間とし，
$\mathcal{U} : X = \bigcup_{i \in I} U_i$ を有限開被覆とする。
$\mathcal{O}_X$-加群の複体 $\mathcal{F}^\bullet$ に対し，
標準写像
$$
\text{Tot}(\check{\mathcal{C}}^\bullet_{alt}(\mathcal{U}, \mathcal{F}^\bullet))
\longrightarrow
R\Gamma(X, \mathcal{F}^\bullet)
$$
が存在する。これは $\mathcal{F}^\bullet$ に関して関手的であり，
(\ref{cohomology-equation-global-sections-to-alternating-cech}) と両立する。
\end{lemma}

\begin{proof}
${\mathcal I}^\bullet$ を，各項が入射的 $\mathcal{O}_X$-加群である
K-入射的複体とする。$\mathcal{I}^\bullet$ に対する写像
(\ref{cohomology-equation-global-sections-to-alternating-cech}) は
$\Gamma(X, {\mathcal I}^\bullet) \to
\text{Tot}(\check{\mathcal{C}}^\bullet_{alt}({\mathcal U},
{\mathcal I}^\bullet))$.
である。Homology, Lemma
\ref{homology-lemma-double-complex-gives-resolution} を二重複体
$\check{\mathcal{C}}^\bullet_{alt}({\mathcal U},
{\mathcal I}^\bullet)$ に適用し，その際補題
\ref{cohomology-lemma-injective-trivial-cech}, \ref{cohomology-lemma-alternating-usual} を用いると，
これはアーベル群の複体の擬同型であることが従う。
${\mathcal F}^\bullet \to {\mathcal I}^\bullet$ を，
${\mathcal F}^\bullet$ から各項が入射的対象である K-入射的複体への
擬同型とする（Injectives, Theorem
\ref{injectives-theorem-K-injective-embedding-grothendieck}).
$R\Gamma(X, {\mathcal F}^\bullet)$ は複体
$\Gamma(X, {\mathcal I}^\bullet)$ によって表現されるので，写像
$$
\text{Tot}(\check{\mathcal{C}}^\bullet_{alt}({\mathcal U},
{\mathcal F}^\bullet))
\longrightarrow
\text{Tot}(\check{\mathcal{C}}^\bullet_{alt}({\mathcal U},
{\mathcal I}^\bullet)).
$$
を用いて補題の写像を得る。関手性と両立性の確認は省略する。
\end{proof}

\begin{lemma}
\label{cohomology-lemma-alternating-cech-complex-complex-ss}
$(X, \mathcal{O}_X)$ を環付き空間とし，
$\mathcal{U} : X = \bigcup_{i \in I} U_i$ を有限開被覆とする。
$\mathcal{F}^\bullet$ を $\mathcal{O}_X$-加群の複体，
$\mathcal{B}$ を $X$ の開部分集合からなる集合とする。次を仮定する。
\begin{enumerate}
\item $X$ の各開集合は，その要素が $\mathcal{B}$ に属する被覆をもつ。
\item すべての $i_0, \ldots, i_p \in I$ に対し
$U_{i_0\ldots i_p} \in \mathcal{B}$ である。
\item すべての $U \in \mathcal{B}$ と $p > 0$ に対し，次が成り立つ。
\begin{enumerate}
\item $H^p(U, \mathcal{F}^q) = 0$ である。
\item $H^p(U, \Coker(\mathcal{F}^{q - 1} \to \mathcal{F}^q)) = 0$ である。
\item $H^p(U, H^q(\mathcal{F})) = 0$ である。
\end{enumerate}
\end{enumerate}
このとき補題 \ref{cohomology-lemma-alternating-cech-complex-complex} の写像
$$
\text{Tot}(\check{\mathcal{C}}^\bullet_{alt}(\mathcal{U}, \mathcal{F}^\bullet))
\longrightarrow
R\Gamma(X, \mathcal{F}^\bullet)
$$
は $D(\textit{Ab})$ における同型である。
\end{lemma}

\begin{proof}
まず $\mathcal{F}^\bullet$ が下に有界であると仮定する。この場合，写像
$$
\text{Tot}(\check{\mathcal{C}}^\bullet_{alt}(\mathcal{U}, \mathcal{F}^\bullet))
\longrightarrow
\text{Tot}(\check{\mathcal{C}}^\bullet(\mathcal{U}, \mathcal{F}^\bullet))
$$
は補題 \ref{cohomology-lemma-alternating-usual} により擬同型である。
実際，二重複体の写像
$\check{\mathcal{C}}^\bullet_{alt}(\mathcal{U}, \mathcal{F}^\bullet) \to
\check{\mathcal{C}}^\bullet(\mathcal{U}, \mathcal{F}^\bullet)$
は，これらの複体に付随する第二スペクトル系列の第一頁の間に同型を誘導する
（Homology, Lemma \ref{homology-lemma-ss-double-complex}）。
また，これらのスペクトル系列は収束する
（Homology, Lemma \ref{homology-lemma-first-quadrant-ss}）。
したがってこの場合の結論は，補題
\ref{cohomology-lemma-cech-complex-complex-computes} と仮定 (3)(a) から従う。

\medskip\noindent
一般の場合，仮定 (3)(c) により，補題 \ref{cohomology-lemma-K-injective}
におけるような分解
$\mathcal{F}^\bullet \to \mathcal{I}^\bullet = \lim \mathcal{I}_n^\bullet$
を選べる。このとき補題の写像は
$$
\lim_n 
\text{Tot}(\check{\mathcal{C}}^\bullet_{alt}(\mathcal{U},
\tau_{\geq -n}\mathcal{F}^\bullet))
\longrightarrow
\Gamma(X, \mathcal{I}^\bullet) =
\lim_n \Gamma(X, \mathcal{I}_n^\bullet)
$$
となる。ここで矢印は導来圏におけるものだが，右辺の等式は複体のレベルで
成り立つ。(3)(b) により
$\tau_{\geq -n}\mathcal{F}^\bullet$ は下に有界で，補題の仮定を満たす
複体であることに注意する。したがって下に有界な複体の場合から，各写像
$$
\text{Tot}(\check{\mathcal{C}}^\bullet_{alt}(\mathcal{U},
\tau_{\geq -n}\mathcal{F}^\bullet))
\longrightarrow
\Gamma(X, \mathcal{I}_n^\bullet)
$$
は擬同型である。左辺の複体のコホモロジーは，各所与の次数において
最終的に定値となる（交代 {\v C}ech 複体は有限だからである）。
したがって右辺についても同じことが成り立つ。ゆえに Homology, Lemma
\ref{homology-lemma-apply-Mittag-Leffler-again}
（または，より強い More on Algebra, Lemma
\ref{more-algebra-lemma-apply-Mittag-Leffler-again}）により，
右辺の極限のコホモロジーはこの定値に等しい。これで補題が従う。
\end{proof}






\section{Hom 複体}
\label{cohomology-section-hom-complexes}

\noindent
$(X, \mathcal{O}_X)$ を環付き空間とし，
$\mathcal{L}^\bullet$ と $\mathcal{M}^\bullet$ を
$\mathcal{O}_X$-加群の二つの複体とする。$\mathcal{O}_X$-加群の複体
$\SheafHom^\bullet(\mathcal{L}^\bullet, \mathcal{M}^\bullet)$.
を次のように構成する。各 $n$ に対して
$$
\SheafHom^n(\mathcal{L}^\bullet, \mathcal{M}^\bullet) =
\prod\nolimits_{n = p + q}
\SheafHom_{\mathcal{O}_X}(\mathcal{L}^{-q}, \mathcal{M}^p)
$$
とおく。$\SheafHom^n$ は，次数付き $\mathcal{O}_X$-加群とみなした
$\mathcal{L}^\bullet$ から $\mathcal{M}^\bullet$ への
次数 $n$ の斉次な $\mathcal{O}_X$-線形写像すべてからなる
$\mathcal{O}_X$-加群の層と考えるとよい。この用語の下で，微分を規則
$$
\text{d}(f) =
\text{d}_\mathcal{M} \circ f - (-1)^n f \circ \text{d}_\mathcal{L}
$$
によって定義する。ただし
$f \in \SheafHom^n_{\mathcal{O}_X}(\mathcal{L}^\bullet, \mathcal{M}^\bullet)$
である。$\text{d}^2 = 0$ の確認は省略する。この構成は
Differential Graded Algebra, Example \ref{dga-example-category-complexes}
の特別な場合である。構成から直ちに，すべての $n \in \mathbf{Z}$ と
すべての開集合 $U \subset X$ に対し
\begin{equation}
\label{cohomology-equation-cohomology-hom-complex}
H^n(\Gamma(U, \SheafHom^\bullet(\mathcal{L}^\bullet, \mathcal{M}^\bullet))) =
\Hom_{K(\mathcal{O}_U)}(\mathcal{L}^\bullet, \mathcal{M}^\bullet[n])
\end{equation}
を得る。

\begin{lemma}
\label{cohomology-lemma-compose}
$(X, \mathcal{O}_X)$ を環付き空間とする。
$\mathcal{O}_X$-加群の複体
$\mathcal{K}^\bullet, \mathcal{L}^\bullet, \mathcal{M}^\bullet$
が与えられると，$\mathcal{K}^\bullet, \mathcal{L}^\bullet,
\mathcal{M}^\bullet$ に関して関手的な $\mathcal{O}_X$-加群の複体の同型
$$
\SheafHom^\bullet(\mathcal{K}^\bullet,
\SheafHom^\bullet(\mathcal{L}^\bullet, \mathcal{M}^\bullet))
=
\SheafHom^\bullet(\text{Tot}(\mathcal{K}^\bullet \otimes_{\mathcal{O}_X}
\mathcal{L}^\bullet), \mathcal{M}^\bullet)
$$
が存在する。
\end{lemma}

\begin{proof}
省略する。ヒント：More on Algebra, Lemma
\ref{more-algebra-lemma-compose} とまったく同様に証明される。
\end{proof}

\begin{lemma}
\label{cohomology-lemma-composition}
$(X, \mathcal{O}_X)$ を環付き空間とする。
$\mathcal{O}_X$-加群の複体
$\mathcal{K}^\bullet, \mathcal{L}^\bullet, \mathcal{M}^\bullet$
が与えられると，$\mathcal{O}_X$-加群の複体の標準射
$$
\text{Tot}\left(
\SheafHom^\bullet(\mathcal{L}^\bullet, \mathcal{M}^\bullet)
\otimes_{\mathcal{O}_X}
\SheafHom^\bullet(\mathcal{K}^\bullet, \mathcal{L}^\bullet)
\right)
\longrightarrow
\SheafHom^\bullet(\mathcal{K}^\bullet, \mathcal{M}^\bullet)
$$
が存在する。
\end{lemma}

\begin{proof}
省略する。ヒント：More on Algebra, Lemma
\ref{more-algebra-lemma-composition} とまったく同様に証明される。
\end{proof}

\begin{lemma}
\label{cohomology-lemma-diagonal-better}
$(X, \mathcal{O}_X)$ を環付き空間とする。
$\mathcal{O}_X$-加群の複体
$\mathcal{K}^\bullet, \mathcal{L}^\bullet, \mathcal{M}^\bullet$
が与えられると，三つの複体すべてに関して関手的な
$\mathcal{O}_X$-加群の複体の標準射
$$
\text{Tot}\left(
\mathcal{K}^\bullet \otimes_{\mathcal{O}_X}
\SheafHom^\bullet(\mathcal{M}^\bullet, \mathcal{L}^\bullet)
\right)
\longrightarrow
\SheafHom^\bullet(\mathcal{M}^\bullet,
\text{Tot}(\mathcal{K}^\bullet \otimes_{\mathcal{O}_X} \mathcal{L}^\bullet))
$$
が存在する。
\end{lemma}

\begin{proof}
省略する。ヒント：More on Algebra, Lemma
\ref{more-algebra-lemma-diagonal-better} とまったく同様に証明される。
\end{proof}

\begin{lemma}
\label{cohomology-lemma-diagonal}
$(X, \mathcal{O}_X)$ を環付き空間とする。
$\mathcal{O}_X$-加群の複体
$\mathcal{K}^\bullet, \mathcal{L}^\bullet$
が与えられると，両方の複体に関して関手的な
$\mathcal{O}_X$-加群の複体の標準射
$$
\mathcal{K}^\bullet
\longrightarrow
\SheafHom^\bullet(\mathcal{L}^\bullet,
\text{Tot}(\mathcal{K}^\bullet \otimes_{\mathcal{O}_X} \mathcal{L}^\bullet))
$$
が存在する。
\end{lemma}

\begin{proof}
省略する。ヒント：More on Algebra, Lemma
\ref{more-algebra-lemma-diagonal} とまったく同様に証明される。
\end{proof}

\begin{lemma}
\label{cohomology-lemma-evaluate-and-more}
$(X, \mathcal{O}_X)$ を環付き空間とする。
$\mathcal{O}_X$-加群の複体
$\mathcal{K}^\bullet, \mathcal{L}^\bullet, \mathcal{M}^\bullet$
が与えられると，三つの複体すべてに関して関手的な
$\mathcal{O}_X$-加群の複体の標準射
$$
\text{Tot}(\SheafHom^\bullet(\mathcal{L}^\bullet,
\mathcal{M}^\bullet) \otimes_{\mathcal{O}_X} \mathcal{K}^\bullet)
\longrightarrow
\SheafHom^\bullet(\SheafHom^\bullet(\mathcal{K}^\bullet,
\mathcal{L}^\bullet), \mathcal{M}^\bullet)
$$
が存在する。
\end{lemma}

\begin{proof}
省略する。ヒント：More on Algebra, Lemma
\ref{more-algebra-lemma-evaluate-and-more} とまったく同様に証明される。
\end{proof}

\begin{lemma}
\label{cohomology-lemma-RHom-into-K-injective}
$(X, \mathcal{O}_X)$ を環付き空間とし，$L,M$ を
$D(\mathcal{O}_X)$ の対象とする。
$\mathcal{I}^\bullet$ を $M$ を表現する $\mathcal{O}_X$-加群の
K-入射的複体とし，$\mathcal{L}^\bullet$ を $L$ を表現する
$\mathcal{O}_X$-加群の複体とする。このとき，すべての開集合
$U \subset X$ に対し
$$
H^0(\Gamma(U, \SheafHom^\bullet(\mathcal{L}^\bullet, \mathcal{I}^\bullet))) =
\Hom_{D(\mathcal{O}_U)}(L|_U, M|_U)
$$
である。
\end{lemma}

\begin{proof}
次が成り立つ。
\begin{align*}
H^0(\Gamma(U, \SheafHom^\bullet(\mathcal{L}^\bullet, \mathcal{I}^\bullet)))
& =
\Hom_{K(\mathcal{O}_U)}(\mathcal{L}^\bullet|_U, \mathcal{I}^\bullet|_U) \\
& =
\Hom_{D(\mathcal{O}_U)}(L|_U, M|_U)
\end{align*}
第一の等式は (\ref{cohomology-equation-cohomology-hom-complex}) である。
第二の等式は，補題 \ref{cohomology-lemma-restrict-K-injective-to-open} により
$\mathcal{I}^\bullet|_U$ が K-入射的であることから成り立つ。
\end{proof}

\begin{lemma}
\label{cohomology-lemma-RHom-well-defined}
$(X, \mathcal{O}_X)$ を環付き空間とする。
$(\mathcal{I}')^\bullet \to \mathcal{I}^\bullet$
を $\mathcal{O}_X$-加群の K-入射的複体の擬同型とし，
$(\mathcal{L}')^\bullet \to \mathcal{L}^\bullet$
を $\mathcal{O}_X$-加群の複体の擬同型とする。このとき
$$
\SheafHom^\bullet(\mathcal{L}^\bullet, (\mathcal{I}')^\bullet)
\longrightarrow
\SheafHom^\bullet((\mathcal{L}')^\bullet, \mathcal{I}^\bullet)
$$
は擬同型である。
\end{lemma}

\begin{proof}
$M$ を $\mathcal{I}^\bullet$ と $(\mathcal{I}')^\bullet$ が表現する
$D(\mathcal{O}_X)$ の対象とし，$L$ を
$\mathcal{L}^\bullet$ と $(\mathcal{L}')^\bullet$ が表現する
$D(\mathcal{O}_X)$ の対象とする。
補題 \ref{cohomology-lemma-RHom-into-K-injective} により，層
$$
H^0(\SheafHom^\bullet(\mathcal{L}^\bullet, (\mathcal{I}')^\bullet))
\quad\text{and}\quad
H^0(\SheafHom^\bullet((\mathcal{L}')^\bullet, \mathcal{I}^\bullet))
$$
はいずれも前層
$$
U \longmapsto \Hom_{D(\mathcal{O}_U)}(L|_U, M|_U)
$$
に付随する層に等しい。したがってこの写像は擬同型である。
\end{proof}

\begin{lemma}
\label{cohomology-lemma-RHom-from-K-flat-into-K-injective}
$(X, \mathcal{O}_X)$ を環付き空間とする。
$\mathcal{I}^\bullet$ を $\mathcal{O}_X$-加群の K-入射的複体，
$\mathcal{L}^\bullet$ を $\mathcal{O}_X$-加群の K-平坦複体とする。
このとき $\SheafHom^\bullet(\mathcal{L}^\bullet, \mathcal{I}^\bullet)$
は $\mathcal{O}_X$-加群の K-入射的複体である。
\end{lemma}

\begin{proof}
実際，$\mathcal{K}^\bullet$ が $\mathcal{O}_X$-加群の
非輪状複体ならば，次が成り立つ。
\begin{align*}
\Hom_{K(\mathcal{O}_X)}(\mathcal{K}^\bullet,
\SheafHom^\bullet(\mathcal{L}^\bullet, \mathcal{I}^\bullet))
& =
H^0(\Gamma(X,
\SheafHom^\bullet(\mathcal{K}^\bullet,
\SheafHom^\bullet(\mathcal{L}^\bullet, \mathcal{I}^\bullet)))) \\
& =
H^0(\Gamma(X, \SheafHom^\bullet(\text{Tot}(
\mathcal{K}^\bullet \otimes_{\mathcal{O}_X} \mathcal{L}^\bullet),
\mathcal{I}^\bullet))) \\
& =
\Hom_{K(\mathcal{O}_X)}(
\text{Tot}(\mathcal{K}^\bullet \otimes_{\mathcal{O}_X} \mathcal{L}^\bullet),
\mathcal{I}^\bullet) \\
& =
0
\end{align*}
第一の等式は (\ref{cohomology-equation-cohomology-hom-complex}) から，
第二の等式は補題 \ref{cohomology-lemma-compose} から，
第三の等式は (\ref{cohomology-equation-cohomology-hom-complex}) から従う。
最後の等式は，$\mathcal{L}^\bullet$ が K-平坦なので
$\text{Tot}(\mathcal{K}^\bullet \otimes_{\mathcal{O}_X} \mathcal{L}^\bullet)$
が非輪状であること（Definition \ref{cohomology-definition-K-flat}）と，
$\mathcal{I}^\bullet$ が K-入射的であることから従う。
\end{proof}








\section{導来圏における内部 Hom}
\label{cohomology-section-internal-hom}

\noindent
$(X, \mathcal{O}_X)$ を環付き空間とし，$L,M$ を
$D(\mathcal{O}_X)$ の対象とする。$D(\mathcal{O}_X)$ の任意の第三の対象
$K$ に対して標準全単射
\begin{equation}
\label{cohomology-equation-internal-hom}
\Hom_{D(\mathcal{O}_X)}(K, R\SheafHom(L, M))
=
\Hom_{D(\mathcal{O}_X)}(K \otimes_{\mathcal{O}_X}^\mathbf{L} L, M)
\end{equation}
が存在するような $D(\mathcal{O}_X)$ の対象
$R\SheafHom(L, M)$ を構成したい。Yoneda の補題
（Categories, Lemma \ref{categories-lemma-yoneda}）により，
この式は $R\SheafHom(L, M)$ を一意な同型を除いて定めることに注意する。

\medskip\noindent
このような対象を構成するため，$M$ を表現する K-入射的複体
$\mathcal{I}^\bullet$ と，$L$ を表現する任意の
$\mathcal{O}_X$-加群の複体 $\mathcal{L}^\bullet$ を選ぶ。そして
$$
R\SheafHom(L, M) = \SheafHom^\bullet(\mathcal{L}^\bullet, \mathcal{I}^\bullet)
$$
とおく。ここで右辺は節 \ref{cohomology-section-hom-complexes} で構成した
$\mathcal{O}_X$-加群の複体である。補題
\ref{cohomology-lemma-RHom-well-defined} により，これは良定義である。関手
$$
D(\mathcal{O}_X)^{opp} \times D(\mathcal{O}_X) \longrightarrow D(\mathcal{O}_X),
\quad
(K, L) \longmapsto R\SheafHom(K, L)
$$
を得る。(\ref{cohomology-equation-internal-hom}) を証明する準備として，
$R\SheafHom(K, L)$ のコホモロジー群を計算する。

\begin{lemma}
\label{cohomology-lemma-section-RHom-over-U}
$(X, \mathcal{O}_X)$ を環付き空間とし，$L,M$ を
$D(\mathcal{O}_X)$ の対象とする。すべての開集合 $U$ に対し
$$
H^0(U, R\SheafHom(L, M)) =
\Hom_{D(\mathcal{O}_U)}(L|_U, M|_U)
$$
であり，特に
$H^0(X, R\SheafHom(L, M)) = \Hom_{D(\mathcal{O}_X)}(L, M)$ である。
\end{lemma}

\begin{proof}
K-入射的な $\mathcal{O}_X$-加群の複体 $\mathcal{I}^\bullet$ で
$M$ を表現するものと，K-平坦複体 $\mathcal{L}^\bullet$ で
$L$ を表現するものを選ぶ。このとき
$\SheafHom^\bullet(\mathcal{L}^\bullet, \mathcal{I}^\bullet)$
は補題 \ref{cohomology-lemma-RHom-from-K-flat-into-K-injective} により K-入射的である。
したがって $U$ 上の切断を取るだけで $U$ 上のコホモロジーを計算でき，
結果は補題 \ref{cohomology-lemma-RHom-into-K-injective} から従う。
\end{proof}

\begin{lemma}
\label{cohomology-lemma-internal-hom}
$(X, \mathcal{O}_X)$ を環付き空間とし，$K,L,M$ を
$D(\mathcal{O}_X)$ の対象とする。上で述べた構成の下で，
$K,L,M$ に関して関手的な $D(\mathcal{O}_X)$ における標準同型
$$
R\SheafHom(K, R\SheafHom(L, M)) =
R\SheafHom(K \otimes_{\mathcal{O}_X}^\mathbf{L} L, M)
$$
が存在する。$H^0(X, -)$ を取ると
(\ref{cohomology-equation-internal-hom}) が復元される。
\end{lemma}

\begin{proof}
$M$ を表現する K-入射的複体 $\mathcal{I}^\bullet$ と，
$L$ を表現する K-平坦な $\mathcal{O}_X$-加群の複体
$\mathcal{L}^\bullet$ を選ぶ。$\mathcal{K}^\bullet$ を
$K$ を表現する任意の $\mathcal{O}_X$-加群の複体とする。このとき
$$
\SheafHom^\bullet(\mathcal{K}^\bullet,
\SheafHom^\bullet(\mathcal{L}^\bullet, \mathcal{I}^\bullet))
=
\SheafHom^\bullet(
\text{Tot}(\mathcal{K}^\bullet \otimes_{\mathcal{O}_X} \mathcal{L}^\bullet),
\mathcal{I}^\bullet)
$$
が補題 \ref{cohomology-lemma-compose} により成り立つ。左辺は
$R\SheafHom(K, R\SheafHom(L, M))$ を表現し（補題
\ref{cohomology-lemma-RHom-from-K-flat-into-K-injective} を用いる），右辺は
$R\SheafHom(K \otimes_{\mathcal{O}_X}^\mathbf{L} L, M)$ を表現することに
注意する。これで補題の表示式が証明された。大域切断を取り，補題
\ref{cohomology-lemma-section-RHom-over-U} を用いると
(\ref{cohomology-equation-internal-hom}) を得る。
\end{proof}

\begin{lemma}
\label{cohomology-lemma-restriction-RHom-to-U}
$(X, \mathcal{O}_X)$ を環付き空間とし，$K,L$ を
$D(\mathcal{O}_X)$ の対象とする。$R\SheafHom(K, L)$ の構成は
開集合への制限と可換する。すなわち，すべての開集合 $U$ に対し
$R\SheafHom(K|_U, L|_U) = R\SheafHom(K, L)|_U$ である。
\end{lemma}

\begin{proof}
構成と補題 \ref{cohomology-lemma-restrict-K-injective-to-open} から明らかである。
\end{proof}

\begin{lemma}
\label{cohomology-lemma-RHom-triangulated}
$(X, \mathcal{O}_X)$ を環付き空間とする。二変数関手
$R\SheafHom(- , -)$ は，どちらの変数についても完全三角を完全三角に写す。
\end{lemma}

\begin{proof}
これは，対応
$$
(\mathcal{L}^\bullet, \mathcal{M}^\bullet) \longmapsto
\SheafHom^\bullet(\mathcal{L}^\bullet, \mathcal{M}^\bullet)
$$
が，いずれかの変数における複体の各次数で分裂する短完全列を，
各次数で分裂する短完全列に写すことから従う。詳細は省略する。
\end{proof}

\begin{lemma}
\label{cohomology-lemma-internal-hom-composition}
$(X, \mathcal{O}_X)$ を環付き空間とする。
$D(\mathcal{O}_X)$ の $K,L,M$ が与えられると，
$K,L,M$ に関して関手的な $D(\mathcal{O}_X)$ における標準射
$$
R\SheafHom(L, M) \otimes_{\mathcal{O}_X}^\mathbf{L} R\SheafHom(K, L)
\longrightarrow R\SheafHom(K, M)
$$
が存在する。
\end{lemma}

\begin{proof}
$M$ を表現する K-入射的複体 $\mathcal{I}^\bullet$，
$L$ を表現する K-入射的複体 $\mathcal{J}^\bullet$，および
$K$ を表現する任意の $\mathcal{O}_X$-加群の複体
$\mathcal{K}^\bullet$ を選ぶ。補題 \ref{cohomology-lemma-composition} により
複体の写像
$$
\text{Tot}\left(
\SheafHom^\bullet(\mathcal{J}^\bullet, \mathcal{I}^\bullet)
\otimes_{\mathcal{O}_X}
\SheafHom^\bullet(\mathcal{K}^\bullet, \mathcal{J}^\bullet)
\right)
\longrightarrow
\SheafHom^\bullet(\mathcal{K}^\bullet, \mathcal{I}^\bullet)
$$
が存在する。$\mathcal{O}_X$-加群の複体
$\SheafHom^\bullet(\mathcal{J}^\bullet, \mathcal{I}^\bullet)$,
$\SheafHom^\bullet(\mathcal{K}^\bullet, \mathcal{J}^\bullet)$，および
$\SheafHom^\bullet(\mathcal{K}^\bullet, \mathcal{I}^\bullet)$
はそれぞれ $R\SheafHom(L, M)$，$R\SheafHom(K, L)$，
$R\SheafHom(K, M)$ を表現する。K-平坦複体
$\mathcal{H}^\bullet$ と擬同型
$\mathcal{H}^\bullet \to
\SheafHom^\bullet(\mathcal{K}^\bullet, \mathcal{J}^\bullet)$,
を選ぶと，写像
$$
\text{Tot}\left(
\SheafHom^\bullet(\mathcal{J}^\bullet, \mathcal{I}^\bullet)
\otimes_{\mathcal{O}_X} \mathcal{H}^\bullet
\right)
\longrightarrow
\text{Tot}\left(
\SheafHom^\bullet(\mathcal{J}^\bullet, \mathcal{I}^\bullet)
\otimes_{\mathcal{O}_X}
\SheafHom^\bullet(\mathcal{K}^\bullet, \mathcal{J}^\bullet)
\right)
$$
が存在し，その始域は
$R\SheafHom(L, M) \otimes_{\mathcal{O}_X}^\mathbf{L} R\SheafHom(K, L)$
を表現する。表示した二つの矢印を合成すると望む写像を得る。
この構成が関手的であることの証明は省略する。
\end{proof}

\begin{lemma}
\label{cohomology-lemma-internal-hom-diagonal-better}
$(X, \mathcal{O}_X)$ を環付き空間とする。
$D(\mathcal{O}_X)$ の $K,L,M$ が与えられると，
$K,L,M$ に関して関手的な $D(\mathcal{O}_X)$ における標準射
$$
K \otimes_{\mathcal{O}_X}^\mathbf{L} R\SheafHom(M, L)
\longrightarrow
R\SheafHom(M, K \otimes_{\mathcal{O}_X}^\mathbf{L} L)
$$
が存在する。
\end{lemma}

\begin{proof}
$K$ を表現する K-平坦複体 $\mathcal{K}^\bullet$ と，
$L$ を表現する K-入射的複体 $\mathcal{I}^\bullet$ を選び，
$M$ を表現する任意の $\mathcal{O}_X$-加群の複体
$\mathcal{M}^\bullet$ を選ぶ。$\mathcal{J}^\bullet$ が K-入射的である
ような擬同型
$\text{Tot}(\mathcal{K}^\bullet \otimes_{\mathcal{O}_X} \mathcal{I}^\bullet)
\to \mathcal{J}^\bullet$
を選ぶ。このとき写像
$$
\text{Tot}\left(
\mathcal{K}^\bullet \otimes_{\mathcal{O}_X}
\SheafHom^\bullet(\mathcal{M}^\bullet, \mathcal{I}^\bullet)
\right)
\to
\SheafHom^\bullet(\mathcal{M}^\bullet,
\text{Tot}(\mathcal{K}^\bullet \otimes_{\mathcal{O}_X} \mathcal{I}^\bullet))
\to
\SheafHom^\bullet(\mathcal{M}^\bullet, \mathcal{J}^\bullet)
$$
を用いる。ここで第一の写像は補題 \ref{cohomology-lemma-diagonal-better}
の写像である。
\end{proof}

\begin{lemma}
\label{cohomology-lemma-internal-hom-diagonal}
$(X, \mathcal{O}_X)$ を環付き空間とする。
$D(\mathcal{O}_X)$ の $K,L$ が与えられると，$K,L$ の両方に関して
関手的な $D(\mathcal{O}_X)$ における標準射
$$
K \longrightarrow R\SheafHom(L, K \otimes_{\mathcal{O}_X}^\mathbf{L} L)
$$
が存在する。
\end{lemma}

\begin{proof}
$K$ を表現する K-平坦複体 $\mathcal{K}^\bullet$ と，
$L$ を表現する任意の $\mathcal{O}_X$-加群の複体
$\mathcal{L}^\bullet$ を選ぶ。K-入射的複体
$\mathcal{J}^\bullet$ と擬同型
$\text{Tot}(\mathcal{K}^\bullet \otimes_{\mathcal{O}_X} \mathcal{L}^\bullet)
\to \mathcal{J}^\bullet$ を選ぶ。このとき
$$
\mathcal{K}^\bullet \to
\SheafHom^\bullet(\mathcal{L}^\bullet,
\text{Tot}(\mathcal{K}^\bullet \otimes_{\mathcal{O}_X} \mathcal{L}^\bullet))
\to
\SheafHom^\bullet(\mathcal{L}^\bullet, \mathcal{J}^\bullet)
$$
を用いる。ここで第一の写像は補題 \ref{cohomology-lemma-diagonal} から生じる。
\end{proof}

\begin{lemma}
\label{cohomology-lemma-dual}
$(X, \mathcal{O}_X)$ を環付き空間とし，$L$ を
$D(\mathcal{O}_X)$ の対象とする。
$L^\vee = R\SheafHom(L, \mathcal{O}_X)$ とおく。
$D(\mathcal{O}_X)$ の $M$ に対して標準写像
\begin{equation}
\label{cohomology-equation-eval}
M \otimes^\mathbf{L}_{\mathcal{O}_X} L^\vee
\longrightarrow
R\SheafHom(L, M)
\end{equation}
が存在し，これは標準写像
$$
H^0(X, M \otimes^\mathbf{L}_{\mathcal{O}_X} L^\vee)
\longrightarrow
\Hom_{D(\mathcal{O}_X)}(L, M)
$$
を誘導する。この写像は $D(\mathcal{O}_X)$ の $M$ に関して関手的である。
\end{lemma}

\begin{proof}
写像 (\ref{cohomology-equation-eval}) は，同一視
$M = R\SheafHom(\mathcal{O}_X, M)$ を用いた補題
\ref{cohomology-lemma-internal-hom-composition} の特別な場合である。
\end{proof}

\begin{lemma}
\label{cohomology-lemma-internal-hom-evaluate}
$(X, \mathcal{O}_X)$ を環付き空間とし，$K,L,M$ を
$D(\mathcal{O}_X)$ の対象とする。$K,L,M$ に関して関手的な
$D(\mathcal{O}_X)$ における標準射
$$
R\SheafHom(L, M) \otimes_{\mathcal{O}_X}^\mathbf{L} K
\longrightarrow
R\SheafHom(R\SheafHom(K, L), M)
$$
が存在する。
\end{lemma}

\begin{proof}
$M$ を表現する K-入射的複体 $\mathcal{I}^\bullet$，
$L$ を表現する K-入射的複体 $\mathcal{J}^\bullet$，および
$K$ を表現する K-平坦複体 $\mathcal{K}^\bullet$ を選ぶ。
写像は，補題 \ref{cohomology-lemma-evaluate-and-more} の写像
$$
\text{Tot}(\SheafHom^\bullet(\mathcal{J}^\bullet,
\mathcal{I}^\bullet) \otimes_{\mathcal{O}_X} \mathcal{K}^\bullet)
\longrightarrow
\SheafHom^\bullet(\SheafHom^\bullet(\mathcal{K}^\bullet,
\mathcal{J}^\bullet), \mathcal{I}^\bullet)
$$
を用いて定義される。ここで選んだ複体により，左辺は
$R\SheafHom(L, M) \otimes_{\mathcal{O}_X}^\mathbf{L} K$
を表現し，右辺は $R\SheafHom(R\SheafHom(K, L), M)$ を表現する。
これが $D(\mathcal{O}_X)$ の三つの対象すべてに関して関手的であることの
証明は省略する。
\end{proof}

\begin{remark}
\label{cohomology-remark-tensor-internal-hom}
$(X, \mathcal{O}_X)$ を環付き空間とする。
$D(\mathcal{O}_X)$ の $K,K',M,M'$ に対して標準写像
$$
R\SheafHom(K, K') \otimes_{\mathcal{O}_X}^\mathbf{L}
R\SheafHom(M, M')
\longrightarrow
R\SheafHom(K \otimes_{\mathcal{O}_X}^\mathbf{L} M,
K' \otimes_{\mathcal{O}_X}^\mathbf{L} M')
$$
が存在する。実際，(\ref{cohomology-equation-internal-hom}) により，これは写像
$$
R\SheafHom(K, K') \otimes_{\mathcal{O}_X}^\mathbf{L}
R\SheafHom(M, M') \otimes_{\mathcal{O}_X}^\mathbf{L}
K \otimes_{\mathcal{O}_X}^\mathbf{L} M
\longrightarrow
K' \otimes_{\mathcal{O}_X}^\mathbf{L} M'
$$
と同じものである。このため，まず中央の二因子を入れ替え
（符号規則は More on Algebra, Section
\ref{more-algebra-section-sign-rules} のとおり），写像
$$
R\SheafHom(K, K') \otimes_{\mathcal{O}_X}^\mathbf{L} K \to K'
\quad\text{and}\quad
R\SheafHom(M, M') \otimes_{\mathcal{O}_X}^\mathbf{L} M \to M'
$$
を用いればよい。これらは $K = R\SheafHom(\mathcal{O}_X, K)$ とみなし，
$K',M,M'$ についても同様にみなしたときの補題
\ref{cohomology-lemma-internal-hom-composition} から得られる。
\end{remark}

\begin{remark}
\label{cohomology-remark-projection-formula-for-internal-hom}
$f : X \to Y$ を環付き空間の射とし，$K,L$ を
$D(\mathcal{O}_X)$ の対象とする。標準写像
$$
Rf_*R\SheafHom(L, K) \longrightarrow R\SheafHom(Rf_*L, Rf_*K)
$$
が存在すると主張する。実際，(\ref{cohomology-equation-internal-hom}) により，
これは写像
$Rf_*R\SheafHom(L, K) \otimes_{\mathcal{O}_Y}^\mathbf{L} Rf_*L \to Rf_*K$.
と同じものである。このために合成
$$
Rf_*R\SheafHom(L, K) \otimes_{\mathcal{O}_Y}^\mathbf{L} Rf_*L \to
Rf_*(R\SheafHom(L, K) \otimes_{\mathcal{O}_X}^\mathbf{L} L) \to
Rf_*K
$$
を用いる。ここで第一の矢印は相対カップ積
（注記 \ref{cohomology-remark-cup-product}）であり，第二の矢印は標準写像
$R\SheafHom(L, K) \otimes_{\mathcal{O}_X}^\mathbf{L} L \to K$
に $Rf_*$ を適用したものである。この標準写像は，一つの位置に
$\mathcal{O}_X$ を入れた補題
\ref{cohomology-lemma-internal-hom-composition} から生じる。
\end{remark}

\begin{remark}
\label{cohomology-remark-relative-cup-and-composition}
$h : X \to Y$ を環付き空間の射とし，$K,M$ を
$D(\mathcal{O}_Y)$ の対象とする。図式
$$
\xymatrix{
Rf_*R\SheafHom_{\mathcal{O}_X}(K, M)
\otimes_{\mathcal{O}_Y}^\mathbf{L} Rf_*K
\ar[r] \ar[d] &
Rf_*\left(R\SheafHom_{\mathcal{O}_X}(K, M)
\otimes_{\mathcal{O}_X}^\mathbf{L} K\right)
\ar[d] \\
R\SheafHom_{\mathcal{O}_Y}(Rf_*K, Rf_*M) \otimes_{\mathcal{O}_Y}^\mathbf{L}
Rf_*K \ar[r] &
Rf_*M
}
$$
は可換である。ここで左の垂直矢印は注記
\ref{cohomology-remark-projection-formula-for-internal-hom} から生じる。
上の水平矢印は注記 \ref{cohomology-remark-cup-product} のものである。
残る二つの矢印は，いずれか一つの項を
$\mathcal{O}_X$ または $\mathcal{O}_Y$ で置き換えた補題
\ref{cohomology-lemma-internal-hom-composition} の写像の例である。
\end{remark}

\begin{remark}
\label{cohomology-remark-prepare-fancy-base-change}
$h : X \to Y$ を環付き空間の射とし，$K,L$ を
$D(\mathcal{O}_Y)$ の対象とする。$D(\mathcal{O}_X)$ における標準写像
$$
Lh^*R\SheafHom(K, L) \longrightarrow R\SheafHom(Lh^*K, Lh^*L)
$$
が存在すると主張する。実際，(\ref{cohomology-equation-internal-hom})
（補題 \ref{cohomology-lemma-internal-hom} で証明した）により，
このような写像は写像
$$
Lh^*R\SheafHom(K, L) \otimes^\mathbf{L} Lh^*K \longrightarrow Lh^*L
$$
と同じものである。この矢印の始域は補題
\ref{cohomology-lemma-pullback-tensor-product} により
$Lh^*(\SheafHom(K, L) \otimes^\mathbf{L} K)$ だから，
標準写像
$$
R\SheafHom(K, L) \otimes^\mathbf{L} K \longrightarrow L.
$$
を構成すれば十分である。このために，(\ref{cohomology-equation-internal-hom})
を介して
$$
\text{id} :
R\SheafHom(K, L)
\longrightarrow
R\SheafHom(K, L)
$$
に対応する矢印を取る。
\end{remark}

\begin{remark}
\label{cohomology-remark-fancy-base-change}
環付き空間の可換図式
$$
\xymatrix{
X' \ar[r]_h \ar[d]_{f'} &
X \ar[d]^f \\
S' \ar[r]^g &
S
}
$$
が与えられたとする。$K,L$ を $D(\mathcal{O}_X)$ の対象とする。
$D(\mathcal{O}_{S'})$ における標準基底変換写像
$$
Lg^*Rf_*R\SheafHom(K, L)
\longrightarrow
R(f')_*R\SheafHom(Lh^*K, Lh^*L)
$$
が存在すると主張する。実際，合成
\begin{align*}
L(f')^*Lg^*Rf_*R\SheafHom(K, L)
& =
Lh^*Lf^*Rf_*R\SheafHom(K, L) \\
& \to
Lh^*R\SheafHom(K, L) \\
& \to
R\SheafHom(Lh^*K, Lh^*L)
\end{align*}
に随伴する写像を取る。ここで第一の矢印は随伴写像
$Lf^*Rf_* \to \text{id}$ を用い，第二の矢印は注記
\ref{cohomology-remark-prepare-fancy-base-change} で構成した標準写像である。
\end{remark}





\section{Ext 層}
\label{cohomology-section-ext}

\noindent
$(X, \mathcal{O}_X)$ を環付き空間とし，
$K, L \in D(\mathcal{O}_X)$ とする。導来圏における内部 Hom の構成を用い，
$D(\mathcal{O}_X)$ の対象 $R\SheafHom(K, L)$ の第 $n$ コホモロジー層を
取ることによって，良定義な $\mathcal{O}_X$-加群の層
$$
\SheafExt^n(K, L) = H^n(R\SheafHom(K, L))
$$
を得る。この対象を $\SheafExt^n_{\mathcal{O}_X}(K, L)$ と書くこともある。
補題 \ref{cohomology-lemma-section-RHom-over-U} により，この
$\SheafExt^n$-層は規則
$$
U \longmapsto \Ext^n_{D(\mathcal{O}_U)}(K|_U, L|_U)
$$
の層化である。例 \ref{cohomology-example-spectral-sequence} により，
常にスペクトル系列
$$
E_2^{p, q} = H^p(X, \SheafExt^q(K, L))
$$
が存在し，適切な条件下では
$\Ext^{p + q}_{D(\mathcal{O}_X)}(K, L)$ に収束する
（例えば $L$ が下に有界で $K$ が上に有界な場合）。





\section{大域導来 Hom}
\label{cohomology-section-global-RHom}

\noindent
$(X, \mathcal{O}_X)$ を環付き空間とし，
$K, L \in D(\mathcal{O}_X)$ とする。導来圏における内部 Hom の構成を
用いると，$D(\Gamma(X, \mathcal{O}_X))$ の良定義な対象
$$
R\Hom_X(K, L) = R\Gamma(X, R\SheafHom(K, L))
$$
を得る。この対象を $R\Hom_{\mathcal{O}_X}(K, L)$ と書くこともある。
補題 \ref{cohomology-lemma-section-RHom-over-U} により
$$
H^0(R\Hom_X(K, L)) = \Hom_{D(\mathcal{O}_X)}(K, L),
\quad
H^p(R\Hom_X(K, L)) = \Ext_{D(\mathcal{O}_X)}^p(K, L)
$$
である。$f : Y \to X$ が環付き空間の射ならば，注記
\ref{cohomology-remark-prepare-fancy-base-change} で定義した写像の大域切断を
取ることにより，$D(\Gamma(X, \mathcal{O}_X))$ における標準写像
$$
R\Hom_X(K, L) \longrightarrow R\Hom_Y(Lf^*K, Lf^*L)
$$
が存在する。








\section{複体の貼り合わせ}
\label{cohomology-section-glueing-complexes}

\noindent
複体は貼り合わせられる！ より正確には，ある状況では，局所的に与えられた
導来圏の対象を大域的な対象へ貼り合わせることができる。
まずいくつかの容易な場合を証明し，次に位相空間と開被覆の設定で，非常に一般的な
\cite[Theorem 3.2.4]{BBD} を証明する。

\begin{lemma}
\label{cohomology-lemma-glue}
$(X, \mathcal{O}_X)$ を環付き空間とし，$X = U \cup V$ を
$X$ の二つの開部分空間の和とする。次が与えられたとする。
\begin{enumerate}
\item $D(\mathcal{O}_U)$ の対象 $A$。
\item $D(\mathcal{O}_V)$ の対象 $B$。
\item 同型 $c : A|_{U \cap V} \to B|_{U \cap V}$。
\end{enumerate}
このとき $D(\mathcal{O}_X)$ の対象 $F$ と，同型
$f : F|_U \to A$，$g : F|_V \to B$ で
$c = g|_{U \cap V} \circ f^{-1}|_{U \cap V}$ を満たすものが存在する。
さらに，次が与えられたとする。
\begin{enumerate}
\item $D(\mathcal{O}_X)$ の対象 $E$。
\item $D(\mathcal{O}_U)$ の射 $a : A \to E|_U$。
\item $D(\mathcal{O}_V)$ の射 $b : B \to E|_V$。
\end{enumerate}
これらが
$$
a|_{U \cap V}  = b|_{U \cap V} \circ c.
$$
を満たすならば，$D(\mathcal{O}_X)$ に射 $F \to E$ が存在し，
その $U$ への制限は $a \circ f$，$V$ への制限は $b \circ g$ である。
\end{lemma}

\begin{proof}
$j_U$, $j_V$, $j_{U \cap V}$ を対応する開埋め込みとする。
完全三角
$$
F \to Rj_{U, *}A \oplus Rj_{V, *}B \to Rj_{U \cap V, *}(B|_{U \cap V})
\to F[1]
$$
を選ぶ。ここで写像
$Rj_{V, *}B \to Rj_{U \cap V, *}(B|_{U \cap V})$ は明らかなものであり，
$Rj_{U, *}A \to Rj_{U \cap V, *}(B|_{U \cap V})$
は
$Rj_{U, *}A \to Rj_{U \cap V, *}(A|_{U \cap V})$
と $Rj_{U \cap V, *}c$ の合成である。$U$ に制限すると
$$
F|_U \to A \oplus (Rj_{V, *}B)|_U \to (Rj_{U \cap V, *}(B|_{U \cap V}))|_U
\to F|_U[1]
$$
を得る。$j : U \cap V \to U$ とおく。開集合への制限とコホモロジーの
両立性により，
$(Rj_{V, *}B)|_U$ および $(Rj_{U \cap V, *}(B|_{U \cap V}))|_U$
はいずれも $Rj_*(B|_{U \cap V})$ と標準同型である。
したがって最後に表示した図式の第二の矢印は切断をもち，射
$F|_U \to A$ が同型であると結論する。同様に射
$F|_V \to B$ も同型である。射 $F \to E$ の存在は
$\Hom$ に対する Mayer--Vietoris 系列から従う。補題
\ref{cohomology-lemma-mayer-vietoris-hom} を参照せよ。
\end{proof}

\begin{lemma}
\label{cohomology-lemma-vanishing-and-glueing}
$f : (X, \mathcal{O}_X) \to (Y, \mathcal{O}_Y)$ を環付き空間の射とし，
$\mathcal{B}$ を $Y$ の位相の基底とする。
\begin{enumerate}
\item $K \in D(\mathcal{O}_X)$ であり，$V \in \mathcal{B}$ に対し
$i < 0$ ならば $H^i(f^{-1}(V), K) = 0$ と仮定する。
このとき $Rf_*K$ の負の次数のコホモロジー層は消滅し，
すべての開集合 $V \subset Y$ と $i < 0$ に対し
$H^i(f^{-1}(V), K) = 0$ であり，規則
$V \mapsto H^0(f^{-1}V, K)$ は $Y$ 上の層である。
\item $K,L \in D(\mathcal{O}_X)$ であり，$V \in \mathcal{B}$ に対し
$i < 0$ ならば
$\Ext^i(K|_{f^{-1}V}, L|_{f^{-1}V}) = 0$ と仮定する。
このとき，すべての開集合 $V \subset Y$ と $i < 0$ に対し
$\Ext^i(K|_{f^{-1}V}, L|_{f^{-1}V}) = 0$ であり，規則
$V \mapsto \Hom(K|_{f^{-1}V}, L|_{f^{-1}V})$ は $Y$ 上の層である。
\end{enumerate}
\end{lemma}

\begin{proof}
補題 \ref{cohomology-lemma-unbounded-describe-higher-direct-images} により，
$H^i(Rf_*K)$ は前層
$V \mapsto H^i(f^{-1}(V), K) = H^i(V, Rf_*K)$.
に付随する層である。(1) の仮定から，$Rf_*K$ の次数 $<0$ の
コホモロジー層は消滅する。したがって任意の開集合 $V \subset Y$ に対し，
コホモロジー群 $H^i(V, Rf_*K)$ は $i < 0$ ならば零であり，
$i = 0$ ならば $H^0(V, H^0(Rf_*K))$ に等しい。これで (1) が証明された。

\medskip\noindent
(2) を証明するには，補題 \ref{cohomology-lemma-section-RHom-over-U} を用いて
読み替え，複体 $R\SheafHom(K, L)$ に (1) を適用する。
\end{proof}

\begin{situation}
\label{cohomology-situation-locally-given}
$(X, \mathcal{O}_X)$ を環付き空間とする。次が与えられている。
\begin{enumerate}
\item $X$ の開集合の族 $\mathcal{B}$。
\item $U \in \mathcal{B}$ に対し，$D(\mathcal{O}_U)$ の対象 $K_U$。
\item $V \subset U$，$V,U \in \mathcal{B}$ に対し，
$D(\mathcal{O}_V)$ における同型
$\rho^U_V : K_U|_V \to K_V$。
\end{enumerate}
これらは，$W \subset V \subset U$ かつ
$U,V,W \in \mathcal{B}$ ならば
$\rho^U_W = \rho^V_W \circ \rho ^U_V|_W$ を満たす。
\end{situation}

\noindent
これ以上の仮定なしには，この状況について何も証明できない。
興味深い一つの場合は，$\mathcal{B}$ が $X$ の位相の基底である場合である。
もう一つは，位相空間の射 $f : X \to Y$ があり，
$\mathcal{B}$ の要素が $Y$ の位相のある基底の要素の逆像である場合である。

\medskip\noindent
状況 \ref{cohomology-situation-locally-given} における{\it 解}とは，
$K$ が $D(\mathcal{O}_X)$ の対象であり，
$\rho_U : K|_U \to K_U$（$U \in \mathcal{B}$）が同型であって，
すべての $V \subset U$，$U,V \in \mathcal{B}$ に対し
$\rho^U_V \circ \rho_U|_V = \rho_V$ を満たすような対
$(K,\rho_U)$ をいう。ある場合には解は一意である。

\begin{lemma}
\label{cohomology-lemma-uniqueness}
状況 \ref{cohomology-situation-locally-given} において，次を仮定する。
\begin{enumerate}
\item $X = \bigcup_{U \in \mathcal{B}} U$ であり，
すべての $U,V \in \mathcal{B}$ に対し
$U \cap V = \bigcup_{W \in \mathcal{B}, W \subset U \cap V} W$,
\item 任意の $U \in \mathcal{B}$ に対し，$i < 0$ ならば
$\Ext^i(K_U, K_U) = 0$ である。
\end{enumerate}
解 $(K,\rho_U)$ が存在するならば，それは一意な同型を除いて一意であり，
さらに $i < 0$ に対し $\Ext^i(K, K) = 0$ である。
\end{lemma}

\begin{proof}
$(K,\rho_U)$ と $(K',\rho'_U)$ を二つの解とする。
$f : X \to Y$ を Topology, Lemma
\ref{topology-lemma-create-map-from-subcollection} で構成された連続写像とし，
$\mathcal{O}_Y = f_*\mathcal{O}_X$ とおく。このとき
$K,K'$ と $\mathcal{B}$ は補題
\ref{cohomology-lemma-vanishing-and-glueing} の (2) の状況にある。
したがって $K$ の負の Ext の消滅を得，規則
$$
V \longmapsto \Hom(K|_{f^{-1}V}, K'|_{f^{-1}V})
$$
が $Y$ 上の層であることが分かる。$(K,\rho_U)$ と $(K',\rho'_U)$ は
いずれも解なので，$U = f^{-1}(f(U))$ 上の写像
$$
(\rho'_U)^{-1} \circ \rho_U : K|_U \longrightarrow K'|_U
$$
は重なり上で一致する。ゆえに上の層の一意な大域切断が得られ，
これは与えられたすべての構造と両立する望ましい同型
$K \to K'$ を定める。
\end{proof}

\begin{remark}
\label{cohomology-remark-uniqueness}
補題 \ref{cohomology-lemma-uniqueness} の記法と仮定を用いる。
$U,V \in \mathcal{B}$ とする。$\mathcal{B}'$ を，
$U \cap V$ に含まれる $\mathcal{B}$ の要素からなる集合とする。このとき
$$
(\{K_{U'}\}_{U' \in \mathcal{B}'},
\{\rho_{V'}^{U'}\}_{V' \subset U'\text{ with }U', V' \in \mathcal{B}'})
$$
は環付き空間 $U \cap V$ 上の系で，補題
\ref{cohomology-lemma-uniqueness} の仮定を満たす。さらに，
$(K_U|_{U \cap V}, \rho^U_{U'})$ と
$(K_V|_{U \cap V}, \rho^V_{U'})$ はいずれもこの系の解である。
補題により，一意な同型
$$
\rho_{U, V} : K_U|_{U \cap V} \longrightarrow K_V|_{U \cap V}
$$
で，すべての $U' \subset U \cap V$，$U' \in \mathcal{B}$ に対して図式
$$
\xymatrix{
K_U|_{U'} \ar[rr]_{\rho_{U, V}|_{U'}} \ar[rd]_{\rho^U_{U'}} & &
K_V|_{U'} \ar[ld]^{\rho^V_{U'}} \\
& K_{U'}
}
$$
を可換にするものが得られる。第三の要素 $W \in \mathcal{B}$ を取る。
同様の性質を満たす同型
$\rho_{U, W} : K_U|_{U \cap W} \to K_W|_{U \cap W}$ および
$\rho_{V, W} : K_U|_{V \cap W} \to K_W|_{V \cap W}$
を得る。最後に
$$
\rho_{U, W}|_{U \cap V \cap W} =
\rho_{V, W}|_{U \cap V \cap W} \circ \rho_{U, V}|_{U \cap V \cap W}
$$
が成り立つ。等式の両辺は，$U'' \subset U \cap V \cap W$，
$U'' \in \mathcal{B}$ に対して写像 $\rho^U_{U''}$ と
$\rho^W_{U''}$ に両立する一意な同型だから，これは補題の一意性から従う。
いくつかの細部は省略する。族 $(K_U,\rho_{U,V})$ は，
$X$ の開被覆 $\mathcal{U} : X = \bigcup_{U \in \mathcal{B}} U$
に関する導来圏の降下データである。この用語では，解の存在を求めることは
「降下データの有効性」を求めることである。
\end{remark}

\begin{lemma}
\label{cohomology-lemma-solution-in-finite-case}
状況 \ref{cohomology-situation-locally-given} において，次を仮定する。
\begin{enumerate}
\item $X = U_1 \cup \ldots \cup U_n$，ただし $U_i \in \mathcal{B}$，
\item $U,V \in \mathcal{B}$ に対し
$U \cap V = \bigcup_{W \in \mathcal{B}, W \subset U \cap V} W$,
\item 任意の $U \in \mathcal{B}$ に対し，$i < 0$ ならば
$\Ext^i(K_U, K_U) = 0$ である。
\end{enumerate}
このとき解が存在し，一意な同型を除いて一意である。
\end{lemma}

\begin{proof}
一意性は補題 \ref{cohomology-lemma-uniqueness} で見た。補題を $n$ に関する
帰納法で証明する。$n = 1$ の場合は直ちに従う。

\medskip\noindent
$n = 2$ の場合。
注記 \ref{cohomology-remark-uniqueness} で構成した同型
$\rho_{U_1, U_2} : K_{U_1}|_{U_1 \cap U_2} \to K_{U_2}|_{U_1 \cap U_2}$
を考える。補題 \ref{cohomology-lemma-glue} により，$D(\mathcal{O}_X)$ の対象
$K$ と，$\rho_{U_1,U_2}$ に両立する同型
$\rho_{U_1} : K|_{U_1} \to K_{U_1}$ および
$\rho_{U_2} : K|_{U_2} \to K_{U_2}$ を得る。
$U \in \mathcal{B}$ を取る。同型 $\rho_U : K|_U \to K_U$ を構成する。
$(K,\rho_U)$ が解であることの確認は読者に委ねる。
$\mathcal{B}'$ を，$U \cap U_1$ または $U \cap U_2$ に含まれる
$\mathcal{B}$ の要素からなる集合とする。このとき
$(K_U,\rho^U_{U'})$ は環付き空間 $U$ 上の系
$(\{K_{U'}\}_{U' \in \mathcal{B}'},
\{\rho_{V'}^{U'}\}_{V' \subset U'\text{ with }U', V' \in \mathcal{B}'})$
の解である。$(K|_U,\tau_{U'})$ も解であると主張する。ここで
$U' \in \mathcal{B}'$ に対する $\tau_{U'}$ は次のように選ぶ。
$U' \subset U_1$ ならば合成
$$
K|_{U'} \xrightarrow{\rho_{U_1}|_{U'}}
K_{U_1}|_{U'} \xrightarrow{\rho^{U_1}_{U'}}
K_{U'}
$$
を取り，$U' \subset U_2$ ならば合成
$$
K|_{U'} \xrightarrow{\rho_{U_2}|_{U'}}
K_{U_2}|_{U'} \xrightarrow{\rho^{U_2}_{U'}}
K_{U'}.
$$
を取る。これが解であることを確認するには，注記
\ref{cohomology-remark-uniqueness} で述べた写像 $\rho_{U_1,U_2}$ の性質と，
$\rho_{U_1},\rho_{U_2}$ が $\rho_{U_1,U_2}$ と両立することを用いる。
以上の下で補題 \ref{cohomology-lemma-uniqueness} を適用すると，
$U' \in \mathcal{B}'$ に対して写像 $\tau_{U'}$ と
$\rho^U_{U'}$ に両立する一意な同型
$K|_{U'} \to K_{U'}$ を得る。

\medskip\noindent
$n > 2$ の場合。開部分空間
$X' = U_1 \cup \ldots \cup U_{n - 1}$ を考え，
$\mathcal{B}'$ を $X'$ に含まれる $\mathcal{B}$ の要素からなる集合とする。
すると環付き空間 $X'$ 上に系
$(\{K_U\}_{U \in \mathcal{B}'}, \{\rho_V^U\}_{U, V \in \mathcal{B}'})$
を得，これに帰納法の仮定を適用できる。解
$(K_{X'},\rho^{X'}_U)$ を得る。次に $X$ の開集合の族
$\mathcal{B}^* = \mathcal{B} \cup \{X'\}$ を考えると，系
$(\{K_U\}_{U \in \mathcal{B}^*},
\{\rho_V^U\}_{V \subset U\text{ with }U, V \in \mathcal{B}^*})$.
を得る。解 $K_{X'}$ に補題 \ref{cohomology-lemma-uniqueness} を適用すると，
この新しい系も条件 (3) を満たすことに注意する。この系について
$X = X' \cup U_n$ である。これで，上で扱った $n = 2$ の場合に帰着する。
\end{proof}

\begin{lemma}
\label{cohomology-lemma-glueing-increasing-union}
$X$ を環付き空間，$E$ を整列集合とし，
$$
X = \bigcup\nolimits_{\alpha \in E} W_\alpha
$$
を開被覆とする。ここで $W_\alpha \subset W_{\alpha+1}$ であり，
$\alpha$ が後続元でなければ
$W_\alpha = \bigcup_{\beta < \alpha} W_\beta$ とする。
$K_\alpha$ を $D(\mathcal{O}_{W_\alpha})$ の対象で，
$i < 0$ に対し $\Ext^i(K_\alpha,K_\alpha)=0$ を満たすものとする。
すべての $\beta < \alpha$ に対し，$D(\mathcal{O}_{W_\beta})$
における同型
$\rho_\beta^\alpha :  K_\alpha|_{W_\beta} \to K_\beta$
が与えられ，
$\rho_\gamma^\alpha = \rho_\gamma^\beta \circ \rho^\alpha_\beta|_{W_\gamma}$
が $\gamma < \beta < \alpha$ に対して成り立つと仮定する。
このとき $D(\mathcal{O}_X)$ の対象 $K$ と同型
$K|_{W_\alpha} \to K_\alpha$（$\alpha \in E$）
が存在し，これらは同型 $\rho_\beta^\alpha$ と両立する。
\end{lemma}

\begin{proof}
この証明では $\alpha,\beta,\gamma,\ldots$ は $E$ の要素を表す。
$K_\alpha$ を表現する $W_\alpha$ 上の K-入射的複体
$I_\alpha^\bullet$ を選ぶ。$\beta < \alpha$ に対し，
$j_{\beta,\alpha}:W_\beta\to W_\alpha$ を包含射とする。
超限再帰を用いて，すべての $\beta<\alpha$ に対し複体の写像
$$
\tau_{\beta, \alpha} :
(j_{\beta, \alpha})_!I_\beta^\bullet
\longrightarrow
I_\alpha^\bullet
$$
で，同型
$\rho^\alpha_\beta : K_\alpha|_{W_\beta} \to K_\beta$.
の逆の随伴を表現するものを構成する。さらに，
$\gamma<\beta<\alpha$ に対して複体の写像として
$$
\tau_{\gamma, \alpha} = \tau_{\beta, \alpha} \circ
(j_{\beta, \alpha})_!\tau_{\gamma, \beta}
$$
となるように構成する。実際，すべての $\gamma<\beta<\alpha$ に対し，
正しく合成する $\tau_{\gamma,\beta}$ がすでに与えられていると仮定する。
$\alpha=\alpha'+1$ が後続元ならば，複体の任意の写像
$$
(j_{\alpha', \alpha})_!I_{\alpha'}^\bullet \to I_\alpha^\bullet
$$
で，同型
$\rho^\alpha_{\alpha'} : K_\alpha|_{W_{\alpha'}} \to K_{\alpha'}$
の逆に随伴するものを選ぶ（$I_\alpha^\bullet$ は K-入射的なので可能である）。
任意の $\beta<\alpha'$ に対し
$$
\tau_{\beta, \alpha} = \tau_{\alpha', \alpha} \circ
(j_{\alpha', \alpha})_!\tau_{\beta, \alpha'}
$$
とおく。$\alpha$ が後続元でなければ，$W_\alpha$ 上の複体
$$
C^\bullet = \colim_{\beta < \alpha} (j_{\beta, \alpha})_!I_\beta^\bullet
$$
を考えられる（各次数ごとの余極限）。ここで系の遷移写像は
$\beta'<\beta<\alpha$ に対する写像 $\tau_{\beta',\beta}$ で与えられる。
$C^\bullet$ は $K_\alpha$ を表現すると主張する。実際，
$\beta<\alpha$ に対して，余極限への標準写像
$(j_{\beta,\alpha})_!I_\beta^\bullet\to C^\bullet$ の制限は写像
$$
\sigma_\beta : I_\beta^\bullet \longrightarrow C^\bullet|_{W_\beta}
$$
を与え，これは擬同型である。実際，$x\in W_\beta$ として茎を見ると
$$
(C^\bullet)_x =
\colim_{\beta' < \alpha}
\left((j_{\beta', \alpha})_!I_{\beta'}^\bullet\right)_x =
\colim_{\beta \leq \beta' < \alpha} (I_{\beta'}^\bullet)_x
\longleftarrow
(I_\beta^\bullet)_x
$$
を得，これは擬同型である。ここでは，茎を取ることが余極限と可換すること，
フィルター余極限が完全であること，および
$\beta\leq\beta'<\alpha$ に対する写像
$(I_\beta^\bullet)_x\to(I_{\beta'}^\bullet)_x$ が擬同型であることを用いた。
したがって $(C^\bullet,\sigma_\beta^{-1})$ は系
$(\{K_\beta\}_{\beta<\alpha},
\{\rho^\beta_{\beta'}\}_{\beta' < \beta < \alpha})$.
の解である。$(K_\alpha,\rho^\alpha_\beta)$ も別の解なので，
すべての写像と両立する $D(\mathcal{O}_{W_\alpha})$ における一意な同型
$\sigma:K_\alpha\to C^\bullet$ を得る。補題
\ref{cohomology-lemma-solution-in-finite-case} を参照せよ
（ここで負の Ext 群の消滅を用いる）。
$\sigma$ を表現する複体の射
$\tau:C^\bullet\to I_\alpha^\bullet$ を選び，
$$
\tau_{\beta, \alpha} = \tau|_{W_\beta} \circ \sigma_\beta
$$
とおけば，望む写像を得る。最後に，$K$ を複体
$$
K^\bullet = \colim_{\alpha \in E} (W_\alpha \to X)_!I_\alpha^\bullet
$$
が表現する導来圏の対象とする。ここで遷移写像は，
$\beta<\alpha$ に対して上で注意深く構成した写像
$\tau_{\beta,\alpha}$ によって与えられる。
上とまったく同じ議論により，すべての $\alpha$ に対して
余極限への標準写像の制限は同型
$$
K|_{W_\alpha} \longrightarrow K_\alpha
$$
を定め，これは与えられた写像 $\rho^\alpha_\beta$ と両立する。
\end{proof}

\noindent
超限帰納法を用いると，一般の場合の結果を証明できる。

\begin{theorem}[BBD 貼り合わせ補題]
\label{cohomology-theorem-glueing-bbd-general}
\begin{reference}
\cite[Theorem 3.2.4]{BBD} の，有界性を仮定しない特別な場合。
\end{reference}
状況 \ref{cohomology-situation-locally-given} において，次を仮定する。
\begin{enumerate}
\item $X = \bigcup_{U \in \mathcal{B}} U$,
\item $U,V \in \mathcal{B}$ に対し
$U \cap V = \bigcup_{W \in \mathcal{B}, W \subset U \cap V} W$,
\item 任意の $U \in \mathcal{B}$ に対し，$i<0$ ならば
$\Ext^i(K_U,K_U)=0$ である。
\end{enumerate}
このとき $D(\mathcal{O}_X)$ の対象 $K$ と，
$U\in\mathcal{B}$ に対する $D(\mathcal{O}_U)$ の同型
$\rho_U:K|_U\to K_U$ が存在し，すべての $V\subset U$，
$U,V\in\mathcal{B}$ に対し
$\rho^U_V\circ\rho_U|_V=\rho_V$ を満たす。
対 $(K,\rho_U)$ は一意な同型を除いて一意である。
\end{theorem}

\begin{proof}
上の本文では，対 $(K,\rho_U)$ を解と呼んだ。一意性は補題
\ref{cohomology-lemma-uniqueness} から従う。$X$ が $\mathcal{B}$ の要素による
有限被覆をもつ場合（例えば $X$ が準コンパクトな場合），定理は補題
\ref{cohomology-lemma-solution-in-finite-case} の帰結である。一般の場合も，
超限帰納法と補題 \ref{cohomology-lemma-glueing-increasing-union} を用いて
まったく同じ仕方で論じる。

\medskip\noindent
まず超限再帰を用いて，任意の順序数 $\alpha$ に対し開集合
$W_\alpha\subset X$ を選ぶ。すなわち $W_0=\emptyset$ とおく。
$\alpha=\beta+1$ が後続順序数ならば，$W_\beta=X$ の場合は
$W_\alpha=X$ とおき，$W_\beta\ne X$ の場合は
$W_\beta$ に含まれない $U_\alpha\in\mathcal{B}$ を取り，
$W_\alpha=W_\beta\cup U_\alpha$ とおく。
$\alpha$ が極限順序数ならば
$W_\alpha=\bigcup_{\beta<\alpha}W_\beta$ とおく。
十分大きい $\alpha$ に対し $W_\alpha=X$ となる。
すべての $\alpha$ に対して開集合 $W_\alpha$ は
$\mathcal{B}$ の要素の和であることに注意する。したがって
$\mathcal{B}_\alpha=\{U\in\mathcal{B},U\subset W_\alpha\}$ とおけば，
$$
S_\alpha = (\{K_U\}_{U \in \mathcal{B}_\alpha},
\{\rho_V^U\}_{V \subset U\text{ with }U, V \in \mathcal{B}_\alpha})
$$
は環付き空間 $W_\alpha$ 上の補題 \ref{cohomology-lemma-uniqueness} のような系である。

\medskip\noindent
すべての $\alpha$ に対し系 $S_\alpha$ が解をもつことを超限帰納法で示す。
十分大きい $\alpha$ に対してこの系は定理で与えられた系なので，
これにより定理が証明される。

\medskip\noindent
$\alpha=\beta+1$ が後続順序数である場合。
（この場合は上の補題の証明ですでに扱ったが，明瞭さのため議論を繰り返す。）
この場合，ある $U_\alpha\in\mathcal{B}$ に対し
$W_\alpha=W_\beta\cup U_\alpha$ であることを思い出そう。
帰納法の仮定により，系 $S_\beta$ の解
$(K_{W_\beta}, \{\rho^{W_\beta}_U\}_{U \in \mathcal{B}_\beta})$
をもつ。次に $W_\alpha$ の開集合の族
$\mathcal{B}_\alpha^*=\mathcal{B}_\alpha\cup\{W_\beta\}$ を考えると，系
$(\{K_U\}_{U \in \mathcal{B}_\alpha^*},
\{\rho_V^U\}_{V \subset U\text{ with }U, V \in \mathcal{B}_\alpha^*})$.
を得る。解 $K_{W_\beta}$ に補題 \ref{cohomology-lemma-uniqueness} を適用すると，
この新しい系も条件 (3) を満たすことに注意する。この系について
$W_\alpha=W_\beta\cup U_\alpha$ である。これで補題
\ref{cohomology-lemma-solution-in-finite-case} で扱った場合に帰着する。

\medskip\noindent
$\alpha$ が極限順序数である場合。この場合
$W_\alpha=\bigcup_{\beta<\alpha}W_\beta$ であることを思い出そう。
$\beta<\alpha$ に対し，
$(K_{W_\beta}, \{\rho^{W_\beta}_U\}_{U \in \mathcal{B}_\beta})$
を $S_\beta$ の解とする。$\gamma<\beta<\alpha$ に対し，
写像 $\rho^{W_\beta}_U$（$U\in\mathcal{B}_\gamma$）を備えた制限
$K_{W_\beta}|_{W_\gamma}$ は $S_\gamma$ の解である。一意性により，
一意な同型
$\rho_{W_\gamma}^{W_\beta} : K_{W_\beta}|_{W_\gamma} \to K_{W_\gamma}$
で，$U\in\mathcal{B}_\gamma$ に対して
$\rho^{W_\beta}_U$ と $\rho^{W_\gamma}_U$ に両立するものを得る。
これらの写像は正しく合成する。すなわち
$\delta<\gamma<\beta<\alpha$ に対し
$\rho_{W_\delta}^{W_\gamma}\circ
\rho_{W_\gamma}^{W_\beta}|_{W_\delta}
=\rho^{W_\delta}_{W_\beta}$ である。
したがって補題 \ref{cohomology-lemma-glueing-increasing-union} を適用できる
（解 $K_{W_\beta}$ に補題 \ref{cohomology-lemma-uniqueness} を適用すると，
$K_{W_\beta}$ の負の Ext が消滅することに注意せよ）。これにより
$K_{W_\alpha}$ と同型
$$
\rho_{W_\beta}^{W_\alpha} :
K_{W_\alpha}|_{W_\beta}
\longrightarrow
K_{W_\beta}
$$
で，$\gamma<\beta<\alpha$ に対して写像
$\rho_{W_\gamma}^{W_\beta}$ と両立するものを得る。

\medskip\noindent
$K_{W_\alpha}$ が解であることを示すには，さらに
$U\in\mathcal{B}_\alpha$ に対し，ある両立性を満たす同型
$\rho_U^{W_\alpha}:K_{W_\alpha}|_U\to K_U$ を構成する必要がある。
$\rho_U^{W_\alpha}$ を，任意の $\beta<\alpha$ と任意の
$V\in\mathcal{B}_\beta$，$V\subset U$ に対して図式
$$
\xymatrix{
K_{W_\alpha}|_V \ar[r]_{\rho_U^{W_\alpha}|_V}
\ar[d]_{\rho_{W_\beta}^{W_\alpha}|_V}
& K_U|_V \ar[d]^{\rho_U^V} \\
K_{W_\beta} \ar[r]^{\rho_V^{W_\beta}}
& K_V
}
$$
を可換にする一意な写像として選ぶ。これは，環付き空間 $U$ 上で
$$
(\{K_V\}_{V \subset U, V \in \mathcal{B}_\beta\text{ for some }\beta < \alpha},
\{\rho_V^{V'}\}_{V \subset V'\text{ with }V, V' \subset U
\text{ and }V, V' \in \mathcal{B}_\beta\text{ for some }\beta < \alpha})
$$
が補題 \ref{cohomology-lemma-uniqueness} のような系であり，かつ
$(K_U,\rho^U_V)$ と
$(K_{W_\alpha}|_U,\rho_V^{W_\beta}\circ
\rho_{W_\beta}^{W_\alpha}|_V)$ がいずれもこの系の解であることから
意味をもつ。これにより存在と一意性を得る。これらの写像が望む両立性を
満たすことの証明は省略する（単なる記帳である）。
\end{proof}



\section{狭義完全複体}
\label{cohomology-section-strictly-perfect}

\noindent
加群の狭義完全複体は，後で擬コヒーレント複体および完全複体の概念を
定義するために用いられる。これは次のように定義される。

\begin{definition}
\label{cohomology-definition-strictly-perfect}
$(X, \mathcal{O}_X)$ を環付き空間とする。
$\mathcal{E}^\bullet$ を $\mathcal{O}_X$-加群の複体とする。
有限個を除くすべての $i$ に対して $\mathcal{E}^i$ が零であり，
すべての $i$ に対して $\mathcal{E}^i$ が有限自由
$\mathcal{O}_X$-加群の直和因子であるとき，
$\mathcal{E}^\bullet$ は {\it 狭義完全} であるという。
\end{definition}

\noindent
注意：$X$ が局所環付き空間であるとは仮定していないので，
有限自由 $\mathcal{O}_X$-加群の直和因子が有限局所自由であるとは
限らない。

\begin{lemma}
\label{cohomology-lemma-cone}
狭義完全複体の射の錐は狭義完全である。
\end{lemma}

\begin{proof}
定義から直ちに従う。
\end{proof}

\begin{lemma}
\label{cohomology-lemma-tensor}
二つの狭義完全複体のテンソル積に付随する全複体は狭義完全である。
\end{lemma}

\begin{proof}
省略する。
\end{proof}

\begin{lemma}
\label{cohomology-lemma-strictly-perfect-pullback}
$f : (X, \mathcal{O}_X) \to (Y, \mathcal{O}_Y)$ を環付き空間の射とする。
$\mathcal{F}^\bullet$ が $\mathcal{O}_Y$-加群の狭義完全複体ならば，
$f^*\mathcal{F}^\bullet$ は $\mathcal{O}_X$-加群の狭義完全複体である。
\end{lemma}

\begin{proof}
有限自由加群の引き戻しは有限自由である。関手 $f^*$ は加法的なので
直和因子を保つ。したがって補題が従う。
\end{proof}

\begin{lemma}
\label{cohomology-lemma-local-lift-map}
$(X, \mathcal{O}_X)$ を環付き空間とする。
$\mathcal{O}_X$-加群の実線図式
$$
\xymatrix{
\mathcal{E} \ar@{..>}[dr] \ar[r] & \mathcal{F} \\
& \mathcal{G} \ar[u]_p
}
$$
が与えられ，$\mathcal{E}$ が有限自由 $\mathcal{O}_X$-加群の直和因子であり，
$p$ が全射であるとする。このとき，図式を可換にする点線の矢印は
$X$ 上局所的に存在する。
\end{lemma}

\begin{proof}
ある $n$ に対して $\mathcal{E} = \mathcal{O}_X^{\oplus n}$ と仮定してよい。
この場合，点線の矢印を見つけることは，基底元の像を
$\Gamma(X, \mathcal{F})$ に持ち上げることと同値である。
層の全射の特徴づけ（Sheaves, Section
\ref{sheaves-section-exactness-points}）により，これは局所的に可能である。
\end{proof}

\begin{lemma}
\label{cohomology-lemma-local-homotopy}
$(X, \mathcal{O}_X)$ を環付き空間とする。
\begin{enumerate}
\item $\alpha : \mathcal{E}^\bullet \to \mathcal{F}^\bullet$ を
$\mathcal{O}_X$-加群の複体の射とし，$\mathcal{E}^\bullet$ は狭義完全，
$\mathcal{F}^\bullet$ は非輪状であるとする。このとき $\alpha$ は
$X$ 上局所的に零とホモトピックである。
\item $\alpha : \mathcal{E}^\bullet \to \mathcal{F}^\bullet$ を
$\mathcal{O}_X$-加群の複体の射とし，$\mathcal{E}^\bullet$ は狭義完全，
$i < a$ に対して $\mathcal{E}^i = 0$，$i \geq a$ に対して
$H^i(\mathcal{F}^\bullet) = 0$ であるとする。このとき $\alpha$ は
$X$ 上局所的に零とホモトピックである。
\end{enumerate}
\end{lemma}

\begin{proof}
第1の主張は第2の主張から従うので，(2) のみを証明する。
複体 $\mathcal{E}^\bullet$ の長さに関する帰納法で証明する。
$\mathcal{E}$ が有限自由 $\mathcal{O}_X$-加群の直和因子であり，
$n \geq a$ が整数であって
$\mathcal{E}^\bullet \cong \mathcal{E}[-n]$ である場合，結論は補題
\ref{cohomology-lemma-local-lift-map} と，仮定した
$H^n(\mathcal{F}^\bullet)$ の消滅により
$\mathcal{F}^{n - 1} \to \Ker(\mathcal{F}^n \to \mathcal{F}^{n + 1})$
が全射であることから従う。
$i \in [a, b]$ を除いて $\mathcal{E}^i$ が零である場合，複体の分裂完全列
$$
0 \to \mathcal{E}^b[-b] \to \mathcal{E}^\bullet \to
\sigma_{\leq b - 1}\mathcal{E}^\bullet \to 0
$$
を得，これは $K(\mathcal{O}_X)$ に完全三角を定める。したがって，
三角圏の公理により完全列
$$
\Hom_{K(\mathcal{O}_X)}(
\sigma_{\leq b - 1}\mathcal{E}^\bullet, \mathcal{F}^\bullet)
\to
\Hom_{K(\mathcal{O}_X)}(\mathcal{E}^\bullet, \mathcal{F}^\bullet)
\to
\Hom_{K(\mathcal{O}_X)}(\mathcal{E}^b[-b], \mathcal{F}^\bullet)
$$
を得る。合成 $\mathcal{E}^b[-b] \to \mathcal{F}^\bullet$ は局所的に
零とホモトピックであるから，考えている写像は上の表示された完全列の
左辺の元から来ると仮定してよい。この元は帰納法の仮定により局所的に零である。
\end{proof}

\begin{lemma}
\label{cohomology-lemma-lift-through-quasi-isomorphism}
$(X, \mathcal{O}_X)$ を環付き空間とする。
$\mathcal{O}_X$-加群の複体の実線図式
$$
\xymatrix{
\mathcal{E}^\bullet \ar@{..>}[dr] \ar[r]_\alpha & \mathcal{F}^\bullet \\
& \mathcal{G}^\bullet \ar[u]_f
}
$$
が与えられ，$\mathcal{E}^\bullet$ は狭義完全で，$j < a$ に対して
$\mathcal{E}^j = 0$，$j > a$ に対して $H^j(f)$ は同型，
$j = a$ に対しては全射であるとする。このとき，ホモトピーを除いて
図式を可換にする点線の矢印が $X$ 上局所的に存在する。
\end{lemma}

\begin{proof}
写像 $f$ に関する仮定から，錐 $C(f)^\bullet$ の次数 $\geq a$ の
コホモロジー層は消滅する。したがって補題
\ref{cohomology-lemma-local-homotopy} により，合成
$\mathcal{E}^\bullet \to \mathcal{F}^\bullet \to C(f)^\bullet$
が各 $U_i$ 上で零とホモトピックになるような開被覆
$X = \bigcup U_i$ が存在する。
$$
\mathcal{G}^\bullet \to \mathcal{F}^\bullet \to C(f)^\bullet \to
\mathcal{G}^\bullet[1]
$$
は $K(\mathcal{O}_{U_i})$ において完全三角へ制限されるので，
$\alpha|_{U_i}$ はホモトピーを除いて写像
$\alpha_i : \mathcal{E}^\bullet|_{U_i} \to \mathcal{G}^\bullet|_{U_i}$
へ持ち上がる。これが望む結論である。
\end{proof}

\begin{lemma}
\label{cohomology-lemma-local-actual}
$(X, \mathcal{O}_X)$ を環付き空間とする。
$\mathcal{E}^\bullet$，$\mathcal{F}^\bullet$ を
$\mathcal{O}_X$-加群の複体とし，$\mathcal{E}^\bullet$ は狭義完全であるとする。
\begin{enumerate}
\item 任意の元
$\alpha \in \Hom_{D(\mathcal{O}_X)}(\mathcal{E}^\bullet, \mathcal{F}^\bullet)$
に対して，$\alpha|_{U_i}$ が複体の射
$\alpha_i : \mathcal{E}^\bullet|_{U_i} \to \mathcal{F}^\bullet|_{U_i}$.
で与えられるような開被覆 $X = \bigcup U_i$ が存在する。
\item 複体の射 $\alpha : \mathcal{E}^\bullet \to \mathcal{F}^\bullet$ が与えられ，
その群
$\Hom_{D(\mathcal{O}_X)}(\mathcal{E}^\bullet, \mathcal{F}^\bullet)$
における像が零であるとする。このとき，$\alpha|_{U_i}$ が
零とホモトピックになるような開被覆 $X = \bigcup U_i$ が存在する。
\end{enumerate}
\end{lemma}

\begin{proof}
(1) を証明する。導来圏の構成により，$\alpha = f^{-1}\beta$ となる
擬同型 $f : \mathcal{F}^\bullet \to \mathcal{G}^\bullet$ と複体の射
$\beta : \mathcal{E}^\bullet \to \mathcal{G}^\bullet$ をとることができる。
したがって結論は補題 \ref{cohomology-lemma-lift-through-quasi-isomorphism} から従う。
(2) の証明は省略する。
\end{proof}

\begin{lemma}
\label{cohomology-lemma-Rhom-strictly-perfect}
$(X, \mathcal{O}_X)$ を環付き空間とする。
$\mathcal{E}^\bullet$，$\mathcal{F}^\bullet$ を
$\mathcal{O}_X$-加群の複体とし，$\mathcal{E}^\bullet$ は狭義完全であるとする。
このとき，内部 Hom
$R\SheafHom(\mathcal{E}^\bullet, \mathcal{F}^\bullet)$ は，各項が
$$
\mathcal{H}^n =
\bigoplus\nolimits_{n = p + q}
\SheafHom_{\mathcal{O}_X}(\mathcal{E}^{-q}, \mathcal{F}^p)
$$
であり，微分が Section \ref{cohomology-section-hom-complexes} で記述したものである
複体 $\mathcal{H}^\bullet$ によって表現される。
\end{lemma}

\begin{proof}
K-入射的複体への擬同型
$\mathcal{F}^\bullet \to \mathcal{I}^\bullet$ を選ぶ。
$(\mathcal{H}')^\bullet$ を，各項が
$$
(\mathcal{H}')^n =
\prod\nolimits_{n = p + q}
\SheafHom_{\mathcal{O}_X}(\mathcal{E}^{-q}, \mathcal{I}^p)
$$
である複体とする。Section \ref{cohomology-section-internal-hom} の構成により，
これは $R\SheafHom(\mathcal{E}^\bullet, \mathcal{F}^\bullet)$ を表現する。
写像
$$
\mathcal{H}^\bullet \longrightarrow (\mathcal{H}')^\bullet
$$
が擬同型であることを示せば十分である。開集合 $U \subset X$ に対して，
直接確かめれば
$$
H^0(\mathcal{H}^\bullet(U)) =
\Hom_{K(\mathcal{O}_U)}(\mathcal{E}^\bullet|_U, \mathcal{I}^\bullet|_U)
\to
H^0((\mathcal{H}')^\bullet(U)) =
\Hom_{D(\mathcal{O}_U)}(\mathcal{E}^\bullet|_U, \mathcal{I}^\bullet|_U)
$$
を得る。補題 \ref{cohomology-lemma-local-actual} により，
$U \mapsto H^0(\mathcal{H}^\bullet(U))$ の層化は
$U \mapsto H^0((\mathcal{H}')^\bullet(U))$ の層化に等しい。
他のコホモロジー層についても同様に論じられる。
したがって $\mathcal{H}^\bullet$ は $(\mathcal{H}')^\bullet$ と
擬同型であり，補題が証明された。
\end{proof}

\begin{lemma}
\label{cohomology-lemma-Rhom-strictly-perfect-K-flat}
補題 \ref{cohomology-lemma-Rhom-strictly-perfect} の状況において，
$\mathcal{F}^\bullet$ が K-平坦ならば，$\mathcal{H}^\bullet$ は K-平坦である。
\end{lemma}

\begin{proof}
狭義完全複体 $\mathcal{E}^\bullet$ の有界性により，Hom 複体の定義に現れる
直積が直和になるので，$\mathcal{H}^\bullet$ は単に Hom 複体
$\SheafHom^\bullet(\mathcal{E}^\bullet, \mathcal{F}^\bullet)$
であることに注意する。$\mathcal{K}^\bullet$ を
$\mathcal{O}_X$-加群の非輪状複体とする。補題
\ref{cohomology-lemma-diagonal-better} の写像
$$
\gamma :
\text{Tot}(\mathcal{K}^\bullet \otimes
\SheafHom^\bullet(\mathcal{E}^\bullet, \mathcal{F}^\bullet))
\longrightarrow
\SheafHom^\bullet(\mathcal{E}^\bullet,
\text{Tot}(\mathcal{K}^\bullet \otimes \mathcal{F}^\bullet))
$$
を考える。$\mathcal{F}^\bullet$ は K-平坦なので，複体
$\text{Tot}(\mathcal{K}^\bullet \otimes \mathcal{F}^\bullet)$ は
非輪状であり，したがって補題 \ref{cohomology-lemma-local-actual}
（あるいは補題 \ref{cohomology-lemma-Rhom-strictly-perfect}）により
$\gamma$ の終域も非輪状である。それゆえ補題を証明するには，
$\gamma$ が複体の同型であることを示せば十分である。
このために，複体 $\mathcal{E}^\bullet$ の長さに関する帰納法を用いる。
長さが $\leq 1$ ならば，ある $n \geq 0$ と $k \in \mathbf{Z}$ に対して
$\mathcal{E}^\bullet$ は $\mathcal{O}_X^{\oplus n}[k]$ の直和因子であり，
この場合は直接確かめられる。長さが $> 1$ ならば，適当な
$a \in \mathbf{Z}$ に対する複体の，各次数で分裂する短完全列
$$
0 \to \sigma_{\geq a + 1} \mathcal{E}^\bullet \to
\mathcal{E}^\bullet \to
\sigma_{\leq a} \mathcal{E}^\bullet \to 0
$$
を考えることで，より短い長さへ帰着する。Homology, Section
\ref{homology-section-truncations} を参照せよ。
このとき $\gamma$ は，複体の，各次数で分裂する短完全列の射に組み込まれる。
帰納法により，$\sigma_{\geq a + 1} \mathcal{E}^\bullet$ と
$\sigma_{\leq a} \mathcal{E}^\bullet$ に対して $\gamma$ は同型であるから，
$\mathcal{E}^\bullet$ に対する結論が従う。細部は省略する。
\end{proof}

\begin{lemma}
\label{cohomology-lemma-Rhom-complex-of-direct-summands-finite-free}
$(X, \mathcal{O}_X)$ を環付き空間とする。
$\mathcal{E}^\bullet$，$\mathcal{F}^\bullet$ を
$\mathcal{O}_X$-加群の複体とし，次を仮定する。
\begin{enumerate}
\item $n \ll 0$ に対して $\mathcal{F}^n = 0$。
\item $n \gg 0$ に対して $\mathcal{E}^n = 0$。
\item $\mathcal{E}^n$ は有限自由 $\mathcal{O}_X$-加群の直和因子と同型である。
\end{enumerate}
このとき，内部 Hom
$R\SheafHom(\mathcal{E}^\bullet, \mathcal{F}^\bullet)$ は，各項が
$$
\mathcal{H}^n =
\bigoplus\nolimits_{n = p + q}
\SheafHom_{\mathcal{O}_X}(\mathcal{E}^{-q}, \mathcal{F}^p)
$$
であり，微分が Section \ref{cohomology-section-internal-hom} で記述したものである
複体 $\mathcal{H}^\bullet$ によって表現される。
\end{lemma}

\begin{proof}
$\mathcal{I}^\bullet$ が入射的加群からなる下に有界な複体となるように，
擬同型 $\mathcal{F}^\bullet \to \mathcal{I}^\bullet$ を選ぶ。
$\mathcal{I}^\bullet$ は K-入射的であることに注意する
（Derived Categories, Lemma
\ref{derived-lemma-bounded-below-injectives-K-injective}）。
したがって Section \ref{cohomology-section-internal-hom} の構成から，
$R\SheafHom(\mathcal{E}^\bullet, \mathcal{F}^\bullet)$ は，各項が
$$
(\mathcal{H}')^n =
\prod\nolimits_{n = p + q}
\SheafHom_{\mathcal{O}_X}(\mathcal{E}^{-q}, \mathcal{I}^p) =
\bigoplus\nolimits_{n = p + q}
\SheafHom_{\mathcal{O}_X}(\mathcal{E}^{-q}, \mathcal{I}^p)
$$
である複体 $(\mathcal{H}')^\bullet$ によって表現される
（非零の項は有限個しかないので等号が成り立つ）。
$\mathcal{H}^\bullet$ は，各項が
$\SheafHom_{\mathcal{O}_X}(\mathcal{E}^{-q}, \mathcal{F}^p)$
である二重複体に付随する全複体であり，
$(\mathcal{H}')^\bullet$ についても同様である。
自然な写像 $(\mathcal{H}')^\bullet \to \mathcal{H}^\bullet$ は
二重複体の写像から来る。したがってこの写像が擬同型であることを示すには，
二重複体のスペクトル系列
（Homology, Lemma \ref{homology-lemma-first-quadrant-ss}）
$$
{}'E_1^{p, q} =
H^p(\SheafHom_{\mathcal{O}_X}(\mathcal{E}^{-q}, \mathcal{F}^\bullet))
$$
を用いればよい。これは $H^{p + q}(\mathcal{H}^\bullet)$ に収束し，
$(\mathcal{H}')^\bullet$ についても同様である。補題の証明を終えるには，
$\mathcal{E}$ が有限自由 $\mathcal{O}_X$-加群の直和因子であるとき，
$\mathcal{F}^\bullet \to \mathcal{I}^\bullet$ がコホモロジー層上で同型
$$
H^p(\SheafHom_{\mathcal{O}_X}(\mathcal{E}, \mathcal{F}^\bullet))
\longrightarrow
H^p(\SheafHom_{\mathcal{O}_X}(\mathcal{E}, \mathcal{I}^\bullet))
$$
を誘導することを示せば十分である。$\mathcal{E}$ が有限自由である場合，
これは明らかなので，結論が従う。
\end{proof}





\section{擬コヒーレント加群}
\label{cohomology-section-pseudo-coherent}

\noindent
この節では擬コヒーレント複体を論じる。

\begin{definition}
\label{cohomology-definition-pseudo-coherent}
$(X, \mathcal{O}_X)$ を環付き空間とし，$\mathcal{E}^\bullet$ を
$\mathcal{O}_X$-加群の複体とする。$m \in \mathbf{Z}$ とする。
\begin{enumerate}
\item 開被覆 $X = \bigcup U_i$ と，各 $i$ に対して複体の射
$\alpha_i : \mathcal{E}_i^\bullet \to \mathcal{E}^\bullet|_{U_i}$
で，$\mathcal{E}_i^\bullet$ は $U_i$ 上狭義完全であり，
$j > m$ に対して $H^j(\alpha_i)$ が同型，$H^m(\alpha_i)$ が
全射となるものが存在するとき，$\mathcal{E}^\bullet$ は
{\it $m$-擬コヒーレント} であるという。
\item すべての $m$ に対して $m$-擬コヒーレントであるとき，
$\mathcal{E}^\bullet$ は {\it 擬コヒーレント} であるという。
\item $D(\mathcal{O}_X)$ の対象 $E$ が {\it $m$-擬コヒーレント}
（それぞれ {\it 擬コヒーレント}）であるとは，$E$ が
$\mathcal{O}_X$-加群の $m$-擬コヒーレント
（それぞれ擬コヒーレント）複体によって表現されることをいう。
\end{enumerate}
\end{definition}

\noindent
$X$ が準コンパクトならば，$D(\mathcal{O}_X)$ の
$m$-擬コヒーレント対象は $D^-(\mathcal{O}_X)$ に属する。
しかし $X$ が準コンパクトでなければ，そうとは限らない。

\begin{lemma}
\label{cohomology-lemma-pseudo-coherent-independent-representative}
$(X, \mathcal{O}_X)$ を環付き空間とし，$E$ を
$D(\mathcal{O}_X)$ の対象とする。
\begin{enumerate}
\item 開被覆 $X = \bigcup U_i$，$U_i$ 上の狭義完全複体
$\mathcal{E}_i^\bullet$，および $D(\mathcal{O}_{U_i})$ における写像
$\alpha_i : \mathcal{E}_i^\bullet \to E|_{U_i}$ で，$j > m$ に対して
$H^j(\alpha_i)$ が同型，$H^m(\alpha_i)$ が全射となるものが存在すれば，
$E$ は $m$-擬コヒーレントである。
\item $E$ が $m$-擬コヒーレントならば，$E$ を表現する任意の複体は
$m$-擬コヒーレントである。
\end{enumerate}
\end{lemma}

\begin{proof}
$\mathcal{F}^\bullet$ を $E$ を表現する任意の複体とし，
$X = \bigcup U_i$ と
$\alpha_i : \mathcal{E}_i^\bullet \to E|_{U_i}$ を (1) のとおりとする。
$\mathcal{F}^\bullet$ が複体として $m$-擬コヒーレントであることを示す。
これで (1) と (2) が同時に証明される。
補題 \ref{cohomology-lemma-local-actual} により，開被覆 $X = \bigcup U_i$ を
細分すれば，写像 $\alpha_i$ を複体の写像
$\alpha_i : \mathcal{E}_i^\bullet \to \mathcal{F}^\bullet|_{U_i}$.
で表現できる。仮定により，$j > m$ に対して $H^j(\alpha_i)$ は同型であり，
$H^m(\alpha_i)$ は全射である。したがって $\mathcal{F}^\bullet$ は
$m$-擬コヒーレントである。
\end{proof}

\begin{lemma}
\label{cohomology-lemma-pseudo-coherent-pullback}
$f : (X, \mathcal{O}_X) \to (Y, \mathcal{O}_Y)$ を環付き空間の射とする。
$E$ を $D(\mathcal{O}_Y)$ の対象とする。$E$ が
$m$-擬コヒーレントならば，$Lf^*E$ は $m$-擬コヒーレントである。
\end{lemma}

\begin{proof}
$E$ を $\mathcal{O}_Y$-加群の複体 $\mathcal{E}^\bullet$ で表現し，
定義 \ref{cohomology-definition-pseudo-coherent} のような開被覆
$Y = \bigcup V_i$ と
$\alpha_i : \mathcal{E}_i^\bullet \to \mathcal{E}^\bullet|_{V_i}$
を選ぶ。$U_i = f^{-1}(V_i)$ とおく。
補題 \ref{cohomology-lemma-pseudo-coherent-independent-representative} により，
$Lf^*\mathcal{E}^\bullet|_{U_i}$ が $m$-擬コヒーレントであることを
示せば十分である。完全三角
$$
\mathcal{E}_i^\bullet \to
\mathcal{E}^\bullet|_{V_i} \to
C \to
\mathcal{E}_i^\bullet[1]
$$
を選ぶ。$\alpha_i$ に関する仮定は，すべての $j \geq m$ に対して
コホモロジー層 $H^j(C)$ が零であることにほかならない。
$f_i : U_i \to V_i$ を $f$ の制限とする。
$Lf^*\mathcal{E}^\bullet|_{U_i} = 
Lf_i^*(\mathcal{E}|_{V_i})$ であることに注意する。
$Lf_i^*$ を適用して完全三角
$$
Lf_i^*\mathcal{E}_i^\bullet \to
Lf_i^*\mathcal{E}|_{V_i} \to
Lf_i^*C \to
Lf_i^*\mathcal{E}_i^\bullet[1]
$$
を得る。左導来関手としての $Lf_i^*$ の構成から，
$j \geq m$ に対して $H^j(Lf_i^*C) = 0$ である
（Derived Categories, Lemma
\ref{derived-lemma-negative-vanishing} の双対）。
したがって $j > m$ に対して $H^j(Lf_i^*\alpha_i)$ は同型であり，
$H^m(Lf^*\alpha_i)$ は全射である。他方，
$Lf_i^*\mathcal{E}_i^\bullet = f_i^*\mathcal{E}_i^\bullet$ は
補題 \ref{cohomology-lemma-strictly-perfect-pullback} により狭義完全である。
これで結論を得る。
\end{proof}

\begin{lemma}
\label{cohomology-lemma-cone-pseudo-coherent}
$(X, \mathcal{O}_X)$ を環付き空間，$m \in \mathbf{Z}$ とする。
$(K, L, M, f, g, h)$ を $D(\mathcal{O}_X)$ の完全三角とする。
\begin{enumerate}
\item $K$ が $(m + 1)$-擬コヒーレントで，$L$ が
$m$-擬コヒーレントならば，$M$ は $m$-擬コヒーレントである。
\item $K$ と $M$ が $m$-擬コヒーレントならば，$L$ は
$m$-擬コヒーレントである。
\item $L$ が $(m + 1)$-擬コヒーレントで，$M$ が
$m$-擬コヒーレントならば，$K$ は $(m + 1)$-擬コヒーレントである。
\end{enumerate}
\end{lemma}

\begin{proof}
(1) を証明する。開被覆 $X = \bigcup U_i$ と
$D(\mathcal{O}_{U_i})$ における写像
$\alpha_i : \mathcal{K}_i^\bullet \to K|_{U_i}$ で，
$\mathcal{K}_i^\bullet$ は狭義完全，$j > m + 1$ に対して
$H^j(\alpha_i)$ は同型，$j = m + 1$ に対して全射となるものを選ぶ。
$\mathcal{K}_i^\bullet$ を
$\sigma_{\geq m + 1}\mathcal{K}_i^\bullet$
で置き換えられるので，$j < m + 1$ に対して
$\mathcal{K}_i^j = 0$ と仮定してよい。
開被覆を細分した後，$D(\mathcal{O}_{U_i})$ における写像
$\beta_i : \mathcal{L}_i^\bullet \to L|_{U_i}$ で，
$\mathcal{L}_i^\bullet$ は狭義完全，$j > m$ に対して
$H^j(\beta)$ は同型，$j = m$ に対して全射となるものを選べる。
補題 \ref{cohomology-lemma-lift-through-quasi-isomorphism} により，被覆をさらに
細分すれば，複体の写像
$\gamma_i : \mathcal{K}^\bullet \to \mathcal{L}^\bullet$
で図式
$$
\xymatrix{
K|_{U_i} \ar[r] & L|_{U_i} \\
\mathcal{K}_i^\bullet \ar[u]^{\alpha_i} \ar[r]^{\gamma_i} &
\mathcal{L}_i^\bullet \ar[u]_{\beta_i}
}
$$
が $D(\mathcal{O}_{U_i})$ で可換になるものを見つけられる
（このためには写像 $\alpha_i$，$\beta_i$ および
$K|_{U_i} \to L|_{U_i}$ を実際の複体の写像で表現する必要がある。
細部は省略する）。錐 $C(\gamma_i)^\bullet$ は狭義完全である
（補題 \ref{cohomology-lemma-cone}）。図式の可換性から，完全三角の射
$$
(\mathcal{K}_i^\bullet, \mathcal{L}_i^\bullet, C(\gamma_i)^\bullet)
\longrightarrow
(K|_{U_i}, L|_{U_i}, M|_{U_i}).
$$
が存在する。誘導されるコホモロジー長完全列の写像と
Homology, Lemmas \ref{homology-lemma-four-lemma} および
\ref{homology-lemma-five-lemma} から，
$C(\gamma_i)^\bullet \to M|_{U_i}$ は次数 $> m$ のコホモロジー上で
同型を，次数 $m$ で全射を誘導する。したがって補題
\ref{cohomology-lemma-pseudo-coherent-independent-representative} により，
$M$ は $m$-擬コヒーレントである。

\medskip\noindent
主張 (2) と (3) は，完全三角を回転して (1) を適用すれば従う。
\end{proof}

\begin{lemma}
\label{cohomology-lemma-tensor-pseudo-coherent}
$(X, \mathcal{O}_X)$ を環付き空間とし，$K,L$ を
$D(\mathcal{O}_X)$ の対象とする。
\begin{enumerate}
\item $K$ が $n$-擬コヒーレントで，$i > a$ に対して
$H^i(K) = 0$ であり，$L$ が $m$-擬コヒーレントで，$j > b$ に対して
$H^j(L) = 0$ ならば，$t = \max(m + a, n + b)$ とおくと
$K \otimes_{\mathcal{O}_X}^\mathbf{L} L$ は $t$-擬コヒーレントである。
\item $K$ と $L$ が擬コヒーレントならば，
$K \otimes_{\mathcal{O}_X}^\mathbf{L} L$ は擬コヒーレントである。
\end{enumerate}
\end{lemma}

\begin{proof}
(1) を証明する。$X$ を開被覆の各要素で置き換えることにより，
狭義完全複体 $\mathcal{K}^\bullet$，$\mathcal{L}^\bullet$ と写像
$\alpha : \mathcal{K}^\bullet \to K$，$\beta : \mathcal{L}^\bullet \to L$
が存在し，$i > n$ に対して $H^i(\alpha)$ が同型，$i = n$ に対して
全射，また $i > m$ に対して $H^i(\beta)$ が同型，$i = m$ に対して
全射であると仮定してよい。このとき写像
$$
\alpha \otimes^\mathbf{L} \beta :
\text{Tot}(\mathcal{K}^\bullet \otimes_{\mathcal{O}_X} \mathcal{L}^\bullet)
\to K \otimes_{\mathcal{O}_X}^\mathbf{L} L
$$
は $i > t$ に対して次数 $i$ のコホモロジー層上で同型を，
$i = t$ に対して全射を誘導する。これは Tor のスペクトル系列から従う
（細部は省略する）。

\medskip\noindent
(2) を証明する。まず $X$ を開被覆の各要素で置き換えることにより，
$K$ と $L$ が上に有界な場合へ帰着できる。このとき主張は (1) から
直ちに従う。
\end{proof}

\begin{lemma}
\label{cohomology-lemma-summands-pseudo-coherent}
$(X, \mathcal{O}_X)$ を環付き空間とし，$m \in \mathbf{Z}$ とする。
$D(\mathcal{O}_X)$ において $K \oplus L$ が $m$-擬コヒーレント
（それぞれ擬コヒーレント）ならば，$K$ と $L$ もそうである。
\end{lemma}

\begin{proof}
$K \oplus L$ が $m$-擬コヒーレントであると仮定する。
$X$ を開被覆の各要素で置き換えれば
$K \oplus L \in D^-(\mathcal{O}_X)$，したがって
$L \in D^-(\mathcal{O}_X)$ と仮定してよい。
完全三角
$$
(K \oplus L, K \oplus L, L \oplus L[1]) =
(K, K, 0) \oplus (L, L, L \oplus L[1])
$$
が存在することに注意する。Derived Categories, Lemma
\ref{derived-lemma-direct-sum-triangles} を参照せよ。
補題 \ref{cohomology-lemma-cone-pseudo-coherent} により，
$L \oplus L[1]$ は $m$-擬コヒーレントである。
したがって $L[1] \oplus L[2]$ も $m$-擬コヒーレントであり，
帰納的に $L[n] \oplus L[n + 1]$ は $m$-擬コヒーレントである。
$L$ は上に有界なので，十分大きな $n$ に対して $L[n]$ は
$m$-擬コヒーレントである。したがって完全三角
$$
(L[n], L[n] \oplus L[n - 1], L[n - 1])
$$
を用いて逆向きに進めると，
$L[n - 1], L[n - 2], \ldots, L$ が $m$-擬コヒーレントであるとわかる。
\end{proof}

\begin{lemma}
\label{cohomology-lemma-complex-pseudo-coherent-modules}
$(X, \mathcal{O}_X)$ を環付き空間とする。$m \in \mathbf{Z}$ とする。
$\mathcal{F}^\bullet$ を $\mathcal{O}_X$-加群の（局所的に）上に有界な
複体とし，すべての $i$ に対して $\mathcal{F}^i$ は
$(m - i)$-擬コヒーレントであるとする。このとき
$\mathcal{F}^\bullet$ は $m$-擬コヒーレントである。
\end{lemma}

\begin{proof}
省略する。ヒント：補題 \ref{cohomology-lemma-cone-pseudo-coherent} と切断を用いる。
これは More on Algebra, Lemma
\ref{more-algebra-lemma-complex-pseudo-coherent-modules} の証明と同様である。
\end{proof}

\begin{lemma}
\label{cohomology-lemma-cohomology-pseudo-coherent}
$(X, \mathcal{O}_X)$ を環付き空間とし，$m \in \mathbf{Z}$ とする。
$E$ を $D(\mathcal{O}_X)$ の対象とする。$E$ が（局所的に）上に有界であり，
すべての $i$ に対して $H^i(E)$ が $(m - i)$-擬コヒーレントならば，
$E$ は $m$-擬コヒーレントである。
\end{lemma}

\begin{proof}
省略する。ヒント：補題 \ref{cohomology-lemma-cone-pseudo-coherent} と切断を用いる。
これは More on Algebra, Lemma
\ref{more-algebra-lemma-cohomology-pseudo-coherent} の証明と同様である。
\end{proof}

\begin{lemma}
\label{cohomology-lemma-finite-cohomology}
$(X, \mathcal{O}_X)$ を環付き空間とする。
$K$ を $D(\mathcal{O}_X)$ の対象とする。$m \in \mathbf{Z}$ とする。
\begin{enumerate}
\item $K$ が $m$-擬コヒーレントで，$i > m$ に対して
$H^i(K) = 0$ ならば，$H^m(K)$ は有限型 $\mathcal{O}_X$-加群である。
\item $K$ が $m$-擬コヒーレントで，$i > m + 1$ に対して
$H^i(K) = 0$ ならば，$H^{m + 1}(K)$ は有限表示
$\mathcal{O}_X$-加群である。
\end{enumerate}
\end{lemma}

\begin{proof}
(1) を証明する。$X$ 上局所的に作業してよい。したがって，
狭義完全複体 $\mathcal{E}^\bullet$ と写像
$\alpha : \mathcal{E}^\bullet \to K$ で，次数 $> m$ の
コホモロジー上で同型を，次数 $m$ で全射を誘導するものが存在すると
仮定してよい。$\mathcal{E}^\bullet$ に対して結論を証明すれば十分である。
$\mathcal{E}^n \not = 0$ となる最大の整数を $n$ とする。
$n = m$ ならば，$H^m(\mathcal{E}^\bullet)$ は
$\mathcal{E}^n$ の商であり，結論は明らかである。
$n > m$ ならば，$H^n(E^\bullet) = 0$ なので
$\mathcal{E}^{n - 1} \to \mathcal{E}^n$ は全射である。
補題 \ref{cohomology-lemma-local-lift-map} により，この全射の切断を局所的に
見つけて
$\mathcal{E}^{n - 1} = \mathcal{E}' \oplus \mathcal{E}^n$.
と書ける。したがって，次数 $n-1$ の項が $\mathcal{E}'$，
次数 $n$ の項が $0$ であることを除いて
$\mathcal{E}^\bullet$ と同じ複体 $(\mathcal{E}')^\bullet$ に対して
結論を証明すれば十分である。$n$ に関する帰納法で結論を得る。

\medskip\noindent
(2) を証明する。$X$ 上局所的に作業してよい。したがって，
狭義完全複体 $\mathcal{E}^\bullet$ と写像
$\alpha : \mathcal{E}^\bullet \to K$ で，次数 $> m$ の
コホモロジー上で同型を，次数 $m$ で全射を誘導するものが存在すると
仮定してよい。(1) の証明と同様に，$i > m + 1$ に対して
$\mathcal{E}^i = 0$ である場合へ帰着できる。このとき
$H^{m + 1}(K) \cong H^{m + 1}(\mathcal{E}^\bullet) =
\Coker(\mathcal{E}^m \to \mathcal{E}^{m + 1})$
であり，これは有限表示である。
\end{proof}

\begin{lemma}
\label{cohomology-lemma-n-pseudo-module}
$(X, \mathcal{O}_X)$ を環付き空間とし，$\mathcal{F}$ を
$\mathcal{O}_X$-加群の層とする。
\begin{enumerate}
\item $D(\mathcal{O}_X)$ の対象とみなした $\mathcal{F}$ が
$0$-擬コヒーレントであるための必要十分条件は，
$\mathcal{F}$ が有限型 $\mathcal{O}_X$-加群であることである。
\item $D(\mathcal{O}_X)$ の対象とみなした $\mathcal{F}$ が
$(-1)$-擬コヒーレントであるための必要十分条件は，
$\mathcal{F}$ が有限表示 $\mathcal{O}_X$-加群であることである。
\end{enumerate}
\end{lemma}

\begin{proof}
一方向の含意には補題 \ref{cohomology-lemma-finite-cohomology} を，
逆方向には補題 \ref{cohomology-lemma-cohomology-pseudo-coherent} を用いる。
\end{proof}





\section{Tor 次元}
\label{cohomology-section-tor}

\noindent
この節では，平坦加群による分解をより詳しく調べる。

\begin{definition}
\label{cohomology-definition-tor-amplitude}
$(X, \mathcal{O}_X)$ を環付き空間とし，$E$ を
$D(\mathcal{O}_X)$ の対象とする。$a \leq b$ なる
$a, b \in \mathbf{Z}$ とする。
\begin{enumerate}
\item すべての $\mathcal{O}_X$-加群 $\mathcal{F}$ と
すべての $i \not \in [a, b]$ に対して
$H^i(E \otimes_{\mathcal{O}_X}^\mathbf{L} \mathcal{F}) = 0$
であるとき，$E$ は {\it $[a,b]$ に Tor 振幅をもつ} という。
\item ある $a,b$ に対して $[a,b]$ に Tor 振幅をもつとき，
$E$ は {\it 有限 Tor 次元をもつ} という。
\item すべての $i$ に対して $E|_{U_i}$ が有限 Tor 次元をもつような
開被覆 $X = \bigcup U_i$ が存在するとき，$E$ は
{\it 局所的に有限 Tor 次元をもつ} という。
\end{enumerate}
$D(\mathcal{O}_X)$ の対象とみなした $\mathcal{F}[0]$ が
$[-d,0]$ に Tor 振幅をもつとき，$\mathcal{O}_X$-加群
$\mathcal{F}$ は {\it Tor 次元 $\leq d$} をもつという。
\end{definition}

\noindent
定義における $E$ が有限 Tor 次元をもつならば，
$E$ は $D^b(\mathcal{O}_X)$ の対象であることに注意する。
これは上の定義で $\mathcal{F} = \mathcal{O}_X$ とすればわかる。

\begin{lemma}
\label{cohomology-lemma-last-one-flat}
$(X, \mathcal{O}_X)$ を環付き空間とする。
$\mathcal{E}^\bullet$ を，$[a,b]$ に Tor 振幅をもつ，
平坦 $\mathcal{O}_X$-加群からなる上に有界な複体とする。
このとき $\Coker(d_{\mathcal{E}^\bullet}^{a - 1})$ は
平坦 $\mathcal{O}_X$-加群である。
\end{lemma}

\begin{proof}
複体 $\mathcal{E}^\bullet$ は平坦加群からなる上に有界な複体なので，
任意の $\mathcal{O}_X$-加群 $\mathcal{F}$ に対して
$\mathcal{E}^\bullet \otimes_{\mathcal{O}_X} \mathcal{F} =
\mathcal{E}^\bullet \otimes_{\mathcal{O}_X}^{\mathbf{L}} \mathcal{F}$
である。したがって，すべての $\mathcal{O}_X$-加群 $\mathcal{F}$ に対して列
$$
\mathcal{E}^{a - 2} \otimes_{\mathcal{O}_X} \mathcal{F} \to
\mathcal{E}^{a - 1} \otimes_{\mathcal{O}_X} \mathcal{F} \to
\mathcal{E}^a \otimes_{\mathcal{O}_X} \mathcal{F}
$$
は中央で完全である。
$\mathcal{E}^{a - 2} \to \mathcal{E}^{a - 1} \to \mathcal{E}^a \to
\Coker(d^{a - 1}) \to 0$
は平坦分解なので，これはすべての $\mathcal{O}_X$-加群
$\mathcal{F}$ に対して
$\text{Tor}_1^{\mathcal{O}_X}(\Coker(d^{a - 1}), \mathcal{F}) = 0$
であることを意味する。ゆえに $\Coker(d^{a - 1})$ は平坦である。
補題 \ref{cohomology-lemma-flat-tor-zero} を参照せよ。
\end{proof}

\begin{lemma}
\label{cohomology-lemma-tor-amplitude}
$(X, \mathcal{O}_X)$ を環付き空間とし，$E$ を
$D(\mathcal{O}_X)$ の対象とする。$a \leq b$ なる
$a,b\in\mathbf{Z}$ とする。次は同値である。
\begin{enumerate}
\item $E$ は $[a,b]$ に Tor 振幅をもつ。
\item $E$ は，平坦 $\mathcal{O}_X$-加群からなる複体
$\mathcal{E}^\bullet$ で，$i \not \in [a,b]$ に対して
$\mathcal{E}^i = 0$ となるものによって表現される。
\end{enumerate}
\end{lemma}

\begin{proof}
(2) が成り立つならば，
$E \otimes_{\mathcal{O}_X}^\mathbf{L} \mathcal{F} =
\mathcal{E}^\bullet \otimes_{\mathcal{O}_X} \mathcal{F}$
と計算でき，(1) が成り立つことは明らかである。

\medskip\noindent
(1) が成り立つと仮定する。$E$ は平坦 $\mathcal{O}_X$-加群からなる
上に有界な複体 $\mathcal{K}^\bullet$ で表現できる。
Section \ref{cohomology-section-flat} を参照せよ。
$\mathcal{K}^n \not = 0$ となる最大の整数を $n$ とする。
$n > b$ ならば，$H^n(\mathcal{K}^\bullet) = 0$ なので
$\mathcal{K}^{n - 1} \to \mathcal{K}^n$ は全射である。
$\mathcal{K}^n$ は平坦だから
$\Ker(\mathcal{K}^{n - 1} \to \mathcal{K}^n)$ は平坦である
（Modules, Lemma \ref{modules-lemma-flat-ses}）。
したがって $\mathcal{K}^\bullet$ を
$\tau_{\leq n - 1}\mathcal{K}^\bullet$ で置き換えられる。
このように $n$ に関する帰納法により，$K^\bullet$ が
平坦 $\mathcal{O}_X$-加群からなる複体で，
$i > b$ に対して $\mathcal{K}^i = 0$ となる場合へ帰着する。

\medskip\noindent
$\mathcal{E}^\bullet = \tau_{\geq a}\mathcal{K}^\bullet$ とおく。
$\mathcal{E}^a$ が平坦であること以外はすべて明らかであり，
その平坦性は補題 \ref{cohomology-lemma-last-one-flat} と定義から直ちに従う。
\end{proof}

\begin{lemma}
\label{cohomology-lemma-tor-amplitude-pullback}
$f : (X, \mathcal{O}_X) \to (Y, \mathcal{O}_Y)$ を環付き空間の射とし，
$E$ を $D(\mathcal{O}_Y)$ の対象とする。$E$ が $[a,b]$ に
Tor 振幅をもつならば，$Lf^*E$ も $[a,b]$ に Tor 振幅をもつ。
\end{lemma}

\begin{proof}
$E$ が $[a,b]$ に Tor 振幅をもつと仮定する。補題
\ref{cohomology-lemma-tor-amplitude} により，$E$ は平坦 $\mathcal{O}$-加群からなる
複体 $\mathcal{E}^\bullet$ で，$i \not \in [a,b]$ に対して
$\mathcal{E}^i = 0$ となるものによって表現できる。このとき
$Lf^*E$ は $f^*\mathcal{E}^\bullet$ によって表現される。
Modules, Lemma \ref{modules-lemma-base-change-flat} により，
加群 $f^*\mathcal{E}^i$ は平坦である。ゆえに補題
\ref{cohomology-lemma-tor-amplitude} により，$Lf^*E$ は $[a,b]$ に
Tor 振幅をもつ。
\end{proof}

\begin{lemma}
\label{cohomology-lemma-tor-amplitude-stalk}
$(X, \mathcal{O}_X)$ を環付き空間とし，$E$ を
$D(\mathcal{O}_X)$ の対象とする。$a \leq b$ なる
$a,b\in\mathbf{Z}$ とする。次は同値である。
\begin{enumerate}
\item $E$ は $[a,b]$ に Tor 振幅をもつ。
\item すべての $x \in X$ に対して，$D(\mathcal{O}_{X,x})$ の対象
$E_x$ は $[a,b]$ に Tor 振幅をもつ。
\end{enumerate}
\end{lemma}

\begin{proof}
$x$ における茎をとることは，環付き空間の射
$(x, \mathcal{O}_{X, x}) \to (X, \mathcal{O}_X)$ による
引き戻しと同じである。したがって含意 (1) $\Rightarrow$ (2) は
補題 \ref{cohomology-lemma-tor-amplitude-pullback} から従う。
逆に，茎をとることはテンソル積と可換する
（Modules, Lemma \ref{modules-lemma-stalk-tensor-product}）。
したがって
$$
(E \otimes_{\mathcal{O}_X}^\mathbf{L} \mathcal{F})_x =
E_x \otimes_{\mathcal{O}_{X, x}}^\mathbf{L} \mathcal{F}_x
$$
である。他方，茎をとることは完全なので，
$$
H^i(E \otimes_{\mathcal{O}_X}^\mathbf{L} \mathcal{F})_x =
H^i((E \otimes_{\mathcal{O}_X}^\mathbf{L} \mathcal{F})_x) =
H^i(E_x \otimes_{\mathcal{O}_{X, x}}^\mathbf{L} \mathcal{F}_x)
$$
であり，$H^i(E \otimes_{\mathcal{O}_X}^\mathbf{L} \mathcal{F})$ が
零かどうかは，そのすべての茎が零かどうかを調べれば判定できる
（Modules, Lemma \ref{modules-lemma-abelian}）。よって (2) は (1) を含意する。
\end{proof}

\begin{lemma}
\label{cohomology-lemma-cone-tor-amplitude}
$(X, \mathcal{O}_X)$ を環付き空間とする。
$(K,L,M,f,g,h)$ を $D(\mathcal{O}_X)$ の完全三角とし，
$a,b\in\mathbf{Z}$ とする。
\begin{enumerate}
\item $K$ が $[a+1,b+1]$ に Tor 振幅をもち，$L$ が $[a,b]$ に
Tor 振幅をもつならば，$M$ は $[a,b]$ に Tor 振幅をもつ。
\item $K$ と $M$ が $[a,b]$ に Tor 振幅をもつならば，
$L$ は $[a,b]$ に Tor 振幅をもつ。
\item $L$ が $[a+1,b+1]$ に Tor 振幅をもち，$M$ が $[a,b]$ に
Tor 振幅をもつならば，$K$ は $[a+1,b+1]$ に Tor 振幅をもつ。
\end{enumerate}
\end{lemma}

\begin{proof}
省略する。ヒント：完全三角に付随するコホモロジー長完全列と，
$- \otimes_{\mathcal{O}_X}^{\mathbf{L}} \mathcal{F}$
が完全三角を保つことから直ちに従う。最も容易なのは (2) であり，
他は平行移動によりこれから従う。
\end{proof}

\begin{lemma}
\label{cohomology-lemma-tensor-tor-amplitude}
$(X, \mathcal{O}_X)$ を環付き空間とし，$K,L$ を
$D(\mathcal{O}_X)$ の対象とする。$K$ が $[a,b]$ に，
$L$ が $[c,d]$ に Tor 振幅をもつならば，
$K \otimes_{\mathcal{O}_X}^\mathbf{L} L$ は
$[a+c,b+d]$ に Tor 振幅をもつ。
\end{lemma}

\begin{proof}
省略する。ヒント：Tor のスペクトル系列を用いる。
\end{proof}

\begin{lemma}
\label{cohomology-lemma-summands-tor-amplitude}
$(X, \mathcal{O}_X)$ を環付き空間とし，$a,b\in\mathbf{Z}$ とする。
$K,L$ を $D(\mathcal{O}_X)$ の対象とする。
$K\oplus L$ が $[a,b]$ に Tor 振幅をもつならば，$K$ と $L$ もそうである。
\end{lemma}

\begin{proof}
Tor 関手が加法的であることから明らかである。
\end{proof}






\section{完全複体}
\label{cohomology-section-perfect}

\noindent
この節では，環付き空間上の完全複体の性質を論じる。

\begin{definition}
\label{cohomology-definition-perfect}
$(X, \mathcal{O}_X)$ を環付き空間とし，$\mathcal{E}^\bullet$ を
$\mathcal{O}_X$-加群の複体とする。開被覆 $X=\bigcup U_i$ が存在し，
各 $i$ に対して複体の射
$\mathcal{E}_i^\bullet \to \mathcal{E}^\bullet|_{U_i}$
が擬同型で，$\mathcal{E}_i^\bullet$ が
$\mathcal{O}_{U_i}$-加群の狭義完全複体となるものが存在するとき，
$\mathcal{E}^\bullet$ は {\it 完全} であるという。
$D(\mathcal{O}_X)$ の対象 $E$ が {\it 完全} であるとは，
$\mathcal{O}_X$-加群の完全複体によって表現されることをいう。
\end{definition}

\noindent
$X$ が準コンパクトならば，$D(\mathcal{O}_X)$ の完全対象は
$D^b(\mathcal{O}_X)$ に属する。しかし $X$ が準コンパクトでなければ，
そうとは限らない。

\begin{lemma}
\label{cohomology-lemma-perfect-independent-representative}
$(X, \mathcal{O}_X)$ を環付き空間とし，$E$ を
$D(\mathcal{O}_X)$ の対象とする。
\begin{enumerate}
\item 開被覆 $X=\bigcup U_i$ と，$D(\mathcal{O}_{U_i})$ において
$E|_{U_i}$ を表現する $U_i$ 上の狭義完全複体
$\mathcal{E}_i^\bullet$ が存在すれば，$E$ は完全である。
\item $E$ が完全ならば，$E$ を表現する任意の複体は完全である。
\end{enumerate}
\end{lemma}

\begin{proof}
補題 \ref{cohomology-lemma-pseudo-coherent-independent-representative} の証明と同じである。
\end{proof}

\begin{lemma}
\label{cohomology-lemma-perfect-on-locally-ringed}
$(X, \mathcal{O}_X)$ を環付き空間とし，$E$ を
$D(\mathcal{O}_X)$ の対象とする。すべての茎
$\mathcal{O}_{X,x}$ が局所環であると仮定する。次は同値である。
\begin{enumerate}
\item $E$ は完全である。
\item 開被覆 $X=\bigcup U_i$ が存在し，$E|_{U_i}$ は
有限局所自由 $\mathcal{O}_{U_i}$-加群からなる有限複体によって表現される。
\item 開被覆 $X=\bigcup U_i$ が存在し，$E|_{U_i}$ は
有限自由 $\mathcal{O}_{U_i}$-加群からなる有限複体によって表現される。
\end{enumerate}
\end{lemma}

\begin{proof}
補題 \ref{cohomology-lemma-perfect-independent-representative} と，$X$ 上では
有限自由加群の任意の直和因子が有限局所自由であることから従う。
Modules, Lemma
\ref{modules-lemma-direct-summand-of-locally-free-is-locally-free} を参照せよ。
\end{proof}

\begin{lemma}
\label{cohomology-lemma-perfect-precise}
$(X, \mathcal{O}_X)$ を環付き空間とし，$E$ を
$D(\mathcal{O}_X)$ の対象とする。$a\leq b$ を整数とする。
$E$ が $[a,b]$ に Tor 振幅をもち，$(a-1)$-擬コヒーレントならば，
$E$ は完全である。
\end{lemma}

\begin{proof}
$X$ を開被覆の各要素で置き換えることにより，狭義完全複体
$\mathcal{E}^\bullet$ と写像 $\alpha:\mathcal{E}^\bullet\to E$ で，
$i\geq a$ に対して $H^i(\alpha)$ が同型となるものが存在すると
仮定してよい。$\mathcal{E}^\bullet$ を
$\sigma_{\geq a - 1}\mathcal{E}^\bullet$ で置き換える。完全三角
$$
\mathcal{E}^\bullet \to E \to C \to \mathcal{E}^\bullet[1]
$$
を選ぶ。$E$ と $\mathcal{E}^\bullet$ のコホモロジー層の消滅，および
$\alpha$ に関する仮定から，
$\mathcal{K}=\Ker(\mathcal{E}^{a - 1}\to\mathcal{E}^a)$ とおけば
$C\cong\mathcal{K}[2-a]$ を得る。
$\mathcal{F}$ を $\mathcal{O}_X$-加群とする。
$-\otimes_{\mathcal{O}_X}^\mathbf{L}\mathcal{F}$ を適用すると，
$E$ が $[a,b]$ に Tor 振幅をもつという仮定により，
$\mathcal{K} \otimes_{\mathcal{O}_X} \mathcal{F} \to
\mathcal{E}^{a - 1} \otimes_{\mathcal{O}_X} \mathcal{F}$ の像は
$\Ker(\mathcal{E}^{a - 1} \otimes_{\mathcal{O}_X} \mathcal{F}
\to \mathcal{E}^a \otimes_{\mathcal{O}_X} \mathcal{F})$ である。
$\mathcal{E}'=\Coker(\mathcal{E}^{a - 1}\to\mathcal{E}^a)$ とおくと，
$\text{Tor}_1^{\mathcal{O}_X}(\mathcal{E}', \mathcal{F}) = 0$
が従う。したがって $\mathcal{E}'$ は平坦である
（補題 \ref{cohomology-lemma-flat-tor-zero}）。
よって Modules, Lemma
\ref{modules-lemma-flat-locally-finite-presentation} により，
$\mathcal{E}'$ は局所的に有限自由加群の直和因子である。
したがって局所的には，複体
$$
\mathcal{E}' \to \mathcal{E}^{a + 1} \to \ldots \to \mathcal{E}^b
$$
は $E$ と擬同型であり，$E$ は完全である。
\end{proof}

\begin{lemma}
\label{cohomology-lemma-perfect}
$(X, \mathcal{O}_X)$ を環付き空間とし，$E$ を
$D(\mathcal{O}_X)$ の対象とする。次は同値である。
\begin{enumerate}
\item $E$ は完全である。
\item $E$ は擬コヒーレントであり，局所的に有限 Tor 次元をもつ。
\end{enumerate}
\end{lemma}

\begin{proof}
(1) を仮定する。定義により，$E|_{U_i}$ が狭義完全複体で表現される
ような開被覆 $X=\bigcup U_i$ が存在する。したがって補題
\ref{cohomology-lemma-pseudo-coherent-independent-representative} により，
$E$ は擬コヒーレント，すなわちすべての $m$ に対して
$m$-擬コヒーレントである。さらに，有限自由加群の直和因子は
平坦なので，補題 \ref{cohomology-lemma-tor-amplitude} により
$E|_{U_i}$ は有限 Tor 次元をもつ。ゆえに (2) が成り立つ。

\medskip\noindent
(2) を仮定する。$X$ を開被覆の各要素で置き換えることにより，
$E$ が $[a,b]$ に Tor 振幅をもつような整数 $a\leq b$ が存在すると
仮定してよい。$E$ はすべての $m$ に対して $m$-擬コヒーレントなので，
補題 \ref{cohomology-lemma-perfect-precise} から結論を得る。
\end{proof}

\begin{lemma}
\label{cohomology-lemma-perfect-pullback}
$f:(X,\mathcal{O}_X)\to(Y,\mathcal{O}_Y)$ を環付き空間の射とし，
$E$ を $D(\mathcal{O}_Y)$ の対象とする。
$E$ が $D(\mathcal{O}_Y)$ において完全ならば，
$Lf^*E$ は $D(\mathcal{O}_X)$ において完全である。
\end{lemma}

\begin{proof}
補題 \ref{cohomology-lemma-perfect}，\ref{cohomology-lemma-tor-amplitude-pullback}，および
\ref{cohomology-lemma-pseudo-coherent-pullback} から従う。
（別の証明として，補題 \ref{cohomology-lemma-pseudo-coherent-pullback} の
証明をそのまま用いてもよい。）
\end{proof}

\begin{lemma}
\label{cohomology-lemma-two-out-of-three-perfect}
$(X,\mathcal{O}_X)$ を環付き空間とし，$(K,L,M,f,g,h)$ を
$D(\mathcal{O}_X)$ の完全三角とする。$K,L,M$ のうち二つが完全ならば，
残る一つも完全である。
\end{lemma}

\begin{proof}
第1の証明：補題 \ref{cohomology-lemma-perfect}，
\ref{cohomology-lemma-cone-pseudo-coherent}，および
\ref{cohomology-lemma-cone-tor-amplitude} を組み合わせる。
第2の証明（概略）：$K$ と $L$ が完全であるとする。
$X$ を開被覆の各要素で置き換えることにより，$K$ と $L$ が
狭義完全複体 $\mathcal{K}^\bullet$ と $\mathcal{L}^\bullet$ で
表現されると仮定してよい。さらに $X$ を開被覆の各要素で
置き換えることにより，写像 $K\to L$ が複体の写像
$\alpha:\mathcal{K}^\bullet\to\mathcal{L}^\bullet$ で与えられると
仮定してよい。補題 \ref{cohomology-lemma-local-actual} を参照せよ。
このとき $M$ は $\alpha$ の錐と同型であり，補題
\ref{cohomology-lemma-cone} によりこの錐は狭義完全である。
\end{proof}

\begin{lemma}
\label{cohomology-lemma-tensor-perfect}
$(X,\mathcal{O}_X)$ を環付き空間とする。
$K,L$ が $D(\mathcal{O}_X)$ の完全対象ならば，
$K\otimes_{\mathcal{O}_X}^\mathbf{L}L$ も完全である。
\end{lemma}

\begin{proof}
補題 \ref{cohomology-lemma-perfect}，\ref{cohomology-lemma-tensor-pseudo-coherent}，および
\ref{cohomology-lemma-tensor-tor-amplitude} から従う。
\end{proof}

\begin{lemma}
\label{cohomology-lemma-summands-perfect}
$(X,\mathcal{O}_X)$ を環付き空間とする。
$K\oplus L$ が $D(\mathcal{O}_X)$ の完全対象ならば，
$K$ と $L$ も完全である。
\end{lemma}

\begin{proof}
補題 \ref{cohomology-lemma-perfect}，\ref{cohomology-lemma-summands-pseudo-coherent}，および
\ref{cohomology-lemma-summands-tor-amplitude} から従う。
\end{proof}

\begin{lemma}
\label{cohomology-lemma-pushforward-perfect}
$(X,\mathcal{O}_X)$ を環付き空間とし，$j:U\to X$ を開部分空間とする。
$E$ を $D(\mathcal{O}_U)$ の完全対象とし，そのコホモロジー層は
閉部分集合 $T\subset U$ に台をもち，$j(T)$ は $X$ で閉であるとする。
このとき $Rj_*E$ は $D(\mathcal{O}_X)$ の完全対象である。
\end{lemma}

\begin{proof}
完全複体であることは $X$ 上局所的な性質である。したがって，
$Rj_*E$ の $U$ および $V=X\setminus j(T)$ への制限が完全であることを
確かめれば十分である。$Rj_*E|_U=E$ は完全である。また
$E|_{U\setminus T}=0$ なので $Rj_*E|_V=0$ である。
\end{proof}

\begin{lemma}
\label{cohomology-lemma-perfect-max-open-coh-loc-free}
$(X,\mathcal{O}_X)$ を環付き空間とし，$D(\mathcal{O}_X)$ の
対象 $E$ は完全であるとする。すべての茎 $\mathcal{O}_{X,x}$ が
局所環であると仮定する。このとき集合
$$
U =
\{x \in X \mid
H^i(E)_x\text{ is a finite free }
\mathcal{O}_{X, x}\text{-module for all }i\in \mathbf{Z}\}
$$
は $X$ の開集合であり，すべての $i\in\mathbf{Z}$ に対して
$H^i(E)|_U$ が有限局所自由となる最大の開集合 $U\subset X$ である。
\end{lemma}

\begin{proof}
すべての $i\in\mathbf{Z}$ に対して $H^i(E)|_V$ が有限局所自由となる
開集合 $V\subset X$ があれば，$V\subset U$ であることに注意する。
$x\in U$ とする。$x$ のある開近傍が $U$ に含まれ，その近傍上で
すべての $i$ に対して $H^i(E)$ が有限局所自由であることを示す。
これで証明が終わる。証明中，$X$ を $x$ の開近傍で
（有限回）置き換えてよい。したがって $E$ が狭義完全複体
$\mathcal{E}^\bullet$ によって表現されると仮定してよい。
$i\not\in[a,b]$ に対して $\mathcal{E}^i=0$ とする。
$b-a$ に関する帰納法で証明する。加群
$H^b(E) = \Coker(d^{b - 1} : \mathcal{E}^{b - 1} \to \mathcal{E}^b)$
は有限表示である。$H^b(E)_x$ は有限自由なので，Modules, Lemma
\ref{modules-lemma-finite-presentation-stalk-free} により，
$x$ のある開近傍上で $H^b(E)$ は有限自由である。
そこで $X$ を（必要ならさらに小さい）開近傍で置き換え，
直和分解 $\mathcal{E}^b=\Im(d^{b - 1})\oplus H^b(E)$ が存在し，
$H^b(E)$ が有限自由であると仮定してよい。補題
\ref{cohomology-lemma-local-lift-map} を参照せよ。
同じ議論をもう一度行えば，
$\mathcal{E}^{b - 1}=\Ker(d^{b - 1})\oplus\Im(d^{b - 1})$
と仮定してよい。複体
$\mathcal{E}^a \to \ldots \to \mathcal{E}^{b - 2} \to \Ker(d^{b - 1})$
は狭義完全であり，$i\not=b$ に対して
$H^i(E)=H^i(E')$ を満たす完全対象 $E'$ を表現する。
したがって帰納法の仮定から結論を得る。
\end{proof}






\section{双対}
\label{cohomology-section-duals}

\noindent
この節では，複体の圏および導来圏の双対可能対象を特徴づける。
特に，$D(\mathcal{O}_X)$ の対象が双対をもつための必要十分条件が
完全であることだとわかる（例 \ref{cohomology-example-dual-derived} と補題
\ref{cohomology-lemma-left-dual-derived} から従う）。

\begin{lemma}
\label{cohomology-lemma-symmetric-monoidal-cat-complexes}
$(X,\mathcal{O}_X)$ を環付き空間とする。
$\mathcal{O}_X$-加群の複体の圏に，テンソル積を
$\mathcal{F}^\bullet \otimes \mathcal{G}^\bullet =
\text{Tot}(\mathcal{F}^\bullet \otimes_{\mathcal{O}_X} \mathcal{G}^\bullet)$
で定めると，対称モノイド圏になる（符号規則については
More on Algebra, Section \ref{more-algebra-section-sign-rules} を参照せよ）。
\end{lemma}

\begin{proof}
省略する。ヒント：単位 $\mathbf{1}$ として，次数 $0$ の項が
$\mathcal{O}_X$ で他の次数では零である複体をとる。このとき明らかな同型
$\text{Tot}(\mathbf{1} \otimes_{\mathcal{O}_X} \mathcal{G}^\bullet) =
\mathcal{G}^\bullet$ および
$\text{Tot}(\mathcal{F}^\bullet \otimes_{\mathcal{O}_X} \mathbf{1}) =
\mathcal{F}^\bullet$.
がある。補題を証明するには種々の図式の可換性を確かめる必要がある。
Categories, Definitions \ref{categories-definition-monoidal-category} および
\ref{categories-definition-symmetric-monoidal-category} を参照せよ。
いずれの確認も直接的である。
\end{proof}

\begin{example}
\label{cohomology-example-dual}
$(X,\mathcal{O}_X)$ を環付き空間とする。$\mathcal{F}^\bullet$ を
$\mathcal{O}_X$-加群の局所有界な複体とし，各 $\mathcal{F}^n$ は
局所的に有限自由 $\mathcal{O}_X$-加群の直和因子であるとする。
言い換えれば，$\mathcal{F}^\bullet|_{U_i}$ が狭義完全となるような
開被覆 $X=\bigcup U_i$ が存在する。Section
\ref{cohomology-section-hom-complexes} のように複体
$$
\mathcal{G}^\bullet = \SheafHom^\bullet(\mathcal{F}^\bullet, \mathcal{O}_X)
$$
を考える。
$$
\eta :
\mathcal{O}_X
\to
\text{Tot}(\mathcal{F}^\bullet \otimes_{\mathcal{O}_X} \mathcal{G}^\bullet)
\quad\text{and}\quad
\epsilon :
\text{Tot}(\mathcal{G}^\bullet \otimes_{\mathcal{O}_X} \mathcal{F}^\bullet)
\to
\mathcal{O}_X
$$
を $\eta=\sum\eta_n$ および $\epsilon=\sum\epsilon_n$ で定める。
ここで $\eta_n : \mathcal{O}_X \to
\mathcal{F}^n \otimes_{\mathcal{O}_X} \mathcal{G}^{-n}$
および
$\epsilon_n : \mathcal{G}^{-n} \otimes_{\mathcal{O}_X} \mathcal{F}^n
\to \mathcal{O}_X$ は Modules, Example \ref{modules-example-dual} の
写像である。このとき Categories, Definition
\ref{categories-definition-dual} の意味で，
$\mathcal{G}^\bullet,\eta,\epsilon$ は $\mathcal{F}^\bullet$ の左双対である。
$(1 \otimes \epsilon) \circ (\eta \otimes 1) = \text{id}_{\mathcal{F}^\bullet}$
および
$(\epsilon \otimes 1) \circ (1 \otimes \eta) =
\text{id}_{\mathcal{G}^\bullet}$ の確認は省略する。
More on Algebra, Lemma \ref{more-algebra-lemma-left-dual-complex} と比較せよ。
\end{example}

\begin{lemma}
\label{cohomology-lemma-left-dual-complex}
$(X,\mathcal{O}_X)$ を環付き空間とし，$\mathcal{F}^\bullet$ を
$\mathcal{O}_X$-加群の複体とする。$\mathcal{F}^\bullet$ が
$\mathcal{O}_X$-加群の複体のモノイド圏において左双対をもつならば
（Categories, Definition \ref{categories-definition-dual}），
$\mathcal{F}^\bullet$ は，各項が局所的に有限自由
$\mathcal{O}_X$-加群の直和因子である局所有界複体であり，
その左双対は例 \ref{cohomology-example-dual} で構成したものである。
\end{lemma}

\begin{proof}
左双対の一意性
（Categories, Remark \ref{categories-remark-left-dual-adjoint}）により，
$\mathcal{F}^\bullet$ が主張どおりであることを示せば最後の主張も従う。
$\mathcal{G}^\bullet,\eta,\epsilon$ を左双対とする。
$\eta=\sum\eta_n$，$\epsilon=\sum\epsilon_n$ と書く。ここで
$\eta_n : \mathcal{O}_X \to
\mathcal{F}^n \otimes_{\mathcal{O}_X} \mathcal{G}^{-n}$
および
$\epsilon_n : \mathcal{G}^{-n} \otimes_{\mathcal{O}_X} \mathcal{F}^n
\to \mathcal{O}_X$ である。
$(1 \otimes \epsilon) \circ (\eta \otimes 1) = \text{id}_{\mathcal{F}^\bullet}$
および
$(\epsilon \otimes 1) \circ (1 \otimes \eta) = \text{id}_{\mathcal{G}^\bullet}$
であるから，Categories, Definition
\ref{categories-definition-dual} により直ちに
$(1 \otimes \epsilon_n) \circ (\eta_n \otimes 1) = \text{id}_{\mathcal{F}^n}$
および
$(\epsilon_n \otimes 1) \circ (1 \otimes \eta_n) =
\text{id}_{\mathcal{G}^{-n}}$ を得る。したがって Modules, Lemma
\ref{modules-lemma-left-dual-module} により，$\mathcal{F}^n$ は
局所的に有限自由 $\mathcal{O}_X$-加群の直和因子である。
和 $\eta=\sum\eta_n$ は局所有限なので，$\mathcal{F}^\bullet$ は
局所有界である。
\end{proof}

\begin{lemma}
\label{cohomology-lemma-internal-hom-evaluate-isom}
$(X,\mathcal{O}_X)$ を環付き空間とし，
$K,L,M\in D(\mathcal{O}_X)$ とする。$K$ が完全ならば，補題
\ref{cohomology-lemma-internal-hom-evaluate} の写像
$$
R\SheafHom(L, M) \otimes_{\mathcal{O}_X}^\mathbf{L} K
\longrightarrow
R\SheafHom(R\SheafHom(K, L), M)
$$
は同型である。
\end{lemma}

\begin{proof}
この写像は大域的に定義され，右辺と左辺の構成は局所化と可換するので
（補題 \ref{cohomology-lemma-restriction-RHom-to-U} を参照せよ），
$X$ 上局所的に証明してよい。したがって $K$ は狭義完全複体
$\mathcal{E}^\bullet$ で表現されると仮定してよい。

\medskip\noindent
$K_1\to K_2\to K_3$ が $D(\mathcal{O}_X)$ の完全三角ならば，
完全三角
$$
R\SheafHom(L, M) \otimes_{\mathcal{O}_X}^\mathbf{L} K_1 \to
R\SheafHom(L, M) \otimes_{\mathcal{O}_X}^\mathbf{L} K_2 \to
R\SheafHom(L, M) \otimes_{\mathcal{O}_X}^\mathbf{L} K_3
$$
および
$$
R\SheafHom(R\SheafHom(K_1, L), M) \to
R\SheafHom(R\SheafHom(K_2, L), M)
R\SheafHom(R\SheafHom(K_3, L), M)
$$
を得る。Section \ref{cohomology-section-flat} および補題
\ref{cohomology-lemma-RHom-triangulated} を参照せよ。
補題 \ref{cohomology-lemma-internal-hom-evaluate} の矢印は $K$ に関して関手的なので，
これらの完全三角の間の射を得る。したがって $K_1$ と $K_3$ に対して
結論が成り立つならば，Derived Categories, Lemma
\ref{derived-lemma-third-isomorphism-triangle} により
$K_2$ に対しても成り立つ。

\medskip\noindent
上の注意と，愚直切断の完全三角
$$
\sigma_{\geq n}\mathcal{E}^\bullet \to \mathcal{E}^\bullet \to
\sigma_{\leq n - 1}\mathcal{E}^\bullet
$$
を組み合わせると，$K$ がある次数に置かれた有限自由
$\mathcal{O}_X$-加群の直和因子からなる場合へ帰着する。
この問題が直和（上で述べたことと同様）およびシフトと明らかに
両立することから，ある整数 $n$ に対して
$K=\mathcal{O}_X^{\oplus n}$ である場合へ帰着する。
この場合は明らかである。
\end{proof}

\begin{lemma}
\label{cohomology-lemma-dual-perfect-complex}
$(X,\mathcal{O}_X)$ を環付き空間とし，$K$ を
$D(\mathcal{O}_X)$ の完全対象とする。このとき
$K^\vee=R\SheafHom(K,\mathcal{O}_X)$ も完全対象であり，
$(K^\vee)^\vee\cong K$ である。$D(\mathcal{O}_X)$ の $M$ に対して，
関手的な同型
$$
M \otimes^\mathbf{L}_{\mathcal{O}_X} K^\vee = R\SheafHom(K, M)
$$
および
$$
H^0(X, M \otimes^\mathbf{L}_{\mathcal{O}_X} K^\vee) =
\Hom_{D(\mathcal{O}_X)}(K, M)
$$
が存在する。
\end{lemma}

\begin{proof}
補題 \ref{cohomology-lemma-internal-hom-evaluate} により標準写像
$$
K = R\SheafHom(\mathcal{O}_X, \mathcal{O}_X)
\otimes_{\mathcal{O}_X}^\mathbf{L} K \longrightarrow
R\SheafHom(R\SheafHom(K, \mathcal{O}_X), \mathcal{O}_X) =
(K^\vee)^\vee
$$
が存在し，これは補題 \ref{cohomology-lemma-internal-hom-evaluate-isom} により同型である。
他の主張を確認する際，内部 Hom の構成が開集合への制限と可換すること
（補題 \ref{cohomology-lemma-restriction-RHom-to-U}）を断りなく用いる。
$K^\vee$ が完全であることは $X$ 上局所的に確認できる。
補題 \ref{cohomology-lemma-dual} により，最後の主張を示すには写像
(\ref{cohomology-equation-eval})
$$
M \otimes^\mathbf{L}_{\mathcal{O}_X} K^\vee
\longrightarrow
R\SheafHom(K, M)
$$
が同型であることを確かめれば十分である。これも $X$ 上局所的な問題である。
したがって，$K$ が狭義完全複体で表現される場合に，これら二つの主張を
証明すれば十分である。

\medskip\noindent
$K$ が狭義完全複体 $\mathcal{E}^\bullet$ で表現されると仮定する。
補題 \ref{cohomology-lemma-Rhom-strictly-perfect} により，$K^\vee$ は
次数 $n$ の項が
$(\mathcal{E}^{-n})^\vee =
\SheafHom_{\mathcal{O}_X}(\mathcal{E}^{-n}, \mathcal{O}_X)$
である複体によって表現される。
$\mathcal{E}^{-n}$ は有限自由 $\mathcal{O}_X$-加群の直和因子なので，
$(\mathcal{E}^{-n})^\vee$ もそうである。したがって $K^\vee$ も
狭義完全複体で表現され，完全である。
(\ref{cohomology-equation-eval}) が同型であることを示すため，
$M$ を複体 $\mathcal{F}^\bullet$ で表現する。補題
\ref{cohomology-lemma-Rhom-strictly-perfect} により，
$R\SheafHom(K,M)$ は各項が
$$
\bigoplus\nolimits_{n = p + q}
\SheafHom_{\mathcal{O}_X}(\mathcal{E}^{-q}, \mathcal{F}^p)
$$
である複体によって表現される。他方，対象
$M\otimes^\mathbf{L}_{\mathcal{O}_X}K^\vee$ は各項が
$$
\bigoplus\nolimits_{n = p + q}
\mathcal{F}^p \otimes_{\mathcal{O}_X} (\mathcal{E}^{-q})^\vee
$$
である複体によって表現される。したがって
(\ref{cohomology-equation-eval}) が同型であるという主張は，標準写像
$$
\mathcal{F}
\otimes_{\mathcal{O}_X}
\SheafHom_{\mathcal{O}_X}(\mathcal{E}, \mathcal{O}_X)
\longrightarrow
\SheafHom_{\mathcal{O}_X}(\mathcal{E}, \mathcal{F})
$$
が，$\mathcal{E}$ が有限自由 $\mathcal{O}_X$-加群の直和因子で，
$\mathcal{F}$ が任意の $\mathcal{O}_X$-加群であるとき同型である，
という主張へ帰着する。これは $\mathcal{E}$ が有限自由である場合の
対応する主張から直ちに従う。
\end{proof}

\begin{lemma}
\label{cohomology-lemma-symmetric-monoidal-derived}
$(X,\mathcal{O}_X)$ を環付き空間とする。導来圏
$D(\mathcal{O}_X)$ は，テンソル積を導来テンソル積とし，
通常の結合制約と可換制約を備えた対称モノイド圏である
（符号規則については More on Algebra, Section
\ref{more-algebra-section-sign-rules} を参照せよ）。
\end{lemma}

\begin{proof}
省略する。補題 \ref{cohomology-lemma-symmetric-monoidal-cat-complexes} と比較せよ。
\end{proof}

\begin{example}
\label{cohomology-example-dual-derived}
$(X,\mathcal{O}_X)$ を環付き空間とし，$K$ を
$D(\mathcal{O}_X)$ の完全対象とする。補題
\ref{cohomology-lemma-dual-perfect-complex} のように
$K^\vee=R\SheafHom(K,\mathcal{O}_X)$ とおく。このとき写像
$$
K \otimes_{\mathcal{O}_X}^\mathbf{L} K^\vee \longrightarrow R\SheafHom(K, K)
$$
は（同補題により）同型である。写像
$$
\eta :
\mathcal{O}_X
\longrightarrow
K \otimes_{\mathcal{O}_X}^\mathbf{L} K^\vee
$$
を，上の同型のもとで $1$ を $\text{id}_K$ に対応する切断へ送る
写像とする。また
$$
\epsilon : 
K^\vee
\otimes_{\mathcal{O}_X}^\mathbf{L} K
\longrightarrow
\mathcal{O}_X
$$
を評価写像とする（たとえばその構成には補題
\ref{cohomology-lemma-internal-hom-composition} を用いられる）。
このとき Categories, Definition \ref{categories-definition-dual} の意味で，
$K^\vee,\eta,\epsilon$ は $K$ の左双対である。
$(1 \otimes \epsilon) \circ (\eta \otimes 1) = \text{id}_K$
および
$(\epsilon \otimes 1) \circ (1 \otimes \eta) =
\text{id}_{K^\vee}$ の確認は省略する。
\end{example}

\begin{lemma}
\label{cohomology-lemma-left-dual-derived}
$(X,\mathcal{O}_X)$ を環付き空間とし，$M$ を
$D(\mathcal{O}_X)$ の対象とする。モノイド圏 $D(\mathcal{O}_X)$ において
$M$ が左双対をもつならば
（Categories, Definition \ref{categories-definition-dual}），
$M$ は完全であり，その左双対は例 \ref{cohomology-example-dual-derived} で
構成したものである。
\end{lemma}

\begin{proof}
$x\in X$ とする。$M$ の $U$ への制限が完全複体になるような
$x$ の開近傍 $U$ を見つければ十分である。したがって証明中，
$X$ を $x$ の開近傍で（有限回）置き換えてよい。
$N,\eta,\epsilon$ を左双対とする。

\medskip\noindent
次の議論を何度か用いる。$M$ を表現する
$\mathcal{O}_X$-加群の任意の複体 $\mathcal{M}^\bullet$ を選ぶ。
$N$ を表現し，各項が平坦 $\mathcal{O}_X$-加群である
K-平坦複体 $\mathcal{N}^\bullet$ を選ぶ。補題
\ref{cohomology-lemma-K-flat-resolution} を参照せよ。写像
$$
\eta : \mathcal{O}_X \to
\text{Tot}(\mathcal{M}^\bullet \otimes_{\mathcal{O}_X} \mathcal{N}^\bullet)
$$
を考える。$X$ を縮小すれば，ある整数 $N$ と，
$i=1,\ldots,N$ に対する整数 $n_i\in\mathbf{Z}$，ならびに
$\mathcal{M}^{n_i}$ と $\mathcal{N}^{-n_i}$ の切断 $f_i$ と $g_i$ で
$$
\eta(1) = \sum\nolimits_i f_i \otimes g_i
$$
を満たすものを見つけられる。$\mathcal{K}^\bullet\subset
\mathcal{M}^\bullet$ を，$i=1,\ldots,N$ に対して切断 $f_i$ を含む，
$\mathcal{O}_X$-加群の任意の部分複体とする。
$\mathcal{N}^n$ の平坦性により
$\text{Tot}(\mathcal{K}^\bullet \otimes_{\mathcal{O}_X} \mathcal{N}^\bullet)
\subset
\text{Tot}(\mathcal{M}^\bullet \otimes_{\mathcal{O}_X} \mathcal{N}^\bullet)$
なので，$\eta$ は
$$
\tilde \eta :
\mathcal{O}_X \to
\text{Tot}(\mathcal{K}^\bullet \otimes_{\mathcal{O}_X} \mathcal{N}^\bullet)
$$
を経由する。$\mathcal{K}^\bullet$ が表現する $D(\mathcal{O}_X)$ の
対象を $K$ とすると，可換図式
$$
\xymatrix{
M \ar[rr]_-{\eta \otimes 1} \ar[rrd]_{\tilde \eta \otimes 1} & &
M \otimes^\mathbf{L} N \otimes^\mathbf{L} M
\ar[r]_-{1 \otimes \epsilon} &
M \\
& &
K \otimes^\mathbf{L} N \otimes^\mathbf{L} M
\ar[u] \ar[r]^-{1 \otimes \epsilon} &
K \ar[u]
}
$$
を得る。上の行の合成は $M$ 上の恒等写像なので，
$D(\mathcal{O}_X)$ において $M$ は $K$ の直和因子である。

\medskip\noindent
上の議論の第1の応用として，すべての $i=1,\ldots,N$ に対して
$a<n_i<b$ となるように部分複体
$\mathcal{K}^\bullet=\sigma_{\geq a}\tau_{\leq b}\mathcal{M}^\bullet$
を選べる。したがって $M$ は $D(\mathcal{O}_X)$ において
有界複体の直和因子であり，$M\in D^b(\mathcal{O}_X)$ と仮定してよい。
（上の手続きでは $X$ を縮小していることを思い出そう。）

\medskip\noindent
$M\in D^b(\mathcal{O}_X)$ なので，$\mathcal{M}^\bullet$ を
平坦加群からなる上に有界な複体に選べる
（Modules, Lemma \ref{modules-lemma-module-quotient-flat} および
Derived Categories, Lemma \ref{derived-lemma-subcategory-left-resolution}）。
すると上の議論で，すべての $i=1,\ldots,N$ に対して $a<n_i$ となるように
$\mathcal{K}^\bullet=\sigma_{\geq a}\mathcal{M}^\bullet$ を選べる。
したがって $M$ は $D(\mathcal{O}_X)$ において，平坦加群からなる
有界複体の直和因子であると仮定してよい。特に $M$ は有限 Tor 振幅をもつ。

\medskip\noindent
$M$ が $[a,b]$ に Tor 振幅をもつとする。$M$ が
$m$-擬コヒーレントであると仮定し，$X$ を縮小した後，
$M$ が $(m-1)$-擬コヒーレントであると仮定できることを示す。
$M$ はいずれにせよ $(b+1)$-擬コヒーレントなので，補題
\ref{cohomology-lemma-perfect-precise} により，これで証明が終わる。
$X$ を縮小すれば，狭義完全複体 $\mathcal{E}^\bullet$ と
$D(\mathcal{O}_X)$ における写像 $\alpha:\mathcal{E}^\bullet\to M$ で，
$i>m$ に対して $H^i(\alpha)$ が同型，$i=m$ に対して全射となるものが
存在すると仮定してよい。また $i<m$ に対して
$\mathcal{E}^i=0$ と仮定する。完全三角
$$
\mathcal{E}^\bullet \to M \to L \to \mathcal{E}^\bullet[1]
$$
を選ぶ。$i\geq m$ に対して $H^i(L)=0$ であることに注意する。
したがって $L$ は，$i\geq m$ に対して $\mathcal{L}^i=0$ となる
複体 $\mathcal{L}^\bullet$ で表現できる。写像
$L\to\mathcal{E}^\bullet[1]$ は複体の写像
$\mathcal{L}^\bullet \to \mathcal{E}^\bullet[1]$
で与えられ，これは次数 $m-1$ を除くすべての次数で零であり，
次数 $m-1$ では写像 $\mathcal{L}^{m - 1}\to\mathcal{E}^m$ を与える。
Derived Categories, Lemma \ref{derived-lemma-negative-exts} を参照せよ。
このとき $M$ は複体
$$
\mathcal{M}^\bullet :
\ldots \to
\mathcal{L}^{m - 2} \to
\mathcal{L}^{m - 1} \to
\mathcal{E}^m \to
\mathcal{E}^{m + 1} \to \ldots
$$
によって表現される。この複体に第2段落の議論を適用し，
$i=1,\ldots,N$ に対する $\mathcal{M}^{n_i}$ の切断 $f_i$ を得る。
$n<m$ に対して，$\mathcal{K}^n\subset\mathcal{L}^n$ を，
$n_i=n$ である $f_i$ と $n_i=n-1$ である $d(f_i)$ により生成される
$\mathcal{O}_X$-部分加群とする。$n\geq m$ に対しては
$\mathcal{K}^n=\mathcal{E}^n$ とおく。明らかに完全三角の射
$$
\xymatrix{
\mathcal{E}^\bullet \ar[r] &
\mathcal{M}^\bullet \ar[r] &
\mathcal{L}^\bullet \ar[r] &
\mathcal{E}^\bullet[1] \\
\mathcal{E}^\bullet \ar[r] \ar[u] &
\mathcal{K}^\bullet \ar[r] \ar[u] &
\sigma_{\leq m - 1}\mathcal{K}^\bullet \ar[r] \ar[u] &
\mathcal{E}^\bullet[1] \ar[u]
}
$$
を得る。すべての射は上で示したものである。
複体 $\mathcal{K}^\bullet$ に対応する $D(\mathcal{O}_X)$ の対象を
$K$ とする。証明の第2段落の議論により，
合成 $M\to K\to M$ が $M$ 上の恒等写像になるような
$D(\mathcal{O}_X)$ の射 $s:M\to K$ を得る。図式
$$
\xymatrix{
\mathcal{E}^\bullet \ar[r] &
\mathcal{K}^\bullet \ar@{=}[r] &
K \\
\mathcal{E}^\bullet \ar[u]^{\text{id}} \ar[r]^i &
\mathcal{M}^\bullet \ar@{=}[r] &
M \ar[u]_s
}
$$
が可換であるかどうかはわからないが，写像 $K\to M$ と合成した後には
可換である。補題 \ref{cohomology-lemma-local-actual} により，$X$ を縮小すれば，
$s\circ i$ が複体の写像
$\sigma:\mathcal{E}^\bullet\to\mathcal{K}^\bullet$ で与えられると
仮定してよい。同じ補題により，$\sigma$ と包含
$\mathcal{K}^\bullet\subset\mathcal{M}^\bullet$ の合成が，
あるホモトピー
$\{h^i:\mathcal{E}^i\to\mathcal{M}^{i-1}\}$ によって
零とホモトピックであると仮定してよい。
そこで $\mathcal{K}^{m-1}$ を
$\mathcal{K}^{m-1}+\Im(h^m)$ で置き換えると
（この置き換え後も $\mathcal{K}^{m-1}$ は有限個の大域切断で
生成されていることに注意する），$\sigma$ 自身が零とホモトピックになる。
これは，可換な実線図式
$$
\xymatrix{
\mathcal{E}^\bullet \ar[r] &
M \ar[r] &
\mathcal{L}^\bullet \ar[r] &
\mathcal{E}^\bullet[1] \\
\mathcal{E}^\bullet \ar[r] \ar[u] &
K \ar[r] \ar[u] &
\sigma_{\leq m - 1}\mathcal{K}^\bullet \ar[r] \ar[u] &
\mathcal{E}^\bullet[1] \ar[u] \\
\mathcal{E}^\bullet \ar[r] \ar[u] &
M \ar[r] \ar[u]^s &
\mathcal{L}^\bullet \ar[r] \ar@{..>}[u] &
\mathcal{E}^\bullet[1] \ar[u]
}
$$
を得ることを意味する。三角圏の公理により，この図式に適合する
点線の矢印を得る。次数 $m-1$ のコホモロジー層を見ると，
$$
\xymatrix{
H^{m - 1}(M) \ar[r] &
H^{m - 1}(\mathcal{L}^\bullet) \ar[r] &
H^m(\mathcal{E}^\bullet) \\
H^{m - 1}(K) \ar[r] \ar[u] &
H^{m - 1}(\sigma_{\leq m - 1}\mathcal{K}^\bullet) \ar[r] \ar[u] &
H^m(\mathcal{E}^\bullet) \ar[u] \\
H^{m - 1}(M) \ar[r] \ar[u] &
H^{m - 1}(\mathcal{L}^\bullet) \ar[r] \ar[u] &
H^m(\mathcal{E}^\bullet) \ar[u]
}
$$
を得る。左列と右列では縦の合成がいずれも恒等写像なので，
中央の縦の合成
$H^{m - 1}(\mathcal{L}^\bullet) \to
H^{m - 1}(\sigma_{\leq m - 1}\mathcal{K}^\bullet) \to
H^{m - 1}(\mathcal{L}^\bullet)$ は全射である。特に
$H^{m - 1}(\sigma_{\leq m - 1}\mathcal{K}^\bullet)\to
H^{m - 1}(\mathcal{L}^\bullet)$ は全射である。
上の完全三角の射がコホモロジー長完全列の間に誘導する写像を用いると，
図式追跡により，$i\geq m$ に対して $H^i(K)\to H^i(M)$ が同型，
$i=m-1$ に対して全射であることが従う。
構成により，ある $r\geq0$ と全射
$\mathcal{O}_X^{\oplus r}\to\mathcal{K}^{m-1}$ を選べる。このとき合成
$$
(\mathcal{O}_X^{\oplus r} \to \mathcal{E}^m \to
\mathcal{E}^{m + 1} \to \ldots ) \longrightarrow
K \longrightarrow M
$$
は次数 $\geq m$ のコホモロジー層上で同型を，次数 $m-1$ で全射を
誘導する。これで証明は完了する。
\end{proof}





\section{その他}
\label{cohomology-section-misc}

\noindent
他のどこにも収まりにくいいくつかの結果を述べる。

\begin{lemma}
\label{cohomology-lemma-colim-and-lim-of-duals}
$(X,\mathcal{O}_X)$ を環付き空間とし，
$(K_n)_{n\in\mathbf{N}}$ を $D(\mathcal{O}_X)$ の完全対象の系とする。
$K=\text{hocolim}K_n$ を導来余極限とする
（Derived Categories, Definition
\ref{derived-definition-derived-colimit}）。このとき
$D(\mathcal{O}_X)$ の任意の対象 $E$ に対して
$$
R\SheafHom(K, E) = R\lim E \otimes^\mathbf{L}_{\mathcal{O}_X} K_n^\vee
$$
である。ここで $(K_n^\vee)$ は双対完全複体の逆系である。
\end{lemma}

\begin{proof}
補題 \ref{cohomology-lemma-dual-perfect-complex} により
$R\lim E \otimes^\mathbf{L}_{\mathcal{O}_X} K_n^\vee =
R\lim R\SheafHom(K_n, E)$
であり，これは完全三角
$$
R\lim R\SheafHom(K_n, E) \to
\prod R\SheafHom(K_n, E) \to
\prod R\SheafHom(K_n, E)
$$
に組み込まれる。同様に $K$ は完全三角
$\bigoplus K_n\to\bigoplus K_n\to K$ に組み込まれるので，
$\prod R\SheafHom(K_n,E)=R\SheafHom(\bigoplus K_n,E)$ を
示せば十分である。これは (\ref{cohomology-equation-internal-hom}) と，
導来テンソル積が直和と可換することから形式的に従う。
\end{proof}

\begin{lemma}
\label{cohomology-lemma-ext-composition-is-cup}
$(X,\mathcal{O}_X)$ を環付き空間とする。$K$ と $E$ を
$D(\mathcal{O}_X)$ の対象とし，$E$ は完全であるとする。図式
$$
\xymatrix{
H^0(X, K \otimes_{\mathcal{O}_X}^\mathbf{L} E^\vee) \times H^0(X, E)
\ar[r] \ar[d] &
H^0(X, K \otimes_{\mathcal{O}_X}^\mathbf{L} E^\vee
\otimes_{\mathcal{O}_X}^\mathbf{L} E) \ar[d] \\
\Hom_X(E, K) \times H^0(X, E) \ar[r] &
H^0(X, K)
}
$$
は可換である。ここで上の水平矢印はカップ積，右の垂直矢印は
$\epsilon : E^\vee \otimes_{\mathcal{O}_X}^\mathbf{L} E \to \mathcal{O}_X$
（例 \ref{cohomology-example-dual-derived}）を用い，左の垂直矢印は補題
\ref{cohomology-lemma-dual-perfect-complex} を用いる。下の水平矢印は明らかなものである。
\end{lemma}

\begin{proof}
$\otimes=\otimes_{\mathcal{O}_X}^\mathbf{L}$ および
$\mathcal{O}=\mathcal{O}_X$ と略記する。$E$ と $K$ をそれぞれ
$R\SheafHom(\mathcal{O},E)$ と $R\SheafHom(\mathcal{O},K)$ と同一視し，
$E^\vee$ を $R\SheafHom(E,\mathcal{O})$ と同一視する。

\medskip\noindent
$\xi\in H^0(X,K\otimes E^\vee)$ および $\eta\in H^0(X,E)$ とする。
$D(\mathcal{O})$ における対応する写像を
$\tilde\xi:\mathcal{O}\to K\otimes E^\vee$ および
$\tilde\eta:\mathcal{O}\to E$ と書く。補題
\ref{cohomology-lemma-second-cup-equals-first} により，カップ積
$\xi\cup\eta$ は
$\tilde\xi\otimes\tilde\eta:\mathcal{O}\to
K\otimes E^\vee\otimes E$ に対応する。

\medskip\noindent
補題 \ref{cohomology-lemma-dual-perfect-complex} によって $\xi$ に対応する写像
$\xi':E\to K$ は合成
$$
E = \mathcal{O} \otimes E
\xrightarrow{\tilde \xi \otimes 1_E}
K \otimes E^\vee \otimes E
\xrightarrow{1_K \otimes \epsilon}
K
$$
であると主張する。補題 \ref{cohomology-lemma-dual-perfect-complex} の構成は
評価写像 (\ref{cohomology-equation-eval}) を用いる。この評価写像は，
$E$ と $R\SheafHom(\mathcal{O},E)$ の同一視，および補題
\ref{cohomology-lemma-internal-hom-composition} で構成した合成
$\underline{\circ}$ を用いて構成される。したがって $\xi'$ は合成
\begin{align*}
E = \mathcal{O} \otimes
R\SheafHom(\mathcal{O}, E)
& \xrightarrow{\tilde \xi \otimes 1}
R\SheafHom(\mathcal{O}, K) \otimes
R\SheafHom(E, \mathcal{O}) \otimes
R\SheafHom(\mathcal{O}, E) \\
& \xrightarrow{\underline{\circ} \otimes 1}
R\SheafHom(E, K) \otimes R\SheafHom(\mathcal{O}, E) \\
& \xrightarrow{\underline{\circ}}
R\SheafHom(\mathcal{O}, K) = K
\end{align*}
である。補題 \ref{cohomology-lemma-internal-hom-composition} で構成した合成
$\underline{\circ}$ が結合的であること
（将来の参照をここに挿入）と，例 \ref{cohomology-example-dual-derived} において
$\epsilon$ が合成
$\underline{\circ}:E^\vee\otimes E\to\mathcal{O}$ として
定義されていることから，主張は直ちに従う。

\medskip\noindent
前の二段落の結果を用いると，補題の主張は
$(1_K \otimes \epsilon) \circ (\tilde \xi \otimes \tilde \eta)$
と
$(1_K \otimes \epsilon) \circ (\tilde \xi \otimes 1_E)
\circ (1_\mathcal{O} \otimes \tilde \eta)$
が等しいということであり，これは明らかである。
\end{proof}

\begin{lemma}
\label{cohomology-lemma-pullback-internal-hom}
$h:X\to Y$ を環付き空間の射とし，$K,M$ を
$D(\mathcal{O}_Y)$ の対象とする。Remark
\ref{cohomology-remark-prepare-fancy-base-change} の標準写像
$$
Lh^*R\SheafHom(K, M) \longrightarrow R\SheafHom(Lh^*K, Lh^*M)
$$
は次の場合に同型である。
\begin{enumerate}
\item $K$ は完全である。
\item $h$ は平坦で，$K$ は擬コヒーレントであり，
$M$ は（局所的に）下に有界である。
\item $\mathcal{O}_X$ は $h^{-1}\mathcal{O}_Y$ 上有限 Tor 次元をもち，
$K$ は擬コヒーレントで，$M$ は（局所的に）下に有界である。
\end{enumerate}
\end{lemma}

\begin{proof}
(1) を証明する。問題は $Y$ 上局所的なので，$K$ が狭義完全複体
$\mathcal{E}^\bullet$ で表現されると仮定してよい。
Section \ref{cohomology-section-perfect} を参照せよ。
$M$ を表現する K-平坦複体 $\mathcal{F}^\bullet$ を選ぶ。
補題 \ref{cohomology-lemma-Rhom-strictly-perfect} により，
$R\SheafHom(K,L)$ は複体
$\mathcal{H}^\bullet =
\SheafHom^\bullet(\mathcal{E}^\bullet, \mathcal{F}^\bullet)$
によって表現される。ここで各項は
$\mathcal{H}^n=\bigoplus\nolimits_{n=p+q}
\SheafHom_{\mathcal{O}_X}(\mathcal{E}^{-q},\mathcal{F}^p)$ である。
Section \ref{cohomology-section-derived-pullback} における $Lh^*$ の構成から，
$Lh^*K$ は狭義完全複体 $h^*\mathcal{E}^\bullet$ で表現される
（補題 \ref{cohomology-lemma-strictly-perfect-pullback}）。
同様に対象 $Lh^*M$ は複体 $h^*\mathcal{F}^\bullet$ で表現される。
最後に，補題 \ref{cohomology-lemma-Rhom-strictly-perfect-K-flat} により
$\mathcal{H}^\bullet$ は K-平坦なので，対象
$Lh^*R\SheafHom(K,M)$ は $h^*\mathcal{H}^\bullet$ で表現される。
したがって証明を終えるには
$h^*\mathcal{H}^\bullet =
\SheafHom^\bullet(h^*\mathcal{E}^\bullet, h^*\mathcal{F}^\bullet)$.
を示せば十分である。そのためには，$\mathcal{E}$ が有限自由
$\mathcal{O}_X$-加群の直和因子であるとき
$h^*\SheafHom(\mathcal{E}, \mathcal{F}) =
\SheafHom(h^*\mathcal{E}, \mathcal{F})$
であることに注意すればよい。

\medskip\noindent
(2) を証明する。$h$ は平坦なので，$\mathcal{O}_Y$-加群の任意の
複体に単に $h^*$ を適用して $Lh^*$ を計算できる。特に，
すべての $i\in\mathbf{Z}$ に対して
$H^i(Lh^*K)=h^*H^i(K)$ である。
$i<a$ に対して $H^i(M)=0$ とする。
$K'\to K$ を $D(\mathcal{O}_Y)$ の射で，すべての $i\geq b$ に対して
同型 $H^i(K')\to H^i(K)$ を定めるものとする。このとき対応する写像
$$
R\SheafHom(K, M) \to R\SheafHom(K', M)
$$
および
$$
R\SheafHom(Lh^*K, Lh^*M) \to R\SheafHom(Lh^*K', Lh^*M)
$$
は次数 $<a-b$ のコホモロジー層上で同型である（細部は省略する）。
したがって補題の主張にある写像が次数 $<a-b$ のコホモロジー層上で
同型を誘導することを示すには，それらの次数で $K'$ に対する結論を
証明すれば十分である。また (1) の証明と同様，問題は $Y$ 上局所的である。
そこで $K$ が狭義完全複体で表現されると仮定してよい。
Section \ref{cohomology-section-pseudo-coherent} を参照せよ。これで (1) に帰着する。

\medskip\noindent
(3) を証明する。望む消滅を得るために，$Lh^*$ のコホモロジー次元が
有界であることを用いる点を除けば，証明は (2) と同じである。
細部は省略する。
\end{proof}

\begin{lemma}
\label{cohomology-lemma-stalk-internal-hom}
$X$ を環付き空間とし，$K,M$ を $D(\mathcal{O}_X)$ の対象とする。
$x\in X$ とする。標準写像
$$
R\SheafHom(K, M)_x \longrightarrow
R\Hom_{\mathcal{O}_{X, x}}(K_x, M_x)
$$
は次の場合に同型である。
\begin{enumerate}
\item $K$ は完全である。
\item $K$ は擬コヒーレントであり，$M$ は（局所的に）下に有界である。
\end{enumerate}
\end{lemma}

\begin{proof}
$Y=\{x\}$ を，構造層が $\mathcal{O}_{X,x}$ で与えられる
一点環付き空間とする。平坦な包含射 $Y\to X$ に補題
\ref{cohomology-lemma-pullback-internal-hom} を適用する。
\end{proof}







\section{導来圏の可逆対象}
\label{cohomology-section-invertible-D-or-R}

\noindent
環付き空間の導来圏における可逆対象を特徴づける
（構造層の茎が局所環である場合と，そうでない場合の両方を扱う）。

\begin{lemma}
\label{cohomology-lemma-category-summands-finite-free}
$(X,\mathcal{O}_X)$ を環付き空間とし，
$R=\Gamma(X,\mathcal{O}_X)$ とおく。有限自由
$\mathcal{O}_X$-加群の直和因子である $\mathcal{O}_X$-加群の圏は，
有限射影 $R$-加群の圏と同値である。
\end{lemma}

\begin{proof}
有限射影 $R$-加群であることは，有限自由 $R$-加群の直和因子で
あることと同値だと注意する。同値は関手
$\mathcal{E}\mapsto\Gamma(X,\mathcal{E})$ によって与えられる。
逆関手は Modules, Lemma
\ref{modules-lemma-construct-quasi-coherent-sheaves} の構成によって与えられる。
\end{proof}

\begin{lemma}
\label{cohomology-lemma-invertible-derived}
$(X,\mathcal{O}_X)$ を環付き空間とし，$M$ を
$D(\mathcal{O}_X)$ の対象とする。次は同値である。
\begin{enumerate}
\item $M$ は $D(\mathcal{O}_X)$ において可逆である。
Categories, Definition \ref{categories-definition-invertible} を参照せよ。
\item 局所有限な直積分解
$$
\mathcal{O}_X = \prod\nolimits_{n \in \mathbf{Z}} \mathcal{O}_n
$$
が存在し，各 $n$ に対して可逆 $\mathcal{O}_n$-加群
$\mathcal{H}^n$ が存在して
（Modules, Definition \ref{modules-definition-invertible}），
$D(\mathcal{O}_X)$ において $M=\bigoplus\mathcal{H}^n[-n]$ である。
\end{enumerate}
(1) と (2) が成り立つならば，$M$ は $D(\mathcal{O}_X)$ の完全対象である。
すべての $x\in X$ に対して $\mathcal{O}_{X,x}$ が局所環ならば，
これらの条件は次の条件とも同値である。
\begin{enumerate}
\item[(3)] 開被覆 $X=\bigcup U_i$ が存在し，各 $i$ に対して，
$M|_{U_i}$ が次数 $n_i$ に置かれた可逆
$\mathcal{O}_{U_i}$-加群で表現されるような整数 $n_i$ が存在する。
\end{enumerate}
\end{lemma}

\begin{proof}
(2) を仮定する。対象 $R\SheafHom(M,\mathcal{O}_X)$ と合成写像
$$
R\SheafHom(M, \mathcal{O}_X) \otimes_{\mathcal{O}_X}^\mathbf{L} M \to
\mathcal{O}_X
$$
を考える。これが同型であることは局所的に証明してよい。したがって
$\mathcal{O}_X=\prod_{a\leq n\leq b}\mathcal{O}_n$ および
$M=\bigoplus_{a\leq n\leq b}\mathcal{H}^n[-n]$ と仮定してよい。
このとき
$$
R\SheafHom(\mathcal{H}^m, \mathcal{O}_X)
\otimes_{\mathcal{O}_X}^\mathbf{L} \mathcal{H}^n
$$
が，$n\not=m$ ならば零，$n=m$ ならば $\mathcal{O}_n$ に等しいことを
示せば十分である。$n\not=m$ の場合は，$\mathcal{O}_n$ と
$\mathcal{O}_m$ が平坦 $\mathcal{O}_X$-代数であり，
$\mathcal{O}_n\otimes_{\mathcal{O}_X}\mathcal{O}_m=0$ であることから従う。
可逆 $\mathcal{O}_X$-加群の局所構造
（Modules, Lemma \ref{modules-lemma-invertible}）を用いて局所的に
作業すれば，$n=m$ の場合の同型も直ちに従う。細部は省略する。
$D(\mathcal{O}_X)$ は対称モノイド圏なので，$M$ は可逆である。

\medskip\noindent
(1) を仮定する。(2) の記述から，$\mathcal{O}_n$ の候補は
$\SheafHom_{\mathcal{O}_X}(H^n(M),H^n(M))$ である。
これが環の層の局所有限な族で，かつ
$\mathcal{O}_X=\prod\mathcal{O}_n$ ならば，直積分解から生じる
$\mathcal{O}_X$ の冪等元を用いて，直ちに直和分解
$M=\bigoplus H^n(M)[-n]$ を得る。このことから，(2) の証明は
$X$ 上局所的に行ってよい。

\medskip\noindent
$D(\mathcal{O}_X)$ の対象 $N$ と同型
$M \otimes_{\mathcal{O}_X}^\mathbf{L} N \cong \mathcal{O}_X$.
を選ぶ。$x\in X$ とする。このとき $N$ はモノイド圏
$D(\mathcal{O}_X)$ における $M$ の左双対であり，補題
\ref{cohomology-lemma-left-dual-derived} により $M$ は完全である。
対称性から $N$ も完全である。$X$ を $x$ の開近傍で置き換えることにより，
$M$ と $N$ が狭義完全複体 $\mathcal{E}^\bullet$ と
$\mathcal{F}^\bullet$ で表現されると仮定してよい。このとき
$M\otimes_{\mathcal{O}_X}^\mathbf{L}N$ は
$\text{Tot}(\mathcal{E}^\bullet\otimes_{\mathcal{O}_X}
\mathcal{F}^\bullet)$ で表現される。$X$ をさらに縮小すれば，
互いに逆な同型
$\mathcal{O}_X \to M \otimes_{\mathcal{O}_X}^\mathbf{L} N$ と
$M \otimes_{\mathcal{O}_X}^\mathbf{L} N \to \mathcal{O}_X$
は複体の写像
$$
\alpha : \mathcal{O}_X \to
\text{Tot}(\mathcal{E}^\bullet \otimes_{\mathcal{O}_X} \mathcal{F}^\bullet)
\quad\text{and}\quad
\beta :
\text{Tot}(\mathcal{E}^\bullet \otimes_{\mathcal{O}_X} \mathcal{F}^\bullet)
\to \mathcal{O}_X
$$
で与えられると仮定してよい。補題 \ref{cohomology-lemma-local-actual} を参照せよ。
このとき複体の写像として $\beta\circ\alpha=1$ であり，
$D(\mathcal{O}_X)$ の射として $\alpha\circ\beta=1$ である。
$X$ を縮小すれば，合成 $\alpha\circ\beta$ は，成分が
$$
\theta^n :
\text{Tot}^n(\mathcal{E}^\bullet \otimes_{\mathcal{O}_X} \mathcal{F}^\bullet)
\to
\text{Tot}^{n - 1}(
\mathcal{E}^\bullet \otimes_{\mathcal{O}_X} \mathcal{F}^\bullet)
$$
であるあるホモトピー $\theta$ によって $1$ とホモトピックであると
仮定してよい。これは先ほどと同じ補題による。
$R=\Gamma(X,\mathcal{O}_X)$ とおく。補題
\ref{cohomology-lemma-category-summands-finite-free} により，次を得る。
\begin{enumerate}
\item $M^\bullet=\Gamma(X,\mathcal{E}^\bullet)$ は有限射影
$R$-加群からなる有界複体である。
\item $N^\bullet=\Gamma(X,\mathcal{F}^\bullet)$ は有限射影
$R$-加群からなる有界複体である。
\item $\alpha$ と $\beta$ は複体の写像
$a : R \to \text{Tot}(M^\bullet \otimes_R N^\bullet)$ と
$b : \text{Tot}(M^\bullet \otimes_R N^\bullet) \to R$ に対応する。
\item $\theta^n$ は写像
$h^n : \text{Tot}^n(M^\bullet \otimes_R N^\bullet) \to
\text{Tot}^{n - 1}(M^\bullet \otimes_R N^\bullet)$ に対応する。
\item $b \circ a = 1$ および $b \circ a - 1 = dh + hd$。
\end{enumerate}
したがって $M^\bullet$ と $N^\bullet$ は $D(R)$ の互いに逆な
対象を定める。More on Algebra, Lemma
\ref{more-algebra-lemma-invertible-derived} により，直積分解
$R=\prod_{a\leq n\leq b}R_n$ と可逆 $R_n$-加群 $H^n$ で，
$M^\bullet\cong\bigoplus_{a\leq n\leq b}H^n[-n]$ を満たすものを得る。
$D(R)$ におけるこの同型は，各 $H^n$ が射影 $R$-加群であるため，
複体の射
$$
\bigoplus H^n[-n] \longrightarrow M^\bullet
$$
へ持ち上げられる。これに対応して，再び補題
\ref{cohomology-lemma-category-summands-finite-free} を用いると，射
$$
\bigoplus H^n \otimes_R \mathcal{O}_X[-n] \to \mathcal{E}^\bullet
$$
を得，これは $D(\mathcal{O}_X)$ における同型である。
$\mathcal{O}_n=R_n\otimes_R\mathcal{O}_X$ とおけば，(2) が従う。

\medskip\noindent
$\mathcal{O}_X$ のすべての茎が局所環ならば，(2) と (3) の同値は
直ちに証明できる。細部は省略する。
\end{proof}






\section{コンパクト対象}
\label{cohomology-section-compact}

\noindent
この節では，環付き空間上の加群の導来圏におけるコンパクト対象を調べる。
コンパクト対象は Derived Categories, Definition
\ref{derived-definition-compact-object} で定義されている。
適切な環付き空間上では，完全対象はコンパクトである。

\begin{lemma}
\label{cohomology-lemma-when-jshriek-compact}
$X$ を環付き空間とし，$j:U\to X$ を開集合の包含とする。
次を満たす整数 $d$ が存在すれば，$\mathcal{O}_X$-加群
$j_!\mathcal{O}_U$ は $D(\mathcal{O}_X)$ のコンパクト対象である。
\begin{enumerate}
\item すべての $p>d$ に対して $H^p(U,\mathcal{F})=0$。
\item 関手 $\mathcal{F}\mapsto H^p(U,\mathcal{F})$ は直和と可換する。
\end{enumerate}
\end{lemma}

\begin{proof}
(1) と (2) を仮定する。Sheaves, Lemma
\ref{sheaves-lemma-j-shriek-modules} により
$\Hom(j_!\mathcal{O}_U, \mathcal{F}) = \mathcal{F}(U)$
なので，$D(\mathcal{O}_X)$ の $K$ に対して
$\Hom(j_!\mathcal{O}_U,K)=R\Gamma(U,K)$ である。
したがって $R\Gamma(U,-)$ が直和と可換することを示せばよい。
第1の仮定は，関手 $F=H^0(U,-)$ のコホモロジー次元が有限であることを
意味する。さらに第2の仮定から，入射的加群の任意の直和は $F$-非輪状である。
$K_i$ を $D(\mathcal{O}_X)$ の対象の族とする。$K_i$ を表現し，
各項が入射的である K-入射的複体 $I_i^\bullet$ を選ぶ。
Injectives, Theorem
\ref{injectives-theorem-K-injective-embedding-grothendieck} を参照せよ。
任意の非輪状対象からなる複体に $F$ を適用して $RF$ を計算でき
（Derived Categories, Lemma
\ref{derived-lemma-unbounded-right-derived}），かつ
$\bigoplus K_i$ は $\bigoplus I_i^\bullet$ によって表現されるので
（Injectives, Lemma \ref{injectives-lemma-derived-products}），
$R\Gamma(U,\bigoplus K_i)$ は
$\bigoplus H^0(U,I_i^\bullet)$ によって表現される。
ゆえに $R\Gamma(U,-)$ は直和と可換する。
\end{proof}

\begin{lemma}
\label{cohomology-lemma-perfect-is-compact}
$X$ を環付き空間とし，その台となる位相空間が次の性質をもつと仮定する。
\begin{enumerate}
\item $X$ は準コンパクトである。
\item 準コンパクト開部分集合からなる基底が存在する。
\item 任意の二つの準コンパクト開集合の共通部分は準コンパクトである。
\end{enumerate}
$K$ を $D(\mathcal{O}_X)$ の完全対象とする。このとき
\begin{enumerate}
\item[(a)] $K$ は次の意味で $D^+(\mathcal{O}_X)$ のコンパクト対象である。
$M=\bigoplus_{i\in I}M_i$ が下に有界ならば，
$\Hom(K,M)=\bigoplus_{i\in I}\Hom(K,M_i)$ である。
\item[(b)] $X$ のコホモロジー次元が有限，すなわち
$i>d$ に対して $H^i(X,\mathcal{F})=0$ となる $d$ が存在するならば，
$K$ は $D(\mathcal{O}_X)$ のコンパクト対象である。
\end{enumerate}
\end{lemma}

\begin{proof}
$K^\vee$ を $K$ の双対とする。補題
\ref{cohomology-lemma-dual-perfect-complex} を参照せよ。このとき
$$
\Hom_{D(\mathcal{O}_X)}(K, M) =
H^0(X, K^\vee \otimes_{\mathcal{O}_X}^\mathbf{L} M)
$$
であり，これは $D(\mathcal{O}_X)$ の $M$ に関して関手的である。
$K^\vee\otimes_{\mathcal{O}_X}^\mathbf{L}-$ は直和と可換するので，
$R\Gamma(X,-)$ が該当する直和と可換することを示せば十分である。

\medskip\noindent
(b) を証明する。$R\Gamma(X,K)=R\Hom(\mathcal{O}_X,K)$ であり，
補題 \ref{cohomology-lemma-quasi-separated-cohomology-colimit} により
$H^p(X,-)$ は直和と可換するので，これは補題
\ref{cohomology-lemma-when-jshriek-compact} の特別な場合である。

\medskip\noindent
(a) を証明する。$\mathcal{I}_i$（$i\in I$）を入射的
$\mathcal{O}_X$-加群の族とする。補題
\ref{cohomology-lemma-quasi-separated-cohomology-colimit} により
$$
H^p(X, \bigoplus\nolimits_{i \in I} \mathcal{I}_i) =
\bigoplus\nolimits_{i \in I} H^p(X, \mathcal{I}_i) = 0
$$
がすべての $p$ に対して成り立つ。(a) のように
$M=\bigoplus M_i$ とすると，$n<a$ に対して $H^n(M_i)=0$ となる
$a\in\mathbf{Z}$ が存在する。したがって，$M_i$ を表現し，
$n<a$ に対して $\mathcal{I}_i^n=0$ となる，入射的
$\mathcal{O}_X$-加群からなる複体 $\mathcal{I}_i^\bullet$ を選べる。
Derived Categories, Lemma
\ref{derived-lemma-injective-resolutions-exist} を参照せよ。
Injectives, Lemma \ref{injectives-lemma-derived-products} により，
直和複体 $\bigoplus\mathcal{I}_i^\bullet$ は $M$ を表現する。
Leray 非輪状性
（Derived Categories, Lemma \ref{derived-lemma-leray-acyclicity}）により
$$
R\Gamma(X, M) = \Gamma(X, \bigoplus \mathcal{I}_i^\bullet) =
\bigoplus \Gamma(X, \bigoplus \mathcal{I}_i^\bullet) =
\bigoplus R\Gamma(X, M_i)
$$
を得る。これが望む結論である。
\end{proof}


















\section{射影公式}
\label{cohomology-section-projection-formula}

\noindent
この節では射影公式のいくつかの変種をまとめる。
最も基本的な形は補題 \ref{cohomology-lemma-projection-formula} である。
これを述べて証明した後，完全複体を含むより一般的な形を論じる。

\begin{lemma}
\label{cohomology-lemma-injective-tensor-finite-locally-free}
$X$ を環付き空間とする。$\mathcal{I}$ を入射的
$\mathcal{O}_X$-加群，$\mathcal{E}$ を $\mathcal{O}_X$-加群とする。
$\mathcal{E}$ は $X$ 上有限局所自由であると仮定する。
Modules, Definition \ref{modules-definition-locally-free} を参照せよ。
このとき $\mathcal{E}\otimes_{\mathcal{O}_X}\mathcal{I}$ は
入射的 $\mathcal{O}_X$-加群である。
\end{lemma}

\begin{proof}
補題の仮定のもとでは
$$
\Hom_{\mathcal{O}_X}(\mathcal{F},
\mathcal{E} \otimes_{\mathcal{O}_X} \mathcal{I})
=
\Hom_{\mathcal{O}_X}(
\mathcal{F} \otimes_{\mathcal{O}_X} \mathcal{E}^\vee, \mathcal{I})
$$
である。ここで
$\mathcal{E}^\vee = \SheafHom_{\mathcal{O}_X}(\mathcal{E}, \mathcal{O}_X)$
は $\mathcal{E}$ の双対であり，これも有限局所自由である。
有限局所自由層とのテンソル積は完全関手なので，Homology, Lemma
\ref{homology-lemma-characterize-injectives} から結論を得る。
\end{proof}

\begin{lemma}
\label{cohomology-lemma-projection-formula}
$f:X\to Y$ を環付き空間の射とする。
$\mathcal{F}$ を $\mathcal{O}_X$-加群，
$\mathcal{E}$ を $\mathcal{O}_Y$-加群とする。
$\mathcal{E}$ は $Y$ 上有限局所自由であると仮定する。
Modules, Definition \ref{modules-definition-locally-free} を参照せよ。
このとき，すべての $q\geq0$ に対して同型
$$
\mathcal{E} \otimes_{\mathcal{O}_Y} R^qf_*\mathcal{F}
\longrightarrow
R^qf_*(f^*\mathcal{E} \otimes_{\mathcal{O}_X} \mathcal{F})
$$
が存在する。実際，$D^{+}(Y)$ に同型
$$
\mathcal{E} \otimes_{\mathcal{O}_Y} Rf_*\mathcal{F}
\longrightarrow
Rf_*(f^*\mathcal{E} \otimes_{\mathcal{O}_X} \mathcal{F})
$$
が存在し，これは $\mathcal{F}$ に関して関手的である。
\end{lemma}

\begin{proof}
$X$ 上の入射分解 $\mathcal{F}\to\mathcal{I}^\bullet$ を選ぶ。
$f^*\mathcal{E}$ も有限局所自由なので，分解
$$
f^*\mathcal{E} \otimes_{\mathcal{O}_X} \mathcal{F}
\longrightarrow
f^*\mathcal{E} \otimes_{\mathcal{O}_X} \mathcal{I}^\bullet
$$
を得，これは補題 \ref{cohomology-lemma-injective-tensor-finite-locally-free} により
入射分解である。$f_*$ を適用すると
$$
Rf_*(f^*\mathcal{E} \otimes_{\mathcal{O}_X} \mathcal{F})
=
f_*(f^*\mathcal{E} \otimes_{\mathcal{O}_X} \mathcal{I}^\bullet).
$$
を得る。したがって，$\mathcal{O}_X$-加群 $\mathcal{F}$ に関して関手的に
$f_*(f^*\mathcal{E}\otimes_{\mathcal{O}_X}\mathcal{F})=
\mathcal{E}\otimes_{\mathcal{O}_Y}f_*(\mathcal{F})$ であることを
示せば補題が従う。$\mathcal{E}=\mathcal{O}_Y^{\oplus n}$ の場合は
明らかであり，一般の場合は $Y$ 上局所的に作業すれば従う。
細部は省略する。
\end{proof}

\noindent
$f:X\to Y$ を環付き空間の射とする。
$E\in D(\mathcal{O}_X)$ および $K\in D(\mathcal{O}_Y)$ とする。
追加の仮定なしに写像
\begin{equation}
\label{cohomology-equation-projection-formula-map}
Rf_*E \otimes^\mathbf{L}_{\mathcal{O}_Y} K
\longrightarrow
Rf_*(E \otimes^\mathbf{L}_{\mathcal{O}_X} Lf^*K)
\end{equation}
が存在する。これは標準写像
$$
Lf^*(Rf_*E \otimes^\mathbf{L}_{\mathcal{O}_Y} K) =
Lf^*Rf_*E \otimes^\mathbf{L}_{\mathcal{O}_X} Lf^*K
\longrightarrow
E \otimes^\mathbf{L}_{\mathcal{O}_X} Lf^*K
$$
の随伴である。ここで標準写像は $Lf^*Rf_*E\to E$ と補題
\ref{cohomology-lemma-pullback-tensor-product} および \ref{cohomology-lemma-adjoint} から生じる。
射影公式のかなり一般的な形は次のとおりである。

\begin{lemma}
\label{cohomology-lemma-projection-formula-perfect}
$f:X\to Y$ を環付き空間の射とする。
$E\in D(\mathcal{O}_X)$ および $K\in D(\mathcal{O}_Y)$ とする。
$K$ が完全ならば
$$
Rf_*E \otimes^\mathbf{L}_{\mathcal{O}_Y} K =
Rf_*(E \otimes^\mathbf{L}_{\mathcal{O}_X} Lf^*K)
$$
が $D(\mathcal{O}_Y)$ で成り立つ。
\end{lemma}

\begin{proof}
(\ref{cohomology-equation-projection-formula-map}) が同型であることは
$Y$ 上局所的に確認してよい。すなわち，写像の $V_j$ への制限が
同型となる開被覆 $\{V_j\to Y\}$ を見つければよい。
完全対象の定義により，$K$ が $\mathcal{O}_Y$-加群の狭義完全複体で
表現されると仮定してよい。一般に，$K=K_1\oplus K_2$ に対して
主張が成り立つための必要十分条件は，$K_1$ と $K_2$ に対して
成り立つことである。したがって $K$ が有限自由
$\mathcal{O}_Y$-加群からなる有限複体であると仮定してよい。
この場合，愚直切断を用いる簡単な議論により，$K$ が有限自由
$\mathcal{O}_Y$-加群で表現される場合へ帰着する。
主張は変数 $K$ の有限直和因子をとることに関して不変なので，
$K=\mathcal{O}_Y[n]$ の場合を証明すれば十分であり，この場合は自明である。
\end{proof}

\noindent
射影公式が完全な一般性のもとで成り立つ場合を次に述べる。

\begin{lemma}
\label{cohomology-lemma-projection-formula-closed-immersion}
$f:X\to Y$ を環付き空間の射とし，$f$ は閉部分集合への同相写像で
あるとする。このとき (\ref{cohomology-equation-projection-formula-map}) は常に同型である。
\end{lemma}

\begin{proof}
$f$ は閉部分集合への同相写像なので，関手 $f_*$ は完全である
（Modules, Lemma \ref{modules-lemma-i-star-exact}）。
したがって任意の代表複体に $f_*$ を適用して $Rf_*$ を計算できる。
$K$ を表現する $\mathcal{O}_Y$-加群の K-平坦複体
$\mathcal{K}^\bullet$ と，$E$ を表現する
$\mathcal{O}_X$-加群の任意の複体 $\mathcal{E}^\bullet$ を選ぶ。
このとき $Lf^*K$ は $f^*\mathcal{K}^\bullet$ で表現され，
これは $\mathcal{O}_X$-加群の K-平坦複体である
（補題 \ref{cohomology-lemma-pullback-K-flat}）。したがって
(\ref{cohomology-equation-projection-formula-map}) の右辺は
$$
f_*\text{Tot}(\mathcal{E}^\bullet
\otimes_{\mathcal{O}_X} f^*\mathcal{K}^\bullet)
$$
で表現される。同じ理由により左辺は
$$
\text{Tot}(f_*\mathcal{E}^\bullet \otimes_{\mathcal{O}_Y} \mathcal{K}^\bullet)
$$
で表現される。$f_*$ は直和と可換するので
（Modules, Lemma \ref{modules-lemma-i-star-right-adjoint}），
$$
f_*(\mathcal{E} \otimes_{\mathcal{O}_X} f^*\mathcal{K}) =
f_*\mathcal{E} \otimes_{\mathcal{O}_Y} \mathcal{K}
$$
を，任意の $\mathcal{O}_X$-加群 $\mathcal{E}$ と
$\mathcal{O}_Y$-加群 $\mathcal{K}$ に対して示せば十分である。
これを茎で確認する。$y\in Y$ とする。$y\not\in f(X)$ ならば，
両辺の茎は零である。$y=f(x)$ ならば
$$
\mathcal{E}_x \otimes_{\mathcal{O}_{X, x}}
(\mathcal{O}_{X, x} \otimes_{\mathcal{O}_{Y, y}} \mathcal{F}_y) =
\mathcal{E}_x \otimes_{\mathcal{O}_{Y, y}} \mathcal{F}_y
$$
を示せばよい（Sheaves, Lemma
\ref{sheaves-lemma-stalks-closed-pushforward} および補題
\ref{sheaves-lemma-stalk-pullback-modules} を用いた）。
この等式は成り立つので，補題が証明された。
\end{proof}

\begin{remark}
\label{cohomology-remark-compatible-with-diagram}
写像 (\ref{cohomology-equation-projection-formula-map}) は，次の意味で Remark
\ref{cohomology-remark-base-change} の基底変換写像と両立する。すなわち，
$$
\xymatrix{
X' \ar[r]_{g'} \ar[d]_{f'} &
X \ar[d]^f \\
Y' \ar[r]^g &
Y
}
$$
が環付き空間の可換図式であるとする。
$E\in D(\mathcal{O}_X)$ および $K\in D(\mathcal{O}_Y)$ とする。
このとき図式
$$
\xymatrix{
Lg^*(Rf_*E \otimes^\mathbf{L}_{\mathcal{O}_Y} K) \ar[r]_p \ar[d]_t &
Lg^*Rf_*(E \otimes^\mathbf{L}_{\mathcal{O}_X} Lf^*K) \ar[d]_b \\
Lg^*Rf_*E \otimes^\mathbf{L}_{\mathcal{O}_{Y'}} Lg^*K \ar[d]_b &
Rf'_*L(g')^*(E \otimes^\mathbf{L}_{\mathcal{O}_X} Lf^*K) \ar[d]_t \\
Rf'_*L(g')^*E \otimes^\mathbf{L}_{\mathcal{O}_{Y'}} Lg^*K \ar[rd]_p &
Rf'_*(L(g')^*E \otimes^\mathbf{L}_{\mathcal{O}_{Y'}} L(g')^*Lf^*K) \ar[d]_c \\
& Rf'_*(L(g')^*E \otimes^\mathbf{L}_{\mathcal{O}_{Y'}} L(f')^*Lg^*K)
}
$$
は可換である。ここでラベル $t$ の矢印は補題
\ref{cohomology-lemma-pullback-tensor-product} の適用，ラベル $b$ の矢印は
Remark \ref{cohomology-remark-base-change} の適用，ラベル $p$ の矢印は
(\ref{cohomology-equation-projection-formula-map}) の適用によって得られ，
$c$ は $L(g')^*\circ Lf^*=L(f')^*\circ Lg^*$ から生じる。
確認は省略する。
\end{remark}















\section{Berthelot と Ogus が導入した作用素}
\label{cohomology-section-eta}

\noindent
この節では More on Algebra, Section
\ref{more-algebra-section-eta} で始めた議論を続ける。
まずその節を読むことを強く勧める。

\begin{lemma}
\label{cohomology-lemma-invertible-ideal-sheaf}
$(X,\mathcal{O}_X)$ を環付き空間とし，
$\mathcal{I}\subset\mathcal{O}_X$ をイデアルの層とする。
次の二条件を考える。
\begin{enumerate}
\item すべての $x\in X$ に対して，$x$ の開近傍 $U\subset X$ と
$f\in\mathcal{I}(U)$ で，$\mathcal{I}|_U=\mathcal{O}_U\cdot f$ かつ
$f:\mathcal{O}_U\to\mathcal{O}_U$ が単射となるものが存在する。
\item $\mathcal{I}$ は $\mathcal{O}_X$-加群として可逆である。
\end{enumerate}
このとき (1) は (2) を含意し，構造層のすべての茎
$\mathcal{O}_{X,x}$ が局所環ならば逆も成り立つ。
\end{lemma}

\begin{proof}
省略する。ヒント：Modules, Lemma
\ref{modules-lemma-invertible-is-locally-free-rank-1} を用いる。
\end{proof}

\begin{situation}
\label{cohomology-situation-eta}
$(X,\mathcal{O}_X)$ を環付き空間とする。
$\mathcal{I}\subset\mathcal{O}_X$ を，補題
\ref{cohomology-lemma-invertible-ideal-sheaf} の条件 (1) を満たすイデアルの層とする
\footnote{この節の議論は，Modules, Section
\ref{modules-section-invertible} で定義されるように，
$\mathcal{I}$ が可逆 $\mathcal{O}_X$-加群であることだけを
仮定する場合へ一般化できる。}。
\end{situation}

\begin{lemma}
\label{cohomology-lemma-I-torsion-free}
状況 \ref{cohomology-situation-eta} において，$\mathcal{F}$ を
$\mathcal{O}_X$-加群とする。次は同値である。
\begin{enumerate}
\item $\mathcal{I}$ によって零化される切断の部分層
$\mathcal{F}[\mathcal{I}]\subset\mathcal{F}$ は零である。
\item すべての $n\geq1$ に対して部分層
$\mathcal{F}[\mathcal{I}^n]$ は零である。
\item 乗法写像
$\mathcal{I} \otimes_{\mathcal{O}_X} \mathcal{F} \to \mathcal{F}$
は単射である。
\item ある $f\in\mathcal{I}(U)$ に対して
$\mathcal{I}|_U=\mathcal{O}_U\cdot f$ となるすべての開集合
$U\subset X$ に対して，写像
$f:\mathcal{F}|_U\to\mathcal{F}|_U$ は単射である。
\item すべての $x\in X$ と，イデアル
$\mathcal{I}_x\subset\mathcal{O}_{X,x}$ の生成元 $f$ に対して，
$f$ は茎 $\mathcal{F}_x$ 上の非零因子である。
\end{enumerate}
\end{lemma}

\begin{proof}
省略する。
\end{proof}

\noindent
状況 \ref{cohomology-situation-eta} において，$\mathcal{F}$ を
$\mathcal{O}_X$-加群とする。補題 \ref{cohomology-lemma-I-torsion-free} の
同値な条件が成り立つとき，$\mathcal{F}$ は
{\it $\mathcal{I}$-捩れなし} であるという。このとき任意の
$i\in\mathbf{Z}$ に対して
$$
\mathcal{I}^i\mathcal{F} =
\mathcal{I}^{\otimes i} \otimes_{\mathcal{O}_X} \mathcal{F}
$$
と書く。すると包含
$$
\ldots \subset
\mathcal{I}^{i + 1}\mathcal{F} \subset
\mathcal{I}^i\mathcal{F} \subset
\mathcal{I}^{i - 1}\mathcal{F} \subset \ldots
$$
を得る。加群 $\mathcal{I}^i\mathcal{F}$ は
$\mathcal{O}_X$-加群として局所的には $\mathcal{F}$ と同型だが，
大域的にはそうとは限らない。

\medskip\noindent
$\mathcal{F}^\bullet$ を $\mathcal{I}$-捩れなし
$\mathcal{O}_X$-加群からなる複体とし，その微分を
$d^i : \mathcal{F}^i \to \mathcal{F}^{i + 1}$ とする。
このとき $\eta_\mathcal{I}\mathcal{F}^\bullet$ を，項が
\begin{align*}
(\eta_\mathcal{I}\mathcal{F})^i
& =
\Ker\left(
d^i, -1 :
\mathcal{I}^i\mathcal{F}^i \oplus \mathcal{I}^{i + 1}\mathcal{F}^{i + 1}
\to
\mathcal{I}^i\mathcal{F}^{i + 1}
\right) \\
& =
\Ker\left(d^i :
\mathcal{I}^i\mathcal{F}^i
\to
\mathcal{I}^i\mathcal{F}^{i + 1}/
\mathcal{I}^{i + 1}\mathcal{F}^{i + 1}
\right)
\end{align*}
で，微分が $d^i$ から誘導される複体として定義する。
言い換えると，$(\eta_\mathcal{I}\mathcal{F})^i$ の局所切断 $s$ とは，
$\mathcal{I}^i\mathcal{F}^i$ の局所切断 $s$ であって，
$\mathcal{I}^i\mathcal{F}^{i + 1}$ におけるその像 $d^i(s)$ が
部分層 $\mathcal{I}^{i + 1}\mathcal{F}^{i + 1}$ に属するものにほかならない。
$\eta_\mathcal{I}\mathcal{F}^\bullet$ もまた
$\mathcal{I}$-捩れなし加群からなる複体であることに注意する。

\medskip\noindent
$a^\bullet : \mathcal{F}^\bullet \to \mathcal{G}^\bullet$ を
$\mathcal{I}$-捩れなし $\mathcal{O}_X$-加群からなる複体の写像とする。
このとき複体の写像
$$
\eta_\mathcal{I} a^\bullet :
\eta_\mathcal{I}\mathcal{F}^\bullet
\longrightarrow
\eta_\mathcal{I}\mathcal{G}^\bullet
$$
を得る。これは写像
$\mathcal{I}^i\mathcal{F}^i \to \mathcal{I}^i\mathcal{G}^i$
から誘導される。これにより
$\mathcal{I}$-捩れなし $\mathcal{O}_X$-加群からなる複体の圏の
自己関手が得られることを読者は確かめられる。

\medskip\noindent
$a^\bullet, b^\bullet : \mathcal{F}^\bullet \to \mathcal{G}^\bullet$ を
$\mathcal{I}$-捩れなし $\mathcal{O}_X$-加群からなる複体の二つの写像とし，
$h = \{h^i : \mathcal{F}^i \to \mathcal{G}^{i - 1}\}$ を
$a^\bullet$ と $b^\bullet$ の間のホモトピーとする。このとき
$\eta_\mathcal{I}h$ を写像族
$(\eta_\mathcal{I}h)^i : (\eta_\mathcal{I}\mathcal{F})^i \to
(\eta_\mathcal{I}\mathcal{G})^{i - 1}$
で，$x$ を $h^i(x)$ に送るものとして定義する。これは確かに意味をもつ。
実際，$x$ が $\mathcal{I}^i\mathcal{F}^i$ の局所切断ならば，
$h^i(x)$ は $\mathcal{I}^i\mathcal{G}^{i - 1}$ の局所切断であり，
これは明らかに $(\eta_\mathcal{I}\mathcal{G})^{i - 1}$ に含まれる。
$\eta_\mathcal{I}h$ が
$\eta_\mathcal{I}a^\bullet$ と $\eta_\mathcal{I}b^\bullet$ の間の
ホモトピーであることを読者は確かめられる。以上から関手
$$
\eta_f :
K(\mathcal{I}\text{-torsion free }\mathcal{O}_X\text{-modules})
\longrightarrow
K(\mathcal{I}\text{-torsion free }\mathcal{O}_X\text{-modules})
$$
が得られる。これは $\mathcal{I}$-捩れなし
$\mathcal{O}_X$-加群の加法圏のホモトピー圏
(Derived Categories, Section \ref{derived-section-homotopy}) 上の関手である。
$\eta_\mathcal{I}$ が三角圏の完全関手となるような意味づけは存在しない。
More on Algebra, Example
\ref{more-algebra-example-eta-not-distinguished} と比較せよ。

\begin{lemma}
\label{cohomology-lemma-eta-stalks}
状況 \ref{cohomology-situation-eta} において，
$\mathcal{F}^\bullet$ を $\mathcal{I}$-捩れなし
$\mathcal{O}_X$-加群からなる複体とする。
$x \in X$ に対して生成元 $f \in \mathcal{I}_x$ を選ぶ。このとき茎
$(\eta_\mathcal{I}\mathcal{F}^\bullet)_x$ は，More on Algebra, Section
\ref{more-algebra-section-eta} で構成した複体
$\eta_f\mathcal{F}^\bullet_x$ と標準的に同型である。
\end{lemma}

\begin{proof}
省略する。
\end{proof}

\begin{lemma}
\label{cohomology-lemma-eta-first-property}
状況 \ref{cohomology-situation-eta} において，
$\mathcal{F}^\bullet$ を $\mathcal{I}$-捩れなし
$\mathcal{O}_X$-加群からなる複体とする。コホモロジー層の標準的同型
$$
\mathcal{I}^{\otimes i} \otimes_{\mathcal{O}_X}
\left(
H^i(\mathcal{F}^\bullet)/H^i(\mathcal{F}^\bullet)[\mathcal{I}]
\right)
\longrightarrow H^i(\eta_\mathcal{I}\mathcal{F}^\bullet)
$$
が存在する。
\end{lemma}

\begin{proof}
写像
$$
\mathcal{I}^{\otimes i} \otimes_{\mathcal{O}_X} H^i(\mathcal{F}^\bullet)
\longrightarrow
H^i(\eta_\mathcal{I}\mathcal{F}^\bullet)
$$
を次のように定義する。$g$ を $\mathcal{I}^{\otimes i}$ の局所切断，
$\overline{s}$ を $H^i(\mathcal{F}^\bullet)$ の局所切断とする。
すると $\overline{s}$ は局所的には
$\Ker(d^i : \mathcal{F}^i \to \mathcal{F}^{i + 1})$
の局所切断 $s$ の類である。そこで $g \otimes \overline{s}$ を
$(\eta_\mathcal{I}\mathcal{F})^i \subset \mathcal{I}^i\mathcal{F}$
の局所切断 $gs$ に送る。もちろん $gs$ は
$\eta_\mathcal{I}\mathcal{F}^\bullet$ 上の $d^i$ の核に属するので，
$H^i(\eta_\mathcal{I}\mathcal{F}^\bullet)$ の局所切断を定める。
これが良定義であることの確認には何の問題もない。
この写像は補題で述べた同型を経由すると主張する。
これは茎上で確かめればよく，補題 \ref{cohomology-lemma-eta-stalks} により
More on Algebra, Lemma
\ref{more-algebra-lemma-eta-first-property} の結果に帰着する。
\end{proof}

\begin{lemma}
\label{cohomology-lemma-eta-qis}
状況 \ref{cohomology-situation-eta} において，
$\mathcal{F}^\bullet \to \mathcal{G}^\bullet$ を
$\mathcal{I}$-捩れなし $\mathcal{O}_X$-加群からなる複体の写像とする。
このとき誘導される写像
$\eta_\mathcal{I}\mathcal{F}^\bullet \to \eta_\mathcal{I}\mathcal{G}^\bullet$
も擬同型である。
\end{lemma}

\begin{proof}
これは補題 \ref{cohomology-lemma-eta-first-property} の同型が
複体の写像と両立することから従う。
\end{proof}

\begin{lemma}
\label{cohomology-lemma-Leta}
状況 \ref{cohomology-situation-eta} において，加法関手\footnote{この関手は
完全ではない，すなわち，特別三角形を特別三角形に移すとは限らないことに注意せよ。}
$L\eta_\mathcal{I} : D(\mathcal{O}_X) \to D(\mathcal{O}_X)$
で，$D(\mathcal{O}_X)$ の対象 $M$ が $\mathcal{I}$-捩れなし
$\mathcal{O}_X$-加群からなる複体 $\mathcal{F}^\bullet$ によって
表されるならば
$L\eta_\mathcal{I}M = \eta_\mathcal{I}\mathcal{F}^\bullet$
となるものが存在する。射についても同様である。
\end{lemma}

\begin{proof}
$\mathcal{I}$-捩れなし $\mathcal{O}_X$-加群からなる充満部分圏を
$\mathcal{T} \subset \textit{Mod}(\mathcal{O}_X)$ と書く。
これに対応して包含
$$
K(\mathcal{T})
\quad\subset\quad
K(\textit{Mod}(\mathcal{O}_X)) = K(\mathcal{O}_X)
$$
があり，$K(\mathcal{T})$ は $K(\mathcal{O}_X)$ の
充満三角部分圏となる。
$S \subset \text{Arrows}(K(\mathcal{T}))$ を擬同型全体とする。
写像
$$
S^{-1}K(\mathcal{T}) \longrightarrow D(\mathcal{O}_X)
$$
が三角圏の同値であることを示すため，Derived Categories, Lemma
\ref{derived-lemma-localization-subcategory} を適用する。
同補題によれば，次を示せば十分である：
$\mathcal{O}_X$-加群からなる複体 $\mathcal{G}^\bullet$ が与えられたとき，
$\mathcal{I}$-捩れなし $\mathcal{O}_X$-加群からなる複体
$\mathcal{F}^\bullet$ と擬同型
$\mathcal{F}^\bullet \to \mathcal{G}^\bullet$ が存在する。
補題 \ref{cohomology-lemma-K-flat-resolution} により，
$\mathcal{F}^\bullet$ が K-平坦であり（このこと自体は用いない），
平坦 $\mathcal{O}_X$-加群 $\mathcal{F}^i$ からなるような擬同型
$\mathcal{F}^\bullet \to \mathcal{G}^\bullet$ を取れる。
補題 \ref{cohomology-lemma-I-torsion-free} の第三の特徴づけにより，
平坦 $\mathcal{O}_X$-加群は $\mathcal{I}$-捩れなし
$\mathcal{O}_X$-加群である。これで証明は完了する。

\medskip\noindent
以上の準備により $L\eta_f$ を定義できる。
実際，この節の補題 \ref{cohomology-lemma-I-torsion-free} に続く議論によって，
すでに良定義な関手
$$
K(\mathcal{T}) \xrightarrow{\eta_f} K(\mathcal{T}) \to
K(\mathcal{O}_X) \to D(\mathcal{O}_X)
$$
を得ており，補題 \ref{cohomology-lemma-eta-qis} によればこれは擬同型を
擬同型に移す。したがって Categories, Lemma
\ref{categories-lemma-properties-left-localization} により，
この関手は $S^{-1}K(\mathcal{T}) = D(\mathcal{O}_X)$ を経由する。
\end{proof}

\noindent
状況 \ref{cohomology-situation-eta} において Bockstein 作用素を構成しよう。
まず可換図式
$$
\xymatrix{
0 \ar[r] &
\mathcal{I}^{i + 1} \ar[r] \ar[d] &
\mathcal{I}^i \ar[r] \ar[d] &
\mathcal{I}^i/\mathcal{I}^{i + 1}
\ar[r] \ar@{=}[d] &
0 \\
0 \ar[r] &
\mathcal{I}^{i + 1}/\mathcal{I}^{i + 2} \ar[r] &
\mathcal{I}^i/\mathcal{I}^{i + 2} \ar[r] &
\mathcal{I}^i/\mathcal{I}^{i + 1} \ar[r] &
0
}
$$
が存在することに注意する。その各行は $\mathcal{O}_X$-加群の
短完全列である。$M$ を $D(\mathcal{O}_X)$ の対象とする。
上の図式を $M$ とテンソル積することにより，
$$
\xymatrix{
M \otimes^\mathbf{L} \mathcal{I}^{i + 1} \ar[r] \ar[d] &
M \otimes^\mathbf{L} \mathcal{I}^i \ar[r] \ar[d] &
M \otimes^\mathbf{L} \mathcal{I}^i/\mathcal{I}^{i + 1} \ar[d]^{\text{id}} \\
M \otimes^\mathbf{L} \mathcal{I}^{i + 1}/\mathcal{I}^{i + 2} \ar[r] &
M \otimes^\mathbf{L} \mathcal{I}^i/\mathcal{I}^{i + 2} \ar[r] &
M \otimes^\mathbf{L} \mathcal{I}^i/\mathcal{I}^{i + 1}
}
$$
という特別三角形の射を得る。特に，下の三角形に付随する
コホモロジー層の長完全列は，すべての $i \in \mathbf{Z}$ に対して
{\it Bockstein 作用素}
$$
\beta = \beta^i :
H^i(M \otimes^\mathbf{L} \mathcal{I}^i/\mathcal{I}^{i + 1})
\longrightarrow
H^{i + 1}(M \otimes^\mathbf{L} \mathcal{I}^{i + 1}/\mathcal{I}^{i + 2})
$$
を定める。後で用いるため，上の可換図式によって
Bockstein 作用素が
\begin{equation}
\label{cohomology-equation-factorization-bockstein}
\vcenter{
\xymatrix{
H^i(M \otimes^\mathbf{L} \mathcal{I}^i/\mathcal{I}^{i + 1})
\ar[r]_\delta \ar[rd]_\beta &
H^{i + 1}(M \otimes^\mathbf{L} \mathcal{I}^{i + 1}) \ar[d] \\
&
H^{i + 1}(M \otimes^\mathbf{L} \mathcal{I}^{i + 1}/\mathcal{I}^{i + 2})
}
}
\end{equation}
と分解することを記しておく。ここで $\delta$ は，上の可換図式の
上段の特別三角形から生じる境界作用素である。複体
\begin{equation}
\label{cohomology-equation-complex-bocksteins}
H^\bullet(M/\mathcal{I}) =
\left[
\begin{matrix}
\ldots \\
\downarrow \\
H^{i - 1}(M \otimes^\mathbf{L} \mathcal{I}^{i - 1}/\mathcal{I}^i) \\
\downarrow \beta \\
H^i(M \otimes^\mathbf{L} \mathcal{I}^i/\mathcal{I}^{i + 1}) \\
\downarrow \beta \\
H^{i + 1}(M \otimes^\mathbf{L} \mathcal{I}^{i + 1}/\mathcal{I}^{i + 2}) \\
\downarrow \\
\ldots
\end{matrix}
\right]
\end{equation}
を得る。すなわち $\beta \circ \beta = 0$ である。
実際，これは茎上で確かめればよく，その場合には
More on Algebra, Section \ref{more-algebra-section-eta} で示した
代数における対応する結果から従う。
別証：短完全列
$0 \to \mathcal{I}^{i + 1}/\mathcal{I}^{i + 2}
\to \mathcal{I}^i/\mathcal{I}^{i + 2}
\to \mathcal{I}^i/\mathcal{I}^{i + 1} \to 0$
は $D(\mathcal{O}_X)$ における写像
$b^i : \mathcal{I}^i/\mathcal{I}^{i + 1} \to
(\mathcal{I}^{i + 1}/\mathcal{I}^{i + 2})[1]$
を定める。これらを $M$ とテンソル積してコホモロジー層を取ると，
上の写像 $\beta$ が誘導される。次に，合成
$b^{i + 1}[1] \circ b^i : \mathcal{I}^i/\mathcal{I}^{i + 1} \to
(\mathcal{I}^{i + 1}/\mathcal{I}^{i + 2})[1] \to
(\mathcal{I}^{i + 2}/\mathcal{I}^{i + 3})[2]$
が $D(\mathcal{O}_X)$ で零であることを示す。
そのために Derived Categories, Lemma
\ref{derived-lemma-cup-ext-1-zero} の判定法と，
加群 $\mathcal{I}^i/\mathcal{I}^{i + 3}$ が
$\mathcal{I}^{i + 1}/\mathcal{I}^{i + 3}$ の
$\mathcal{I}^i/\mathcal{I}^{i + 1}$ による拡大であることを用いる。

\begin{lemma}
\label{cohomology-lemma-eta-second-property}
状況 \ref{cohomology-situation-eta} において，
$M$ を $D(\mathcal{O}_X)$ の対象とする。標準的同型
$$
L\eta_\mathcal{I}M \otimes^\mathbf{L} \mathcal{O}_X/\mathcal{I}
\longrightarrow
H^\bullet(M/\mathcal{I})
$$
が $D(\mathcal{O}_X)$ において存在する。ここで右辺は複体
(\ref{cohomology-equation-complex-bocksteins}) である。
\end{lemma}

\begin{proof}
補題 \ref{cohomology-lemma-eta-qis} における $L\eta_\mathcal{I}$ の構成により，
$M$ は $\mathcal{I}$-捩れなし $\mathcal{O}_X$-加群からなる複体
$\mathcal{F}^\bullet$ によって表されると仮定してよい。
このとき $L\eta_\mathcal{I}M$ は複体
$\eta_\mathcal{I}\mathcal{F}^\bullet$ によって表され，これも
$\mathcal{I}$-捩れなし $\mathcal{O}_X$-加群からなる複体である。
したがって
$L\eta_\mathcal{I}M \otimes^\mathbf{L} \mathcal{O}_X/\mathcal{I}$
は複体
$\eta_\mathcal{I}\mathcal{F}^\bullet \otimes \mathcal{O}_X/\mathcal{I}$
によって表される。同様に，複体 $H^\bullet(M/\mathcal{I})$ の項は
$H^i(\mathcal{F}^\bullet \otimes \mathcal{I}^i/\mathcal{I}^{i + 1})$
である。

\medskip\noindent
$f$ を $\mathcal{I}$ の局所生成元とする。
$s$ を $(\eta_\mathcal{I}\mathcal{F})^i$ の局所切断とする。
このとき $\mathcal{F}^i$ の局所切断 $s'$ を用いて
$s = f^is'$ と書け，同様に $\mathcal{F}^{i + 1}$ の局所切断 $t$ を
用いて $d^i(s) = f^{i + 1}t$ と書ける。したがって $d^i$ は
$f^is'$ を
$\mathcal{F}^{i + 1} \otimes \mathcal{I}^i/\mathcal{I}^{i + 1}$
の零に写す。それゆえ $s$ を
$H^i(\mathcal{F}^\bullet \otimes \mathcal{I}^i/\mathcal{I}^{i + 1})$
における $f^is'$ の類へ写せる。この規則は
$$
(\eta_\mathcal{I}\mathcal{F})^i \otimes \mathcal{O}_X/\mathcal{I}
\longrightarrow
H^i(\mathcal{F}^\bullet \otimes \mathcal{I}^i/\mathcal{I}^{i + 1})
$$
という $\mathcal{O}_X$-加群の写像を定める。計算により，これらの写像は
微分と両立することが分かる（本質的には $\beta$ が $f^is'$ の類を
$f^{i + 1}t$ の類に写すからである）。したがって，補題の主張にある射を
表す複体の写像を得る。

\medskip\noindent
証明を終えるため，前段落の構成が茎上では More on Algebra, Lemma
\ref{more-algebra-lemma-eta-second-property}
で構成した写像と一致することに注意すればよい。これで結論が従う。
\end{proof}

\begin{lemma}
\label{cohomology-lemma-eta-tensor-invertible}
状況 \ref{cohomology-situation-eta} において，
$\mathcal{F}^\bullet$ を $\mathcal{I}$-捩れなし
$\mathcal{O}_X$-加群からなる複体とする。
$\mathcal{L}$ を可逆 $\mathcal{O}_X$-加群とする。このとき
$\eta_\mathcal{I}(\mathcal{F}^\bullet \otimes \mathcal{L}) =
(\eta_\mathcal{I}\mathcal{F}^\bullet) \otimes \mathcal{L}$.
\end{lemma}

\begin{proof}
構成から直ちに従う。
\end{proof}

\begin{lemma}
\label{cohomology-lemma-eta-cohomology-locally-free}
状況 \ref{cohomology-situation-eta} において，
$M$ を $D(\mathcal{O}_X)$ の対象とする。
$x \in X$ とし，$\mathcal{O}_{X, x}$ は零でないとする。
$H^i(M)_x$ が $\mathcal{O}_{X, x}$ 上有限自由ならば，
$H^i(L\eta_\mathcal{I}M)_x$ も $\mathcal{O}_{X, x}$ 上有限自由で，
その階数は同じである。
\end{lemma}

\begin{proof}
実際，$f \in \mathcal{O}_{X, x}$ が茎 $\mathcal{I}_x$ を生成するとする。
このとき $f$ は $\mathcal{O}_{X, x}$ の非零因子であり，したがって
$H^i(M)_x[f] = 0$ である。それゆえ補題
\ref{cohomology-lemma-eta-first-property} により，
$H^i(L\eta_\mathcal{I}M)_x$ は
$\mathcal{I}^i_x \otimes_{\mathcal{O}_{X, x}} H^i(M)_x$
と同型である。これは所望のとおり同じ階数の自由加群である。
\end{proof}
